% !TeX program = xelatex
% Standalone English volume. Use XeLaTeX until the references stabilize.
\documentclass[a4paper,11pt]{article}
\usepackage[T1]{fontenc}
\usepackage{lmodern}
\usepackage[english]{babel}
\usepackage[left=26mm,right=26mm,top=25mm,bottom=25mm,headheight=15pt,headsep=7mm,footskip=11mm]{geometry}
\usepackage{amsmath,amssymb,mathtools,amsthm}
\newtheorem{proposition}{Proposition}[section]
\newtheorem{lemma}[proposition]{Lemma}
\newtheorem{theorem}[proposition]{Theorem}
\newtheorem{corollary}[proposition]{Corollary}
\newtheorem{maintheorem}[proposition]{Main Theorem}
\theoremstyle{definition}
\newtheorem{definition}[proposition]{Definition}
\theoremstyle{remark}
\newtheorem{remark}[proposition]{Remark}
\makeatletter
\renewcommand*\l@section{\@dottedtocline{1}{0em}{2.4em}}
\renewcommand*\l@subsection{\@dottedtocline{2}{1.45em}{3.45em}}
\renewcommand*\l@subsubsection{\@dottedtocline{3}{3.8em}{4.55em}}
\makeatother
\usepackage{xcolor,fancyhdr,titlesec,needspace,xurl}
\usepackage{longtable,multicol}
\usepackage{hyperref}
\definecolor{ink}{HTML}{1A3542}
\definecolor{muted}{HTML}{4F606A}
\hypersetup{unicode=true,colorlinks=true,linkcolor=ink,filecolor=ink,urlcolor=ink,citecolor=ink,bookmarksnumbered=true,
 pdfdisplaydoctitle=true,pdflang={en},
 pdftitle={The generalized Riemann hypothesis for Dirichlet L-functions. Edition 930, Volume XIII},
 pdfsubject={A natural-language reconstruction of the proposed proof},
 pdfauthor={A mathematical reconstruction}}
\titleformat{\section}{\Large\bfseries\color{ink}}{\thesection}{.65em}{}
\titlespacing*{\section}{0pt}{.8em}{.65em}
\setlength{\parindent}{1.2em}
\linespread{1.045}
\setlength{\parskip}{.22em plus .05em minus .03em}
\setlength{\emergencystretch}{2em}
\setcounter{topnumber}{2}
\setcounter{tocdepth}{3}\setcounter{secnumdepth}{3}
\widowpenalty=4500
\clubpenalty=4500
\pagestyle{fancy}\fancyhf{}
\fancyhead[L]{\small\color{muted}Dirichlet GRH \textperiodcentered\ Volume I}
\fancyhead[R]{\small\color{muted}Dirichlet zeros and power intervals}
\fancyfoot[C]{\small\thepage}
\renewcommand{\headrulewidth}{.3pt}
\newcommand{\lead}[1]{\par\smallskip\noindent{\bfseries\color{ink}#1.}\enspace}
\newcommand{\sour}[2]{\unskip\penalty150\hspace{.35em}{\small\hyperlink{ref#1}{[#2]}}}
\titleformat{\subsection}{\large\bfseries}{\thesubsection}{.6em}{}
\titlespacing*{\subsection}{0pt}{1.15em plus .15em}{.4em}
\setlength{\abovedisplayskip}{9pt plus 2pt minus 3pt}
\setlength{\belowdisplayskip}{9pt plus 2pt minus 3pt}
\newcommand{\RR}{\mathbb R}\newcommand{\CC}{\mathbb C}\newcommand{\QQ}{\mathbb Q}
\newcommand{\ZZ}{\mathbb Z}\newcommand{\ord}{\operatorname{ord}}
\newcommand{\Frob}{\operatorname{Frob}}\newcommand{\Gal}{\operatorname{Gal}}
\newcommand{\len}{\operatorname{length}}\newcommand{\Res}{\operatorname*{Res}}
\newcommand{\wmean}{\operatorname{av}}\newcommand{\Ker}{\operatorname{ker}}
\newcommand{\Lh}{\mathcal L}\newcommand{\Gg}{\mathcal G}
\newcommand{\real}{\operatorname{Re}}\newcommand{\imag}{\operatorname{Im}}
\newcommand{\abs}[1]{\left|#1\right|}\newcommand{\eps}{\varepsilon}
\newcommand{\anchoritem}[2]{\par\Needspace{2.8\baselineskip}\hypertarget{ref#1}{}\noindent{\bfseries\color{ink}#2.}\enspace}
\newcommand{\srcloc}[2]{{\small\nolinkurl{#1.lean}}, line~#2}
\numberwithin{equation}{section}
\raggedbottom

\newcommand{\Sel}{\mathcal S}
\newcommand{\Cl}{\operatorname{Cl}}
\newcommand{\Aut}{\operatorname{Aut}}
\newcommand{\GL}{\operatorname{GL}}
\newcommand{\OO}{\mathcal O}
\hypersetup{bookmarksdepth=3}


\addto\captionsenglish{\renewcommand{\partname}{Volume}}


\makeatletter
\expandafter\def\csname r@cor:original200\endcsname{{83.7}{57}{}{proposition.83.7}{Dirichlet-GRH-v930-Volume-XVI.pdf}}
\expandafter\def\csname r@cor:shortcenter460\endcsname{{9.27}{130}{}{proposition.9.27}{Dirichlet-GRH-v930-Volume-I.pdf}}
\expandafter\def\csname r@eq:Gchi\endcsname{{1.1}{3}{}{equation.1.1}{Dirichlet-GRH-v930-Volume-I.pdf}}
\expandafter\def\csname r@eq:gap\endcsname{{6.3}{26}{}{equation.6.3}{Dirichlet-GRH-v930-Volume-I.pdf}}
\expandafter\def\csname r@eq:letters111\endcsname{{9.1}{117}{}{equation.9.1}{Dirichlet-GRH-v930-Volume-I.pdf}}
\expandafter\def\csname r@lem:balance82\endcsname{{94.1}{37}{}{proposition.94.1}{Dirichlet-GRH-v930-Volume-XVIII.pdf}}
\expandafter\def\csname r@lem:degree181\endcsname{{104.6}{6}{}{proposition.104.6}{Dirichlet-GRH-v930-Volume-XX.pdf}}
\expandafter\def\csname r@lem:far14\endcsname{{7.1}{88}{}{proposition.7.1}{Dirichlet-GRH-v930-Volume-I.pdf}}
\expandafter\def\csname r@lem:flagdegree456\endcsname{{120.16}{27}{}{proposition.120.16}{Dirichlet-GRH-v930-Volume-XXIII.pdf}}
\expandafter\def\csname r@lem:good175\endcsname{{72.18}{60}{}{proposition.72.18}{Dirichlet-GRH-v930-Volume-XIV.pdf}}
\expandafter\def\csname r@lem:mass32\endcsname{{121.17}{48}{}{proposition.121.17}{Dirichlet-GRH-v930-Volume-XXIII.pdf}}
\expandafter\def\csname r@lem:periodinteger463\endcsname{{141.26}{25}{}{proposition.141.26}{Dirichlet-GRH-v930-Volume-XXVII.pdf}}
\expandafter\def\csname r@lem:rankoneshift449\endcsname{{120.11}{22}{}{proposition.120.11}{Dirichlet-GRH-v930-Volume-XXIII.pdf}}
\expandafter\def\csname r@lem:signalsupport452\endcsname{{141.14}{15}{}{proposition.141.14}{Dirichlet-GRH-v930-Volume-XXVII.pdf}}
\expandafter\def\csname r@lem:terminalfourier454\endcsname{{141.20}{20}{}{proposition.141.20}{Dirichlet-GRH-v930-Volume-XXVII.pdf}}
\expandafter\def\csname r@lem:unit179\endcsname{{104.1}{3}{}{proposition.104.1}{Dirichlet-GRH-v930-Volume-XX.pdf}}
\expandafter\def\csname r@lem:wordfirstrow450\endcsname{{141.11}{12}{}{proposition.141.11}{Dirichlet-GRH-v930-Volume-XXVII.pdf}}
\expandafter\def\csname r@nav:I-23\endcsname{{6.1}{24}{}{subsection.6.1}{Dirichlet-GRH-v930-Volume-I.pdf}}
\expandafter\def\csname r@nav:I-28\endcsname{{6.6}{32}{}{subsection.6.6}{Dirichlet-GRH-v930-Volume-I.pdf}}
\expandafter\def\csname r@prop:cartan244\endcsname{{95.2}{66}{}{proposition.95.2}{Dirichlet-GRH-v930-Volume-XVIII.pdf}}
\expandafter\def\csname r@prop:desc50\endcsname{{6.28}{61}{}{proposition.6.28}{Dirichlet-GRH-v930-Volume-I.pdf}}
\expandafter\def\csname r@prop:field12\endcsname{{6.4}{27}{}{proposition.6.4}{Dirichlet-GRH-v930-Volume-I.pdf}}
\expandafter\def\csname r@prop:minor181\endcsname{{104.8}{7}{}{proposition.104.8}{Dirichlet-GRH-v930-Volume-XX.pdf}}
\expandafter\def\csname r@prop:ranks47\endcsname{{92.1}{3}{}{proposition.92.1}{Dirichlet-GRH-v930-Volume-XVIII.pdf}}
\expandafter\def\csname r@prop:res23\endcsname{{6.11}{32}{}{proposition.6.11}{Dirichlet-GRH-v930-Volume-I.pdf}}
\expandafter\def\csname r@prop:short107\endcsname{{9.9}{114}{}{proposition.9.9}{Dirichlet-GRH-v930-Volume-I.pdf}}
\expandafter\def\csname r@prop:weights14\endcsname{{7.2}{89}{}{proposition.7.2}{Dirichlet-GRH-v930-Volume-I.pdf}}
\expandafter\def\csname r@sec:abel\endcsname{{4.1}{17}{}{subsection.4.1}{Dirichlet-GRH-v930-Volume-I.pdf}}
\expandafter\def\csname r@sec:diagonal32\endcsname{{121.15}{50}{}{subsection.121.15}{Dirichlet-GRH-v930-Volume-XXIII.pdf}}
\expandafter\def\csname r@sec:missing\endcsname{{6.6}{32}{}{subsection.6.6}{Dirichlet-GRH-v930-Volume-I.pdf}}
\expandafter\def\csname r@sec:ranks47\endcsname{{92.1}{3}{}{subsection.92.1}{Dirichlet-GRH-v930-Volume-XVIII.pdf}}
\expandafter\def\csname r@seriesguide\endcsname{{1.1}{4}{}{subsection.1.1}{Dirichlet-GRH-v930-Volume-I.pdf}}
\expandafter\def\csname r@thm:classden312\endcsname{{122.13}{71}{}{proposition.122.13}{Dirichlet-GRH-v930-Volume-XXIII.pdf}}
\expandafter\def\csname r@thm:common233\endcsname{{21.5}{42}{}{proposition.21.5}{Dirichlet-GRH-v930-Volume-IV.pdf}}
\expandafter\def\csname r@thm:conductor232\endcsname{{19.15}{12}{}{proposition.19.15}{Dirichlet-GRH-v930-Volume-IV.pdf}}
\expandafter\def\csname r@thm:cutoff233\endcsname{{21.4}{41}{}{proposition.21.4}{Dirichlet-GRH-v930-Volume-IV.pdf}}
\expandafter\def\csname r@thm:dyad234\endcsname{{110.4}{40}{}{proposition.110.4}{Dirichlet-GRH-v930-Volume-XXI.pdf}}
\expandafter\def\csname r@thm:even233\endcsname{{94.18}{52}{}{proposition.94.18}{Dirichlet-GRH-v930-Volume-XVIII.pdf}}
\expandafter\def\csname r@thm:first72\endcsname{{32.15}{40}{}{proposition.32.15}{Dirichlet-GRH-v930-Volume-VI.pdf}}
\expandafter\def\csname r@thm:fixed243\endcsname{{88.2}{54}{}{proposition.88.2}{Dirichlet-GRH-v930-Volume-XVII.pdf}}
\expandafter\def\csname r@thm:fullsignal452\endcsname{{141.15}{16}{}{proposition.141.15}{Dirichlet-GRH-v930-Volume-XXVII.pdf}}
\expandafter\def\csname r@thm:globalperiodcontent448\endcsname{{141.9}{10}{}{proposition.141.9}{Dirichlet-GRH-v930-Volume-XXVII.pdf}}
\expandafter\def\csname r@thm:harmonic188\endcsname{{104.28}{24}{}{proposition.104.28}{Dirichlet-GRH-v930-Volume-XX.pdf}}
\expandafter\def\csname r@thm:height42\endcsname{{6.18}{52}{}{proposition.6.18}{Dirichlet-GRH-v930-Volume-I.pdf}}
\expandafter\def\csname r@thm:integral181\endcsname{{104.9}{8}{}{proposition.104.9}{Dirichlet-GRH-v930-Volume-XX.pdf}}
\expandafter\def\csname r@thm:mixedfourier334\endcsname{{83.19}{68}{}{proposition.83.19}{Dirichlet-GRH-v930-Volume-XVI.pdf}}
\expandafter\def\csname r@thm:orientation46\endcsname{{6.23}{57}{}{proposition.6.23}{Dirichlet-GRH-v930-Volume-I.pdf}}
\expandafter\def\csname r@thm:outsidedifferent453\endcsname{{141.18}{18}{}{proposition.141.18}{Dirichlet-GRH-v930-Volume-XXVII.pdf}}
\expandafter\def\csname r@thm:polycenter460\endcsname{{9.26}{129}{}{proposition.9.26}{Dirichlet-GRH-v930-Volume-I.pdf}}
\expandafter\def\csname r@thm:project53\endcsname{{6.36}{68}{}{proposition.6.36}{Dirichlet-GRH-v930-Volume-I.pdf}}
\expandafter\def\csname r@thm:projectedmatrices489\endcsname{{143.10}{57}{}{proposition.143.10}{Dirichlet-GRH-v930-Volume-XXVII.pdf}}
\expandafter\def\csname r@thm:rangenew50\endcsname{{6.29}{62}{}{proposition.6.29}{Dirichlet-GRH-v930-Volume-I.pdf}}
\expandafter\def\csname r@thm:scale232\endcsname{{110.18}{53}{}{proposition.110.18}{Dirichlet-GRH-v930-Volume-XXI.pdf}}
\expandafter\def\csname r@thm:sector234\endcsname{{9.2}{106}{}{proposition.9.2}{Dirichlet-GRH-v930-Volume-I.pdf}}
\expandafter\def\csname r@thm:signedcovariance478\endcsname{{143.5}{53}{}{proposition.143.5}{Dirichlet-GRH-v930-Volume-XXVII.pdf}}
\expandafter\def\csname r@thm:single82\endcsname{{94.2}{37}{}{proposition.94.2}{Dirichlet-GRH-v930-Volume-XVIII.pdf}}
\expandafter\def\csname r@thm:tower170\endcsname{{72.26}{68}{}{proposition.72.26}{Dirichlet-GRH-v930-Volume-XIV.pdf}}
\expandafter\def\csname r@thm:transport233\endcsname{{94.16}{50}{}{proposition.94.16}{Dirichlet-GRH-v930-Volume-XVIII.pdf}}
\expandafter\def\csname r@thm:twists\endcsname{{5.1}{19}{}{proposition.5.1}{Dirichlet-GRH-v930-Volume-I.pdf}}
\expandafter\def\csname r@thm:weylrowjet449\endcsname{{120.12}{23}{}{proposition.120.12}{Dirichlet-GRH-v930-Volume-XXIII.pdf}}
\makeatother
\begin{document}
\setcounter{part}{12}
\fancyhead[L]{\small\color{muted}Dirichlet GRH \textperiodcentered\ Volume XIII}
\fancyhead[R]{\small\color{muted}CM theta derivatives}
\thispagestyle{plain}\begin{center}
{\LARGE\bfseries The generalized Riemann hypothesis\\[2pt]for Dirichlet \(L\)-functions}\\[6pt]
{\small 27 November 2026 \quad\textperiodcentered\quad Edition 930}\end{center}
\part{CM theta derivatives and Frobenius translates}\label{part:XIII}
\begin{abstract}\noindent We construct the CM theta derivatives, compute local divisibility and follow Frobenius through the prescribed rows. Interpolation at split primes gives integral and archimedean estimates for the marked elliptic jets.\end{abstract}
\setcounter{section}{63}
\noindent The \hyperref[seriesguide]{guide in Volume~I} lists all twenty-seven volumes.
\begingroup\small\setlength{\parskip}{0pt}\tableofcontents\endgroup
\clearpage

\section{Complex multiplication and integral sections}\label{nav:g-V-0}\label{nav:volume-XIII}

\subsection{Complex multiplication and the local prime-power count}\label{nav:V-00}
\label{sec:cm37}

The first binary character of the field in Proposition \ref{prop:field12} cuts out \(K=\QQ(\sqrt{-q})\), where \(q\equiv3\pmod4\). The three splitting behaviors in this one quadratic field determine which prime powers can contribute at a place. We first prove a bound for derivatives of integral functions on a formal group with complex multiplication by \(K\). It determines the same three local densities, including the factor of one half at \(q\).

\subsubsection{An integral bound for an invariant derivative}\label{nav:V-01}
Let \(p\) be a rational prime, let \(R\) be the ring of integers of a finite extension of \(\QQ_p\), and use the valuation \(v_p(p)=1\). Suppose a place \(v\) of \(K\) over \(p\) is nonsplit, so \(K_v\) is quadratic. Write \(e\) for its ramification index, \(Q\) for the size of its residue field, and \(\varpi\) for a uniformizer of \(\OO_{K_v}\). Assume \(R\) contains the image of \(\OO_{K_v}\) under the chosen tangent embedding. Let \(\mathcal F\) be a smooth commutative one-dimensional formal group over \(R\) with an action of \(\OO_{K_v}\) by formal endomorphisms whose tangent action is that embedding. For \(a\ge1\) define
\[
 L_{p,K}(a)=v_p(a!)-\frac1e
                        \left\lfloor\frac{a-1}{Q-1}\right\rfloor .
\]

\begin{theorem}[Integral derivatives on a product of CM formal groups]\label{thm:cm37}
Let \(g\ge1\), let \(D\) be any invariant derivation defined by an integral tangent vector on \(\mathcal F^g\), and let \(f\) be any integral formal function at the identity. Then
\begin{equation}
              v_p\big((D^a f)(0)\big)\ \ge\ L_{p,K}(a).        \label{eq:jet37}
\end{equation}
Inert primes have \((e,Q)=(1,p^2)\), and the ramified prime has \((e,Q)=(2,p)\). At a split prime integrality always gives \(v_p((D^a f)(0))\ge0\). A zero value has valuation \(+\infty\).
\end{theorem}
\begin{proof}
Start with any parameter \(t\) on the one-dimensional group. For the Teichm\"uller units \(\zeta\in\mu_{Q-1}\subset\OO_{K_v}\), the coordinate
\[
       T=\frac1{Q-1}\sum_{\zeta\in\mu_{Q-1}}\zeta^{-1}[\zeta](t)
\]
has linear coefficient one and is integral. Here a scalar acting on tangent coordinates is read through the given embedding. Reindexing the sum gives \([\zeta]^*T=\zeta T\). In this coordinate the endomorphism \([\varpi]\) has expansion \(P(T)=\varpi T+\sum_{i\ge2} b_iT^i\) with integral coefficients. Commutativity with the Teichm\"uller action implies \(b_i=0\) unless \(i\equiv1\pmod{Q-1}\). Thus every nonlinear term has degree at least \(Q\).

Over the fraction field the formal group admits a logarithm with linear coefficient one. In dimension one, integrate the invariant differential with value \(dT\) at zero, taking constant term zero. Translation invariance gives the logarithm addition identity, by differentiating in the second variable and then setting it to zero. Its inverse \(E(Z)=Z+\sum_{n\ge2}c_nZ^n\) is the exponential. The action in logarithm coordinates is its linear part since an additive power series in characteristic zero is linear. Consequently
\[
  (\varpi^n-\varpi)c_n
   =\sum_{i=Q}^{n} b_i
        \sum_{\substack{k_1+\cdots+k_i=n\\ k_j\ge1}}
                            c_{k_1}\cdots c_{k_i},\qquad c_1=1 .
\]
For \(n>1\), the factor on the left has valuation \(1/e\). Induction gives
\(v_p(c_n)\ge-e^{-1}\lfloor(n-1)/(Q-1)\rfloor\). In fact for \(i\ge Q\) the sum of the \(i\) floors for the indices \(k_j\), plus one, is at most \(\lfloor(n-1)/(Q-1)\rfloor\). This accounts for the one division by \(\varpi\) at each induction step, also over ramified coefficient rings.

In the product coordinates write the integral tangent vector as \((u_1,\ldots,u_g)\). Pull back the function along \(Z\mapsto(E(u_1Z),\ldots,E(u_g Z))\). Differentiation in \(Z\) on this formal one-parameter subgroup is the invariant derivation \(D\), so its \(a\)-th derivative at zero is \(a!\) times the coefficient of \(Z^a\). Every term of positive order contributing to this coefficient is a product of at least one exponential coefficient and integral scalars. For indices \(k_1+\cdots+k_i=a\) its denominator valuation is at most \(e^{-1}\sum_j\lfloor(k_j-1)/(Q-1)\rfloor\le e^{-1}\lfloor(a-1)/(Q-1)\rfloor\). A finite sum of such terms obeys the same bound. This proves \eqref{eq:jet37}. In all cases an invariant integral derivation on the formal group preserves its integral functions, proving the final assertion.
\end{proof}

The theorem applies also to a section of a line bundle regular in an integral frame near the identity, by applying it to its scalar in that frame. This is a statement about a scalar Taylor coefficient. When lower Taylor coefficients vanish, the first nonzero coefficient changes by a unit under a change of integral frame.

\subsubsection{The same quadratic field and the same weights}\label{nav:V-02}
Fix a common interval and the fields \(E_x=E_0(\sqrt{-q},\sqrt{-d})\), with \(q,d\) fixed throughout the interval. Write \(B=\Gal(E_x/\QQ)\), \(\Phi=\Gal(E_x/K)\), and use additive notation. The cosets and the bijection \(p_t:B\setminus\Phi\to\Phi\), \(g\mapsto(1,1,t)+g\), are those of Volume~I. For a lift \(F_p\) of Frobenius and the inertia subgroup \(I_p\), the contribution of \(p\) to a single count is
\begin{equation}
 C_p(t,a,w)=\frac{2\log p}{|I_p|}
 \sum_{j\ge1}\ \sum_{g\in(jF_p+I_p)\setminus\Phi}
                         (a/p^j-w_{p_t(g)})_+ .                \label{eq:ct37}
\end{equation}
The sum over all primes is \(C(t,a,w)\). This formula retains the inertia average even at primes ramified in \(E_0\).

Take \(b=\exp(x^\eps)\), \(a=\lfloor bx\rfloor\), and define
\[
                W=2\big(1+\lceil\log_2(a+1)\rceil\big).
\]
Give every distinguished position weight \(\lceil b\rceil\) and every other position weight \(W\). For large \(x\), \(\lceil b\rceil\ge W\) and \(a/W\ge\exp(x^{\eps/3})/3\). Indeed \(\log W=O_\eps(\log(2x))\), while \(\log a\ge x^\eps+\log x-o(1)\). These choices satisfy the general integral-weight estimate of Proposition \ref{prop:weights14}. In particular they keep the distinguished prime sums and the error budget for discarded characters on the same intervals.

\begin{theorem}[The local count at every prime and translate]\label{thm:count37}
For all these weights, at every \(x\) on a sufficiently distant common interval and for every \(t,u\) and stage, the following holds. At a prime nonsplit in \(K\), and for every integral CM jet as in Theorem~\ref{thm:cm37},
\[
             \frac{C_p(t,a,w)}{2\log p}
                    \ \le\ v_p\big((D^a f)(0)\big).
\]
At a split prime \(C_p=0\). These assertions hold also when one replaces the weights by any list of real numbers at least \(W\).
\end{theorem}
\begin{proof}
At an unramified prime of \(K\), inertia lies in \(\Phi\). At a split prime so does Frobenius, hence the inner sum is empty. At an inert prime that entire coset lies outside \(\Phi\) for odd \(j\), and lies inside for even \(j\). If \(p\) ramifies in \(K\), then \(I_p\to B/\Phi\) is onto and exactly half the elements of every Frobenius coset are outside \(\Phi\). Thus the normalized count is at most
\[
 U_p=\begin{cases}
       \displaystyle\sum_{j\ge1\ {\rm odd}}(a/p^j-W)_+,&p\text{ inert},\\
       \displaystyle\frac12\sum_{j\ge1}(a/p^j-W)_+,&p\text{ ramified}.
      \end{cases}
\]
Let \(S_p\) denote the corresponding sum of \(a/p^j\) without the cutoff and weights, retaining the factor one half in the ramified case. The geometric series give \(S_p=ap/(p^2-1)\) and \(S_p=a/(2(p-1))\), respectively. Legendre's formula \(v_p(a!)=(a-s_p(a))/(p-1)\), with \(s_p(a)\) the sum of the base-\(p\) digits, gives in both cases
\[
                 L_{p,K}(a)\ge S_p-s_p(a)/(p-1).
\]
Here \(s_p(a)/(p-1)\le1+\lfloor\log_p a\rfloor\le W/2\). If \(U_p>0\), at least one active term loses \(W\) at an inert prime and \(W/2\) at the ramified prime. All other terms are at most their terms in \(S_p\). Thus \(U_p\le S_p-W/2\le L_{p,K}(a)\). If \(U_p=0\), integrality of the derivative suffices. No bound depends on the translate, the size of the inertia group in \(E_x\), or the number of copies of the formal group.
\end{proof}

\subsubsection{Algebraic jets over a common number field}\label{nav:V-03}
For clarity we realize the local formal groups on the fields of the counts. For each \(K=\QQ(\sqrt{-q})\) take a CM elliptic curve \(C\) with endomorphism ring \(\OO_K\), over a number field containing \(K\) over which all its endomorphisms are defined. The classical algebraicity and integrality of the singular modulus of \(\CC/\OO_K\) give such a curve over \(\overline\QQ\), with integral \(j\)-invariant. Its finitely many required endomorphisms are defined over a finite extension. The potential good reduction criterion for elliptic curves then gives good reduction everywhere after a further finite extension. One can take a single extension by first adjoining full torsion of order \(3\) and \(4\). Indeed potential good reduction makes the inertia image on a prime-to-residue-characteristic Tate module finite. At each prime use the \(3\)-adic module if the residue characteristic is different from \(3\) and the \(2\)-adic module otherwise. A finite subgroup of the kernel modulo \(3\), respectively modulo \(4\), is trivial by the binomial expansion of a matrix of prime order congruent to the identity. The criterion of N\'eron, Ogg and Shafarevich gives good reduction. Endomorphisms extend to the abelian scheme by the N\'eron mapping property. These standard facts are recalled in \cite{XIII-X-IV-cm37,XIII-X-IV-silver37}.

Choose one such number field \(J\) containing every \(E_x\) on the interval. This is possible because \(q,d,K_\psi\) are fixed there and the real cyclotomic layers are nested with a largest index on the bounded interval. One may use \(C^g\) directly. The CM type given by the embeddings of \(E_x\) above the chosen embedding of \(K\) also has the integral model
\[
                \mathcal A=\OO_{E_x}\otimes_{\OO_K} C .
\]
Here the notation denotes the direct factor of \(C^n\) cut out by an idempotent representing the finite projective \(\OO_K\)-module \(\OO_{E_x}\) as a summand of \(\OO_K^n\). The idempotent and its complement add to the identity and their images give the product decomposition. This constructs the abelian scheme, with the indicated action of \(\OO_{E_x}\). Its tangent space in characteristic zero is \(E_x\otimes_K J\) tensored with the tangent line of \(C\), so it has the specified embeddings each with multiplicity one.

At a nonsplit \(p\) the action of \(\OO_K\) on the formal group of \(C\) extends to its \(p\)-adic completion \(\OO_{K_v}\). To see continuity, multiplication by \(p\) has linear term \(pT\) and all other terms have order at least two. Iterating it sends \(T\) into \((p,T)^{n+1}\) after \(n\) steps. Endomorphisms whose scalars differ by a multiple of \(p^n\) therefore agree modulo \((p,T)^{n+1}\) in formal coordinates, by the formal addition law. Completeness gives the extension. The tangent action is the corresponding embedding of \(K_v\). Over \(\OO_{K_v}\) a basis adapted to the projective summand conjugates the idempotent to a diagonal projection. This basis change acts integrally on the power of the formal group. Consequently the formal group of \(\mathcal A\) at every place above \(p\) is a power of the formal group in Theorem~\ref{thm:cm37}.

\begin{corollary}[An ideal dividing every integral CM jet]\label{cor:ideal37}
Use an integral invariant tangent vector on \(\mathcal A\) and a scalar jet of order \(a\) at its identity defined integrally over \(\OO_J\). Let \(z\in J\) be its \(a\)-th derivative at the identity. For each finite place \(w\mid p\) write \(e(w/p)\) for its ramification index. The integral ideal
\[
 \mathfrak c(t,a,w_{\bullet})=
 \prod_{p}\prod_{w\mid p}\mathfrak p_w^{
           \left\lceil e(w/p)\,C_p(t,a,w_{\bullet})/(2\log p)\right\rceil}
\]
contains \(z\). If \(z\ne0\), then
\begin{equation}
            C(t,a,w_{\bullet})\le
      \frac{2}{[J:\QQ]}\sum_{\sigma:J\hookrightarrow\CC}\log|\sigma z|.
                                                               \label{eq:globaljet37}
\end{equation}
Products of jets give sums of the corresponding counts whenever the product is nonzero.
\end{corollary}
\begin{proof}
Taylor evaluation to order \(a\) depends only on the jet modulo terms of order \(a+1\), so it may be computed in any local integral formal frame for the scalar jet. The two preceding theorems give \(v_w(z)\ge e(w/p)C_p/(2\log p)\). At a split prime integrality gives the same inequality. The integral valuations give the ceilings, and only finitely many primes contribute. For a nonzero algebraic integer, the sum of its finite valuations times log residue norms equals the sum of its log absolute values over all embeddings. For each \(p\), \(\sum_{w\mid p} e(w/p)[\kappa(w):\mathbf F_p]=[J:\QQ]\). Summing therefore gives \eqref{eq:globaljet37}. Valuations and log absolute values add under products.
\end{proof}

For example the canonical theta section of \(\mathcal O_C([0])\) vanishes simply at the identity. In an integral local frame its \(a\)-th tensor power has leading term \(uT^a\) with \(u\) a unit. The corresponding derivative against a local tangent generator is \(a!u\). This gives an actual nonzero theta jet to which the bound applies. The scalar periods of Section~\ref{sec:diagonal32} integrate a chain of Weil kernels and a Bessel test. Using \eqref{eq:globaljet37} on those periods requires an integral identification with the jets, including the factors introduced by projections and normalization. It also requires an archimedean bound after the main term of the count is subtracted.

\lead{Locations in the accompanying source} The note's second dated entry on page 3 concerns the odd prime-power series at an inert prime. Related formal coefficients occur in \srcloc{Development/AbelianVarietyTheory/Proofs/DecoratedLatticeProofs/HeightTwoHondaTypeFirstCoeffNormEqOneStatement}{6}. Iteration of an invariant derivation occurs in \srcloc{Development/AbelianVarietyTheory/Proofs/DivisorProofs/ConstantCoeffIterateInvariantDerivationCongr}{52}.

\subsection{The exact ideal of scalar CM derivatives}\label{nav:V-04}
\label{sec:exact41}

Let \(K=\QQ(\sqrt{-q})\) with the primes \(q\equiv3\pmod8\) of the common intervals. Its discriminant is \(-q\). For \(Q=p^f\) and an integer \(n\ge0\) put
\[
       F_Q(n)=\sum_{j\ge1}\lfloor n/Q^j\rfloor
             =\frac{n-s_Q(n)}{Q-1},
\]
where \(s_Q(n)\) is the sum of the base-\(Q\) digits. Define an element of \(K\) by
\begin{equation}
 H_K(a)=\frac{a!}{
  \displaystyle\prod_{p\ {\rm split}}p^{F_p(a)}
  \prod_{p\ {\rm inert}}p^{F_{p^2}(a)}
  (\sqrt{-q})^{F_q(a)}} .                                    \label{eq:generator41}
\end{equation}
Only primes at most \(a\) occur. The two choices of square root give the same principal ideal. We shall compute the image of a derivative map over a common coefficient field. Its domain consists of all integral scalar jets along the identity of the CM abelian scheme of Section~\ref{sec:cm37}.

\begin{lemma}[A factorial bound for the formal exponential]\label{lem:factorial41}
Let \(v_p(p)=1\), let \(e\) be the ramification index of a local scalar field over \(\QQ_p\), let \(\pi\) be its uniformizer and let \(Q=p^f\) be its residue cardinality. In a formal coordinate over the integers of a finite coefficient extension suppose that
\[
      [\pi](T)=\pi T+\sum_{i\ge Q}b_iT^i .
\]
For the formal exponential \(E(Z)=\sum_{n\ge1}c_n Z^n\) with \(c_1=1\), one has \(v_p(c_n)\ge-F_Q(n)/e\). For any integral scalar jet on a power of this group and any integral invariant tangent vector \(D\),
\begin{equation}
             v_p((D^a f)(0))\ \ge\ v_p(a!)-F_Q(a)/e .
                                                                   \label{eq:sharp41}
\end{equation}
\end{lemma}
\begin{proof}
Normalize valuations temporarily by \(v(\pi)=1\). Use \([\pi](E(Z))=E(\pi Z)\) as in Section~\ref{sec:cm37}. In the coefficient of \(E^i\), group the selections by the multiplicities \(m_j\) of indices \(j\), so \(\sum m_j=i\), \(\sum j m_j=n\). The denominator exponent before division by \(\pi^n-\pi\) is at most
\[
       \sum_j m_jF_Q(j)-
                 e v_p\!\left(i!\big/\textstyle\prod_j m_j!\right)
\]
by induction. Always \(F_Q(n)\ge\sum m_jF_Q(j)\) by superadditivity of the floors. We claim that this inequality is strict whenever \(i\ge Q\) and the multinomial is a \(p\)-adic unit. Indeed the unit condition, by the digit formula for \(v_p(i!)\), means that adding the \(m_j\) in base \(p\) produces no carries. If every \(m_j<Q=p^f\), the sum with no carries has no digit in position \(f\) or higher. This contradicts \(i\ge Q\). Thus some index \(j\ge1\) occurs at least \(Q\) times. Adding \(Q\) copies of that positive index in base \(Q\) produces a carry, and adding the remaining positive indices cannot undo the loss in digit sum. The digit formula for \(F_Q\) proves the strict inequality. If the multinomial is a nonunit it supplies at least one unit of the normalized valuation instead. In both cases there is room for division by \(\pi^n-\pi\), of valuation one for \(n>1\). Every index on the right is smaller than \(n\). This proves the coefficient bound by induction.

On a product use the parameterization \((E(u_1Z),\ldots,E(u_gZ))\) for the integral coordinates of \(D\). Extract the coefficient of \(Z^a\), multiply by \(a!\), and apply superadditivity of \(F_Q\) to each product of coefficients. Returning to \(v_p(p)=1\) proves \eqref{eq:sharp41}.
\end{proof}

\begin{lemma}[The first term after rescaling]\label{lem:unit41}
At a place above \(p\) the formal group of \(\mathcal A=\OO_{E_x}\otimes_{\OO_K} C\) is a power of the elliptic formal group of \(C\). Use the scalar completion of \(K\) determined by its tangent action, with index \(e\) and residue cardinality \(Q\). There is an integral coordinate as in Lemma~\ref{lem:factorial41} with \(b_Q\) a unit. In an extension choose \(\lambda^{Q-1}=\pi\), and put \(H(Z)=\lambda^{-1}E(\lambda Z)\). This series is integral and its reduction has the form
\[
                    \overline H(Z)=\sum_{j\ge0} h_j Z^{Q^j},
                \qquad h_0=1,\quad h_j\ne0\quad(j\ge0).
\]
\end{lemma}
\begin{proof}
The product description at nonsplit places was proved in Section~\ref{sec:cm37}. At a split prime the continuous action extends to \(\OO_K\otimes\ZZ_p=\ZZ_p\times\ZZ_p\) by the same continuity argument. Over the fraction field the logarithm identifies each idempotent with its linear action on the tangent line. One acts by one and the other by zero, and the identities persist integrally since the coefficient ring is torsion-free. The projective module \(\OO_{E_x}\) has rank \([E_x:K]\) at both factors and is free over the product of these local rings. A change of basis in the idempotent construction of Section~\ref{sec:cm37} gives the claimed power also at split primes. The coordinate averaged over the Teichm\"uller units in the scalar completion gives the displayed form of \([\pi]\) at every prime.

It remains to verify its first nonlinear unit coefficient. An endomorphism of the good elliptic curve extends to the abelian scheme; the degree of a nonzero CM endomorphism with scalar \(\alpha\) on every fiber is \(|N_{K/\QQ}\alpha|\). For example the degree is the degree of the pullback of the origin divisor. This line bundle has constant fiber degree in a smooth proper family, equal in characteristic zero to the lattice index. At an inert prime take \(\pi=p\), of degree \(p^2=Q\). At the ramified prime take \(\pi=\sqrt{-q}\), of degree \(q=Q\). In both cases the reduction has zero linear coefficient. Its local degree at the origin is at most the total degree \(Q\); it is the order of the pulled-back local parameter on the special fiber. Since degrees below \(Q\) have no nonzero reduced coefficient, \(b_Q\) is a unit.

At a split prime take \(\pi=p\). For a positive class number multiple \(d\), choose \(\alpha\) generating the \(d\)-th power of the prime ideal of \(K\) opposite to the tangent place. It is a unit in the tangent completion, and \(\alpha\bar\alpha=p^d\). The reduction of \([\alpha]\) is separable of degree \(p^d\) since its differential is nonzero. Thus the separable degree of \([p^d]\) is at least \(p^d\). The separable degree of \([p]\) is at most \(p\), since its differential is zero and its degree is \(p^2\). By multiplicativity, the separable degree is exactly \(p\). All points of the geometric kernel have the same local multiplicity by translation, which therefore equals \(p\). The local expansion of \([p]\) has its first nonzero reduced term in degree \(Q=p\), as required.

The bound on \(c_n\) gives \(v_p(c_n\lambda^{n-1})\ge0\), since \(F_Q(n)\le(n-1)/(Q-1)\) for \(n\ge1\). Substitute \(T=\lambda H(Z)\) in the exponential identity and divide by \(\pi\lambda\). Modulo the maximal ideal the result is
\[
                Z=\overline H+\bar b_Q\overline H^{\,Q}.
\]
On the left \(H(\pi Z)/\pi\) reduces to \(Z\); on the right every term of degree index \(i>Q\) has a factor \(\lambda^{i-Q}\). This equation determines a formal solution uniquely coefficient by coefficient. In characteristic \(p\) the displayed series is its solution with \(h_j=-\bar b_Q h_{j-1}^Q\) for \(j\ge1\), so every \(h_j\) is nonzero.
\end{proof}

\begin{theorem}[A global scalar jet with the minimum at every prime]\label{thm:exact41}
Fix a sufficiently distant common interval. There is a single finite extension \(J\) for that interval such that for each of its fields \(E=E_x\) one can fix one integral tangent vector \(D_E\) on \(\mathcal A_E/\OO_J\) with the following property. For every order \(a=\lfloor x\exp(x^\eps)\rfloor\) used on that field, the image of
\[
 \{\text{integral scalar jets to order }a\text{ at the identity}\}
    \longrightarrow \OO_J,\qquad f\longmapsto (D_E^a f)(0),
\]
is exactly \(H_K(a)\OO_J\). Every integral invariant tangent vector has image contained in this ideal. In particular some single integral jet has derivative \(H_K(a)\), simultaneously at every finite and infinite place.
\end{theorem}
\begin{proof}
Section~\ref{sec:cm37} supplies an initial common field and abelian schemes. The closed infinitesimal neighborhood of order \(a\) of the identity has functions forming a finite locally free \(\OO_J\)-module. Indeed its filtration by powers of the ideal of the identity has graded pieces the symmetric powers of the relative cotangent space, by smoothness. The derivative kills terms of degree above \(a\) upon evaluation at zero. It gives an integral linear map on this module and commutes with localization and extension of scalars. Locally the module is the truncated formal power series ring. The valuation of \eqref{eq:generator41} over \(p\) is precisely the right-hand side of \eqref{eq:sharp41}, using \(v_q(\sqrt{-q})=1/2\). These valuations are nonnegative. This proves the universal inclusion and integrality of the generator.

We prove that the bound can be attained locally with a choice of tangent vector. Write \(a=\sum a_jQ^j\) with \(0\le a_j<Q\), and set \(k=\sum a_j\). Use one distinct factor of the formal group for each nonzero digit. There are enough factors since the rank \(g=[E_x:K]\) is at least the degree \(2^l\ge x^{2+8\eps}\) of its totally real cyclotomic layer, whereas \(1+\log_2 a=O(x^\eps+\log x)\). Use the monomial \(\prod_j T_j^{a_j}\) in these coordinates and tangent components \(u_j\). The coefficient of \(Z^a\) after pullback is \(\lambda^{k-a}\) times the coefficient of \(Z^a\) in \(\prod_j H(u_jZ)^{a_j}\). The latter reduces to a homogeneous polynomial of degree \(a\) in the components. It is nonzero. Indeed the coefficient of \(\prod_j u_j^{a_j Q^j}\) is \(\prod_j h_j^{a_j}\): the only way to express \(a_jQ^j\) as a sum of precisely \(a_j\) powers of \(Q\) is to use \(a_j\) copies of \(Q^j\). Any carry would lower the digit sum below \(a_j\). At a point where this polynomial is nonzero, the derivative valuation is therefore
\[
          v_p(a!)+(k-a)v_p(\lambda)
                      =v_p(a!)-F_Q(a)/e .
\]
The reduction calculation uses a root \(\lambda\) to detect units; the coefficients \(h_j\) lie in the residue field of the original local model by their recurrence. The jet and its tangent vector can be taken in the original formal coordinates after a finite residue extension.

We next make the same tangent vector work globally and at all the orders for each field. Each bounded interval uses finitely many fields and integer orders. Let \(A\) be the largest order and fix one of the fields \(E\). Its integral tangent module \(M\) is projective of rank \(g\ge2\) over the ring of integers of the initial coefficient field. The structure theorem for projective modules over a Dedekind domain writes it as a sum of \(g-1\) free lines and an invertible ideal. Choose a surjective linear form \(L:M\to\OO_J\) on one free line and a vector \(v_0\) with \(L(v_0)=1\). At each place above \(p\le A\), take the product of the preceding nonzero homogeneous polynomials for the relevant orders. They may be written in the same coordinates on the residue tangent space. Their product times the reduction of \(L\) remains a nonzero polynomial. It has a nonvanishing value over a finite extension of the residue field. Scale the resulting vector by a nonzero scalar so that \(L=1\); homogeneity preserves every nonvanishing condition.

All needed finite residue extensions can be included at once in a finite extension of the coefficient field. For each rational \(p\le A\) one may adjoin the roots of unity of order \(p^b-1\) with \(b\) a common multiple of the finitely many required residue degrees over \(\mathbf F_p\). At every place above \(p\) these roots reduce distinctly and supply \(\mathbf F_{p^b}\). Take the compositum for all the primes and fields. Over the resulting \(J\), choose the residue directions at every place above \(p\le A\); under further extension the local polynomial nonvanishings persist. Subtract \(v_0\) from the chosen directions and apply the Chinese remainder theorem to the projective module \(\ker L\). Adding \(v_0\) back gives a global vector \(D_E\) with \(L(D_E)=1\) meeting all the prescribed reductions.

For each order and each \(p\le A\), the corresponding local monomial now attains equality. Computing in a completion is sufficient, since the image ideal of a linear map of finite modules commutes with flat base change to a completed discrete valuation ring. For \(p>A\), the reduction of \(D_E\) is also nonzero because \(L(D_E)=1\). A local parameter with unit derivative in that direction, raised to the \(a\)-th power, has derivative \(a!\) times a unit. It is a unit there since \(p>a\). The two ideals thus have identical valuation at every prime of \(J\). They are equal. The image is the image of a linear map on the global module over the affine base \(\operatorname{Spec}\OO_J\), so its generator itself is attained by a global jet. This proves the assertion for each field over the same \(J\).
\end{proof}

In particular the set of derivatives for the chosen vector, after division by \(H_K(a)\), is the whole coefficient ring, and the value one occurs. This computes the full common divisor of scalar jets. The theta periods in Section~\ref{sec:diagonal32} are specified by a kernel, a projection and an integral. A comparison for those prescribed values must retain restrictions on the jet or on the map to the period.

\subsection{Realizing the derivative by global theta sections}\label{nav:V-05}
\label{sec:theta43}

We now impose that a scalar jet be the expansion of a global algebraic theta section on the CM abelian scheme. We specify its line bundle, its order, its frame at the identity and its behavior under inversion.

\begin{lemma}[Jets from a very ample linear system]\label{lem:very43}
Let \(R\) be a Dedekind domain, let \(X/R\) be smooth and projective with section \(e\), and let \(L\) be an invertible sheaf together with an identification \(e^*L=R\). Suppose a finite projective module \(V\) of sections defines a closed immersion \(X\hookrightarrow\mathbf P_R(V)\) with tautological bundle \(L\), using the convention that a point is a rank-one quotient. There exists \(s_0\in V\) with \(s_0(e)=1\). For every \(a\ge1\), all integral scalar jets of order \(a\) at \(e\) are expansions of \(s/s_0^a\) for sections \(s\in H^0(X,L^{\otimes a})\). One may use the sections generated by degree-\(a\) products from \(V\).
\end{lemma}
\begin{proof}
Evaluation at \(e\) is a surjection \(V\to R\), so choose \(s_0\) above one and put \(V_0=\ker(V\to R)\). This splits \(V=Rs_0\oplus V_0\). The open set where \(s_0\) is invertible contains the whole section \(e\). In the corresponding affine chart of the projective bundle, its coordinate algebra is \(\operatorname{Sym}_R V_0\); the linear coordinate belonging to \(s\in V_0\) restricts on \(X\) to \(s/s_0\). The ideal of the origin in this chart is the positive-degree ideal. Since \(X\) is closed in the chart and contains the origin section, these linear coordinates generate the ideal of \(e\) in the induced affine open of \(X\). Its order-\(a\) neighborhood is contained in the chart because it has the same underlying space as \(e\). Every element of its coordinate algebra is consequently represented by a polynomial of degree at most \(a\): all monomials of strictly larger degree vanish modulo the \((a+1)\)-st power of the ideal. A monomial of degree \(j\le a\) homogenizes to the product of its \(j\) sections and \(a-j\) copies of \(s_0\). Adding these products gives the asserted global section over \(R\).
\end{proof}

\begin{theorem}[Theta sections attaining the minimum]\label{thm:globaltheta43}
On the common fields and intervals of Theorem~\ref{thm:exact41}, give each CM abelian scheme \(\mathcal A_E\) any relatively ample symmetric rigidified theta line bundle \(P\) over the coefficient ring. There exists a frame \(s_0\), supplied by a global section of \(P^3\) which equals one at the identity, such that for the vector \(D_E\) from that theorem and every required order \(a\),
\begin{equation}
 \left\{D_E^a(s/s_0^a)(0):s\in H^0(\mathcal A_E,P^{3a})\right\}
                      =H_K(a)\OO_J .                          \label{eq:theta43}
\end{equation}
A product principal polarization \(P\) exists over the common field used there. With that choice the positive theta bundle \(P^{3a}\) has type \((3a,\ldots,3a)\) at every complex embedding. In particular the smallest nonzero normalized log norm is the quantity computed in Theorem \ref{thm:height42}. It is attained by one global theta section in the indicated frame.
\end{theorem}
\begin{proof}
On an abelian variety over a field the cube of an ample line bundle is very ample, by the classical very ampleness theorem for abelian varieties. This theorem holds also in positive characteristic. We recall why these fiberwise embeddings give a relative linear system over the entire coefficient ring. For a symmetric ample line bundle \(L\) on a geometric fiber and an integer \(N\) prime to its characteristic, multiplication by \(N\) is finite \'etale of degree \(N^{2g}\). Pullback on \(H^i(L)\) is split injective by the trace divided by this degree. By the theorem of the cube \([N]^*L\simeq L^{N^2}\). Serre vanishing for large such \(N\) gives \(H^i(L)=0\) for \(i>0\). Apply this to \(L=P^3\) on every fiber. The Euler characteristic is locally constant in this proper flat family, so \(h^0(L)\) is constant. Cohomology and base change for the projective abelian scheme gives \(V=H^0(\mathcal A_E,P^3)\) finite projective with its sections spanning the complete system on every fiber. Evaluation is surjective by Nakayama's lemma. The morphism to the projective bundle is a closed immersion: it is proper with finite geometric fibers, hence finite, and the cokernel of the finite algebra map from the structure sheaf of the projective bundle vanishes on every fiber by the fieldwise closed immersion. Nakayama gives surjectivity. These classical results on linear systems of abelian varieties are given in \cite{XIII-X-IV-av43}.

Apply Lemma~\ref{lem:very43} to \(L=P^3\), rigidified using \(P\). The full domain of the derivative map in Theorem~\ref{thm:exact41} is in the image of the global theta sections. The integral inclusion and attainment of the generator from that theorem prove \eqref{eq:theta43} for the same vector at every order. A nonzero integral multiple of the generator has norm at least that of the generator, proving the norm assertion.

For the asserted principal polarization we can construct the bundle directly. As a projective module over the Dedekind domain \(\OO_K\), \(\OO_E\) is isomorphic to \(\OO_K^{g-1}\oplus I\) for an invertible ideal \(I\). The construction by idempotents in Section~\ref{sec:cm37} respects direct sums and isomorphisms. It identifies \(\mathcal A_E\) with \(C^{g-1}\times C_I\) where \(C_I=I\otimes_{\OO_K} C\) is an abelian scheme of relative dimension one. Indeed it is a direct factor of a power of \(C\), and its generic tangent space has dimension one. The factors are thus elliptic schemes. On each use the line bundle of its origin, tensoring by a bundle from the base to rigidify it. It is symmetric and induces the principal polarization on each fiber. Their exterior product gives \(P\); its \(3a\)-th power has the stated type. This product choice uses an \(\OO_K\)-module decomposition of \(\OO_E\). The equality \eqref{eq:theta43} also applies to a different integral polarization prescribed in advance whenever it satisfies the hypotheses.
\end{proof}

There is a direct analytic description of the last construction. At any complex embedding write each elliptic factor as \(\CC/(\ZZ+\tau_j\ZZ)\) with \(\operatorname{Im}\tau_j>0\). The odd Jacobi theta function \(\vartheta_1(\pi z_j,\tau_j)\) vanishes simply at the lattice and has no other zeros. Its multipliers for \(z_j\mapsto z_j+1\) and \(z_j\mapsto z_j+\tau_j\) are \(-1\) and \(-\exp(-\pi i\tau_j-2\pi iz_j)\), respectively. These identities identify the associated bundle with the bundle of the origin. Therefore a section of \(P^{3a}\) lifts to a holomorphic function on \(\CC^g\) with these multipliers raised to the power \(3a\) in coordinate \(j\). This is the usual theta transformation law for the product polarization. The ratios in \eqref{eq:theta43} are computed in the common frame fixed once over \(J\). Changing that frame can change a higher derivative when lower derivatives are present. \cite{XIII-X-IV-nist43}

\begin{corollary}[Theta sections of prescribed parity]\label{cor:parity43}
Write \(\iota=[-1]\) and let \(s_*=s_0\,\iota^*s_0\) in \(P^6\), using the symmetry fixing the rigidification at zero. Use the frame \(s_*^a\) for \(P^{6a}\). For global theta sections of parity \((-1)^a\), the image ideal of the order-\(a\) derivative is between \(2H_K(a)\OO_J\) and \(H_K(a)\OO_J\), inclusive. Consequently it is the full scalar derivative ideal at every place of odd residue characteristic. There is one global section of that parity with derivative exactly \(2H_K(a)\).
\end{corollary}
\begin{proof}
Rigidified symmetry is an involution: its square is an automorphism of the line bundle, hence a scalar from the base since the fibers are proper and geometrically connected, and that scalar is one at the identity. Thus \(s_*\) is invariant and has value one. Given the section \(s\) with derivative \(H_K(a)\) from Theorem~\ref{thm:globaltheta43}, put \(t=s(\iota^*s_0)^a\). Then \(t/s_*^a=s/s_0^a\) at the identity. The integral section \(t+(-1)^a\iota^*t\) has the stated parity. Inversion acts by \(-1\) on the tangent space and sends the formal one-parameter subgroup for \(D_E\) to that for \(-D_E\). It therefore multiplies the \(a\)-th derivative at zero by \((-1)^a\) in the invariant frame. The derivative of the sum is \(2H_K(a)\). The universal divisibility proved in Theorem~\ref{thm:exact41} gives the other inclusion.
\end{proof}

The theta sections constructed here have an integral line bundle and a prescribed derivative at every place in one frame. Their coefficients and the frame can depend on the field. The Bessel periods of Section~\ref{sec:diagonal32} are obtained from specified Weil kernels, choices of splittings, common test vectors and a projection. Comparing those periods to sections subject to these restrictions requires further input; the minimum over all global theta sections alone is given by \eqref{eq:theta43}.

\subsection{The trace polarization and integral Fourier components}\label{nav:V-07}
\label{sec:pol49}

Fix a common interval \(I=[V,V^{3/2}]\), and take the fields \(E=E_x\)
of Volume~I. Let \(K=\QQ(\sqrt{-d})\) with \(\sqrt{-d}\mapsto i\sqrt d\).
The source uses the inverse of one factor of the relative different
and a purely imaginary scalar to define its CM lattice. We construct
this polarization on the same fields. We then compute the denominators
of translations in a linear system pulled back by a split CM isogeny.
The ideal and scalar occur in
\srcloc{Formalism/LocalRingCardinality/DataHost}{149}.

\begin{theorem}[A principally polarized CM lattice over the interval]\label{thm:pol49}
Let \(\mathfrak c\) be a nonzero integral ideal of \(K\) and
\(u=N_{K/\QQ}\mathfrak c\), the absolute ideal norm. On each \(E\)
one can choose an integral ideal \(\mathfrak b\), stable under
\(\Gal(E/K)\), with
\(\mathfrak b\overline{\mathfrak b}=\mathfrak D_{E/K}\) and
\(\mathfrak b+\overline{\mathfrak b}=\OO_E\). Set
\[
 M=\mathfrak c\OO_E\,\mathfrak b^{-1},\qquad
 \xi=-1/(u\sqrt{-d}),\qquad
 e(x,y)=\operatorname{Tr}_{E/\QQ}(\xi x\bar y).
\]
The form \(e\) gives a principal polarization on the complex torus
of \(M\) with CM embeddings extending the chosen embedding of \(K\).
Its Rosati involution restricts to conjugation on \(E\).
For one choice of \(\mathfrak c\) per field all the resulting tori
and polarizations have abelian scheme models and relatively ample
symmetric rigidified line bundles over the integers of one number
field \(J\) on \(I\).
\end{theorem}
\begin{proof}
Outside primes dividing \(2Fqd\), the field \(E\) is unramified over
\(\QQ\), where \(F\) is the conductor of \(\psi\). At \(d\), the
extension obtained by adjoining \(K_\psi,C_l,\sqrt{-q}\) to \(\QQ\)
is unramified. Base change of its local unramified algebras to \(K_d\)
(here the completion at \(d\)) is unramified. Every remaining rational
prime in the support of the relative different splits in \(K\).
Indeed \(-d\equiv1\bmod 8\), and \(-d\equiv1\bmod p\) for \(p=q\)
and for every odd \(p\mid F\) on sufficiently distant intervals.
For each such split \(p\), choose one of its two primes in \(K\) and
put into \(\mathfrak b\) the full part of \(\mathfrak D_{E/K}\) over
that prime. Conjugation exchanges the two parts and \(\Gal(E/K)\)
preserves each. This proves all the assertions about \(\mathfrak b\).

Transitivity of the different and the quadratic discriminant give
\(\mathfrak D_{E/\QQ}=\mathfrak D_{E/K}(\sqrt{-d})\).
Also \(\mathfrak c\overline{\mathfrak c}=(u)\), hence
\(\xi M\overline M=\mathfrak D_{E/\QQ}^{-1}\). By the defining
trace-duality property of the inverse different,
\(\{x\in E:e(x,M)\subset\ZZ\}=M\). The trace form is alternating
since \(\bar\xi=-\xi\). In the coordinates of the CM type the real
extension of \(e\) satisfies
\[
 e(v,iv)=\frac{2}{u\sqrt d}\sum_{\tau|K}|v_\tau|^2>0\quad(v\ne0),
 \qquad e(iv,iw)=e(v,w).
\]
The Riemann bilinear relations consequently give a polarization, which
is principal by trace duality. The identity \(e(ax,y)=e(x,\bar a y)\)
gives the asserted involution.~\cite{XIII-X-IV-cm37}

Choose the elliptic scheme \(C_d\) after finite extension as in
Section~\ref{sec:cm37}, with the complex uniformization
\(\CC/\OO_K\). The finite projective \(\OO_K\)-module \(M\) gives
the abelian scheme \(\mathcal A=M\otimes_{\OO_K}C_d\). Its complex
lattice and tangent action are exactly those above. The complex
polarizing homomorphism is defined over a number field. For example,
after isogenies the source and target are powers of \(C_d\), and their
rational homomorphisms are matrices in \(K\); homomorphisms descend
under algebraically closed extension in characteristic zero. A line
bundle representing the polarization can be chosen over
\(\overline{\QQ}\) in the corresponding component of the Picard
scheme. It can be made symmetric by twisting with a degree-zero
bundle, since multiplication by two on the dual is surjective.
These finite data descend after a finite extension.

Extend a divisor of this bundle by closure on the regular abelian
scheme over the integer ring, and rigidify the resulting bundle \(P\)
along the identity. A line bundle trivial on the generic fiber is
represented by a vertical divisor. All fibers are smooth and
geometrically integral, so such divisors are pulled back from the
base. Rigidification thus gives the symmetry on the integral model.
To see ampleness on every fiber, diagonalize the positive definite
\(K\)-Hermitian form \(u^{-1}\operatorname{Tr}_{E/K}(x\bar y)\)
by orthogonalization. Its diagonal entries are positive rational
numbers. After clearing denominators in the coordinate maps, for
some positive integer \(D\) there are independent maps
\(\ell_i:M\to\OO_K\) and positive integers \(c_i\) with
\[
 D u^{-1}\operatorname{Tr}_{E/K}(x\bar y)
             =\sum_i c_i\ell_i(x)\overline{\ell_i(y)} .
\]
The map \(\ell=(\ell_i)\) induces an isogeny to a power of \(C_d\)
on every fiber: a positive integer multiple of its inverse on the
rational lattices maps the integral lattices back. Thus its kernel
is contained in the torsion of that integer and its image has full
dimension. Pull back the product of the elliptic origin bundles with
exponents \(c_i\). Its polarization homomorphism agrees with that of
\(P^D\) on the generic fiber by the displayed equality and the
elliptic trace form \(-1/\sqrt{-d}\). Equality extends over the base
by separatedness. Fiberwise their quotient is algebraically trivial,
so \(P^D\) is ample on every fiber. The fiberwise criterion for relative
ampleness now applies. Taking a compositum works for all the finitely
many fields on \(I\).
\end{proof}

\begin{lemma}[Two principal split primes of controlled size]\label{lem:aux49}
For \(T\ge\max(d,3)\) there exist primes \(T<g_1<g_2\)
and \(\pi_i\in\OO_K\) with
\(\pi_i\bar\pi_i=g_i,\ (\pi_i)+(\bar\pi_i)=\OO_K\), and
\(\log(g_1g_2)\ll\log(2T)\) with an absolute implied constant.
On \(I\) both primes can be fixed above every integer order
\(\lceil x e^{x^\eps}\rceil\), all rational primes ramified in any
\(E_x\), and all the dimensions \(2(2n)(r+t_{\rm r})+1\) from
Section~\ref{sec:ranks47}. Then
\(\log(g_1g_2)\ll_{\psi,\eps}x^{3\eps/2}+\log(2x)\) for every \(x\in I\).
\end{lemma}
\begin{proof}
Let \(\omega=(1+\sqrt{-d})/2\), and take
\(c=2\lceil\sqrt{T/d}\rceil+2\). The binary integral form
\[
 N_{K/\QQ}(u+cv\omega)=(u+cv/2)^2+dc^2v^2/4
\]
is primitive, positive definite, and has discriminant \(-dc^2\).
Thorner and Zaman's unconditional theorem gives a prime represented
by it of size at most an absolute constant times \((dc^2)^{694}\).
In any such representation \(v\ne0\) since a square cannot be prime.
Hence \(g\ge dc^2/4>T\), while \(dc^2\ll T\).
The representing element has norm \(g\); as \(g>d\) it is unramified
and its two conjugate prime ideals are distinct. Apply this argument
first at \(T\) and then at \(g_1\).~\cite[Thm.~1.2]{XIII-X-IV-tz49}

For the interval take \(T\) larger by one than the largest bound
specified in the statement. The ranks are polynomial in \(V\), the
largest order has logarithm \(O(V^{3\eps/2}+\log(2V))\), and the
quadratic prime selection gives
\(\log(qd)\ll_{\psi,\eps} V^{\eps/5}(\log(2V))^2\).
The remaining ramified primes divide \(2F\). Since \(x\ge V\), this
gives the required bound at every point on the interval.
\end{proof}

\begin{theorem}[Integral Fourier components of split CM translations]\label{thm:fourier49}
For one of these primes write \(g=\pi\bar\pi\). Take a product
\(\mathcal V\) of copies of \(\mathcal A\) with the indicated
\(\OO_K\)-action on some factors and the conjugate action on the others.
Let \(G_+,G_-\) be the sums of their dimensions and \(G=G_++G_-\).
Write \(f=[\pi]\), \(h=[\bar\pi]\), \(B=\ker f\). After a common
finite extension all generic \(g\)-torsion points and roots of unity
of order \(g\) are defined. For any relatively ample integral bundle
\(P_0\) on \(\mathcal V\), consider global sections of \(f^*P_0^N\).
Translations by \(B\) act using \(f\circ t_z=f\).

At a place \(w\mid g\) normalize \(v_w(g)=1\), and let \(G_{\rm et}\)
be \(G_+\) if \(\pi\) is a unit at \(w\), and \(G_-\) otherwise.
For every character \(\lambda\) of \(B(J)\), the Fourier projector
\[
       e_\lambda= g^{-G}\sum_{z\in B(J)}
                         \lambda(z)^{-1}t_z^*
\]
preserves the integral lattice of sections after multiplication by
\(g^{G_{\rm et}}\). The same bound holds for
\(\sum_\lambda u_\lambda e_\lambda\) for any \(w\)-integral
\(u_\lambda\). On the full lattice of sections, the denominator for
each single projector is sharp for all sufficiently large \(N\).
Away from \(g\) all these single projectors preserve the integral lattice.
\end{theorem}
\begin{proof}
On a single factor the kernel of \([\pi]\) has generic order
\(g^{\dim\mathcal A}\) by its index on \(M\). It is killed by \(g\);
both \([\pi]\) and its conjugate are isogenies of abelian schemes,
since their composite is \([g]\) and they have finite kernels and
full image on every fiber. When the tangent scalar is a unit the
isogeny is \'etale at the identity and hence everywhere by translation.
At the other places over \(g\), its dual is \'etale because the Rosati
involution from Theorem~\ref{thm:pol49} exchanges \(\pi,\bar\pi\).
The kernel of the dual isogeny is the Cartier dual of the original
kernel. The same argument applies with the two roles exchanged on a
factor with conjugate action.~\cite{XIII-X-IV-dual49}

Over the completed integer ring \(R_w\), the finite \'etale kernel
factors and the finite \'etale dual factors are constant since all
their generic geometric points are rational. Indeed each point extends
by properness and the closures are disjoint sections of a finite
\'etale scheme over a normal local domain with split generic algebra.
Thus \(B=B_{\rm et}\times B_{\rm mult}\), where
\(B_{\rm et}\) is constant of order \(g^{G_{\rm et}}\), and
\(B_{\rm mult}\) is the diagonalizable group of the character group
\(\Gamma\simeq(\ZZ/g)^{G-G_{\rm et}}\).

The identity \(f\circ t=f\) on \(B\times\mathcal V\) linearizes
\(f^*P_0^N\) integrally; the group law gives the compatibility of
successive translations. For the diagonalizable factor, write a
section's coaction as \(\sum_{\gamma\in\Gamma}s_\gamma\otimes X^\gamma\)
in the group algebra \(R_w[\Gamma]\). Coassociativity and the counit
give respectively that \(s_\gamma\) has weight \(\gamma\) and that
\(\sum s_\gamma=s\). Each \(s_\gamma\) is integral since \(X^\gamma\)
is an integral basis. On the generic fiber this decomposition equals
ordinary finite Fourier projection by character orthogonality.
The remaining projection averages over the constant group of order
\(g^{G_{\rm et}}\) with coefficients roots of unity. This proves the
claimed exponent, the bound for linear combinations, and integrality
away from \(g\).

For sharpness restrict to the finite flat fiber \(B=f^{-1}(0)\),
where the pullback bundle has one frame from the origin. Its functions
are \(\operatorname{Map}(B_{\rm et}(J),R_w)\otimes R_w[\Gamma]\).
For the given \(\lambda\) take the point mass at the identity of the
first factor times the corresponding weight character of the second.
Its projected value at the identity is \(1/g^{G_{\rm et}}\).
For \(N\) sufficiently large this finite-fiber section extends to a
global section over \(R_w\). In fact \(f^*P_0\) is relatively ample,
and relative Serre vanishing for the ideal sheaf of \(B\) tensored
with \((f^*P_0)^N\) gives the surjectivity of restriction over the
affine base. It follows that the denominator cannot be decreased on
that lattice. On the bounded interval all fields and \(g_i\) used
here form finite lists, so a common finite coefficient extension
suffices.
\end{proof}

\begin{corollary}[Pullback of the projected integral jet]\label{cor:jet49}
Use a frame of \(P_0^N\) near the origin and the frames it induces under
\(f\) and \(fh\). Let \(D\) be a product of integral invariant
derivations in the two groups of factors, of orders \(a_+,a_-\).
Write \(J_D\) for evaluation of the resulting scalar derivative at
zero and \(e'_\lambda\) for the \(B\)-projector on \((fh)^*P_0^N\).
For integral \(s\in H^0(\mathcal V,f^*P_0^N)\) one has
\[
 J_D(e'_\lambda h^*s)=
   \bar\pi^{\,a_+}\pi^{a_-} J_D(e_{\lambda\circ h^{-1}}s).
\]
Here \(h^{-1}\) in the character denotes the inverse on \(B(J)\).
Consequently at \(w\mid g\),
\[
 v_w\bigl(J_D(e'_\lambda h^*s)\bigr)\ge
       a_+v_w(\bar\pi)+a_-v_w(\pi)-G_{\rm et}.
\]
For \(r(r+t_{\rm r})\) copies and total dimension
\(G=2n\,r(r+t_{\rm r})\), the logarithm of the full worst-case
denominator at the two primes is bounded by
\(G\log(g_1g_2)=o(b/x)\), uniformly at every \(x\) on \(I\).
\end{corollary}
\begin{proof}
Since \((\pi,\bar\pi)=1\), a B\'ezout identity in \(\OO_K\) gives an
inverse to \(h\) on \(B\). For every \(z\in B(J)\) one has
\(t_z^*h^*=h^*t_{h(z)}^*\) with the canonical pullback actions
and frames. Reindexing the Fourier average gives
\(e'_\lambda h^*=h^*e_{\lambda\circ h^{-1}}\).
Differentiating \(h\circ t_v=t_{h(v)}\circ h\) at the identity
gives the chain rule for an invariant direction everywhere.
Its scalar is \(\bar\pi\) on the first factors and \(\pi\) on the
conjugate factors. Iterating yields the formula with no lower-order
term because the frame is pulled back as well. Integral invariant
derivations preserve integral functions, so Theorem~\ref{thm:fourier49}
gives the valuation bound. Finally \(G=O(n^3)\) by the rank choice,
and Lemma~\ref{lem:aux49} bounds the logarithms by a fixed power of
\(x\). This proves the asymptotic claim.
\end{proof}

The dimension in the corollary gives exactly the group order
\(g^{f r n_V}\) specified in
\srcloc{Formalism/SharpRestriction/Divisibility/CanonicalTranslateJets}{84},
where \(f=2n\) and \(n_V=r+t_{\rm r}\). Lines 176--188 and 198--210
of that file specify further identities between translated sections,
automorphic evaluations and jets in matched frames. The theorem gives
the integral Fourier and isogeny calculation for the explicit polarized
CM lattice. Identifying an individual restricted automorphic lift with
one of these sections, with the same integral frame and a bound for
its projected global norm, remains open in this manuscript.

\subsection{Frobenius and mixed derivatives in CM eigendirections}\label{nav:V-08}
\label{sec:mixed52}

The minimum for a single scalar derivative in Section~\ref{sec:exact41}
depends on the quadratic splitting type of the prime. For mixed
derivatives the individual CM embeddings can affect the answer.
We work on the principally polarized abelian scheme of
Theorem~\ref{thm:pol49}, with \(\mathfrak c=\OO_K\),
\(K=\QQ(\sqrt{-d})\), and with the fields of
Proposition \ref{prop:desc50}.
Fix a rational prime \(p\) inert in \(K\) and unramified in \(E\).
Let \(R\) be a completed integer ring of a common coefficient field
containing \(E\), with good reduction as above. Normalize \(v_p(p)=1\).
Write \(Q=p^2\), \(\Sigma=\operatorname{Hom}_K(E,J)\), and
\(\sigma\tau=\tau\circ\Frob_p^2\) for \(\tau\in\Sigma\), with arithmetic
Frobenius in \(\Gal(E/\QQ)\). For nonnegative integers \(a_\tau\) put
\[
  A=\sum_\tau a_\tau,\qquad
  B_p(\boldsymbol a)=\sum_\tau\bigl(v_p(a_\tau!)-F_Q(a_\tau)\bigr).
\]
In \(a_\tau=\sum_{j\ge0}(l_{\tau j}+p h_{\tau j})Q^j\), both
\(l_{\tau j}\) and \(h_{\tau j}\) are between zero and \(p-1\). Define
\[
      L_\nu=\sum_{j\ge0}l_{\sigma^{-j}\nu,j},\qquad
      H_\nu=\sum_{j\ge0}h_{\sigma^{-j}\nu,j}\qquad(\nu\in\Sigma).
\]

\begin{theorem}[An orbit-dependent divisibility bound]\label{thm:gain52}
Let \(D_\tau\) be invariant derivations generating the integral
tangent eigenlines of \(\OO_E\) at \(p\).
For every integral scalar jet \(f\) through order \(A\) at the identity,
\begin{equation}
   v_p\!\left(\left(\prod_\tau D_\tau^{a_\tau} f\right)(0)\right)
        \ \ge\ B_p(\boldsymbol a)+\sum_{\nu\in\Sigma}\lfloor L_\nu/p\rfloor .
                                                              \label{eq:gain52}
\end{equation}
If \(L_\nu<p\) and \(H_\nu<p\) for every \(\nu\), some integral
jet attains exactly \(B_p(\boldsymbol a)\). In the frame of
Theorem~\ref{thm:globaltheta43}, this minimum is attained by a global
section of \(P^{3A}\) at the chosen place, for \(A\ge1\).
\end{theorem}
\begin{proof}
Locally the CM ideal \(M\) is free of rank one over the semilocal
algebra \(\OO_E\otimes_{\OO_K}\OO_{K,p}\), which is finite \'etale
over \(\OO_{K,p}\). Choose its generator and a basis of this algebra
over \(\OO_{K,p}\). The formal group is then a power of the elliptic
formal group used in Section~\ref{sec:exact41}. After base change to
\(R\) the primitive idempotents of the \'etale algebra give tangent
vectors \(v_\tau\), the columns of an invertible matrix \(V\).
These are integral generators of the tangent eigenlines in the product
coordinates of the elliptic group. Their entries lie in an unramified
extension of \(\OO_{K,p}\), since they are the columns of the inverse
of the evaluation matrix of the chosen algebra basis. In particular,

\begin{equation}
                        v_\tau^{[Q^j]}\equiv v_{\sigma^j\tau}\pmod p,
                                                              \label{eq:fr52}
\end{equation}
where brackets raise each coordinate to that power. To see this,
the row at \(\mu\) of the evaluation matrix times \(v_\tau\) is
\(\mathbf1_{\mu=\tau}\). Raising this equality to \(Q\) on the
residue field replaces row \(\mu\) by row \(\sigma\mu\), and the
matrix remains invertible there. The congruence holds modulo \(p\) in the
unramified coefficient ring, hence also in \(pR\) if \(R\) itself is
ramified. Iteration proves \eqref{eq:fr52}. We may prove the result
using these generators, since rescaling any one by a unit only rescales
the final derivative by a unit.

Let \(\mathcal E(Z)=\sum_{s\ge1}c_s Z^s\), \(c_1=1\), be the formal
exponential in the elliptic coordinate of Lemma~\ref{lem:unit41}.
The normalized invariant operators
\[
          T_{\tau,b}=\frac{p^{F_Q(b)}}{b!}D_\tau^b
\]
preserve integral formal functions. At the identity this follows
by extracting the coefficient of degree \(b\) after substituting
\(\mathcal E(v_{i\tau}Z)\) for the \(i\)-th coordinate and using
Lemma~\ref{lem:factorial41}. For the operator statement apply the
same extraction to \(f(X+_{\mathcal A}Y)\) in \(Y\), with the
coefficients now in \(R[[X]]\). Translation invariance identifies
the result with \(T_{\tau,b}f(X)\).

For \(b=Q^j\), every positive splitting \(b=s_1+\cdots+s_k\) with
\(k\ge2\) satisfies \(F_Q(b)>\sum_iF_Q(s_i)\). Indeed the base-\(Q\)
digit sum of \(b\) is one and the sum of the digit sums on the
right is at least two. Their difference is a positive multiple of
\(Q-1\). Thus, after multiplying by \(p^{F_Q(b)}\), each contribution
from a monomial of degree at least two is divisible by \(p\).
The linear term contributes \(u_j v_\tau^{[Q^j]}\) with
\(u_j=p^{F_Q(Q^j)}c_{Q^j}\in R^\times\), a unit by
Lemma~\ref{lem:unit41} and \(F_Q(Q^j)=(Q^j-1)/(Q-1)\).
Using \eqref{eq:fr52} and translating once more gives a congruence
of operators on integral formal functions,
\begin{equation}
                T_{\tau,Q^j}\equiv u_j D_{\sigma^j\tau}\pmod {p}.
                                                               \label{eq:op52}
\end{equation}

Factor \(T_{\tau,a_\tau}\) as a unit times the product of
\(l_{\tau j}\) copies of \(T_{\tau,Q^j}\) and \(h_{\tau j}\)
copies of \(T_{\tau,pQ^j}\) at each \(j\). The exponents \(F_Q\)
add exactly by the digit formula. The remaining factor is a unit
because the multinomial coefficient that combines these base-\(p\)
digits has valuation zero, by Legendre's formula for factorials.
After this factorization, group the first kind of operator according
to \(\nu=\sigma^j\tau\). Any \(p\) in a group have product congruent
modulo \(p\) to a unit times \(D_\nu^p\) by \eqref{eq:op52}.
Since \(p<Q\), the integral operator \(T_{\nu,p}\) equals
\(D_\nu^p/p!\). Thus each such product sends integral functions
into \(pR[[X]]\). These constant multiples of invariant derivatives
commute. After dividing each disjoint group of \(p\) operators by \(p\),
their composition and all unused factors still preserve integral
functions. Restoring the factorials proves \eqref{eq:gain52}.

For attainment, use the integral linear coordinates \(T'_\nu\)
obtained from the elliptic parameters \(T_i\) by \(V^{-1}\).
Let \(d_{\tau j}=l_{\tau j}+p h_{\tau j}\),
\(k_\nu=\sum_j d_{\sigma^{-j}\nu,j}\), and \(k=\sum_\nu k_\nu\).
Test the jet of the monomial \(\prod_\nu(T'_\nu)^{k_\nu}\).
Introduce independent variables \(Z_\tau\) in the tangent space and
adjoin temporarily \(\lambda^{Q-1}=p\). The reduction of
\(\lambda^{-1}\mathcal E(\lambda Z)\) is
\(\sum_{j\ge0}h_j Z^{Q^j}\) with every \(h_j\ne0\), by
Lemma~\ref{lem:unit41}. After inserting
\(\lambda\sum_\tau v_{i\tau} Z_\tau\) in coordinate \(i\) and applying
\(V^{-1}\), the reduction of the resulting \(T'_\nu/\lambda\) is
\[
                \sum_{j\ge0} h_j Z_{\sigma^{-j}\nu}^{Q^j}.
\]
To obtain \(\prod Z_\tau^{a_\tau}\) from the \(k\) linear-coordinate
factors after this substitution, exactly \(d_{\tau j}\) choices of
\(Z_\tau^{Q^j}\) must occur. Indeed the least number of powers of
\(Q\) summing to \(a_\tau\) is its base-\(Q\) digit sum. The sum of
these minima is already \(k\), so no carries in base \(Q\) are possible.
The indicated coefficient is the product of nonzero powers of \(h_j\)
times the multinomials at each \(\nu\) for the groups of sizes
\(d_{\sigma^{-j}\nu,j}\). The assumptions \(L_\nu,H_\nu<p\) say
exactly that addition within these groups has no base-\(p\) carry.
Every multinomial is therefore a unit. Undoing the rescaling and
multiplying by \(\prod_\tau a_\tau!\) gives valuation
\(\sum_\tau v_p(a_\tau!)+(k-A)/(Q-1)=B_p(\boldsymbol a)\).

The bundle \(P\) of Theorem~\ref{thm:pol49} is relatively ample,
symmetric and rigidified. The argument of Theorem~\ref{thm:globaltheta43}
therefore gives a relatively very ample \(P^3\) and an integral
section \(s_0\) of \(P^3\) with value one at the origin on this
\(K\)-scheme. Lemma~\ref{lem:very43} gives a surjection from the
degree-\(A\) products of its global sections onto the integral jets
in the frame \(s_0^A\). Its
base change to \(R\) is also onto. If every global section had
valuation strictly greater than this minimum then so would the
whole image of their \(R\)-span, contrary to the monomial just
constructed. This proves the last assertion.
\end{proof}

\begin{corollary}[Two mixed derivatives of the order in the prime count]\label{cor:separate52}
At every sufficiently large \(x\) in the common intervals there is
an unramified prime \(p\), split in \(K_q=\QQ(\sqrt{-q})\) and inert
in \(K\), and there are two lists of exponents on the same polarized
scheme with the same multiset of entries and total order
\(A=\lceil x e^{x^\eps}\rceil\), with the following property. At the
chosen place over \(p\), for the first list some global theta section
of \(P^{3A}\) has derivative valuation \(B_p\). For the second list
every integral section has derivative valuation at least \(B_p+1\).
Here \(B_p\) has the same value for both lists, and the tangent
directions and frame are fixed for the comparison.
\end{corollary}
\begin{proof}
In the Galois coordinates \((i,j,z)\) of Volume~I choose
\(g=(0,1,z)\) with \(z\) of order four in \(A_0\). This exists at
large \(x\) since the cyclic order \(|A_0|\) tends to infinity.
The restriction map from the units modulo
\(N=\operatorname{lcm}(F,2^{l'+2})qd\) to \(\Gal(E/\QQ)\) is
onto. Use Linnik's theorem in any preimage residue class to take a
prime \(p\ll N^{C}\) with Frobenius \(g\). It is unramified in \(E\)
and has the stated two signs. Since \(\log N=o(x^\eps)\) uniformly,
\(p^2+p-1<A\) for large \(x\). The permutation
\(\sigma=g^2\) on \(\Sigma\) consists of disjoint pairs, since it
is a nonidentity translation of order two on the torsor of embeddings
over \(K\).~\cite[Thm.~1.1]{XIII-X-IV-tz49}

Fix an embedding \(\nu\) for the exponent \(p-1\). For the first
list give the exponent \(p^2\) to an embedding in a different
\(\sigma\)-orbit. For the second move only this exponent to
\(\sigma^{-1}\nu\). To complete both lists write
\[
                 A-(p^2+p-1)=\sum_{j\ge0}c_j Q^j,\quad 0\le c_j<Q .
\]
For each nonzero \(c_j\), give the exponent \(c_j Q^j\) to an embedding
in a fresh orbit, identical for the two lists. At most
\(1+\log_2 A\) such orbits are needed and there are
\(|\Sigma|/2=n\gg1+\log_2 A\). Give every other embedding exponent
zero. In the first list no two contributions from different
embeddings reach the same \(\nu'\) in the digit sums of the theorem.
There is no carry in either kind of digit, hence the minimum is
attained. In the second list the \(p-1\) low digits at \(\nu\) and the
one low digit from the \(p^2\)-term now give \(L_\nu=p\). Apply
\eqref{eq:gain52}. The list of exponents has only been permuted,
which leaves \(B_p\) unchanged.
\end{proof}

This comparison detects an individual Frobenius action in the
divisibility of genuine global theta derivatives at the required total
order. The two multidegrees in the corollary are chosen to exhibit
that action. To apply it to the automorphic periods one must determine
the multidegrees of their prescribed harmonic coefficients, control
the sum over the period classes in the same integral frame and bound
the projected global norms. Section \ref{nav:I-28} still requires
these estimates and the remaining odd-primary cohomology bound.

\subsection{Reduction of polynomial differential operators at an inert prime}\label{nav:V-09}
\label{sec:algebra54}

The polynomials occurring in a theta evaluation can have several
monomials with the same total order. To decide whether a power of a
prime divides every value, one must also detect cancellations among
these monomials. We compute the algebra of the normalized invariant
differential operators modulo the maximal ideal. Its multiplication
keeps track of the permutation of the CM embeddings by Frobenius.

Use the scheme and prime \(p\) of Section~\ref{sec:mixed52}. More
generally, we allow any finite product of copies of that scheme.
In this case \(\Sigma\) is the disjoint union of the sets of tangent
eigendirections of the factors. A factor with conjugate CM action
has its embeddings relabeled by complex conjugation, which commutes
with \(\sigma=\Frob_p^2\). Frobenius acts within each factor. Let
\(R\) be the completed integer ring at the chosen place, with
residue field \(k\), and put \(Q=p^2\). We use the integral local
generators of the tangent lines constructed in the proof of
Theorem~\ref{thm:gain52}. Write
\[
 T_{\tau,c}=\frac{p^{F_Q(c)}}{c!}D_\tau^c,\qquad
 X_\tau=D_\tau,\qquad Y_\tau=D_\tau^p/p! .
\]
These are commuting operators on integral formal functions. A bar
denotes their action after reduction to \(k\). The same notation for a
scalar denotes its image in \(k\).

\begin{theorem}[The algebra of the reduced operators]\label{thm:algebra54}
The \(k\)-algebra generated by all \(\overline T_{\tau,c}\) is
\begin{equation}
 k[y_\tau:\tau\in\Sigma]\big/(y_\tau^{p^2}:\tau\in\Sigma),
                   \qquad y_\tau\longmapsto\overline Y_\tau .
                                                           \label{eq:algebra54}
\end{equation}
In this presentation \(\overline Y_\tau^{\,p}
=\overline w_\tau\,\overline X_{\sigma\tau}\) for a unit \(w_\tau\).
For the digits and the integers \(L_\nu,H_\nu\) of
Section~\ref{sec:mixed52}, the image of
\(\prod_\tau T_{\tau,a_\tau}\) is the monomial
\[
         \prod_{\nu\in\Sigma} y_\nu^{\,H_\nu+pL_{\sigma\nu}}
\]
times a unit. The proof gives the unit for each monomial.
\end{theorem}
\begin{proof}
By \eqref{eq:op52} are units \(u_{\tau,j}\) such that
\[
            T_{\tau,Q^j}\equiv u_{\tau,j}X_{\sigma^j\tau}
                                \pmod {p},\qquad u_{\tau,0}=1 .
\]
In the product coordinates of that proof one can take the coefficient
\(p^{F_Q(Q^j)}[Z^{Q^j}]\mathcal E(Z)\) of the elliptic exponential
on the factor containing \(\tau\). The congruence means that the
difference maps integral formal functions into \(pR[[T]]\), even
when \(R\) is ramified. Define the rational unit at \(p\)
\[
            c_j=\frac{p!\,((Q^j)!)^p}{(pQ^j)!}.
\]
Legendre's formula gives
\(v_p((pQ^j)!)=1+p\,v_p((Q^j)!)\). Also
\(F_Q(pQ^j)=pF_Q(Q^j)\), so the exact identity in characteristic
zero is \(T_{\tau,pQ^j}=c_j T_{\tau,Q^j}^{\,p}/p!\).
Write \(T_{\tau,Q^j}=u_{\tau,j}X_{\sigma^j\tau}+pB\), with \(B\)
integral. The two terms commute since \(B\) is a linear combination
of commuting invariant operators. In the binomial expansion of
their \(p\)-th power every term containing \(B\) is divisible as
an operator by \(p^2\). For the last term this uses \(p^p\),
including when \(p=2\), and for an intermediate term it uses the
binomial coefficient and the factor \(p\). Since \(v_p(p!)=1\),
\begin{equation}
            T_{\tau,pQ^j}\equiv
            c_j u_{\tau,j}^{\,p}Y_{\sigma^j\tau}\pmod p .
                                                            \label{eq:high54}
\end{equation}
On the other hand
\[
        Y_\tau^p=\frac{(p^2)!}{p(p!)^p}\,T_{\tau,Q}.
\]
The rational factor is a unit. Thus we may take
\(w_\tau=((p^2)!/(p(p!)^p))u_{\tau,1}\).
The exact equality \(X_\tau^p=p!Y_\tau\) and the preceding
congruence give \(\overline Y_\tau^{\,p^2}=0\).

For each \(a_\tau=\sum_j(l_{\tau j}+p h_{\tau j})Q^j\), factor
\(T_{\tau,a_\tau}\) into the \(Q^j\) and \(pQ^j\) operators with
the indicated multiplicities, as in Theorem~\ref{thm:gain52}.
The unit multiplying their product is
\[
      \frac{\prod_j ((Q^j)!)^{l_{\tau j}}
                          ((pQ^j)!)^{h_{\tau j}}}{a_\tau!}.
\]
Apply the two operator congruences and then replace
\(\overline X_\nu\) by
\(\overline w_{\sigma^{-1}\nu}^{-1}\overline Y_{\sigma^{-1}\nu}^{\,p}\).
These substitutions give the stated exponents and all scalar factors.
They show that \eqref{eq:algebra54} maps onto the operator algebra.

To prove injectivity use the monomials
\(\prod_\nu\overline X_\nu^{\,b_\nu}\overline Y_\nu^{\,d_\nu}\),
where \(0\le b_\nu,d_\nu<p\). By the relation for
\(\overline Y_\nu^{\,p}\), these monomials differ by units from
the proposed basis of monomials in the \(y_\nu\) with exponents below
\(p^2\), each occurring once. Put \(e_\nu=b_\nu+pd_\nu\) and
\(s=\sum_\nu e_\nu\). Choose integral linear formal coordinates
\(T'_\nu\) dual to the \(D_\nu\) in the tangent space at zero.
For this operator, its value on a monomial of total degree \(s\) at
zero comes only from applying each differentiation to a coordinate.
All terms where a derivative acts on the coefficients of a vector
field have order below \(s\) on the test monomial, so they vanish.
As a result the value on \(\prod_\nu(T'_\nu)^{e_\nu}\) is
\(\prod_\nu e_\nu!/(p!)^{d_\nu}\), a unit since
\(e_\nu<p^2\) and \(v_p(e_\nu!)=d_\nu\). It is zero on every other
monomial of that same total degree. These assertions follow in
characteristic zero before reduction, where division by the factorials
is legitimate. An operator of lower order vanishes on every
monomial of degree \(s\) at zero as well. In any proposed nonzero
linear relation, choose a term with \(s\) maximal among those with
nonzero coefficient. Its test monomial therefore isolates a nonzero
value in \(k\), a contradiction. This proves independence, including
independence of the resulting distributions evaluated at zero.
\end{proof}

\begin{corollary}[An exact test for additional divisibility]\label{cor:carry54}
For any list of orders on these factors and any integral scalar
jet through their total order one has
\begin{equation}
 v_p\!\left(\left(\prod_\tau D_\tau^{a_\tau}F\right)(0)\right)
 \ \ge\ B_p(\boldsymbol a)+
       \sum_{\nu\in\Sigma}\left\lfloor
             \frac{H_\nu+pL_{\sigma\nu}}{p^2}\right\rfloor .
                                                             \label{eq:carry54}
\end{equation}
The minimum over integral jets equals \(B_p(\boldsymbol a)\) if
and only if every exponent \(H_\nu+pL_{\sigma\nu}\) is below \(p^2\).
\end{corollary}
\begin{proof}
Group the factors from \eqref{eq:high54} for a common \(\nu\).
Any \(p\) of them together are congruent modulo \(p\) to a unit
times \(X_{\sigma\nu}\), by the computation of \(Y_\nu^p\).
At a position \(\mu\) this gives
\(\lfloor H_{\sigma^{-1}\mu}/p\rfloor\) disjoint groups congruent
to unit multiples of \(X_\mu\), in addition to the \(L_\mu\)
original factors of that kind. Any \(p\) such groups together
are divisible by \(p\) as operators because \(X_\mu^p=p!Y_\mu\).
All the groups are disjoint and commute. Dividing their products by
the powers of \(p\) so obtained leaves integral operators. The number
of powers obtained at \(\mu\) is
\(\lfloor(H_{\sigma^{-1}\mu}+pL_\mu)/p^2\rfloor\).
Restoring the original factorials proves \eqref{eq:carry54}.
If all these terms are zero, the reduced monomial of
Theorem~\ref{thm:algebra54} is nonzero and its test monomial gives
the scalar minimum. Its degree is at most \(\sum a_\tau\).
Indeed replacing \(p\) occurrences of \(Y_\nu\) by
\(X_{\sigma\nu}\) lowers the order, while the initial replacement
of each \(T_{\tau,Q^j}\) or \(T_{\tau,pQ^j}\) has order at most
that of the replaced factor. If some term is positive,
\eqref{eq:carry54} excludes the scalar minimum.
\end{proof}

For example, fix \(\nu\) and prescribe one high digit at each
\((\sigma^{-j}\nu,j)\) for \(0\le j<p\), all other digits zero.
Then \(H_\nu=p\) gives \(y_\nu^p\ne0\), and the scalar minimum
is attained. Now add \(p-1\) low units at \((\sigma\nu,0)\).
The same exponent becomes \(p+p(p-1)=p^2\), so every integral
section acquires at least one extra power of \(p\).

\begin{theorem}[Polynomial coefficients in a projected theta section]\label{thm:poly54}
Use both auxiliary primes and the common interval of
Theorem \ref{thm:project53}.
Keep \(p\) unramified in \(E\), inert in \(K\) and distinct
from those two primes, and put \(a=\lceil x\exp(x^\eps)\rceil\). Let
\(\mathcal P(\boldsymbol D)=\sum_{\boldsymbol a}c_{\boldsymbol a}
\prod_\tau D_\tau^{a_\tau}\) be a nonzero polynomial over
\(\operatorname{Frac}R\), of total degree at most \(a\).
For its nonzero coefficients form
\[
       b_{\boldsymbol a}
         =c_{\boldsymbol a}\prod_\tau
                     \frac{a_\tau!}{p^{F_Q(a_\tau)}} .
\]
Choose \(\beta\) of smallest valuation among these numbers. In the
algebra \eqref{eq:algebra54} form
\[
      W_{\mathcal P}
       =\sum_{\boldsymbol a}
          \overline{b_{\boldsymbol a}/\beta}\,
                          \prod_\tau\overline T_{\tau,a_\tau}.
\]
In every translation character the minimum valuation at zero of
\(\mathcal P(\boldsymbol D)\) on the integral projected sections in
the frame of Theorem \ref{thm:project53} equals \(v_p(\beta)\) precisely when
\(W_{\mathcal P}\ne0\). When \(W_{\mathcal P}=0\) every value has
valuation at least \(v_p(\beta)+1/e_w\), where \(e_w\) is the absolute
ramification index of \(R\). This test applies to any integral
harmonic polynomial expressed in the chosen directions.
\end{theorem}
\begin{proof}
Dividing the operator by \(\beta\) gives exactly a sum of the integral
normalized operators with coefficients \(b_{\boldsymbol a}/\beta\).
If \(W_{\mathcal P}=0\), this operator vanishes after reduction, so
its image lies in the maximal ideal of \(R\). If \(W_{\mathcal P}\ne0\),
collect its monomials using the units in Theorem~\ref{thm:algebra54}.
The maximal-order test in that proof gives an integral jet with
unit value. Its required degree is at most \(a\) by the order
calculation in Corollary~\ref{cor:carry54}. Theorem \ref{thm:project53}
shows that each \(\lambda\)-component supplies every such jet in
the specified frame. This argument also holds for the product of
factors in the present section. In fact the product isogeny is
still finite \'etale at \(p\), its character lines on the target
are torsion, and the cube of its ample symmetric polarization is
relatively very ample, exactly the hypotheses used in that proof.
Flat base change identifies the local sections with the span
of the global integral sections. If all projected global sections
gave values in \(\beta\) times the maximal ideal, their \(R\)-span
would too. Thus when \(W_{\mathcal P}\ne0\) some global section
attains the minimum in that character at the chosen place.
\end{proof}

\begin{corollary}[The full local image ideal]\label{cor:all54}
In the setting of Theorem~\ref{thm:poly54}, choose elliptic formal
parameters \(T_i\) on the product and write \(v_{i\tau}\) for the
coordinates of \(D_\tau\) on the tangent space. Let
\(\mathcal E_i(Z)\) be the elliptic formal exponential in coordinate
\(i\), with linear term \(Z\). For \(0\le|\boldsymbol k|\le a\) set
\[
 j_{\boldsymbol k}
 =\sum_{\boldsymbol a} c_{\boldsymbol a}
       \left(\prod_\tau a_\tau!\right)
       [\boldsymbol Z^{\boldsymbol a}]\,
        \prod_i\mathcal E_i\!\left(\sum_\tau v_{i\tau}Z_\tau\right)^{k_i}.
\]
For each \(\lambda\), the image of evaluation at zero on all
\(R\)-integral projected sections is the nonzero fractional ideal
\(\sum_{|\boldsymbol k|\le a}Rj_{\boldsymbol k}\). In particular
the smallest valuation is \(\min_{|\boldsymbol k|\le a}v_p(j_{\boldsymbol k})\),
also when the first reduction \(W_{\mathcal P}\) vanishes.
\end{corollary}
\begin{proof}
In characteristic zero the formal exponential takes addition on
the tangent space to the formal group law. Differentiating the
displayed product in the variables \(Z_\tau\) thus evaluates
\(\mathcal P(\boldsymbol D)\) on the monomial
\(\prod_iT_i^{k_i}\) at the origin. These monomials form an \(R\)-basis
of the relative jets through order \(a\). Their evaluations therefore
generate exactly the asserted image by the surjectivity for each
character. The polynomial is nonzero in characteristic zero. Among
its terms of maximum degree, a monomial in dual tangent coordinates
with the matching exponents pairs nontrivially with one term and
annihilates the others at zero. Thus at least one generator is nonzero.

Only finitely many coefficients are needed in this formula. In an
elliptic coordinate of Section~\ref{sec:exact41}, write
\([p](T)=pT+\sum_{m\ge Q}d_mT^m\) and
\(\mathcal E_i(Z)=\sum_{s\ge1}h_sZ^s\). Starting with \(h_1=1\),
comparison of \([p](\mathcal E_i(Z))\) and
\(\mathcal E_i(pZ)\) gives for \(2\le s\le a\)
\[
        (p^s-p)h_s=\sum_{m=Q}^s d_m
           \sum_{\substack{j_1+\cdots+j_m=s\\j_i\ge1}}
                         h_{j_1}\cdots h_{j_m}.
\]
Every index on the right is smaller than \(s\) and the scalar on
the left is nonzero. This supplies the required coefficients from
the group law through order \(a\).
\end{proof}

In the source, the homogeneous harmonic coefficient and the
evaluation on a translated section appear in
\srcloc{Formalism/SharpRestriction/Divisibility/CanonicalTranslateJets}{165}.
The criterion computes the first possible local valuation once the
coefficient is specified. The isomorphism \eqref{eq:algebra54}
determines cancellation between its monomials. To bound the
arithmetic periods of Section \ref{nav:I-28}, the coefficients
selected by the global lifts, their sum over the period classes
and their archimedean normalization still require a comparison.

\section{Elliptic derivatives on definite classes}\label{nav:g-V-1}

\subsection{The norm of the prescribed products of maximal minors}\label{nav:V-10}
\label{sec:min55}

The auxiliary norm in the source uses an ordered product of maximal
minors at every CM embedding. Its definition fixes both the
polynomials and the position of the class sum. We compute the norm
with these choices. Let \(\Sigma\) be the CM type over
\(K=\QQ(\sqrt{-d})\) on the fields of
Section \ref{nav:I-28},
and write \(f=|\Sigma|=[E:K]\).
Take ranks \(1\le r\le s\), weights \(b_\tau\in\ZZ_{\ge0}\), and
write \(Z=(Z_\tau)_{\tau\in\Sigma}\) for \(s\)-by-\(r\) complex
matrix variables. Set \(k_\tau=rb_\tau\), \(k=\sum_\tau k_\tau\)
and \(D=frs\). For an \(r\)-element subset \(I\) of row indices,
write \(\Delta_I(Z_\tau)\) for its minor in increasing row order.
For every ordered choice \(T=(I_{\tau j}:1\le j\le b_\tau)\) put
\[
              P_T(Z)=\prod_{\tau,j}\Delta_{I_{\tau j}}(Z_\tau).
\]
Empty products are one. Let \(c,h\) range over the source and target
classes for the two compact levels, with weights \(m_c,m'_h\) equal
to the reciprocals of the respective stabilizer orders. Write
\(F_{ch}(Z)\) for the holomorphic theta function at that pair of
representatives, including the prescribed finite Weil action on the
chosen vector. The quantity computed here at a complex embedding is
\begin{equation}
 \mathcal N=\sum_h m'_h\sum_T
        \left|\sum_c m_c P_T(\partial)F_{ch}(0)\right|^2 .
                                                    \label{eq:pair55}
\end{equation}
These are the minors, derivatives and two class sums of
\srcloc{Development/AbelianVarietyTheory/Definitions/Foundations/NormSumPairRankZeroOfCurveCoordNorm}{105}
and lines 153--162 there. The ordinary derivative convention is
specified in
\srcloc{Formalism/SharpRestriction/GlobalSchurBounds/GL3Theta}{533}.

On polynomials in the matrix entries give distinct monomials
orthogonal Fischer inner products and let
\(\|Z^\alpha\|_{\mathrm F}^2=\alpha!=\prod_i\alpha_i!\).
For a germ \(F\) let \(J_{\boldsymbol k}F\) be its Taylor component
homogeneous of degree \(k_\tau\) in each block. Denote by \(\Pi\)
the orthogonal projection onto the span of all \(P_T\) in that
component.

\begin{theorem}[An exact archimedean norm for the class sums]\leavevmode\label{thm:frame55}
Define
\[
          C_{r,b}=\prod_{j=0}^{r-1}\frac{(b+j)!}{j!}.
\]
For the functions in \eqref{eq:pair55} one has
\begin{equation}
 \mathcal N=\left(\prod_{\tau}C_{r,b_\tau}\right)
       \sum_h m'_h
         \left\|\Pi J_{\boldsymbol k}\!\left(\sum_c m_c F_{ch}\right)
                                              \right\|_{\mathrm F}^2 .
                                                      \label{eq:fisch55}
\end{equation}
In particular this identity keeps the cross terms between distinct
source classes mapping to the same target class.
\end{theorem}
\begin{proof}
First use one matrix variable and one integer \(b\). The
Cauchy--Binet determinant identity, raised to the \(b\)-th power,
gives
\[
        \det(W^*Z)^b=\sum_T P_T(Z)\overline{P_T(W)} .
\]
Ordered choices in this sum count each term with its multinomial
multiplicity. Take the Fischer inner product conjugate-linear in the
first entry. Define \(S H=\sum_T P_T\langle P_T,H\rangle_{\mathrm F}\).
This is self-adjoint and positive. Its range is the span \(V_b\)
of the \(P_T\), since its kernel is their common orthogonal complement.
It commutes with unitary changes of the \(s\) rows. Indeed
the displayed kernel is unchanged by the same unitary change of
rows in both matrices, and the Fischer inner product is invariant.
For the latter assertion its integral representation is the ordinary
Gaussian inner product for independent entries with density
\(\pi^{-1}e^{-|z|^2}\). Integration of the monomials in polar
coordinates verifies this representation. The kernel determines
\(S\) by pairing its second polynomial variable with \(H\).

Let \(P_0\) be the \(b\)-th power of the minor on the first \(r\)
rows. If \(b>0\), any of the other indexed polynomials uses a row
outside the first \(r\), with positive degree in that row in every
monomial. Thus \(S P_0=\|P_0\|_{\mathrm F}^2P_0\). For \(b=0\)
this also holds since the sole index gives the constant one.
The unitary row translates of \(P_0\) span \(V_b\). To see this,
the kernels \(\det(W^*Z)^b\) as \(W\) varies span \(V_b\).
Orthogonality to all of them implies that the polynomial \(S H\)
vanishes at every \(W\). For full-rank \(W\), orthonormalizing its
columns writes this kernel as a scalar multiple of a unitary
translate of \(P_0\). For deficient rank the kernel is zero when
\(b>0\). Hence \(S=\|P_0\|_{\mathrm F}^2\Pi\).

It remains to compute the scalar. For an \(r\)-by-\(r\) matrix of
independent standard complex Gaussian entries it is
\(\mathbb E|\det Z|^{2b}\). Orthonormalize the columns in order.
The squared determinant is the product of the squared lengths of
their components normal to the preceding columns. Given these
columns, in complex codimension \(j\) the squared length of the
next component is a sum of \(j\) independent variables of density
\(e^{-t}\) on \(t\ge0\), by unitary invariance. Its \(b\)-th
moment is \((b+j-1)!/(j-1)!\), obtained by integrating the gamma
density \(t^{j-1}e^{-t}/(j-1)!\). Successive conditional expectation
gives \(C_{r,b}\). For several embeddings the operators, their spans,
and their inner products tensor. Finally \(P_T\) has real coefficients,
so \(P_T(\partial)F(0)=\langle P_T,J_{\boldsymbol k}F\rangle_{\mathrm F}\).
Apply the tensor identity to the sum of the \(F_{ch}\) at each \(h\).
\end{proof}

\begin{proposition}[The Gaussian bound after summing all minors]\label{prop:gauss55}
Write \(|Z|^2=\sum_\tau\operatorname{tr}(Z_\tau^*Z_\tau)\),
\(M_{\rm in}=\sum_c m_c\), \(M_{\rm out}=\sum_hm'_h\).
If \(|F_{ch}(Z)|^2\le S\,\exp(\pi |Z|^2/A)\) for every \(c,h,Z\)
with \(A>0\), then
\begin{equation}
    \mathcal N\le M_{\rm out}M_{\rm in}^2\, S\,
            2^D(2\pi/A)^k\prod_\tau C_{r,b_\tau}.       \label{eq:gauss55}
\end{equation}
The factorial in this bound satisfies
\[
 0\le\log C_{r,b}-r\log(b!)
      =\sum_{j=0}^{r-1}\log\binom{b+j}{j}
      \le\tfrac12r(r-1)\log(1+b+r).
\]
\end{proposition}
\begin{proof}
Set \(G_h=\sum_c m_c F_{ch}\). The triangle inequality gives
\[
              |G_h(Z)|^2\le M_{\rm in}^2S\exp(\pi|Z|^2/A).
\]
Integrate with the probability
density \((t/\pi)^D\exp(-t|Z|^2)\), where \(t=2\pi/A\).
The integral is at most \(2^D M_{\rm in}^2S\). Independent scalar
rotations of the matrix blocks act unitarily for this Gaussian. The
Fourier projection under these rotations onto multidegree
\(\boldsymbol k\) is \(J_{\boldsymbol k}\) by the Taylor expansion,
and is a contraction by Cauchy--Schwarz. Gaussian scaling gives
\(\int|J_{\boldsymbol k}G_h|^2=t^{-k}\|J_{\boldsymbol k}G_h\|_{\rm F}^2\).
The further projection \(\Pi\) is also a contraction. Substitute in
\eqref{eq:fisch55}. The last inequality follows from
\(\binom{b+j}{j}\le(1+b+r)^j\).
\end{proof}

\begin{theorem}[A class-independent bound for the CM theta metric]\label{thm:lattice55}
For a class pair let \(L,L'\) be self-dual Hermitian
\(\mathcal O_E\)-lattices in the definite standard spaces of ranks
\(r,s\). Take \(\mathfrak b\) from Theorem~\ref{thm:pol49}, and
write \(\Phi(X)=(\tau X)_{\tau\in\Sigma}\).
For a nonzero scalar \(\delta\in\CC\), set
\[
 \Lambda=\delta\Phi\!\left(\mathfrak b^{-1}
             \operatorname{Hom}_{\mathcal O_E}(L,L')\right),
 \qquad A=|\delta|^2\sqrt d/2,\qquad H(Z,Z)=|Z|^2/A .
\]
In this metric the lattice has real covolume one and each nonzero
vector has squared length at least \(2f/\operatorname{rd}(E)\),
where \(\operatorname{rd}(E)=|\operatorname{disc}E|^{1/[E:\QQ]}\).
If \(\Theta\) is a holomorphic theta section for this lattice in
the Hermitian frame, let \(\|\Theta\|_2^2\) be the average of
\(|\Theta(Z)|^2e^{-\pi H(Z,Z)}\) on the torus. Then
\begin{equation}
 \sup_Z|\Theta(Z)|^2 e^{-\pi H(Z,Z)}
 \le \left[e\max\left(1,\frac{2rs\,\operatorname{rd}(E)}{\pi}\right)
      \right]^D\|\Theta\|_2^2.                          \label{eq:sup55}
\end{equation}
\end{theorem}
\begin{proof}
The module \(M_0=\operatorname{Hom}_{\mathcal O_E}(L,L')\) is
self-dual for \(h(X,Y)=\operatorname{tr}(Y^*X)\). Indeed it is
the tensor product of the dual source lattice with the target
lattice, and the tensor product of their perfect Hermitian
pairings is exactly \(h\). This can be checked on free bases after
localization. On \(\mathfrak b^{-1}M_0\) the trace form
\(\operatorname{Tr}_{E/\QQ}(-h(X,Y)/\sqrt{-d})\) is unimodular
by the inverse different as in Theorem~\ref{thm:pol49}. On
\(\Lambda\) it is \(-\operatorname{Im}H\). Its top exterior
power divided by \(D!\) is the Euclidean volume form for \(H\),
up to sign. A symplectic lattice basis therefore has volume one.

For nonzero \(X\in\mathfrak b^{-1}M_0\), \(v=h(X,X)\) is totally
positive and belongs to \(\mathfrak D_{E/K}^{-1}\), since
\(\mathfrak b\bar{\mathfrak b}=\mathfrak D_{E/K}\).
Thus \(|N_{E/\QQ}v|\ge1/N_{E/\QQ}\mathfrak D_{E/K}\).
The latter ideal norm is \(|\operatorname{disc}E|/d^f\).
The arithmetic--geometric mean inequality over the \(f\) places
gives
\[
   \sum_\tau\operatorname{tr}((\tau X)^*(\tau X))
       \ge f\,|N_{E/\QQ}v|^{1/(2f)}
       \ge f\sqrt d/\operatorname{rd}(E).
\]
Multiplication by \(|\delta|^2/A\) gives the asserted separation.

For a point \(z\), the translated holomorphic function
\[
 U_z\Theta(u)=\Theta(u+z)\exp\{-\pi H(u,z)-\pi H(z,z)/2\}
\]
has metric squared magnitude equal at \(u\) to the original one
at \(u+z\). Here \(H\) is complex-linear in its first variable.
Apply the mean-value inequality for the squared absolute value of
this function on the open \(H\)-ball of radius
\(\rho\), where \(\rho^2=\min(D/\pi,f/(2\operatorname{rd}(E)))\).
The ball projects injectively up to its boundary by the separation
just proved. Inside it the unweighted integral is at most
\(e^{\pi\rho^2}\|\Theta\|_2^2\), since the metric density descends
and the torus has volume one. The ball has volume
\(\pi^D\rho^{2D}/D!\). Using \(D!\le D^D\) proves \eqref{eq:sup55}.
\end{proof}

\begin{corollary}[Translation and transport of the vacuum]\label{cor:move55}
Suppose the functions entering \eqref{eq:pair55} are the finite
Heisenberg translates of a transported theta vacuum, with
\[
         F_{ch}=\sum_i q_{i,ch} U_{t_{i,ch}}\Theta_{ch},
            \qquad \left(\sum_i|q_{i,ch}|\right)^2\le Q^{rs}.
\]
Unitary character factors may be included in the coefficients.
If each \(\Theta_{ch}\) is obtained from a vacuum of squared norm
\(\nu_0\) by normalized finite lattice intertwiners for the common
Hermitian form, then \eqref{eq:gauss55} holds with
\[
           S=Q^{rs}
            [e\max(1,2rs\,\operatorname{rd}(E)/\pi)]^D\nu_0.
\]
This applies for every list of the prescribed minor weights at the
same \(x\), including the ranks in Theorem \ref{thm:rangenew50}.
\end{corollary}
\begin{proof}
For two commensurable self-dual period lattices \(\Lambda_1,\Lambda_2\),
put \(B=\Lambda_1\cap\Lambda_2\), \(q=[\Lambda_2:B]\).
For compatible semicharacters a normalized intertwiner is
\(q^{-1/2}\sum_{t\in\Lambda_2/B}\epsilon_t U_t\) with
\(|\epsilon_t|=1\), the phases chosen for the target automorphy.
This is the normalization in
\srcloc{Formalism/LocalToGlobal/KernelLattices/ColumnData}{117}.
Each translate has the same squared norm on the torus for \(B\),
using average Haar measure throughout. Distinct translates are
orthogonal there. Indeed for \(t-t'\notin\Lambda_1\), self-duality
gives \(w\in\Lambda_1\) such that
\(\operatorname{Im}H(w,t-t')\notin\ZZ\). The unitary operator
\(U_w\) on that common automorphy space has distinct eigenvalues
on \(U_t\Theta\) and \(U_{t'}\Theta\), since its commutator with
\(U_t\) is \(e^{-2\pi i\operatorname{Im}H(w,t)}\).
The norm squared of the normalized sum is therefore the original
squared norm. For a function with target automorphy its average
metric norm is the same over \(B\) and \(\Lambda_2\).
Iteration proves \(\|\Theta_{ch}\|_2^2=\nu_0\). Finally the metric
supremum is preserved under any \(U_t\). The triangle inequality and
\eqref{eq:sup55} give the common envelope used in \eqref{eq:gauss55}.
\end{proof}

The last corollary uses the finite translation presentation of the
vector. The squared norm \(\nu_0\) records its normalization at the
complex embedding in question. For an arithmetic application one
also needs the normalization at the conjugate embeddings, the
integral sections realizing the periods, and the Bessel projection
on the row tower. The identity \eqref{eq:fisch55} concerns the
auxiliary norm of the prescribed minors before this last comparison.

\subsection{Arithmetic of the normalized elliptic theta derivatives}\label{nav:V-11}
\label{sec:ell57}

Use the discriminant-one curve \(C=\CC/\Lambda\), the differential
\(dz=dx/y\), and a nonzero point \(t\in C[2]\) from
Section \ref{nav:I-28}.
Write \(A\) for the area of \(\Lambda\). For a lift \(w\in\CC\) put
\[
    U_wF(z)=e^{-\pi z\bar w/A-\pi|w|^2/(2A)}F(z+w).
\]
Let \(\Theta_t\) be the normalized even theta in
Proposition \ref{prop:ellnorm56}. We use its ordinary complex derivatives
in the fixed differential. They determine the Taylor coefficients
by division by the corresponding factorials. This convention is
specified in
\srcloc{Formalism/SharpRestriction/GlobalSchurBounds/StandardLatticesGL3}{45}.

\begin{theorem}[An algebraic formula for every derivative]\label{thm:germ57}
For \(p\) a rational prime write \(v_p(p)=1\) at each place over
\(p\). All the numbers \(H_n=\Theta_t^{(n)}(0)\), \(n\ge0\),
belong to a single number field. At every complex embedding their
images are exactly the derivatives of the normalized theta on the
conjugate curve, with the conjugate mark and differential.
For \(p\ge5\) they satisfy
\[
       v_p(H_n)\ge
          \begin{cases}
             0 & p\ne d,\\
             -\tfrac12\lfloor n/2\rfloor & p=d .
          \end{cases}                                      \label{eq:ellval57}
\]
\end{theorem}
\begin{proof}
Let \(\sigma,\zeta,\wp\) be the Weierstrass functions of
\(\Lambda\), so that \((\log\sigma)'=\zeta\), \(\zeta'=-\wp\),
\(\wp=x\), and \(\wp'=y\). The increments
\(\eta_\Lambda(\lambda)=\zeta(z+\lambda)-\zeta(z)\) are additive
on the lattice. Their real-linear extension to \(\CC\) has the
form \(u_\Lambda z+v_\Lambda\bar z\). If
\(\lambda_1,\lambda_2\) is a positively oriented basis, the
residue of \(\zeta\) at its single pole in a period parallelogram
gives
\(\eta_\Lambda(\lambda_1)\lambda_2
-\eta_\Lambda(\lambda_2)\lambda_1=2\pi i\).
Since \(\bar\lambda_1\lambda_2-\bar\lambda_2\lambda_1=2iA\),
one gets \(v_\Lambda=\pi/A\). These definitions and
quasiperiod relations are also given in
\href{https://dlmf.nist.gov/23.2}{NIST, Weierstrass functions}.

The remaining complex number \(u_\Lambda\) has an algebraic
description. Let \(\alpha\in\OO_K\setminus\RR\) and sum over the
kernel \(C[\alpha]\) of its action. Comparing poles in the
elliptic functions \(\sum_P\wp(z+P)\) and
\(\alpha^2\wp(\alpha z)\) gives
\[
     \sum_{P\in C[\alpha]}\wp(z+P)=\alpha^2\wp(\alpha z)+c_\alpha,
             \qquad c_\alpha=\sum_{P\ne0}x(P).
\]
Their poles have the same principal parts and their difference at
zero has the stated constant term. Integrate using \(\zeta'=-\wp\)
and then take the increment under \(z\mapsto z+\lambda\).
The additive integration constant drops out and we obtain
\[
 N_{K/\QQ}(\alpha)\eta_\Lambda(\lambda)
          -\alpha\eta_\Lambda(\alpha\lambda)=-c_\alpha\lambda.
\]
The antiholomorphic terms cancel. With \(\alpha=\sqrt{-d}\) this
gives
\begin{equation}
  u_\Lambda=\frac{c_\alpha}{\alpha(\alpha-\bar\alpha)}
    =-\frac1{2d}\sum_{0\ne P\in C[\sqrt{-d}]}x(P).
                                                        \label{eq:quasi57}
\end{equation}
Multiplication by elements of \(\OO_K\) on this algebraic CM
curve is algebraic over \(\overline{\QQ}\). For instance a torsion
kernel is algebraic, the quotient isogeny is algebraic, and
isomorphisms of elliptic curves with prescribed origin between the
quotient and \(C\) are defined over \(\overline{\QQ}\).
The derivative of the transported endomorphism is the image of
\(\alpha\) on the transported differential, so its kernel is
the corresponding kernel there. In particular replacing
\(\sqrt{-d}\) by its negative leaves the kernel in
\eqref{eq:quasi57} fixed. The formula proves both algebraicity
and transport of \(u_\Lambda\).

The odd Hermitian theta with derivative one is
\(e^{-u_\Lambda z^2/2}\sigma(z)\). Indeed the Weierstrass sigma
transformation law and the formula for \(\eta_\Lambda\) give its
unitary multipliers, with sign \(+1\) on \(2\Lambda\) and \(-1\)
on the other lattice vectors. It has its sole simple zero at the
origin. Translate by a lift of \(t\) and normalize the value at
zero. Oddness of \(\zeta\) and its increment under \(2t\) give
\(\zeta(t)=\eta_\Lambda(2t)/2\). Hence
\[
  \Theta_t(z)=e^{-u_\Lambda z^2/2-\zeta(t)z}
                      \frac{\sigma(z+t)}{\sigma(t)},\qquad
       (\log\Theta_t)''(z)=-u_\Lambda-\wp(z+t).
\]
The first logarithmic derivative at zero is zero. If \(L_i\)
denotes the \(i\)-th logarithmic derivative there, we have
\[
 L_2=-u_\Lambda-x(t),\qquad L_i=-(\partial^{\,i-2}x)(t)
          \quad(i\ge3),\qquad
 \partial x=y,\quad \partial y=6x^2-g_2/2 .
\]
These identities follow by differentiating the Weierstrass
equation. They show that every \(L_i\) is algebraic in the same
number field and commutes with each embedding. The identity
\(\Theta'_t=\Theta_t(\log\Theta_t)'\) and Leibniz's rule give
\[
 H_0=1,\qquad H_n=\sum_{i=1}^n\binom{n-1}{i-1}L_i H_{n-i}.
\]
This proves the assertions about the \(H_n\) over number fields.

At \(p\ge5\) the equations \(g_2^3=j/1728\) and
\(27g_3^2=g_2^3-1\) imply that \(g_2,g_3\) are integral since
\(j\) is integral. The unit discriminant gives good reduction in
this model. A nonzero point of order prime to \(p\) cannot reduce
to zero, since prime-to-\(p\) torsion in good reduction specializes
injectively. Its affine coordinates are therefore integral. At
\(p\ne d\) use \eqref{eq:quasi57}. All its points have order
dividing \(d\), giving \(v_p(u_\Lambda)\ge0\). At \(p=d\) use the
first expression of \eqref{eq:quasi57} with
\(\alpha=(1+\sqrt{-d})/2\). Its norm \((1+d)/4\) is prime to \(d\)
and \(v_d(\alpha(\alpha-\bar\alpha))=1/2\). Thus
\(v_d(u_\Lambda)\ge-1/2\). The mark \(t\) is prime-to-\(p\)
torsion in both cases, so \(L_i\) is integral for \(i\ge3\).
Only \(L_2\) can cost a factor of valuation \(-1/2\), and every
occurrence uses two orders in the recurrence. Induction gives
the stated local bounds.
\end{proof}

\begin{theorem}[Algebraic torsion translates in the Hermitian frame]\label{thm:trans57}
Let \(w\in\CC\) satisfy \(m w\in\Lambda\) for an integer \(m\ge1\).
The derivatives at zero of \(U_w\Theta_t\) lie in one finite
extension of the field in Theorem~\ref{thm:germ57}. At a complex
embedding they become the corresponding derivatives of
\(U_{w'}\Theta_{t'}\) times a single root of unity, independent
of the order of differentiation. Here the marked curve is
transported and \(w'\) lifts the transported point \([w]\).
At every \(p\ge5\) prime to \(m\) the same valuation bounds as in
Theorem~\ref{thm:germ57} hold for the translated derivatives.
\end{theorem}
\begin{proof}
In this proof write \(P=[w]\) and \(R_w(z)=U_w\Theta_t(z)/\Theta_t(z)\).
Under a period translation \(\lambda\in\Lambda\), the function
\(R_w\) is multiplied by a root of unity of order dividing \(m\).
This follows by dividing the transformation laws of the numerator
and denominator. Their ratio of multipliers is
\(\exp(\pm2\pi i\,\imag(\lambda\bar w)/A)\), whose \(m\)-th power
is one by integrality of the alternating form on \(\Lambda\).
In particular \(R_w^m\) is an elliptic function, with divisor
\(m([t-P]-[t])\). This degree-zero divisor has trivial class
in \(\mathrm{Pic}^0(C)=C\) since \(mP=0\). Choose
\(h\in\overline{\QQ}(C)^\times\) with that divisor. Then
\(R_w^m=ch\) for a nonzero complex constant \(c\).

Successive translations along the same vector have no commutator
phase, so \(U_w^m=U_{mw}\). More explicitly, the product of the
exponential factors at \(z+jw\), \(0\le j<m\), is the exponential
factor for \(mw\). Consequently
\[
 \prod_{j=0}^{m-1}R_w(z+jw)=\epsilon_t(mw),\qquad
 c^m\prod_{j=0}^{m-1}h(z+jP)=\epsilon_t(mw)^m .
\]
The multiplier \(\epsilon_t(mw)\) is \(1\) or \(-1\) because the
characteristic is a half-characteristic. The product of the
translates of \(h\) is a nonzero algebraic constant. Its divisor
telescopes to zero. Thus \(c\) is algebraic.

More precisely \(R_w(mz)\) is the uniformization of a rational
function on \(C\) over \(\overline{\QQ}\). It is periodic by the
multiplier computation and its divisor is
\([m]^*([t-P]-[t])\). This divisor is principal over
\(\overline{\QQ}\) because pullback by \([m]\) on \(\mathrm{Pic}^0\)
is multiplication by \(m\). Divide \(R_w(mz)\) by an algebraic
function with that divisor. The quotient is constant, and its
\(m\)-th power is algebraic by \(R_w(mz)^m=c\,h([m]z)\).
This also gives a finite field for the rational function.
For the formal parameter \(T=-2x/y\) at zero we have \(T=z+O(z^2)\).
The expansions \(T(z)\) and its formal inverse have coefficients
in \(\QQ(g_2,g_3)\), by formal integration of \(dx/y\) in \(T\)
and reversion of its series with linear term one. They commute
with field embeddings. Substituting \(z/m\) gives a single finite
coefficient field for the germ \(R_w(z)\). Multiplying by
\(\Theta_t\) gives one for the translated theta.

At a conjugate embedding transport \(h\) and the product formula
above. For the genuine translated analytic theta on the new curve
use the same transported \(h\) and the chosen lift \(w'\).
The ratio of the two constants \(c\) in these equations has
\(m\)-th power a sign. Thus the two germs for \(R_w\) have
quotient a root of unity of order dividing \(2m^2\). This follows
also at a zero of either germ by removing the common order of
vanishing first. Theorem~\ref{thm:germ57} gives exact transport
for \(\Theta_t\). This proves the asserted single phase for the
whole translated germ.

For the local assertion take a completed integer ring containing
the torsion and fields under consideration and the good Weierstrass
model at \(p\). Scale \(h\) locally so that its divisor on the
smooth proper model consists of its horizontal divisor only.
This is possible because the special fiber is geometrically
integral, so every vertical divisor is an integer multiple of
that fiber. Translation by each torsion section is an automorphism
of the model. The product of translates of the rescaled \(h\) is
a constant of valuation zero by its vertical divisor. Thus its
associated \(c\) is a unit by the product formula. The pole
section \(t\) misses the origin in the special fiber. The section
\(t-P\) also misses it unless \(P=t\), by injectivity of
prime-to-\(p\) specialization. It follows that in the completed
local ring at zero one has \(c h=T^{m e}b(T)\) with \(b\) an
integral unit and \(e=1\) if \(P=t\), \(e=0\) otherwise.
Here the case \(P=0\) simply has zero divisor. After adjoining a
chosen root of \(b(0)\), the equation for the root of \(b(T)\)
has a unique power series solution with that constant term by
Hensel's lemma, since \(m\) is a unit. In particular \(R_w\) is
integral in the formal parameter \(T\). The invariant vector
field dual to \(dx/y\) acts on this completed ring by an integral
unit times \(d/dT\). Every undivided derivative of \(R_w\) at
zero is therefore integral. Apply Leibniz's rule and the bounds
on \(H_n\) to \(R_w\Theta_t\) to finish.
\end{proof}

\begin{corollary}[Simultaneous finite translates on a product]\leavevmode\label{cor:prod57}
Let \(a\ge1\), and on \(C^a\) let \(\Theta=\prod_{j=1}^a\Theta_t(z_j)\).
Choose finitely many torsion vectors \(w_i\) killed by \(m\) and
algebraic weights \(b_i\). The derivatives of
\(F=\sum_i b_i U_{w_i}\Theta\) lie in one number field. At each
complex embedding the Taylor series is the Taylor series of
\(\sum_i \rho(b_i)\xi_{i,\rho}U_{w'_{i}}\Theta'\) for roots of
unity \(\xi_{i,\rho}\). Consequently
\[
       \|F^\rho\|_2^2
        \le \Big(\sum_i|\rho(b_i)|\Big)^2\nu_\rho^a ,
\]
where the norm uses a common period sublattice and normalized Haar
measure. At every \(p\ge5\) prime to \(m\) where all weights are
integral, \(v_p(\partial^{\boldsymbol n}F(0))\ge0\) for \(p\ne d\).
At \(p=d\) it is at least
\(-\tfrac12\sum_j\lfloor n_j/2\rfloor\) under the same hypotheses.
\end{corollary}
\begin{proof}
The Hermitian translation factors over the product. Apply
Theorem~\ref{thm:trans57} to every component and take the
compositum of the finitely many fields. The phases multiply.
The metric norm of each translated product over the common
sublattice is the product of the one-factor norms, since its
metric magnitude is a torus translate of the original product.
The triangle inequality proves the bound. At a finite place each
mixed derivative of a product is the product of its indicated
one-variable derivatives. The valuation bounds follow term by term.
\end{proof}

These results identify the analytic elliptic normalization and its
arithmetic derivatives for the actual function used in the source.
The torsion level matters in the local statement. At a prime dividing
that level the Hensel argument does not apply. For an iterated theta
lift the effects of passage between lattices and of the class and
Bessel sums must also be accounted for.

\subsection{Conjugate class sums for the normalized theta vacua}\label{nav:V-12}
\label{sec:arch59}

We apply Sections~\ref{sec:min55} and~\ref{sec:ell57} at all
embeddings of one coefficient field. Throughout \(C=\CC/(\delta
\OO_K)\) is the discriminant-one curve, \(t_*\) is the two-torsion
mark on which both idempotents act nontrivially, and \(E/K\) is
the field belonging to \(x\) in the common interval. Put
\(W=\operatorname{Mat}_{s\times r}(E)\), \(f=[E:K]\),
\(D=frs\), and write \(\Phi\) for the CM embedding of \(W\) into
the \(D\) matrix entries of Section~\ref{sec:min55}.
The trace Hermitian form on \(W\) is
\[
          h_K(X,Y)=\operatorname{Tr}_{E/K}\operatorname{tr}(Y^*X).
\]
For a pair of Hermitian lattice classes set
\(M_0=\mathfrak b^{-1}\operatorname{Hom}_{\OO_E}(L,L')\).
This is self-dual for \(h_K\) by the inverse relative different
and the perfect Hermitian pairing on \(\operatorname{Hom}(L,L')\).
In particular \(h_K(M_0,M_0)\subset\OO_K\).

\begin{lemma}[The theta construction in an orthonormal basis]\label{lem:basis59}
Suppose \(rs\) is even. The space \(W\) has a \(K\)-orthonormal
basis \(v_1,\ldots,v_D\). Let \(V\) have columns \(\Phi(v_i)\);
it is unitary with algebraic entries. In coordinates \(z=V^*Z\)
let \(\Gamma=\delta\OO_K^D\) and
\(\Gamma_M=\delta M\), where \(M\subset K^D\) is \(M_0\) written
in this basis. These lattices admit compatible theta multipliers.
On the product lattice the section with this multiplier is
\(\Theta(z)=\prod_j\Theta_{t_*}(z_j)\), up to scalar.
The normalized intertwiner from \(\Gamma\) to \(\Gamma_M\)
sends this section to a nonzero vacuum.
\end{lemma}
\begin{proof}
Diagonalize the positive form \(\operatorname{Tr}_{E/K}(u\bar v)\)
on \(E\) over \(K\) by Gram--Schmidt. The diagonal entries are
positive rational numbers, since they are fixed by conjugation
in \(K\). Each occurs \(rs\) times on \(W\). For a diagonal pair
\(a,a\) one can choose \(\lambda,\mu\in K\) with
\(|\lambda|^2+|\mu|^2=1/a\). Indeed Meyer's theorem that
every indefinite rational quadratic form in at least five
variables has a nontrivial rational zero applies to
\(X^2+dY^2+U^2+dT^2-S^2/a\). At such a zero \(S\ne0\) by
positivity of the first four terms. Division by \(S^2\) gives the
two required norms. The two columns
\((\lambda,\mu)\) and \((-\bar\mu,\bar\lambda)\) are orthonormal
for the diagonal pair. This proves existence of the basis.
For the even-by-even ranks of the column tower there is also a
direct construction in groups of four repeated entries. Write
\(1/a\) as a sum of four rational squares by Lagrange's
four-square theorem applied after multiplying by a denominator
squared. Left multiplication by the resulting rational quaternion
has transpose times itself equal to \(a^{-1}I_4\).
Finally \(V^*V=I_D\) by the trace formula for \(h_K\).

On either lattice define a sign at the point with coordinates
\(\delta b\) by \((-1)^{\sum_i b_i\bar b_i}\).
The exponent is an integer because both lattices before multiplication
by \(\delta\) are integral for the \(K\)-Hermitian form.
For \(h\in\OO_K\) the integers \(2\real h\) and
\(2\imag h/\sqrt d\) have the same parity, as follows on writing
\(h=u+v(1+\sqrt{-d})/2\). Expanding the exponent of the sign on
a sum therefore gives precisely the sign required by the unitary
translation law for the alternating form \(\imag H\). The two
signs coincide on the intersection of the lattices by their
definition. In one product coordinate \(t_*=[\delta/2]\), which
corresponds to \(1\in\OO_K/2\). The odd theta has sign \(-1\)
on the two periods \(\delta,\delta(1+\sqrt{-d})/2\). Translation
by \(\delta/2\) preserves the first sign and reverses the second,
by the commutator formula in Corollary~\ref{cor:move55}. The signs
thus obtained agree with \((-1)^{b\bar b}\) on these generators
since \((1+d)/4\) is even. The multiplication law extends agreement
to all periods. The product has the asserted multiplier.
The space of sections for the product of the degree-one bundles
is one-dimensional by the product formula for sections.

Both lattices are self-dual for the alternating form of the metric,
by Theorem~\ref{thm:lattice55} for \(\Gamma_M\) and directly in
each product coordinate for \(\Gamma\). Multiply each \(U_w\Theta\)
by the sign of \(w\in\Gamma_M\), sum over
\(\Gamma_M/(\Gamma\cap\Gamma_M)\) and divide by the square root
of its order. The summands are independent of the chosen lifts
after multiplying by the sign, by the multiplication law and
the compatibility on the intersection. Translation by any member
of \(\Gamma_M\) permutes the sum with the required sign.
Corollary~\ref{cor:move55} shows that its norm is that of \(\Theta\),
which is positive.
\end{proof}

This construction is the finite expression in
\srcloc{Formalism/SharpRestriction/GlobalSchurBounds/StandardLatticesGL3}{197}.
We allow successive normalized intertwiners through any finite list
of such trace lattices, and multiplication by algebraic scalars all
of whose complex conjugates have magnitude one. All statements below
apply to a single intertwiner as well.

\begin{theorem}[Metric and lattice after arithmetic conjugation]\label{thm:transport59}
Let \(G\) be a vacuum constructed in Lemma~\ref{lem:basis59},
including the allowed successive transports. Its Taylor coefficients
in \(Z\) belong to one number field. At an embedding \(\rho\)
the transported Taylor series represents an entire function
\(G^\rho\). Write \(\Lambda_\rho=\delta_\rho\mathfrak c\) for
the period lattice of the conjugate elliptic differential, where
\(\mathfrak c\) is an integral ideal of \(K\). In product
coordinates for \(\rho(V)\) the metric magnitude of \(G^\rho\)
descends to the lattice
\[
                    \Gamma_{M,\rho}=\delta_\rho\mathfrak c\,\rho(M).
\]
Here \(\rho(M)\) applies the restriction of \(\rho\) to \(K\)
coordinatewise. With the Hermitian metric of scale
\(A_\rho=\operatorname{area}(\Lambda_\rho)\), one has
\begin{align}
      \|G^\rho\|_2^2&=\nu_\rho^D, \notag\\
      \sup_Z |G^\rho(Z)|^2e^{-\pi |Z|^2/A_\rho}
          &\le [e\max(1,2rs\,\operatorname{rd}(E)/\pi)]^D\,\nu_\rho^D .
                                                       \label{eq:vac59}
\end{align}
Finite torsion translates of \(G\) conjugate to the corresponding
translates of \(G^\rho\) up to roots of unity.
\end{theorem}
\begin{proof}
Fix an integer \(n\) such that every period lattice in the
construction lies between \(n\Gamma\) and \(n^{-1}\Gamma\);
enlarge it to put all translation vectors under consideration
in \(n^{-1}\Gamma\). Work first in product coordinates.
For \(w\in n^{-1}\Gamma\) write \(R_w=U_w\Theta/\Theta\).
The proof of Theorem~\ref{thm:trans57} shows in each coordinate
that \(R_w(nz)\) is an algebraic rational function on the
elliptic curve. Thus the product is rational on \(C^D\) over a
number field. The calculation also covers a larger multiple \(n\)
of the torsion order by composition with integer multiplication.
Every finite expression in intertwiners consequently has the form
\[
                 G(Vz)=\Theta(z)\,S([z/n]),\qquad
                 S\in\overline{\QQ}(C^D).
\]
Here \([z/n]\) denotes the point of \(C^D\) represented by \(z/n\).
For example, for a translate of this expression its ratio to
\(\Theta\) is \(R_w(z) S([z/n]+[w/n])\), again rational on the
same cover. This proves the assertion by induction, including
the normalization factors. Finitely many rational functions
and torsion points suffice. Their fields together with the field
of Theorem~\ref{thm:germ57} and the entries of \(V\) lie in one
finite Galois extension \(J/\QQ\). One may include the full
\(n^2\)-torsion and the CM endomorphisms. In a finite sum,
products of the \(U_w\) may equally be combined into single
translations. Their commutator and composition phases are roots
of unity since the alternating pairing is rational on
\(\delta K^D\). Corollary~\ref{cor:prod57} then proves the entire
convergence of every conjugate series, even at a possible pole of
the displayed ratio. Entries of \(V\) lie in the abelian CM
field \(E\) as embedded in \(\CC\). Its automorphisms commute
with conjugation, so \(\rho(V)\) is unitary and conjugation
transports the linear change of variables.

More precisely, select \(w'\) so that \([w'/n]=\rho([w/n])\) on
the conjugate product curve. By the rational-function description,
conjugation sends \(R_w(z)\) to \(R'_{w'}(z)\) up to the constant
root of unity in Theorem~\ref{thm:trans57}. It carries translation
of \(S\) by \([w/n]\) to translation of \(S^\rho\) by \([w'/n]\).
This proves the last assertion of the theorem, and proves that a
unitary multiplier which is a root of unity at \(w\) transports
to a unitary multiplier at \(w'\). For any further translate one
may enlarge \(n\) and \(J\) for this assertion.

We identify all the transported periods. In the quotient
\(n^{-1}\Gamma/n\Gamma\) the map \(w\mapsto[w/n]\) identifies
the group with \(C[n^2]^D\). The scalar torsion
is a free rank-one \(\OO_K/n^2\OO_K\)-module. Transport by
\(\rho\) is semilinear for \(\rho|_K\) because the scalar CM
endomorphisms transport with their tangent multipliers. Hence
the generator \([\delta/n^2]\) is sent to
\([\delta_\rho a/n^2]\), where
\(a\in\mathfrak c\) generates
\(\mathfrak c/n^2\mathfrak c\) over \(\OO_K/n^2\OO_K\).
For clarity, the module representing \(\Gamma_M/n\Gamma\) in
\((\OO_K/n^2\OO_K)^D\) is \(nM/n^2\OO_K^D\).
Its image is
\[
       n\mathfrak c\,\rho(M)/n^2\mathfrak c^D
                       \ \subset\ (\mathfrak c/n^2\mathfrak c)^D .
\]
Indeed at each prime over \(n\), \(a\) generates
\(\mathfrak c\) locally by Nakayama's lemma. At every other
prime these quotient modules vanish. This proves the equality
and gives the displayed lattice in the theorem. The lattice
\(\delta_\rho\mathfrak c\) exists since the conjugate curve has
CM by the same maximal order, so its period lattice is homothetic
to an invertible ideal. On \(n\Lambda_\rho^D\) the periodicity of
the ratio and the theta law already give a unitary multiplier.
Hence \(\Gamma_{M,\rho}\) is a lattice of metric periodicity.

For the first normalized intertwiner its conjugated sum consists
of translates of the conjugate product section by distinct
elements modulo the conjugate product lattice. Its coefficients
have magnitude equal to the inverse square root of their number.
All translations have norm squared \(\nu_\rho^D\) on the common
period lattice, with Haar probability measure. They are orthogonal
by the commutator proof of Corollary~\ref{cor:move55}, giving
the first equality of \eqref{eq:vac59}. To iterate the argument
we verify self-duality and separation for each intermediate
lattice. For the trace lattice with the ideal \(\mathfrak c\)
inserted, \(A_\rho=|\delta_\rho|^2 N\mathfrak c\sqrt d/2\).
The trace pairings used in Theorem~\ref{thm:lattice55} are divided
by \(N\mathfrak c\), and
\(\mathfrak c\overline{\mathfrak c}=(N\mathfrak c)\) exactly
cancels this divisor. Thus the real alternating form is still
unimodular. For a nonzero matrix in
\(\mathfrak c\,\mathfrak b^{-1}\operatorname{Hom}(L,L')\), the
quantity \(\operatorname{tr}(X^*X)/N\mathfrak c\) is totally
positive in \(\mathfrak D_{E/K}^{-1}\). The same statements
hold in the transported coordinates, since the trace and all
ideal relations are preserved under field embeddings. Applying
the arithmetic--geometric mean argument of that theorem gives
the same squared separation \(2f/\operatorname{rd}(E)\).
Consequently every subsequent normalized sum again preserves
the norm. Indeed the preceding section now has unitary
multipliers on its transported self-dual lattice, and the
transported cosets are distinct modulo that lattice by the
torsion-group identification above. This is exactly the
orthogonality argument for the first intertwiner. The
mean-value proof of Theorem~\ref{thm:lattice55} applies with
this separation and covolume one, giving the supremum bound.
\end{proof}

\begin{theorem}[The field norm of the auxiliary class sum]\label{thm:arch59}
At each class pair \(c,h\) of Section~\ref{sec:min55}, take a
vacuum \(G_{ch}\) as just constructed and a finite section
\[
      F_{ch}=\sum_i a_{i,ch}U_{w_{i,ch}}G_{ch},\qquad
      \left(\sum_i|\rho(a_{i,ch})|\right)^2\le Q^{rs}
                         \quad\text{at every embedding }\rho .
\]
The \(w_{i,ch}\) are rational torsion vectors and the coefficients
are algebraic. Let \(J/\QQ\) be a finite Galois field containing
the coefficients of these data and their expansions, including
the fields used in Theorem~\ref{thm:transport59}. The number
\(\mathcal N\) of \eqref{eq:pair55} belongs to \(J\). If one
class coordinate is nonzero, set
\(\mathfrak h=[J:\QQ]^{-1}\log|N_{J/\QQ}\mathcal N|\).
Use the weights, \(M_{\rm in},M_{\rm out},k,D\)
from Proposition~\ref{prop:gauss55}. Then
\begin{align}
 \mathfrak h\le {}&
     \log(M_{\rm out}M_{\rm in}^2)+rs\log Q
              +D\log\!\big(2e\max(1,2rs\,\operatorname{rd}(E)/\pi)\big)
                      \notag\\
    &+\sum_{\tau}\log C_{r,b_\tau}
       +k\big(\log(\sqrt d/\pi)-\gamma+D_{\xi_d}\big) \notag\\
    &+D\big(\tfrac14\log d-\tfrac12\gamma+\tfrac12D_{\xi_d}\big).
                                                        \label{eq:field59}
\end{align}
The bound is independent of the indices of the intertwining
lattices, the choice of an orthonormal basis, and the auxiliary
coefficient field. It holds at every \(x\) on the common
interval with the same choice of \(d\).
\end{theorem}
\begin{proof}
Write \(v_{h,T}=\sum_c m_cP_T(\partial)F_{ch}(0)\).
These coordinates lie in \(J\) by Theorem~\ref{thm:transport59}.
The rational mass weights and the stability of \(J\) under complex
conjugation put \(\mathcal N=\sum_{h,T}m'_h v_{h,T}\bar v_{h,T}\)
in \(J\). At the original embedding it is positive under the
stated nonvanishing assumption. Define
\[
              B_\rho=\sum_{h,T}m'_h|\rho(v_{h,T})|^2 .
\]
This is positive at every embedding. At an arbitrary embedding
complex conjugation on \(J\) need not commute with \(\rho\).
Writing \(\iota\) for the original complex conjugation on \(J\),
Cauchy--Schwarz gives
\[
          |\rho(\mathcal N)|\le \sqrt{B_\rho B_{\rho\iota}}.
\]
The map \(\rho\mapsto\rho\iota\) permutes the embeddings. Hence
\(\mathfrak h\le[J:\QQ]^{-1}\sum_\rho\log B_\rho\).

The minors have integer coefficients, so \(B_\rho\) is the
sum of squares of exactly the prescribed derivatives of the
conjugated sections with exactly the given weights. Each
translated section has the same metric supremum as its vacuum.
Theorem~\ref{thm:transport59} makes the possible conjugation
phases immaterial in the triangle inequality for \(F_{ch}\).
The envelope in Proposition~\ref{prop:gauss55} at this embedding
can therefore be taken as
\[
 S_\rho=Q^{rs}[e\max(1,2rs\,\operatorname{rd}(E)/\pi)]^D\nu_\rho^D .
\]
Insert this into \eqref{eq:gauss55}, take logarithms, and average.
The exact mean of \(k\log(2\pi/A_\rho)+D\log\nu_\rho\) is
equation \eqref{eq:archmain58} of
Theorem \ref{thm:norm58}.
This proves \eqref{eq:field59}.
\end{proof}

The coefficient bound used here is the arithmetic finite
translation bound in
\srcloc{Development/AbelianVarietyTheory/Definitions/Foundations/NormSumPairRankZeroOfCurveCoordNorm}{114}.
Formula~\eqref{eq:field59} applies to the explicit vacuum and
class-section presentations in its hypotheses. In the iterated
automorphic lift, one must identify the Weil translates at the
class representatives with such presentations, prove the required
finite-place divisibilities, and identify a nonzero global Bessel
projection. The field norm estimate alone does not provide those
identifications or the odd-primary strict cohomology bound.

\subsection{Split representatives and the finite Weil action}\label{nav:V-13}
\label{sec:split61}

We now start with one finite vector in a theta lift and follow it
through the two sets of class representatives. The local calculation
can be performed at split primes outside the level of that vector.
Put \(k=E^+\) and let \(U_n\) be the unitary group of the standard
positive definite Hermitian space \(E^n\) over \(k\). All adelic
groups in this section use the finite adeles.

\begin{theorem}[Class representatives supported at split primes]\label{thm:class61}
Let \(J\subset U_n(\mathbb A_{k,f})\) be an open compact subgroup
and let \(S\) be any finite set of rational primes. Every class in
\[
                 U_n(k)\backslash U_n(\mathbb A_{k,f})/J
\]
has a representative whose components are the identity except
over at most two rational primes that split completely in \(E\)
and lie outside \(S\). For \(n\ge2\) one of the two primes can be
fixed for all the classes. For \(n=1\) one prime suffices.
\end{theorem}
\begin{proof}
Shrink \(J\) to a product of local subgroups which is full integral
at almost all places. A representative for a class at the finer
level also represents its image at the stated level. For \(T=U_1\)
the diagonal first entry splits the determinant \(U_n\to T\).
For \(g\in U_n(\mathbb A_{k,f})\) put \(t=\det g\). Write
\(t=z/\bar z\) for a finite idele \(z\) of \(E\) by local Hilbert
90. At almost every place the extension is split or unramified
and \(t\) is a unit. In the split case take the pair \((t_1,1)\).
In the unramified field case a Hilbert 90 solution can be
multiplied by a power of a uniformizer of \(k_v\) to have valuation
zero. These choices give an idele.

Choose a positive integer \(M\) such that for all integral unit
ideles \(u\) congruent to \(1\) modulo \(M\) at primes over \(M\)
we have \(\operatorname{diag}(u/\bar u,1,\ldots,1)\in J\).
Such an \(M\) exists by openness and the full level away from
finitely many primes. The ray class of \(z\) for \(M\OO_E\)
contains a prime \(\mathfrak P\) over a rational prime \(p\)
splitting completely in \(E\), outside any prescribed finite set.
Indeed let \(H\) be the ray class field and take its normal closure
over \(\QQ\). Extend the desired ray automorphism to an element of
the normal closure fixing \(E\). Chebotarev over \(\QQ\) supplies
a rational prime and a prime of the normal closure with precisely
this Frobenius. Its restriction to \(E\) is trivial and \(E/\QQ\)
is Galois. The Frobenius over \(E\) is therefore the same element,
and its restriction to \(H\) gives the chosen class by reciprocity.
We use the convention in which uniformizers give arithmetic
Frobenius. Thus \(z=a z_{\mathfrak P}u\) with \(a\in E^\times\),
\(z_{\mathfrak P}\) supported at \(\mathfrak P\), and \(u\) in
the indicated ray unit group. Removing \(\operatorname{diag}(a/\bar a)\)
on the left and \(\operatorname{diag}(u/\bar u)\) on the right
leaves a representative \(g_1\) with determinant supported over
\(p\). In these diagonal matrices all other entries on the
diagonal are one. This proves the result for \(n=1\).

For \(n\ge2\) fix another rational prime \(p_0\notin S\)
splitting completely in \(E\), excluding it when choosing \(p\).
Fix a place \(v_0\) of \(k\) above \(p_0\).
The group \(SU_n\) is simply connected and absolutely almost
simple over \(k\), since it becomes \(\mathrm{SL}_n\) over \(E\).
Its points over \(k_{v_0}\) are \(\mathrm{SL}_n(k_{v_0})\)
and are noncompact. The strong approximation theorem for simply
connected almost simple groups therefore makes \(SU_n(k)\) dense
in its finite adeles away from \(v_0\). Apply it to
\(g_1\operatorname{diag}(\det(g_1)^{-1},1,\ldots,1)\in
SU_n(\mathbb A_{k,f})\). At every place away from \(p,p_0\),
choose approximation modulo \(J_v\cap SU_n(k_v)\).
This specifies an open set in the restricted product since
the subgroup is full integral almost everywhere. The rational
element so obtained, acting inversely on the left of \(g_1\),
puts all these components in \(J_v\). Remove them on the right
using the product subgroup. Only the two split primes remain.
\end{proof}

\begin{lemma}[The split oscillator on a lattice vector]\label{lem:split61}
Let \(X\) be a finite-dimensional vector space over a nonarchimedean
local field, \(X^\vee\) its frequency space with a perfect additive
pairing, and \(U\subset X\) a compact open lattice. Its annihilator
is denoted by \(U^\circ\subset X^\vee\). In the position
realization take the unit vector
\(\varphi_U=\operatorname{vol}(U)^{-1/2}{\bf1}_U\).
For \(A\in\mathrm{GL}(X)\) the normalized change of variables
\[
       f(y)\longmapsto|\det A|^{-1/2}f(A^{-1}y)
\]
sends it to \(\varphi_{AU}\). In a unitary lattice realization of
the same Heisenberg representation this change on the vacuum is
the positive square-root normalized finite intertwiner from
\(U\oplus U^\circ\) to \(AU\oplus(AU)^\circ\), with trivial
vacuum characters in the position realization.
\end{lemma}
\begin{proof}
The first identity follows by transformation of Haar measure.
Position translations by \(U\) and frequency modulations by
\(U^\circ\) fix \(\varphi_U\). Their common invariant line consists
of functions supported on \(U\) and constant there. The first
statement about support follows from the fact that \(U\) is the
annihilator of \(U^\circ\), using all its modulations. The same
description holds for \(AU\). Let
\(b=[U:U\cap AU]\) and \(c=[AU:U\cap AU]\).
The target period lattice modulo its intersection with the
source period lattice has order \(bc\). Indeed the index in
the frequency component is \(b\) by annihilator duality.
The normalized finite sum over these \(bc\) cosets is on the
target line, has norm one, and has inner product \(1/\sqrt{bc}\)
with the source vacuum. For the inner product only the zero coset
contributes by commutator orthogonality, exactly as in
Corollary~\ref{cor:move55}. Direct integration also gives
\[
       \langle\varphi_U,\varphi_{AU}\rangle
        =\frac{\operatorname{vol}(U\cap AU)}
                {\sqrt{\operatorname{vol}(U)\operatorname{vol}(AU)}}
        =(bc)^{-1/2}.
\]
The nonzero positive inner product on this line fixes the scalar.
This proves the asserted compatibility of normalizations. A unitary
splitting character multiplying the change of variables multiplies
the answer by its own scalar.
\end{proof}

We specify the global finite Heisenberg realization used in applying
the lemma. Identify the holomorphic vacuum \(G_0\) on the standard
matrix lattice with the finite vacuum. Here all vector equalities
are transported with this fixed choice of scale. For
\(a\in W=\operatorname{Mat}_{s\times r}(E)\) write
\(q(a)=h_K(a,a)\in\QQ\) and
\[
       W(a)=e^{\pi i q(a)}U_{\delta\Phi(a)} .
\]
These operators fix \(G_0\) on the integral period module
\(M=\mathfrak b^{-1}\operatorname{Mat}_{s\times r}(\OO_E)\),
by the sign computed in Lemma~\ref{lem:basis59}. For
\(z\in K\) define the rational functional
\(\lambda(z)=\real z+\imag z/\sqrt d\) in the chosen embedding.
Multiplying the exponentials in the definition gives
\[
 W(a)W(b)=e^{-2\pi i\lambda(h_K(a,b))}\,W(a+b).
\]
Moreover \(\lambda(\OO_K)\subset\ZZ\). Extend this cocycle to
finite adeles by the product of the local additive characters
which on rational arguments gives \(e^{-2\pi i u}\). It is
trivial on \(\widehat M\times\widehat M\). This describes the
usual lattice realization without choosing complex
representatives for finite adelic vectors. Indeed
\[
             W/M \ \simeq\ (W\otimes_\QQ\mathbb A_{\QQ,f})/\widehat M .
\]
To check the isomorphism, choose a \(\ZZ\)-basis of \(M\).
After clearing a common denominator surjectivity is the Chinese
remainder theorem, and a rational number integral everywhere
is an integer. For one rational lift of each coset the vectors
\(W(a)G_0\) are orthogonal of the same norm by the commutator
argument. Changing the lift multiplies by the cocycle with a
period, a root of unity. On this basis adelic translations act
by changing cosets with precisely the cocycle just displayed.
For instance if \(y=a+u\), \(u\in\widehat M\), the action on
the vacuum is the inverse of the cocycle of \((a,u)\) times
\(W(a)G_0\). Independence of the lift follows from the cocycle
identity and its triviality on periods. The same construction
works with the period module of any class vacuum.

The description factors into local descriptions since the cosets
are a restricted product of local cosets and the cocycle factors.
At a prime completely split in \(E\) outside \(2\operatorname{disc}(E)\),
the two conjugate components of the matrix space are a position
space \(X\) and its frequency space \(X^\vee\). Each is totally
isotropic. The standard matrix module at this prime is
\(U\oplus U^\circ\), by self-duality and the splitting.
In particular the local vacuum above is identified with
\(\varphi_U\). For detail, translate \(\varphi_U\) by the local
position and frequency vectors. Position cosets give disjoint
supports, and distinct frequency characters on a given coset
are orthogonal. Fourier expansion on the compact group \(U\)
makes their span all locally constant compactly supported
functions. Taking the phase convention of \(W(a)\) amounts to
multiplying the product of a position and a frequency operator by
the scalar determined by the displayed cocycle. The scalar is one
for integral components and thus fixes the same vacuum. These
formulas give the isometry on the coset bases and identify the
operators. The split pair action is
\(A:y\mapsto g' y g^{-1}\) on the position component and the
contragredient action on its frequency component. It preserves
\(h_K\), hence preserves the chosen phase convention.
This is the change-of-variables action of the split oscillator,
with the prescribed splitting character if one is used.
See also the position formulas in \cite{XIII-X-IV-gan,XIII-X-IV-kudla29}.

\begin{theorem}[The single-step class sections from one finite vector]\label{thm:one61}
Use \(rs\) even, the standard normalized vacuum \(G_0\), and
begin with a single vector of the form
\[
       F=\sum_i a_i W(b_i)G_0,\qquad Q b_i\in M,\qquad
                  \left(\sum_i|\rho(a_i)|\right)^2\le Q^{rs}
\]
at all embeddings. Take the usual finite oscillator action of
the unitary pair, with normalized change of variables at split
places. One may multiply that action there by unitary algebraic
determinant characters of the norm-one torus, or by the unitary
Hecke factors on finite ideles of \(E\) with integral circle
weights in the source splitting. At both levels choose representatives by
Theorem~\ref{thm:class61}, avoiding the primes dividing
\(2Q\operatorname{disc}(E)\). For every such pair of representatives
the finite Weil translate of \(F\) has a presentation as in
Theorem~\ref{thm:arch59}, with the same bound \(Q^{rs}\).
Its translation vectors can be taken to have order dividing
\(Q\) relative to the new period lattice. Consequently the
field norm estimates of Section~\ref{sec:missing} of
Theorem \ref{thm:opt60} and Corollary \ref{cor:mass60} apply to the nonzero auxiliary class sum
formed from these representatives.
\end{theorem}
\begin{proof}
Let \(g,g'\) be the chosen pair and write \(B(y)=g'y g^{-1}\)
on the adelic matrix space. The old adelic period module
\(\widehat M\) moves to \(B\widehat M\). It is the completion
of the rational period module
\(M'=\mathfrak b^{-1}\operatorname{Hom}_{\OO_E}(L,L')\) for
the two moved lattices, by localizing the definition of
\(\operatorname{Hom}\). At the moving places Lemma~\ref{lem:split61}
identifies the new vacuum with the normalized lattice intertwiner
up to the indicated character factors. At all other places the
action is the identity. The finite intertwining sum globally is
the product of these local sums. Indeed its coset quotient is
the product of the local quotients, its square-root normalization
is the product of their positive square-root normalizations, and
in the \(W\)-convention each summand is translation by a new
period. The old and new vacua have trivial corrected-period
characters. For rational periods this last statement follows
from \(e^{\pi i q(a)}=(-1)^{q(a)}\) and extends to their finite
completions in the coset model. Hence this is exactly a vacuum
allowed by Theorem~\ref{thm:transport59}.

The permitted character factors have magnitude one under every
complex embedding. To check this, on the finite norm-one ideles
the character has an open unit kernel and on the archimedean
circles it has integral weights. The norm-one torus class quotient
at this kernel is finite, for instance by the ray-class argument
in the first proof. Some positive power of a given finite
idele therefore equals a rational norm-one element times a member
of the kernel. The corresponding power of its character value is
the inverse archimedean character of that rational element by
automorphy. Each conjugate of an element \(u\in E^\times\)
with \(u\bar u=1\) has absolute value one since \(E\) is CM.
Its products to integral powers have the same property. Taking
roots proves both algebraicity and the required magnitudes.
For a Hecke character on the ideles of \(E\) as specified in
\srcloc{Formalism/SharpRestriction/GlobalSchurBounds/StandardLatticesGL3}{409},
use the finite ray class group of \(E\) instead. A finite idele
raised to some positive power is a principal idele \(a\) times
a ray unit in the kernel. The square of its archimedean character
at \(a\) is a product of integral powers of
\(\tau(a)/\overline{\tau(a)}\), by the integral circle weights.
Again all conjugates have absolute value one. This applies in
particular to the determinant lifts at split places in lines
463--470 of the same source.

Conjugating \(W(b_i)\) by the split change of variables translates
by \(B b_i\) in the new model. This follows on position and frequency
operators separately. The correcting phase is preserved because
the pair preserves \(h_K\). As \(Q B b_i\in B\widehat M\), the
adelic coset of \(B b_i\) has order dividing \(Q\) modulo
\(\widehat {M'}\). Choose its rational representative \(b_i'\)
by the quotient isomorphism above, now with lattice \(M'\).
Then \(Q b_i'\in M'\). Acting on the moved vacuum changes this
adelic translation to \(W(b_i')\) times a root of unity, since
the additive cocycle at finite places has finite-order values.
Finally \(W(b_i')\) and the Hermitian translation \(U_{\delta\Phi(b_i')}\)
differ by \(e^{\pi i q(b_i')}\), also a root of unity. Absorb
these scalars and the common character factor in \(a_i\). All
the absolute coefficient sums over embeddings are preserved.
This gives exactly the hypotheses of Theorem~\ref{thm:arch59}
at every pair of representatives and proves the consequence.
\end{proof}

The theorem follows a prescribed single-step vector through the
finite class sums. For the iterated lift the successive inputs,
their finite integral divisibility, and the nonzero Bessel period
extracted at the final rank still require a comparison on the
same interval.

\subsection{An inert-prime congruence for the prescribed class sums}\label{nav:V-14}
\label{sec:period63}

Fix \(x\) on the common interval, \(r\le s\) with \(rs\) even,
and the single-step construction of Theorem~\ref{thm:one61}.
Keep the same \(\delta\), \(K\), \(E\), \(M\), and standard
elliptic mark \(t_*\). Let \(S\) be any finite set of primes
\(p\ge5\) inert in \(K\), unramified in \(E\), and prime to \(Q\).
We choose one vacuum for the entire set. The archimedean estimate
of Theorem \ref{thm:opt60}
applies to this choice. At each \(p\in S\) and each coefficient
place above it use the integer ring \(R\) of a sufficiently large
finite extension. Congruences modulo \(pR\) are normalized by
\(v_p(p)=1\).

\begin{lemma}[One orthonormal basis for the specified primes]\leavevmode\label{lem:basis63}
One can choose the \(K\)-orthonormal basis of Lemma~\ref{lem:basis59}
so that it is an \(\mathcal O_K\otimes\ZZ_p\)-basis of
\(M\otimes\ZZ_p\) simultaneously at every \(p\in S\).
At all places above these primes its CM embedding matrix \(V\)
and its inverse have integral entries.
\end{lemma}
\begin{proof}
At \(p\in S\), \(K_p/\QQ_p\) is the unramified quadratic
extension. The trace Hermitian lattice \(M_p\) is self-dual.
Its reduction modulo \(p\) is therefore a nondegenerate Hermitian
space over \(\mathbb F_{p^2}\). It has a vector of nonzero norm.
Indeed if all norms vanished, polarization of \(h(v+aw,v+aw)\)
for every \(a\in\mathbb F_{p^2}\) would give
\(\operatorname{Tr}_{\mathbb F_{p^2}/\mathbb F_p}(a h(w,v))=0\)
for every \(a,v,w\), contrary to nondegeneracy. A lift has unit
norm. The norm on units of \(\mathcal O_{K_p}\) is onto
\(\ZZ_p^\times\). It is onto on residue fields by cyclicity,
and a unit in \(1+p\ZZ_p\) has a square root in \(\ZZ_p\) by
Hensel's lemma, which is also a norm from \(K_p\).
Rescale the lift to length one, split off its orthogonal complement
integrally by its pairing, and repeat. This gives a local
orthonormal basis at every selected prime.

Relative to any initial rational orthonormal basis, the local bases
are matrices in the unitary group of the identity form over \(K_p\).
The condition that such a matrix give a basis of \(M_p\) is open.
The Cayley transform \(X\mapsto(1-X)(1+X)^{-1}\) parametrizes
the unitary matrices away from eigenvalue \(-1\) by an open part
of the rational vector space of skew-Hermitian matrices. Its
inverse has the same formula. To reach this open part within
each local condition, multiply the selected matrix by a central
norm-one unit arbitrarily close to one, avoiding the finitely
many forbidden eigenvalues. The values
\((1+c\sqrt{-d})/(1-c\sqrt{-d})\) with \(c\in\QQ_p\) close to
zero supply such units. Approximate the inverse Cayley parameters
at all \(p\in S\) by a single rational point of the vector space,
using weak approximation on its rational coordinates. Its
transform exists and lies in all the specified open sets.
For the empty set simply use the initial basis.

The primes in \(S\) do not divide the relative different, so
\(\mathfrak b\) is a unit there. Thus every entry of every \(v_i\)
is integral at all places of \(E\) above \(S\). The same is
true for its conjugates, and \(V^{-1}=V^*\) gives the assertion
for the inverse.
\end{proof}

To construct the vacuum of Lemma~\ref{lem:basis59}, start in the
basis of Lemma~\ref{lem:basis63} with the product of the marked
elliptic sections, then use the normalized intertwiner
to the standard trace lattice, with scalar one. For the input
weights impose the source condition that \(Q^{rs}a_i\) is an
algebraic integer. Choose the split representatives of
Theorem~\ref{thm:one61}. A matrix-coordinate index
\(\nu=(\tau,i,j)\) specifies a CM embedding, row, and column.
Let \(c_E\) denote complex conjugation on \(E\) and \(\sigma_p\)
the arithmetic Frobenius of \(E/\QQ\) at \(p\). Its value is
independent of the place because \(E/\QQ\) is abelian. Define
\[
    \pi_p(\tau,i,j)=(\sigma_p c_E\circ\tau,i,j).
\]
Here each image of \(E\) in \(\CC\) is identified with the same
normal field. The new embedding again belongs to the CM type
since both \(\sigma_p\) and \(c_E\) induce conjugation on \(K\).
Write \(\kappa=-(\partial^{p-1}x)(t_*)\) for the elliptic unit
of Lemma \ref{lem:kappa62}, with \(x\) the Weierstrass function.

\begin{theorem}[Reduction of the actual class-section derivatives]\label{thm:jet63}
For each \(F_{ch}\) obtained from the single input vector and
the indicated representatives, all undivided derivatives in the
matrix coordinates are integral at every place over \(S\).
At \(p\in S\) set \(H_\alpha=\partial^\alpha F_{ch}(0)\).
Writing \(e_\nu\) for the unit multi-index, one has
\begin{equation}
 H_{\alpha+p e_\nu}
      \equiv\kappa\,\alpha_{\pi_p(\nu)}
                H_{\alpha-e_{\pi_p(\nu)}}\pmod{pR}.
                                                     \label{eq:class63}
\end{equation}
The right side is zero if \(\alpha_{\pi_p(\nu)}=0\). Also
\(H_{\alpha+p^2 e_\nu}\equiv0\bmod pR\) for every \(\alpha,\nu\).
The same assertions hold for all classes using the single basis
of the lemma.
\end{theorem}
\begin{proof}
At each prime in \(S\), the product period lattice and the
standard trace lattice agree in the coordinates of the lemma.
The class representatives move lattices only at split primes,
so all the moved trace lattices also agree with them locally at
\(p\). Every index in the two normalized intertwiners is
therefore prime to \(p\), as can be checked on the finite quotient
of each target lattice by its intersection with the source.
Their square-root normalization factors are units. In producing
the finite sums of translates of the original elliptic product,
each translation has order prime to \(p\) in the product torus.
Indeed all representatives of the intertwining cosets belong
locally to the product lattice, and the further vectors have
order dividing \(Q\) in the new period lattice. The two lattices
again agree locally at \(p\). Composing translations only introduces
roots of unity. The weights \(a_i\) are integral here by the
assumption on \(Q^{rs}a_i\).

The determinant factors in Theorem~\ref{thm:one61} are units at
all places above \(S\). In its ray class proof take the finite
idele lifts to be supported at the moving split primes. Whenever
a power of a lift equals a principal idele times units, the
principal generator is a unit at all primes above \(S\).
The formula for a power of the character value as a product of
conjugates and their inverses then proves the assertion, also for
the norm-one torus. We have thus expressed \(F_{ch}\), in product
coordinates, as a sum of products of prime-to-\(p\) translates of
\(\Theta_{t_*}\) with integral coefficients. All their undivided
derivatives are integral by Corollary~\ref{cor:prod57}.

In the ring of undivided derivative sequences use the
multivariable product rule, with a binomial coefficient in each
coordinate. For a product coordinate \(z_j\) denote its shift
derivation by \(D_j\). Theorem \ref{thm:hurwitz62}, gives
\(D_j^p F=\kappa z_j F\) in this ring over \(R/pR\) for every
one of the translated products, hence also for the sum.
The matrix coordinates are related by
\(z_j=\sum_\mu\overline V_{\mu j}Z_\mu\).
Since \(V^*\) is integral, the chain rule preserves integrality
of all undivided derivatives. More explicitly, each derivative
in a matrix direction is a linear combination of the \(D_j\)
with integral coefficients, and the same holds on iterating.
Coefficientwise substitution of the linear forms respects the
sequence product by the chain and product rules.

The commuting derivations give
\(\partial_\nu^p=\sum_j\overline V_{\nu j}^{\,p}D_j^p\)
after this substitution modulo \(pR\). Since \(E/\QQ\) is
unramified at \(p\), every \(p\)-integral entry satisfies
\(\overline V_{\nu j}^{\,p}\equiv V_{\pi_p(\nu),j}\bmod pR\).
This congruence first holds in the unramified integer ring at
each place of \(E\) and remains valid after extension to \(R\).
Consequently
\[
 \partial_\nu^p F_{ch}
 \equiv\kappa\sum_{\mu,j}V_{\pi_p(\nu),j}
                         \overline V_{\mu j}\,Z_\mu F_{ch}
        =\kappa Z_{\pi_p(\nu)}F_{ch}\pmod{pR},
\]
where the last equality is \(VV^*=1\). Taking the
\(\alpha\)-th entry proves \eqref{eq:class63}.
The derivation \(\partial_\nu^p\) kills the linear coordinate
\(Z_{\pi_p(\nu)}\). Its \(p\)-fold iteration therefore
multiplies \(F_{ch}\) by \(\kappa^pZ_{\pi_p(\nu)}^p=0\) in
the sequence ring. This gives the remaining congruence.
\end{proof}

\begin{corollary}[Two divisibility tests on the prescribed minor weights]\label{cor:min63}
Fix \(p\in S\). Suppose that at least one embedding \(\tau\)
satisfies one of the following two conditions.
\begin{align*}
    b_\tau&>s(p^2-1),\\
    \pi_p(\tau)&\ne\tau,\qquad
    b_\tau>p b_{\pi_p(\tau)}+s(p-1).
\end{align*}
Here \(\pi_p\) on embeddings omits the row and column, which it
fixes. Then every prescribed \(P_T(\partial)F_{ch}(0)\) belongs
to \(pR\), at every place over \(p\).
\end{corollary}
\begin{proof}
Each monomial of the determinant product has exponent sum \(b_\tau\)
in each column at \(\tau\). Under the first condition some entry
in a column has exponent at least \(p^2\), and Theorem~\ref{thm:jet63}
applies. For the second condition, fix a column and denote the
row exponents at \(\tau\) by \(a_i\), and those at \(\pi_p(\tau)\)
by \(a_i'\). Apply \eqref{eq:class63} in row \(i\) at \(\tau\)
\(\lfloor a_i/p\rfloor\) times. Each application consumes one
exponent from the other embedding in that same row, or gives zero.
The two embeddings are distinct so reducing the second set of
exponents leaves the first set untouched except for the prescribed
subtractions of \(p\). Nonzero reduction would require
\(\sum_i\lfloor a_i/p\rfloor\le\sum_i a_i'=b_{\pi_p(\tau)}\).
But \(\sum_i a_i\le p\sum_i\lfloor a_i/p\rfloor+s(p-1)\).
The assumed strict inequality contradicts this requirement.
Thus each monomial gives zero modulo \(pR\), and the determinant
coefficients are integers. This proves the claim.
\end{proof}

\begin{theorem}[Finite divisibility and the archimedean estimate for one class sum]\label{thm:joint63}
Let \(S_*\subset S\) be the primes passing either test of
Corollary~\ref{cor:min63}. Put \(e=[E:\QQ]\),
\(A_r=|\mathrm{GL}_{re}(\mathbb F_3)|\),
\(A_s=|\mathrm{GL}_{se}(\mathbb F_3)|\), and
\[
        Y=A_sA_r^2\,\mathcal N,\qquad
        \mathcal N=\sum_{h,T}m'_h
           \left|\sum_c m_c P_T(\partial)F_{ch}(0)\right|^2 .
\]
At every place above every \(p\in S_*\), \(v_p(Y)\ge2\).
If a class coordinate is nonzero, define
\[
  B_-(Y)=\frac1{[J:\QQ]}
    \sum_{p\notin S_*}\log p\sum_{v\mid p}[J_v:\QQ_p]
                                            \max(0,-v_p(Y)).
\]
Only finitely many primes occur in this denominator term. With the
same coefficients, class representatives and number field \(J\)
in both bounds,
\begin{equation}
 2\sum_{p\in S_*}\log p-B_-(Y)
     \le \log(A_s A_r^2)+\frac{\log|N_{J/\QQ}\mathcal N|}{[J:\QQ]} .
                                                      \label{eq:joint63}
\end{equation}
The second term on the right has the upper bound of Theorem \ref{thm:opt60}, on the same \(E=E_x\) and power interval.
\end{theorem}
\begin{proof}
Every source stabilizer acts faithfully on its integral lattice
of \(\ZZ\)-rank \(re\), and the target stabilizers act on lattices
of rank \(se\). They are finite by definiteness. Reduction modulo
three is injective on every finite subgroup of
\(\mathrm{GL}_n(\ZZ)\). Indeed a nontrivial element in the
congruence kernel of prime order would have the form \(1+3^jB\)
with some unit entry of \(B\) and \(j\ge1\). Raising to a prime
other than three preserves this first valuation, and raising
to three increases it by exactly one, by the binomial expansion.
Thus the reciprocal stabilizer orders have common denominators
\(A_r\) and \(A_s\) respectively.

Write \(v_{h,T}\) for the inner class sum. By
Corollary~\ref{cor:min63}, \(A_r v_{h,T}\) lies in \(pR\)
for every place over any \(p\in S_*\). Its complex conjugate
satisfies the same condition since conjugation permutes the
places above \(p\) in the common Galois field. As \(A_sm'_h\)
is integral, every term of \(Y\) has valuation at least two
there. When one class coordinate is nonzero, \(Y\ne0\) since
the original archimedean sum has positive weights.
For a nonzero element of a number field the logarithm of the
absolute rational norm, divided by the degree, is the sum of
\(v_p(Y)[J_v:\QQ_p]\log p/[J:\QQ]\) over all finite places.
It has finite support. Retain the positive lower bounds above
\(S_*\) and the negative contributions elsewhere. This gives
\eqref{eq:joint63}. All the presentations used above have the
archimedean coefficient bounds in Theorem~\ref{thm:one61}, which
proves the final assertion.
\end{proof}

The finite-place bound here records one factor of \(p\) in every
class coordinate at primes satisfying either test. The weighted
prime count in Volume~I requires higher valuations, control of
the denominator term, and a nonzero Bessel projection of the
iterated lifts. These requirements and the odd-primary strict
cohomology estimate remain open.

\subsection{Higher powers in the prescribed class sums}\label{nav:V-15}
\label{sec:powers65}

Fix the fields \(K,E=E_x\), the ranks \(r,s\), the weights
\(b_\tau\), the level \(Q\) and the interval of
Section~\ref{sec:period63}. Write \(f=[E:K]\) and \(D=frs\).
Choose any one finite set \(S\) as there and the one basis
of Lemma~\ref{lem:basis63} for all its primes. We keep the same
input, class representatives, and archimedean normalization as in
Theorem~\ref{thm:joint63}. At each \(p\in S\) use
\[
   F_p(n)=v_p(n!),\qquad \omega_p=p/(p^2-1),\qquad
   \alpha!=\prod_{\nu=1}^{D}\alpha_\nu!,\qquad
   |\alpha|=\sum_{\nu=1}^{D}\alpha_\nu .
\]

\begin{theorem}[A bound for every matrix derivative]\label{thm:matrix65}
For every pair of classes \(c,h\), every place over \(p\in S\),
and every multi-index \(\alpha\), one has
\begin{equation}
 v_p\big(\partial^\alpha F_{ch}(0)\big)\ \ge\ 
 \max\left\{0,\sum_\nu F_p(\alpha_\nu)
                         -\lfloor\omega_p|\alpha|\rfloor\right\}.
                                                        \label{eq:matrix65}
\end{equation}
\end{theorem}
\begin{proof}
For a power series in any finite number of variables over the
local coefficient field, impose the condition that its ordinary
coefficient of total degree \(n\) at each monomial belongs to
\(p^{-\lfloor\omega_p n\rfloor}R\). Sums with integral
coefficients and products preserve this condition. For a product,
the two degrees add and the sum of their floors is at most the
floor of their sum. An integral linear change of variables also
preserves it. Indeed the expansion of a monomial of degree \(n\)
has integral coefficients in that degree.

In product coordinates every local elliptic factor
\(U_w\Theta_{t_*}\) satisfies the condition by the ordinary
coefficient estimate proved in
Theorem \ref{thm:powers64}.
The expression for \(F_{ch}\) in the proof of Theorem~\ref{thm:jet63}
is a finite sum with integral coefficients of products of these
factors. All its translates have order prime to the selected \(p\).
To recover matrix coordinates one substitutes the linear forms
given by the integral matrix \(V^*\). Thus the stated condition
holds for \(F_{ch}\) in those coordinates. Multiplication by
\(\alpha!\) takes its coefficient at \(Z^\alpha\) to
\(\partial^\alpha F_{ch}(0)\). Theorem~\ref{thm:jet63} supplies
integrality of the latter, also when the second bound is negative.
\end{proof}

For a determinant term, the available information in a column is
its total degree. The cheapest possible factorial valuation for
that column has an exact expression.

\begin{lemma}[Distribution of a column degree]\label{lem:digit65}
For integers \(b\ge0\), \(s\ge1\), define
\[
    M_p(b,s)=\sum_{j\ge1}\max\!\left(0,
                  \left\lceil\frac{b-s(p^j-1)}{p^j}\right\rceil
                                 \right).
\]
The sum is finite. It equals the minimum of
\(\sum_{i=1}^{s}F_p(n_i)\) over all tuples of nonnegative
integers with \(\sum_i n_i=b\).
\end{lemma}
\begin{proof}
Write each \(n_i\) as a multiple of \(p^j\) plus its remainder.
The sum of the remainders is at most \(s(p^j-1)\), so each sum
of quotients is at least the indicated summand. Summing gives
the lower bound. To attain the bound at all levels together, start
with \(b_0=b\) and set
\[
 b_{j+1}=\max\{0,\lceil(b_j-s(p-1))/p\rceil\},\qquad
 d_j=b_j-pb_{j+1}.
\]
Here \(0\le d_j\le s(p-1)\). In fact, if \(b_j>s(p-1)\),
then \(s(p-1)-p<d_j\le s(p-1)\); in the other case
\(d_j=b_j\). Whenever \(b_j>0\) the next quotient is smaller,
so this process terminates. Independently at each position \(j\),
distribute \(d_j\) as a sum of \(s\) digits between zero and
\(p-1\), by filling successive slots up to \(p-1\).
Use these digits as the base-\(p\) digits of \(n_1,\ldots,n_s\).
They then sum to \(b\), and the sum of their quotients on division
by \(p^j\) is exactly \(b_j\). To evaluate the recursion, the
integer \(b_j\) is the least nonnegative integer \(q\) such that
\(b\le p^j q+s(p^j-1)\). This holds at \(j=0\).
At the next step the requirement is \(b_j\le p q+s(p-1)\),
which is equivalent to
\(b\le p^{j+1}q+s(p^{j+1}-1)\) by the induction hypothesis.
This gives precisely the ceilings in the assertion.
\end{proof}

\begin{theorem}[Higher finite valuations on the common interval]\label{thm:joint65}
Let \(k=r\sum_\tau b_\tau\), \(B=\max_\tau b_\tau\), and let
\(\delta_p\) be one if either condition in
Corollary~\ref{cor:min63} holds and zero otherwise. Put
\begin{equation}
 m_p=\max\left\{\delta_p,\ r\sum_\tau M_p(b_\tau,s)
                                -\lfloor\omega_p k\rfloor\right\}.
                                                     \label{eq:mixed65}
\end{equation}
At every place above every \(p\in S\), the integer \(p^{m_p}\)
divides every local scalar \(P_T(\partial)F_{ch}(0)\) in \(R\)
and every cleared coordinate \(A_r v_{h,T}\). With the normalization
of Theorem~\ref{thm:joint63}, \(v_p(Y)\ge 2m_p\) there.
If at least one coordinate is nonzero, then, with
\(B_-^{(S)}(Y)\) the normalized sum of the negative valuations
outside \(S\) as in that theorem,
\begin{equation}
 2\sum_{p\in S}m_p\log p-B_-^{(S)}(Y)
    \ \le\ \log(A_s A_r^2)+
                     \frac{\log|N_{J/\QQ}\mathcal N|}{[J:\QQ]} .
                                                       \label{eq:joint65}
\end{equation}
For fixed \(p\), the explicit exponent also satisfies
\[
       m_p\ \ge\ \frac{k}{p^2-1}
                        -D\big(1+\log_p(1+B)\big).
\]
\end{theorem}
\begin{proof}
Every monomial of \(P_T\) uses one matrix entry per column per
minor. At each embedding and in each of the \(r\) columns the
\(s\) exponents therefore sum to \(b_\tau\). Apply
Lemma~\ref{lem:digit65} to each column and insert the result in
Theorem~\ref{thm:matrix65}. Every monomial has degree \(k\) and
every determinant coefficient is integral. The other bound
\(\delta_p\) is Corollary~\ref{cor:min63}. Multiplication of a
class sum by \(A_r\) clears all its stabilizer weights.
The bounds hold at every place above \(p\), and hence also for
the conjugate class coordinates. Clearing the outer weights by
\(A_s\) proves the bound on
\(Y=\sum_{h,T} A_s m'_h |A_r v_{h,T}|^2\). At the distinguished
embedding the sum is positive when some coordinate survives.
Apply the rational norm to this nonzero \(Y\) and take logarithms,
as in Theorem~\ref{thm:joint63}, to obtain \eqref{eq:joint65}.

Finally \(F_p(n)=(n-s_p(n))/(p-1)\). For \(0\le n\le B\) the
digit sum is at most \((p-1)(1+\log_p(1+B))\). Every admissible
column distribution thus has factorial valuation at least
\(b_\tau/(p-1)-s(1+\log_p(1+B))\).
Sum these inequalities and use
\(1/(p-1)-\omega_p=1/(p^2-1)\).
\end{proof}

The right side of \eqref{eq:joint65} has the upper bound in
Theorem \ref{thm:opt60} and Corollary \ref{cor:mass60}, with the same \(Y\), \(E_x\), \(Q\) and classes.
The bound on \(m_p\) is uniform for that choice of data. Its
leading coefficient \(1/(p^2-1)\) has a convergent sum against
\(\log p\) over the primes, by comparison with
\(\sum_{n\ge2}\log n/(n^2-1)\). This part of the estimate alone
therefore supplies only a bounded multiple of \(k\) in the
prime sum. The factorial terms at infinity can grow as
\(r\sum_\tau b_\tau\log b_\tau\). Recovering the weights of the
distinguished Frobenius classes, controlling the remaining
denominator and passing to the nonzero iterated Bessel period
are still necessary for the proposed proof.

\subsection{The elliptic normalization on the common interval}\label{nav:V-20}
Keep \(d\), \(K=\QQ(\sqrt{-d})\), the CM field \(E\) over \(K\),
and the interval \(I\) of Section~\ref{sec:missing}. Take
\(\Lambda=\delta\mathcal O_K\) with
\(g_2(\Lambda)^3-27g_3(\Lambda)^2=1\). Such a \(\delta\) exists
by the scaling law of weight twelve for the lattice discriminant.
Write \(C=\CC/\Lambda\), with differential \(dz\) and model
\[
                y^2=4x^3-g_2x-g_3,\qquad dz=dx/y.
\]
Here \(j=1728g_2^3\). The classical theorem of complex
multiplication for the maximal order says that \(j\) is an
algebraic integer and its distinct conjugates over \(\QQ\) are the
\(j\)-invariants of the ideal classes of \(K\), each occurring once
in its minimal polynomial. In particular \(g_2,g_3\) are algebraic.
We use the theorem precisely in this form. For the modular
identities below we use the usual normalization \(q=e^{2\pi i\tau}\).
\href{https://www.jmilne.org/math/CourseNotes/CM.pdf}{Milne, Complex Multiplication}.

Let \(J\) be any number field containing \(g_2,g_3\) and the
coordinates of a chosen nonzero two-torsion point \(t\).
At each embedding \(\rho:J\hookrightarrow\CC\) take the period
lattice of \(dx/y\) on the conjugate model. Write it as
\(\Lambda_\rho=\omega_\rho(\ZZ+\tau_\rho\ZZ)\) with
\(\tau_\rho\) in the usual closed modular fundamental domain.
Put \(y_\rho=\imag\tau_\rho\), \(A_\rho=|\omega_\rho|^2 y_\rho\)
and \(t_\rho=\rho(t)\) as an elliptic curve point. Boundary choices
in the fundamental domain have the same height.

\begin{proposition}[The metric norm of the marked elliptic theta]\label{prop:ellnorm56}
In the Hermitian frame of metric factor \(e^{-\pi|z|^2/A_\rho}\),
let \(\Theta_\rho\) be the theta section with simple zero at
\(t_\rho\) normalized by \(\Theta_\rho(0)=1\). Its squared norm
\(\nu_\rho\), taken with probability Haar measure, satisfies
\begin{equation}
 A_\rho=4\pi^2 y_\rho|\eta(\tau_\rho)|^4,\qquad
 \nu_\rho=\frac{1}{\sqrt{2y_\rho}\,
            |\vartheta[\epsilon,\beta](0,\tau_\rho)|^2}.
                                                      \label{eq:ellnorm56}
\end{equation}
The characteristic is one of the three even characteristics and
is determined by \(t_\rho\). For an absolute constant \(c\),
\[
       \log^+(A_\rho^{-1})+\log^+\nu_\rho\le c(1+y_\rho),
             \qquad \log^+u=\max(0,\log u).
\]
This includes the point used to define the normalized vacuum in
\srcloc{Formalism/SharpRestriction/GlobalSchurBounds/GL3Theta}{198}.
\end{proposition}
\begin{proof}
Suppress the subscript \(\rho\). The modular discriminant formula
for the lattice is \(g_2^3-27g_3^2=(2\pi)^{12}\omega^{-12}
\eta(\tau)^{24}\), where
\(\eta(\tau)=q^{1/24}\prod_{n\ge1}(1-q^n)\).
Taking absolute values proves the area formula.
\href{https://dlmf.nist.gov/23.15}{NIST, Modular Functions}.

For \(w=z/\omega\) use the convention
\[
 \vartheta[\epsilon,\beta](w,\tau)
   =\sum_{m\in\ZZ} e^{\pi i(m+\epsilon)^2\tau+
                                      2\pi i(m+\epsilon)(w+\beta)} .
\]
The odd characteristic has one zero, at the origin modulo periods.
Indeed pairing opposite half-integers gives the zero, and the
argument principle on a translated period parallelogram gives one
zero counting multiplicity. For the count, the logarithmic
derivative is periodic under \(w\mapsto w+1\) and decreases by
\(2\pi i\) under \(w\mapsto w+\tau\), as follows by reindexing the
series. Translation by the three nonzero half-periods therefore
gives the three even characteristics, with simple zero at the
corresponding marked point. For the appropriate one the Hermitian
theta is
\[
      e^{\pi w^2/(2y)}
        \frac{\vartheta[\epsilon,\beta](w,\tau)}
             {\vartheta[\epsilon,\beta](0,\tau)} .
\]
Its metric magnitude is periodic by the same transformation laws.
It agrees with the unitary translate of the odd Hermitian theta,
divided by its value at zero, as in the source definition. Indeed the
ratio has no zeros or poles, and the logarithm of its magnitude is
a periodic harmonic function, hence constant. The normalization
fixes the constant ratio.

Integrate over \(w=u+v\tau\), \(0\le u,v<1\), with measure
\(du\,dv\). The exponential factor and the metric contribute
\(e^{-2\pi yv^2}\) to the squared magnitude. Integration in \(u\)
kills the distinct Fourier modes. The remaining numerator is
\[
  \sum_{m\in\ZZ}\int_0^1 e^{-2\pi y(m+\epsilon+v)^2}\,dv
                 =\int_\RR e^{-2\pi yv^2}\,dv=(2y)^{-1/2}.
\]
Absolute convergence on the period parallelogram justifies the
calculation. This proves the formula for \(\nu\).

Since \(y\ge\sqrt3/2\), \(|q|\le e^{-\pi\sqrt3}<1\).
Jacobi's product formulas show that the two even constants with
\(\epsilon=0\) are bounded above and below by positive absolute
constants in magnitude. The remaining one has magnitude comparable
to \(e^{-\pi y/4}\). In each product, after extracting the leading
factor, the upper and lower bounds follow from convergence of the
products of \(1+|q|^{n/2}\) and \(1-|q|^{n/2}\) uniformly in this
range. The product for \(\eta\) gives
\(\log|\eta(\tau)|=-\pi y/12+O(1)\) by the same argument.
Substitution yields the asserted bounds.
\href{https://dlmf.nist.gov/20.5}{NIST, Jacobi products}.
\end{proof}

\begin{theorem}[The mean area in the discriminant-one normalization]\label{thm:mean56}
For the quadratic character \(\xi_d\) of \(K\), let
\(D_{\xi_d}=L'(1,\xi_d)/L(1,\xi_d)\).
Taking all embeddings of \(J\) with equal weight gives
\begin{equation}
 \frac{1}{[J:\QQ]}\sum_\rho\log(A_\rho/\pi)
       =\gamma+\log(2\pi/\sqrt d)-D_{\xi_d}.           \label{eq:mean56}
\end{equation}
Moreover the embedding means of \(y_\rho\), \(\log^+(A_\rho^{-1})\)
and \(\log^+\nu_\rho\) are each
\(O(1+\log d+|D_{\xi_d}|)\), with an absolute implied constant.
\end{theorem}
\begin{proof}
Complex multiplication and the equal number of extensions of every
embedding of \(\QQ(j)\) to \(J\) identify the mean of every
function of the reduced lattice shape with its ideal class mean.
The discriminant-one condition fixes the area within a similarity
class by Proposition~\ref{prop:ellnorm56}. The mark can vary above
the same \(j\). The bounds of that proposition hold for all three
choices.

For \(\tau=u+iy\), define the full Epstein sum
\[
       Z(\tau,s)=\sum_{(m,n)\ne(0,0)}y^s|m+n\tau|^{-2s},
                     \qquad \real s>1.
\]
For an integral ideal \(\mathfrak a\) of reduced shape \(\tau\),
its area in \(\CC\) is \(\sqrt d\,N\mathfrak a/2\). Sending
\(z\in\mathfrak a\setminus\{0\}\) to
\((z)\mathfrak a^{-1}\) enumerates the integral ideals in the
inverse class, with multiplicity \(w_K=|\mathcal O_K^\times|\).
Since \(N((z)\mathfrak a^{-1})=|z|^2/N\mathfrak a\), summing
over the \(h\) ideal classes proves
\[
         \frac1h\sum_{[\mathfrak a]} Z(\tau_{\mathfrak a},s)
       =\frac{w_K}{h}(\sqrt d/2)^s\zeta(s)L(s,\xi_d).
\]
Here the factorization follows either by comparing the three
splitting types at each rational prime or from
Section~\ref{sec:abel}.

We compute the constant term of \(Z\) at one. The terms with
\(n=0\) give \(2y^s\zeta(2s)\). For \(n\ne0\) apply the Mellin
integral for the inverse power and Poisson summation in \(m\)
to the Gaussian centered at \(-nu\). The zero Fourier mode
summed over \(n\) is
\[
  2\sqrt\pi\,\frac{\Gamma(s-\tfrac12)}{\Gamma(s)}
                   y^{1-s}\zeta(2s-1).
\]
Its expansion is
\(\pi/(s-1)+2\pi\gamma-2\pi\log2-\pi\log y+O(s-1)\).
This follows from the Laurent expansion of \(\zeta\) and the
gamma duplication formula at one.

The mode with \(n k\ne0\) has Mellin integral proportional to
\[
  e^{2\pi i knu}\int_0^\infty
   v^{s-3/2}e^{-(ny)^2v-\pi^2 k^2/v}\,dv .
\]
The double sum is locally normally convergent near \(s=1\).
Indeed substituting \(v=\pi|k|t/(|n|y)\) leaves a power factor
times \(e^{-\pi |nk|y(t+1/t)}\). One half of this exponential
gives \(e^{-\pi|nk|y}\), the other makes the integral uniformly
finite on compact sets of \(s\). At \(s=1\) the Fourier
transform of \(1/(v^2+a^2)\) at \(k\) is
\((\pi/a)e^{-2\pi |k|a}\), by a residue calculation. With the
factor \(y\) and both signs of \(n,k\), the nonzero modes sum to
\[
  4\pi\sum_{n,k\ge1}\frac{e^{-2\pi nky}\cos(2\pi nku)}{n}
        =-4\pi\sum_{k\ge1}\log|1-e^{2\pi i k\tau}|.
\]
Adding \(2y\zeta(2)=\pi^2y/3\) from \(n=0\) proves
\[
 Z(\tau,s)=\frac{\pi}{s-1}+2\pi(\gamma-\log2)
                    -\pi\log(y|\eta(\tau)|^4)+O_\tau(s-1).
\]
The residue in the ideal class identity fixes the scalar on its
right to have residue \(\pi\). Its constant term divided by \(\pi\)
is therefore \(\gamma+D_{\xi_d}+\log(\sqrt d/2)\).
Comparing constants and using \eqref{eq:ellnorm56} gives
\eqref{eq:mean56}.

Finally the eta product gives
\(-\log(y|\eta|^4)=\pi y/3-\log y+O(1)\ge c_1y-c_2\)
with \(c_1>0\) absolute. Average this inequality and apply the
identity just proved. The bounds for the other two means follow
from Proposition~\ref{prop:ellnorm56}.
\end{proof}

\begin{corollary}[The elliptic factors at every point of the same interval]\label{cor:interval56}
Take the \(d\) selected for \(I\) in Theorem~\ref{thm:twists}.
For all \(x\in I\) the embedding mean of
\(\log^+(A_\rho^{-1})+\log^+\nu_\rho\) is
\(O_{\psi,\eps}(x^{\eps/4})\). For integers \(k,m\ge0\), even
when they vary with \(x\), one has
\[
  \frac1{[J:\QQ]}\sum_\rho
      \log^+\!\left((2\pi/A_\rho)^k\nu_\rho^m\right)
                    \ll_{\psi,\eps}(k+m)x^{\eps/4}.
\]
\end{corollary}
\begin{proof}
The selected character satisfies
\(|D_{\xi_d}|\le3\log x\) and
\(\log d\le Y^{\eps/5}(\log Y)^2\) with \(Y\le x\).
For fixed \(\eps>0\) both quantities are
\(O_\eps(x^{\eps/4})\). Apply Theorem~\ref{thm:mean56} and use
subadditivity of \(\log^+\) on products.
\end{proof}

The source prescribes a vacuum of squared norm
\(\nu_\rho^{[E:K]rs}\) at ranks \(r,s\), as recorded in
\srcloc{Formalism/SharpRestriction/GlobalSchurBounds/StandardLatticesGL3}{496}.
Equation~\eqref{eq:mean56} proves the elliptic area identity
appearing in
\srcloc{Development/AbelianVarietyTheory/Proofs/FoundationProofs/MeanLogPeriodConjugatesEqChowlaSelbergStatement}{8}.
The corollary controls the cost of those two elliptic factors in
the estimate of Section \ref{nav:V-10}. To use that estimate on an algebraic global period
one must still identify the conjugates of its section in the
indicated normalization and compare its class sum with the
specified Bessel projection.

\subsection{The two-torsion point prescribed by the split prime}\label{nav:V-21}
There is an exact identity for the second elliptic factor of
Corollary~\ref{cor:interval56} when the mark is the one used in
the source. Write \(\omega_K=(1+\sqrt{-d})/2\). Since
\(d\equiv7\bmod16\), the ring
\(\mathcal O_K/2\mathcal O_K=\mathbb F_2\times\mathbb F_2\)
has complementary idempotents given by \(\omega_K,\bar\omega_K\).
Indeed their sum is one and their product \((1+d)/4\) is even.
On \(C[2]\) the action is the regular module for this ring, as
is seen from \((\tfrac12\mathcal O_K)/\mathcal O_K\).
There is exactly one point \(t_*\) on which both idempotents act
nontrivially. It is the sum of the two points generating the
distinct summands. This characterizes the point under every
embedding, including those which exchange the summands.

\begin{lemma}[The difference of the other two roots]\label{lem:weber58}
Let \(P,Q\) be the two nonzero points of \(C[2]\) different from
\(t_*\), and define the nonzero algebraic number
\[
                  B_*=(x(P)-x(Q))^2
\]
in the discriminant-one model of Proposition~\ref{prop:ellnorm56}.
At every finite place over a rational prime \(p\), normalized by
\(v_p(p)=1\), one has
\[
                      v_p(B_*)=
            \begin{cases}-4 & p=2,\\ 0 & p\ \text{odd}.\end{cases}
\]
At every complex embedding \(\rho\) the squared metric norm of the
elliptic vacuum with mark \(t_*\) is
\begin{equation}
            \nu_\rho=\frac{\pi}{\sqrt{2A_\rho}\,
                                               |\rho(B_*)|^{1/4}}.
                                                    \label{eq:weber58}
\end{equation}
\end{lemma}
\begin{proof}
The CM \(j\)-invariant is integral. By the criterion for potentially
good reduction of elliptic curves, after finite extension at any
chosen place take an integral Weierstrass equation with unit
discriminant. Enlarge the field to contain \(C[2]\) and the CM
endomorphisms. In comparing \(x\) with the affine coordinate
\(X\) of this good model, the coordinate change has the form
\(x=u^2X+r\). The original equation in the variables \((x,y/2)\)
also has discriminant one. The change scales the discriminant by
\(u^{12}\), so \(u\) is a unit even if \(r\) has a denominator.
It follows that the valuation of \(x(P)-x(Q)\) is that of
\(X(P)-X(Q)\).

For \(p\ne2\) the distinct two-torsion points have distinct
reductions by good reduction. The two nonzero points have integral
affine coordinates. On a smooth Weierstrass cubic two points of
the same \(X\)-coordinate coincide or are negatives, since the
vertical line meets the cubic in that pair and the origin.
Each of the points is its own negative. Hence their reduced
\(X\)-coordinates differ and their difference is a unit.

For \(p=2\), the action of
\(\mathcal O_K\otimes\ZZ_2=\ZZ_2\times\ZZ_2\) on the
\(2\)-divisible group of the good model splits it into two factors.
The endomorphisms extend to the good elliptic scheme by the
extension property for abelian schemes over a discrete valuation
ring. Each factor has height one, as can be read from the order
\(2^n\) of its generic \(2^n\)-torsion using the CM lattice.
The ranks are preserved in the special fiber by finite flatness.
In the special fiber the two dimensions add to one. The
dimension-zero factor is \'{e}tale, and the dimension-one factor
of height one is connected. Indeed its connected formal group
has positive height, leaving height zero for an \'{e}tale quotient.
Thus exactly one of \(P,Q\) reduces to the origin. Write it as
\(P\), so that \(X(Q)\) is integral.

In a parameter \(T\) of the good formal group, multiplication by
two is \(2T+aT^2+O(T^3)\) with \(a\) a unit. To see the unit,
reduce modulo the maximal ideal. The kernel at the origin in the
special fiber has length two by the connected height-one factor.
This length is exactly the order of the first nonzero term of
\([2](T)\) there. Now \(v_2(T(P))>0\) and \([2](T(P))=0\).
If \(v_2(T(P))<1\) the quadratic term has strictly least valuation.
If \(v_2(T(P))>1\) the linear term does. Thus \(v_2(T(P))=1\).
For the Weierstrass parameter \(T=-X/Y\), the expansion of \(X\)
is \(T^{-2}\) times an integral unit series with constant one.
It follows that \(v_2(X(P))=-2\). Since \(X(Q)\) is integral,
squaring the difference gives the asserted \(-4\).

At a complex place write the first full period as \(\omega\).
The classical half-period identities express the square of the
even theta constant in Proposition~\ref{prop:ellnorm56} by
\[
  |\vartheta[\epsilon,\beta](0,\tau)|^4
          =|\omega/\pi|^2\,|x(P)-x(Q)|.
\]
For the mark \(\omega/2\), the constant is \(\vartheta_2(0,\tau)\)
and the difference is \(e_2-e_3\), where
\(e_1=\wp(\omega/2)\), \(e_2=\wp(\omega(1+\tau)/2)\),
\(e_3=\wp(\omega\tau/2)\). For the other two marks the
\(\vartheta_3,\vartheta_4\) identities give the differences of
the other respective pair. The coefficient uses the full period
\(\omega\), twice the half-period \(\omega_1\) in the classical
notation. Substitute this identity into \eqref{eq:ellnorm56}
and use \(A=|\omega|^2\imag\tau\). This gives \eqref{eq:weber58}.
\href{https://dlmf.nist.gov/23.6}{NIST, Weierstrass roots and Jacobi constants}.
\end{proof}

\begin{theorem}[Exact mean norm of the prescribed elliptic vacuum]\label{thm:norm58}
For the paired fields and the unique mark \(t_*\), one has
\begin{equation}
   \frac1{[J:\QQ]}\sum_\rho\log\nu_\rho
       =\frac14\log d-\frac{\gamma}{2}+\frac12D_{\xi_d}.
                                                       \label{eq:norm58}
\end{equation}
In particular, with \(k,m\ge0\) and the same embeddings,
\begin{align}
 \frac1{[J:\QQ]}\sum_\rho\log\!\left((2\pi/A_\rho)^k\nu_\rho^m\right)
   &= k\big(\log(\sqrt d/\pi)-\gamma+D_{\xi_d}\big)       \notag\\
   &\quad +m\big(\tfrac14\log d-\tfrac12\gamma
                                        +\tfrac12D_{\xi_d}\big).
                                                       \label{eq:archmain58}
\end{align}
\end{theorem}
\begin{proof}
One may enlarge \(J\) to contain all the points above since this
repeats each embedding the same number of times. Lemma~\ref{lem:weber58}
implies that \(16B_*\) has valuation zero at every finite place.
It is an algebraic unit. Its field norm has absolute value one,
and therefore \([J:\QQ]^{-1}\sum_\rho\log|\rho(B_*)|=-4\log2\).
Average the logarithm of \eqref{eq:weber58}. This gives
\[
    \frac1{[J:\QQ]}\sum_\rho\log\nu_\rho
      =\log(\sqrt2\,\pi)-\frac1{2[J:\QQ]}\sum_\rho\log A_\rho.
\]
Theorem~\ref{thm:mean56} now gives \eqref{eq:norm58}. The same
theorem gives the mean of \(\log(2\pi/A_\rho)\). Linearity gives
\eqref{eq:archmain58}.
\end{proof}

This evaluates the actual elliptic metric factors with their signs
in the product formula. A source statement on means for the marked
vacuum is
\srcloc{Development/AbelianVarietyTheory/Proofs/DecoratedLatticeProofs/ThetaZeroMeanLogNormConjugatesOfWeberStatement}{13}.
The proof here uses only the classical half-period identities, the
valuation of the difference of two roots, and the class mean in
Theorem~\ref{thm:mean56}. Applying the identity to a higher-rank
period requires an arithmetic identification of its transported
section at the same embeddings.

\subsection{A Gram matrix comparison for the analytic theta section}\label{nav:V-16}
\label{sec:gram66}

Fix \(p\in S\), let \(O=\mathcal O_{K_p}\), and use the elliptic
curve and nonzero two-torsion mark of Section~\ref{sec:period63}.
Write \(\vartheta=\Theta_{t_*}\), with \(\vartheta(0)=1\). Let
\(T\) be a good parameter at the origin, with linear term \(Z\)
in the additive coordinate. By an integral germ in \(n\) elliptic
parameters we mean a member of \(R[[T_1,\ldots,T_n]]\). We view
this ring inside the characteristic-zero Taylor series by substituting
\(T_i=T(u_i)\), where \(u_i\) are additive coordinates.
A row \(a\) of length \(n\) gives the linear form \(a u\).
For a matrix \(A\) with \(n\) columns and entries in \(O\) put
\[
                  G_A(u)=\prod_i\vartheta((Au)_i).
\]
Conjugation on \(K_p\) is its nontrivial automorphism over
\(\QQ_p\), and \(A^*\) denotes its conjugate transpose.

\begin{lemma}[Cancellation of equal integral Hermitian forms]\label{lem:rational66}
Let \(a_j\in\mathcal O_K^n\) be finitely many rows and let
\(d_j\in\ZZ\) satisfy \(\sum_j d_j a_j^*a_j=0\).
The product \(H(u)=\prod_j\vartheta(a_j u)^{d_j}\) is a unit
in the ring of integral elliptic germs, with constant coefficient one.
\end{lemma}
\begin{proof}
A zero row contributes one. For every other row the map
\(C^n\to C\) given by \(a_j\) is a homomorphism of CM abelian
varieties. Under translation by \(\lambda\in\Lambda^n\) the
elliptic theta transformation law multiplies its pullback by
a sign times
\[
  \exp\!\left(\frac{\pi}{A_0}(a_j u)\overline{a_j\lambda}
                  +\frac{\pi}{2A_0}|a_j\lambda|^2\right),
\]
where \(A_0\) is the area of \(\Lambda\). The sign comes from
the two-torsion characteristic. The Gram identity cancels both
exponents in the product. Hence \(H^2\) is a meromorphic
function on the projective complex variety \(C^n\), with divisor
\(2\sum_j d_j a_j^*[t_*]\); here the superscript on a divisor
means pullback by the row homomorphism. By algebraicity of
meromorphic functions on a projective complex variety it is
rational. The divisor and the origin are defined over a number
field. A rational function with that divisor and value one at
the origin descends to its field of definition after a finite
extension. Indeed the Riemann--Roch space for the negative of
this principal divisor is one dimensional after base change
to \(\CC\), with generator \(H^2\). Global sections of the
corresponding line bundle commute with field extension.
Choose a generator over the field of definition and normalize
its value at the origin. It has the prescribed divisor since
its base change is a scalar multiple of \(H^2\).

Pass to the good model over \(R\), enlarging the coefficient field
if necessary. All the row morphisms extend to its elliptic scheme,
by the extension property for homomorphisms of abelian schemes
over a discrete valuation ring with good reduction. Every such
homomorphism sends the origin to the origin, and the reduction of
\(t_*\) is a distinct point since two is a unit. Thus the closures
of the horizontal zeros and poles avoid the closed origin.
The special fiber of the product is geometrically integral.
In the regular local ring at that origin the divisor of \(H^2\)
therefore consists only of an integer multiple of the uniformizer
divisor. After removing that power of the uniformizer, a rational
function with neither zeros nor poles in codimension one is a unit,
by normality. Evaluating along the zero section gives a unit of
\(R\) times the removed power. The normalized value one forces
the power to be zero. Thus \(H^2\) is a unit in the completed
local ring with value one. Since two is a unit, it has a unique
formal square root with constant coefficient one, obtained
recursively by degree. This is \(H\).
\end{proof}

\begin{theorem}[Dependence on the Gram matrix over the local CM integers]\label{thm:gram66}
If \(A,B\) have entries in \(O\), the same number of columns,
and \(A^*A=B^*B\), then \(G_A/G_B\) is an integral elliptic
unit. More generally, a product of such functions with integer
exponents and zero signed Gram matrix is an integral elliptic unit.
In particular:
\begin{enumerate}
\item If \(A^*A=0\), then \(G_A\) is an integral elliptic unit.
\item If \(A^*A=\left(\begin{smallmatrix}0&I_n\\I_n&0\end{smallmatrix}\right)\)
and the coordinates are \(u_1,\ldots,u_n,v_1,\ldots,v_n\), then
\[
       G_A(u,v)\prod_{j=1}^n
           \frac{\vartheta(u_j)\vartheta(v_j)}{\vartheta(u_j+v_j)}
\]
is an integral elliptic unit.
\end{enumerate}
\end{theorem}
\begin{proof}
The equation for the Gram matrix over \(O\) need not hold for
arbitrary integral approximations to the rows. We account for their
Gram matrices before passing to the limit. Choose
\(\lambda\in\mathcal O_K\) with
\(\mathcal O_K=\ZZ+\ZZ\lambda\).
For a fixed number \(d\) of columns take the rows
\[
 e_i\ (1\le i\le d),\qquad
 e_i+e_j,\ e_i+\lambda e_j\ (1\le i<j\le d).
\]
Their rank-one matrices \(a^*a\) form a \(\ZZ\)-basis of
the Hermitian matrices over \(\mathcal O_K\). In the \((i,j)\)
position with \(i<j\), the last two rows contribute \(1,\lambda\);
having prescribed all such positions, one corrects the diagonal
by integer multiples of the first rows. This argument also proves
uniqueness. After completion the same matrices form a
\(\ZZ_p\)-basis of the Hermitian matrices over \(O\), since
\(p\) is inert.

For a finite signed list of local rows approximate every entry
by elements of \(\mathcal O_K\). This is possible since its
completion at \(p\) is \(O\). At each stage expand the signed
Gram matrix of the approximating list in the displayed basis,
with coefficients \(h_a\in\ZZ\). Multiply the approximating
product by \(\prod_a\vartheta(a u)^{-h_a}\).
The resulting signed Gram matrix vanishes exactly. By
Lemma~\ref{lem:rational66} the corrected product is an integral
elliptic unit. The finitely many \(h_a\) converge to zero in
\(\ZZ_p\), by the assumption on the original Gram matrix and
continuity of its basis coordinates.

For each fixed degree the coefficients converge in the local
field. Substitution of a linear form is polynomial in its entries
in that degree. For any series \(1+U\) with zero constant term
in \(U\), the same degree in \((1+U)^m\) involves only finitely
many polynomials \(\binom mj\) in the exponent \(m\). These
statements remain true after the fixed substitution
\(u_i=\log_C(T_i)\), inverse to \(T(u_i)\). Thus the correcting
factors tend coefficientwise to one. The valuation ring is
closed and constant coefficients are one. This proves the general
assertion. For the hyperbolic Gram matrix, the row \(e_i+e_{n+i}\)
with exponent one and the two rows \(e_i,e_{n+i}\) with
exponent minus one give precisely its two off-diagonal unit
entries. Sum over \(i\) and apply the signed assertion to
the resulting ratio.
\end{proof}

\begin{corollary}[Restriction of each class section]\label{cor:class66}
In product coordinates \(z=V^*Z\), write \(F_{ch}\) for the
class section of Theorem~\ref{thm:jet63}. For every matrix \(A\)
with entries in \(O\), the germ \(F_{ch}(Au)/G_A(u)\) is integral
in the elliptic parameters at every place over the chosen prime.
Consequently the restriction \(F_{ch}(Au)\) to \(A^*A=0\) is
integral in those parameters. In a hyperbolic basis as in
Theorem~\ref{thm:gram66}, multiplying \(F_{ch}(A(u,v))\) by
the same displayed theta ratios gives an integral elliptic germ.
\end{corollary}
\begin{proof}
The proof of Theorem~\ref{thm:jet63} expresses each class section
in the coordinates \(z\) as a sum with integral coefficients
of translates of the marked product. Dividing by its common
factor \(\prod_i\vartheta(z_i)\) leaves sums of products of the
integral elliptic germs \(R_w\) from Theorem~\ref{thm:trans57}.
Substitution \(z_i=(Au)_i\) preserves this elliptic integrality.
For integral global rows it is pullback by a morphism of the good
elliptic schemes. For a row over \(O\), use global approximations.
In each fixed degree the series \(T(a u)\), after
\(u_j=\log_C(T_j)\), depends continuously on the entries of
\(a\). Its coefficients are therefore still integral.
Substituting these into an integral germ at the origin is valid
since all their constant terms are zero. Multiplication by the
unit of Theorem~\ref{thm:gram66} proves the two consequences.
\end{proof}

The factor in the hyperbolic case is explicit in terms of the
same marked analytic theta used at infinity. On either of the two
isotropic halves it is identically one. This permits the derivative
in the other directions to be counted before an integral restriction
is made. The source introduces an isotropic condition in
\srcloc{Formalism/SharpRestriction/GlobalSchurBounds/GL3Theta}{241}
and a hyperbolic basis in the same file at line~296.

\section{Frobenius translates and higher derivatives}\label{nav:g-V-2}

\subsection{The Frobenius count in a single list of minor degrees}\label{nav:V-17}
\label{sec:count73}

We compare the \(K_d\)-oriented prime count with the
derivatives of Section~\ref{sec:period63}. Use one of the common
intervals \([V,V^{3/2}]\) in
Section \ref{nav:I-28}.
At each \(x\) write
\[
 E=K_\psi C_l(\sqrt{-q},\sqrt{-d}),\quad
 K=\QQ(\sqrt{-d}),\quad C_l=\QQ(\zeta_{2^{l+2}})^+ .
\]
Use the index of the chosen fields in Volume~I, including its
shortened index when that family is used. In both cases
\(2^l\gg_\psi x^{2+8\eps}\) and \([E:\QQ]\ll_\psi x^{2+8\eps}\).
Identify the CM embeddings over \(K\) with
\(\Phi=\Gal(E/K)\) by the base embedding, and use additive
notation for the abelian Galois group. Complex conjugation is
\(c_E\). Let \(t\in\Phi\) be the position corresponding to a
translate in the count. Set
\[
 b=e^{x^\eps},\qquad a=\lceil xb\rceil,\qquad
 W=\lceil\sqrt b\rceil,\qquad B'=\lceil b\rceil,\qquad X=a/W .
\]
The list of integer thresholds \(w_\tau\), \(\tau\in\Phi\), may
be any list in \([W,B']\). In particular it may be the
distinguished-set list of Volume~I for any of the translates.
The translate of a Frobenius element outside \(\Phi\) is
\(h\mapsto t+c_E+h\). Write \(C_d(t,a,w)=\sum_p C_{d,p}\)
for that oriented count, with the inertia average of Section \ref{nav:I-23}. At an unramified inert prime, with Frobenius
\(\sigma_p\), the exact term is
\begin{equation}
 C_{d,p}=2\log p\sum_{j\ge1,\ j\ {\rm odd}}
       [a/p^j-w_{t+j\delta_p}]_+,\qquad \delta_p=c_E+\sigma_p\in\Phi .
                                                        \label{eq:cp73}
\end{equation}
Indeed even powers restrict trivially to \(K\); at the odd
powers one has \(c_E+j\sigma_p=j\delta_p\).
Use fixed ranks \(r\le s\) of a definite single step with
\(rs\) even, allowing these ranks to vary polynomially with
\(x\). Fix the single input, level \(Q\), and class sections
\(F_{ch}\) of Theorem~\ref{thm:jet63}.

\begin{lemma}[Odd multiples in the CM type]\label{lem:cycles73}
Decompose \(\Phi=P\oplus P'\) into its 2-primary and odd-order
subgroups. For \(\delta=g+y\) with \(0\ne g\in P\) and
\(y\in P'\), the odd multiples of \(\delta\) form the set
\[
               \operatorname{Gen}(\langle g\rangle)
                                   +\langle y\rangle .
\]
Here \(\operatorname{Gen}(C)\) denotes the generators of
\(C\). The number of nontrivial cyclic subgroups of \(P\)
is \(O_\psi(1+\log x)\).
\end{lemma}
\begin{proof}
An odd residue modulo the order of \(g\) runs through exactly
the units modulo that 2-power. The Chinese remainder theorem
allows any such residue together with any residue modulo the
odd order of \(y\), giving the stated set. Put
\(A_0=\Gal(K_\psi C_l/K_\psi)\), regarded as a subgroup of
\(\Phi\) fixing the two quadratic factors. Its restriction
to \(C_l\) is injective, so it is a cyclic 2-group.
Its index in \(\Phi\) is \(2o\), with
\(o=[K_\psi:\QQ]\). For \(i\ge1\) there are at most
\(2o\cdot2^i\) elements of \(P\) killed by \(2^i\):
each fiber in \(P/A_0\) has at most \(|A_0[2^i]|\)
such elements. Every cyclic subgroup of order \(2^i\) has
\(2^{i-1}\) generators, and their sets are disjoint. Thus
there are at most \(4o\) per order. The degree bound on \(E\)
bounds the number of possible orders by \(O_\psi(1+\log x)\).
\end{proof}

For each nontrivial cyclic subgroup \(C\) of \(P\), choose
one generator \(g_C\) and set
\[
 u_{C,t}=\min_{g\in\operatorname{Gen}(C),\,y\in P'}
                    w_{t+g+y}.
\]
Use the following minor degrees in \eqref{eq:pair55}:
\begin{equation}
    b_t=a,\qquad b_{t+g_C}=u_{C,t}\quad(C\ne0),\qquad
    b_\tau=0\quad\hbox{at all remaining embeddings}.
                                                        \label{eq:degrees73}
\end{equation}
The positions in this prescription are all distinct. The list
is chosen once for the whole prime sum. Let \(S\) contain all
primes \(p\le X\) which are inert in \(K\), unramified in
\(E\), prime to \(Q\), and at least five. We may take exactly
this finite set. Use the common basis of Lemma~\ref{lem:basis63}
for every prime in \(S\). On a bounded common interval one
can also take the union of the relevant good primes for each
fixed field, ranks and level, since the integer parameters range over
finite sets, and Lemma~\ref{lem:basis63} applies to their
union at fixed field and lattices.

\begin{theorem}[Divisibility for the actual weighted count]\label{thm:count73}
For every \(p\in S\) with nonzero 2-primary component of
\(\delta_p\), at every place above \(p\), all the derivatives
\(P_T(\partial)F_{ch}(0)\) for \eqref{eq:degrees73} and all their
cleared class coordinates \(A_r v_{h,T}\) have valuation at least
\[
      r\sum_{j\ge1,\ j\ {\rm odd}}
                  [d_{p^j}(a,s)-w_{t+j\delta_p}]_+ ,\qquad
      d_L(a,s)=\max(0,\lceil(a-s(L-1))/L\rceil).
\]
Their common guaranteed lower bound is at least
\[
       r\left(\frac{C_{d,p}}{2\log p}
            -s\,\#\{j\ge1:j\ {\rm odd},\ p^j\le X\}\right).
\]
If the thresholds factor through a quotient of \(\Phi\) of
exponent two, the low half of the cycle in this proof has total
degree exactly \(w_{t+j\delta_p}\) for every odd \(j\).
This includes the dense distinguished set \(T_2=A_0\) for
the trivial character and every quadratic \(\psi\).
\end{theorem}
\begin{proof}
In the Frobenius cycle through \(t\), choose as its high half
the even translates of \(t\) by \(\delta_p\). The order of
\(\delta_p\) is even by hypothesis. Write
\(\delta_p=g+y\) as in Lemma~\ref{lem:cycles73}, and set
\(C=\langle g\rangle\). The other half is
\(t+\operatorname{Gen}(C)+\langle y\rangle\).
It contains the position \(t+g_C\) and no other selected
position of \eqref{eq:degrees73}. Indeed a different selected
generator belongs to a different cyclic subgroup; the position
\(t\) is on the even side. The degree sum on the low half is
therefore exactly \(u_{C,t}\). This is at most
\(w_{t+j\delta_p}\) at every odd \(j\).
If the thresholds factor through an exponent-two quotient,
every generator of \(C\) has the same image there and \(P'\)
has zero image. This proves the asserted equality in that case.
For \(T_2=A_0\) the list factors through \(\Phi/A_0\),
which in the coordinates of Volume~I is a product of a
group of order two with \(\Gal(K_\psi/\QQ)\).

The contribution of the high half to the sum of the
\(d_{p^j}(b_\tau,s)\) is at least \(d_{p^j}(a,s)\).
Apply Corollary \ref{cor:minor67} to this cycle. Every other cycle has
nonnegative contribution. That corollary estimates the actual
determinant monomials before summing the classes and remains
valid for the class coordinates after clearing their stabilizer
weights. This proves the first bound. Since
\(d_L(a,s)\ge a/L-s\), replacing it by \(a/L\) costs at
most \(s\) at each active level, by the Lipschitz bound for
the positive part. A term with \(p^j>X\) is zero in
\eqref{eq:cp73}. The second bound follows.
\end{proof}

\begin{theorem}[A single-step norm over the full prime range]\label{thm:norm73}
Put \(Y_t=A_s A_r^2\mathcal N_t\) for the squared class sum
\eqref{eq:pair55} with degrees \eqref{eq:degrees73}.
The divisibility in Theorem~\ref{thm:count73} holds at twice
its stated exponents for \(Y_t\). If some scalar class
coordinate survives, \(Y_t\ne0\). In a common finite Galois
field \(J\) of its coefficients one then has
\begin{align}
 C_d(t,a,w)\le{}&
    \frac{[J:\QQ]^{-1}\log|N_{J/\QQ}Y_t|
                     +B_-^{(S)}(Y_t)}{r}
          +E_Q+O_{\psi,\eps}(b/x),                         \label{eq:norm73}\\
 E_Q={}&2a\!\sum_{\substack{p\mid Q,\ p\le X\\p\ {\rm inert\ in}\ K}}
                       \frac{\log p}{p-1}.                 \notag
\end{align}
Here the error is uniform for ranks bounded by a fixed power
of \(x\), with threshold allowed to depend on that power.
The denominator term is the normalized sum of negative
valuations outside \(S\) as in Theorem~\ref{thm:joint63}.
For \(Q\) supported on fixed primes one has \(E_Q=0\) on
every sufficiently distant common interval.
\end{theorem}
\begin{proof}
The cleared squared sum has terms
\(A_s m'_h (A_r v_{h,T})\overline{(A_r v_{h,T})}\) with
integral outer weights. Both scalar factors satisfy the lower
bound at all places over \(p\), since complex conjugation
permutes these places. At the distinguished complex embedding
the sum is strictly positive if any coordinate is nonzero.
All primes of \(S\) have integral coordinates even when the
2-primary component under consideration is zero.

We bound the contributions to the count at primes missed by
Theorem~\ref{thm:count73}. If \(\delta_p\) has odd order then
it acts trivially on the real 2-power cyclotomic layer
\(C_l\). Complex conjugation also acts trivially there, so
\(p\equiv\pm1\pmod M\) with \(M=2^{l+2}\).
At an unramified prime the entire count is at most
\(2a\log p/(p-1)\). Among integers at least two in the
two residue classes, all are of the form \(jM\pm1\)
with \(j\ge1\). Summing the reciprocal denominators up to
\(X\), and bounding each logarithm by \(\log(2X)\), gives
\[
      \sum_{\substack{p\le X\\p\equiv\pm1\bmod M}} C_{d,p}
      \ll a M^{-1}\log^2(2X)
      \ll_{\psi,\eps} b x^{-1-6\eps}.
\]
The same pointwise bound holds at ramified primes by the
normalized inertia average. The size estimates for \(M\) and
\(\log X=O_\eps(x^\eps)\) prove the last inequality.

A prime split in \(K\) contributes zero, including when it
ramifies in \(E/K\), since its entire inertia-averaged
Frobenius coset fixes \(K\). The primes dividing \(2\operatorname{cond}(\psi)\),
as well as 3, split in \(K\) on distant intervals by Theorem \ref{thm:twists}. The only other possible ramified primes are
\(q,d\). Both grow faster than every fixed power of \(x\)
there, and their total contribution is at most
\(2a\sum_{p=q,d}\log p/(p-1)=o(b/x)\).
The remaining omitted inert primes dividing the level are
covered by \(E_Q\). The same splitting argument shows that
this term vanishes for fixed prime support.

Finally the accumulated rounding error in twice the local
estimate after dividing by \(r\) is at most
\[
         2s\sum_{p^j\le X}\log p
             \le 2s X\log(2X)=o(b/x).
\]
The inequality uses uniqueness of the prime base of a prime
power, and \(X\asymp x\sqrt b\) gives the last estimate.
The field norm formula for \(Y_t\) gives an upper bound for
the sum of its finite valuations over \(S\) after adding all
negative contributions outside \(S\) to the normalized log
norm. Combining these inequalities proves \eqref{eq:norm73}.
\end{proof}

\begin{proposition}[The additional degrees in the archimedean estimate]\leavevmode\label{prop:cost73}
Uniformly on the same family,
\[
  0\le\sum_{\tau\ne t} b_\tau\ll_\psi b(1+\log x),\qquad
  0\le\frac1r\sum_{\tau\ne t}\log C_{r,b_\tau}
                     \ll_{\psi,\eps} b x^\eps(1+\log x).
\]
In particular the second expression is \(o(b\sqrt x)\).
The estimate concerns the explicit factorials in the
archimedean norm. Its other factor is the norm of the projected vector.
\end{proposition}
\begin{proof}
There is at most one additional nonzero degree per cyclic
subgroup in Lemma~\ref{lem:cycles73}, each at most \(B'\).
For a degree \(m>0\), the product defining \(C_{r,m}\) in
Theorem~\ref{thm:frame55} has \(rm\) factors at most \(m+r\).
Since \(\log(B'+r)=O_\eps(x^\eps)\) for polynomial ranks
on distant intervals, this proves both assertions.
\end{proof}

The scalar \(Y_t\) in \eqref{eq:norm73} uses a single definite
theta step and the displayed auxiliary choice of degrees. Its
field norm has the archimedean upper bound of
Theorem \ref{thm:opt60}.
To deduce the centered count estimate needed for Dirichlet GRH
one still needs a surviving projection, control of
\(B_-^{(S)}(Y_t)\), and an upper estimate in the normalization
of the full Bessel period. The simultaneous large-prime local
cohomology chart on these same fields is in
Theorem \ref{thm:first72}.

\subsection{Nonzero translated class sums at one definite step}\label{nav:V-18}
\label{sec:mean74}

The estimate in Theorem~\ref{thm:norm73} involves a norm of a
nonzero sum. We now construct such sums at a definite step, for
the positive normalization of the split oscillator on the fixed
class representatives. We keep all the stabilizer weights. The
calculation uses the freedom to translate the finite input.

Let \(z=(z_1,\ldots,z_D)\) be the matrix coordinates at all
CM embeddings, orthonormal for the Euclidean form used in
Proposition~\ref{prop:gauss55}. Put \(c_0=\pi/A\), with
\(A\) the area of the marked elliptic period lattice.
The Hermitian metric on functions in this trivialization is
\(|F(z)|^2e^{-c_0|z|^2}\). Suppose a full lattice
\(\Gamma\subset\CC^D\) and a common unitary multiplier give
all the functions under consideration the transformation law
\[
 F(z+\gamma)=\epsilon(\gamma)
   \exp\big(c_0 z\cdot\bar\gamma+c_0|\gamma|^2/2\big)F(z)
           \qquad(\gamma\in\Gamma).
\]
Use Haar probability on the torus \(T_\Gamma=\CC^D/\Gamma\)
and the inner product from this metric, conjugate-linear in
the first variable. Write \(\|F\|_\Gamma\) for its norm.

\begin{lemma}[The translation average of a theta derivative]\label{lem:mean74}
For every holomorphic section \(F\) with this transformation
law and every homogeneous polynomial \(P\) of degree \(k\)
in the matrix coordinates,
\begin{equation}
 \int_{T_\Gamma}
       |P(\partial_z)(U_wF)(0)|^2\,dw
     =c_0^k\|P\|_{\rm F}^2\|F\|_\Gamma^2,
                                                        \label{eq:mean74}
\end{equation}
where \(\|P\|_{\rm F}^2=\sum_\alpha\alpha!|P_\alpha|^2\) and
\[
          (U_wF)(z)=
          e^{-c_0z\cdot\bar w-c_0|w|^2/2}F(z+w).
\]
\end{lemma}
\begin{proof}
On smooth sections put
\(L_j=\partial/\partial\bar w_j\) and
\(D_j=\partial/\partial w_j-c_0\bar w_j\).
Both preserve the transformation law by direct differentiation.
Integration by parts with the metric gives \(L_j^*=-D_j\).
Indeed the derivative of the metric produces the term
\(-c_0\bar w_j\) in \(D_j\), and opposite faces of a fundamental
parallelepiped cancel since the metric pairings are periodic.
The mixed commutator is \([L_i,-D_j]=c_0\mathbf1_{i=j}\).
The \(L_i\)'s commute among themselves, as do the \(D_j\)'s,
and all \(L_i F\) vanish.

For \(|\alpha|=|\beta|=k\) move the creation operators
\((-D)^\alpha\) in the first argument of the inner product to
the second using this adjoint identity. Each \(L_i\) then
commutes to the right, picking up \(c_0\beta_i\) when it removes
one \(-D_i\). A surviving \(L_i\) at the far right kills \(F\).
This gives
\[
    \langle D^\alpha F,D^\beta F\rangle_\Gamma
       =\mathbf1_{\alpha=\beta}\,c_0^k\alpha!\|F\|_\Gamma^2.
\]
The two signs \((-1)^k\) cancel. Finally differentiating the
translation formula at \(z=0\) yields
\(e^{-c_0|w|^2/2}P(D)F(w)\), since differentiating in \(w\)
treats \(\bar w\) as constant. The last two displays prove
\eqref{eq:mean74}, including the periodicity of its integrand.
\end{proof}

Use now the standard normalized vacuum \(G_0\) on the rational
matrix lattice \(M\) of Section~\ref{sec:period63}. At a pair
of class representatives let \(G_{ch}\) be its finite Weil
transform, using the positive split changes of variables of
Lemma~\ref{lem:split61} at the moving places. We take the
optional determinant factors at those places to be trivial.
Write \(M_{ch}\) for the moved rational lattice and set
\[
          n_{ch}=[M_{ch}:M_{ch}\cap M],\qquad
          H_h=\sum_c m_c G_{ch}.
\]
Take \(\Gamma\) to be the analytic period image of \(mM\)
for a positive integer \(m\) with \(mM\subseteq M_{ch}\)
at all pairs. The corrected periods \(W(\gamma)\) of Section~\ref{nav:V-13}
fix each corresponding vacuum on its lattice, hence the
multipliers for \(U_\gamma\) agree on this common sublattice.
It also serves for \(G_0\).

\begin{proposition}[Positive overlap before differentiation]\label{prop:overlap74}
At every target class \(h\),
\[
  \langle G_0,H_h\rangle_\Gamma
      =\|G_0\|_\Gamma^2\sum_c m_c n_{ch}^{-1/2}>0 .
\]
For the fixed representatives, define
\[
       \mathcal N_t(w)=\sum_hm'_h\sum_T
                  |P_T(\partial)(U_wH_h)(0)|^2
\]
using the degrees of \eqref{eq:degrees73}, or any one prescribed
list of nonnegative minor degrees. Its mean is exactly
\begin{equation}
 \int_{T_\Gamma}\mathcal N_t(w)\,dw
       = c_0^k\left(\sum_T\|P_T\|_{\rm F}^2\right)
                       \sum_h m'_h\|H_h\|_\Gamma^2 >0,
             \qquad k=r\sum_\tau b_\tau .                 \label{eq:positive74}
\end{equation}
\end{proposition}
\begin{proof}
The positive finite intertwiner taking \(G_0\) to \(G_{ch}\)
has one translate per coset of \(M_{ch}\cap M\) in \(M_{ch}\)
and a coefficient \(n_{ch}^{-1/2}\) on the zero coset.
This follows by multiplying the split local intertwiners in
Lemma~\ref{lem:split61}. Distinct cosets give translates
orthogonal to \(G_0\) except at zero. In fact the self-duality
of the source period lattice supplies a source period with
nontrivial commutator against each nonzero coset. Its unitary
translation preserves the common section space on a small enough
common period lattice and acts on the two sections by different
unitary eigenvalues. Decreasing the common lattice does not
affect this inner product, since its periodic integrand has the
same average on any finite cover. Equivalently the orthogonality
is the computation of Corollary~\ref{cor:move55}.
The zero coset gives the stated inner product. Positivity follows
from \(m_c>0\), \(n_{ch}>0\), and \(G_0\ne0\).
Apply Lemma~\ref{lem:mean74} to \(H_h\) and sum the finite
identities. There is at least one product of minors and every
such product is a nonzero polynomial since \(r\le s\). This
proves the strict inequality.
\end{proof}

\begin{theorem}[One rational torsion input for all the indicated derivatives]\label{thm:nonzero74}
Fix a rational prime \(P_0\ge5\) away from the fixed conductor
and choose all the class representatives to be the identity
there. On sufficiently distant intervals \(P_0\) splits in
\(K_d\). At each fixed \(x,E,r,s\) there is a single finite
vector \(W(v)G_0\) with \(v\in mM[1/P_0]\) for which all the
fixed-representative sums \(\mathcal N_t\) in Section
\ref{sec:count73} are nonzero simultaneously. One may include
finitely many additional lists of minor degrees in this choice.
For the resulting algebraic numbers \(Y_t=A_sA_r^2\mathcal N_t\)
the inequality \eqref{eq:norm73} holds with \(E_Q=0\).
\end{theorem}
\begin{proof}
Take one product of minors for each of the finitely many lists
and one target class \(h\). By Lemma~\ref{lem:mean74} and
Proposition~\ref{prop:overlap74} its squared derivative on
\(U_wH_h\) is a nonnegative real analytic function on the
connected torus, with strictly positive mean. Its zero set has
empty interior by uniqueness of real analytic continuation.
The intersection of the finitely many nonvanishing sets is
open and dense. The images of \(mM[1/P_0]\) are dense as
well, since in a basis of \(mM\) they give all grids with mesh
\(P_0^{-e}\), \(e\ge0\). Choose one such vector \(v\) in the
open intersection, and let \(w\) be its analytic coordinate.
The scalar relating \(W(v)\) and \(U_w\) is the common root
of unity \(e^{\pi i q(v)}\) from Section~\ref{nav:V-13}.

We verify that this translation can be made before the split
Weil transforms at the chosen representatives. Call the finite
adelic transformation at a pair \(B_{ch}\). The oscillator
action conjugates \(W(v)\) to \(W(B_{ch}v)\), by the position
and frequency operators of Section~\ref{nav:V-13}. At a moving place,
which lies away from \(P_0\), the component of \(v\) belongs
to \(mM\) locally. Both \(v\) and \(B_{ch}v\) therefore belong to
the target lattice there, since
\(B_{ch}M=M_{ch}\) locally. The two corrected local period
operators fix the target vacuum. At every other place
\(B_{ch}\) is the identity. The finite product gives
\[
              \omega(B_{ch})W(v)G_0=W(v)G_{ch}
\]
in the common realization of Section~\ref{nav:V-13}, with the same scalar
relating its right side to \(U_wG_{ch}\) at all pairs.
Linearity now gives the stated nonzero class sums for the
single input. The existence of representatives avoiding a
prescribed finite set is Theorem~\ref{thm:class61}.
Theorem \ref{thm:twists}, gives \(d\equiv-1\pmod{P_0}\) once its growing cutoff exceeds \(P_0\). Thus the fixed prime splits in \(K_d\).

Some \(P_0^e v\) belongs to \(mM\), so take \(Q=P_0^e\).
This input has a single coefficient one in the finite
translation basis. Theorem~\ref{thm:transport59} and
Theorem~\ref{thm:trans57} make its transformed derivatives
algebraic. At every selected inert prime the additional
translation has order prime to that prime. The reduction,
Gram comparison, and divisibility arguments of Sections
\ref{sec:period63}, \ref{sec:gram66} and \ref{sec:count73}
therefore apply with the same single basis. The field norm
can now be taken at all the chosen translates since their
squared sums are positive at the distinguished embedding.
The support of \(Q\) splits in \(K_d\), giving \(E_Q=0\).
\end{proof}

The choice of torsion point in this theorem is existential.
For the fixed representatives the one coefficient has magnitude
one at every complex embedding, so the bound for the translated
vacuum uses that magnitude in place of \(Q^{rs/2}\).
Its arithmetic denominator at other places, including \(P_0\),
still enters \(B_-^{(S)}(Y_t)\). The theorem evaluates the
finite sums at the representatives specified before the
translation. Identifying these evaluations with an automorphic
integral at the originally prescribed levels requires the
corresponding level invariance and determinant characters.
The nonzero projection of the full iterated row and column
lifts and its quantitative normalization remain to be proved.

\subsection{The cost of the prescribed derivatives}\label{nav:V-22}
Use the class sections constructed and transported in
Theorem \ref{thm:transport59}.
The scale in their pointwise estimate is \(a_\rho=\pi/A_\rho\).
We extract each homogeneous component on a sphere of its own
radius. This determines the coefficient of the total derivative
order in the averaged estimate.

\begin{lemma}[A sphere estimate for a homogeneous Taylor component]\label{lem:sphere60}
Let \(G\) be entire on a product of \(f\) copies of \(\CC^l\)
and suppose \(|G(Z)|^2\le S\exp(a\sum_{\tau=1}^f|Z_\tau|^2)\),
where \(a>0\). Write \(G_{\boldsymbol j}\) for the Taylor
component of degree \(j_\tau\ge0\) in block \(\tau\) and
\(j=\sum_\tau j_\tau\). In the Fischer norm, in which
\(\|Z^\alpha\|^2=\alpha!\) and monomials are orthogonal,
\begin{equation}
   \|G_{\boldsymbol j}\|^2\le S\,a^j\prod_\tau H_l(j_\tau),
    \quad H_l(u)=\frac{(u+l-1)!}{(l-1)!}\left(\frac eu\right)^u
                       \ (u\ge1),\quad H_l(0)=1.           \label{eq:sphere60}
\end{equation}
For every \(u\ge1\),
\[
 \log H_l(u)\le l\log(1+u/l)+u\log(1+l/u)
                         \le l\big(1+\log(1+u/l)\big).
\]
\end{lemma}
\begin{proof}
Set the zero-degree blocks to zero throughout. For a homogeneous
polynomial \(P\) of degree \(u\) on \(\CC^l\), the standard
complex Gaussian integral of \(|P|^2\) equals \(\|P\|^2\).
Indeed the coordinate integrals give the factorials. The squared
radius of the same Gaussian has density
\(t^{l-1}e^{-t}/(l-1)!\) by polar integration, and its direction
has the uniform probability measure on the unit sphere, independently
of the radius. Taking the \(u\)-th moment of the squared radius gives
\[
       \int_{|\omega|=1}|P(\omega)|^2\,d\omega
                        =\frac{(l-1)!}{(u+l-1)!}\,\|P\|^2 .
\]
Apply the product version to \(G_{\boldsymbol j}\) on spheres of
radii \(R_\tau>0\). The integral of its square is bounded by the
integral of \(|G|^2\) there. In fact the homogeneous component
is the Fourier projection obtained by rotating each block
independently by \(e^{i\theta_\tau}\). Orthogonal projection is
a contraction on the product of the spheres, whose measure is
rotation invariant. Uniform Taylor convergence on this compact
product justifies taking that Fourier projection. The pointwise
estimate thus gives the factor
\(\exp(aR_\tau^2)R_\tau^{-2j_\tau}\) from block \(\tau\) after
rescaling to unit spheres. Its minimum occurs at
\(R_\tau^2=j_\tau/a\), giving \eqref{eq:sphere60}.
For the bound on \(H_l\) apply
\(\sum_{i=0}^{u-1}\log(l+i)\le\int_0^u\log(l+t)\,dt\)
to the factorial ratio and use \(\log(1+l/u)\le l/u\).
\end{proof}

In the following theorem \(r\le s\) are the two ranks with
\(rs\) even, \(l=rs,\ D=fl,\ f=[E:K]\), \(b_\tau\) are the
prescribed minor weights, \(k_\tau=rb_\tau\), \(k=\sum_\tau k_\tau\),
and \(B=\max_\tau b_\tau\). Let \(m_c,m'_h\) be the rational
stabilizer weights, with sums \(M_{\rm in},M_{\rm out}\).
Write \(P_T\) for an ordered product of \(b_\tau\) maximal
minors at each \(\tau\). For the prescribed family of sections put
\[
 v_{h,T}=\sum_c m_c P_T(\partial)F_{ch}(0),\qquad
 \mathcal N=\sum_{h,T}m'_h|v_{h,T}|^2.
\]
Precisely, \(F_{ch}\) is a finite sum of torsion translates of
a normalized vacuum as in Theorem \ref{thm:arch59}. Its
coefficients \(a_{i,ch}\) obey
\((\sum_i|\rho(a_{i,ch})|)^2\le Q^{rs}\) at each embedding.
The estimates use the trace lattice for the
particular pair \(c,h\). If some \(v_{h,T}\ne0\), write
\(\mathfrak h=[J:\QQ]^{-1}\log|N_{J/\QQ}\mathcal N|\) for a
common Galois coefficient field \(J\).

\begin{theorem}[The signed upper bound at the prescribed weights]\label{thm:opt60}
Define
\[
       L_d=\log(\sqrt d/(2\pi))-\gamma+D_{\xi_d},\qquad
       V_d=\tfrac14\log d-\tfrac12\gamma+\tfrac12D_{\xi_d}.
\]
For every \(x\) in the common interval and all the class sections
above one has
\begin{align}
 \mathfrak h\le {}&
       \log(M_{\rm out}M_{\rm in}^2)+rs\log Q
                              + k L_d+D V_d
                              +r\sum_\tau\log(b_\tau!)          \notag\\
  &+D\log\!\left[e\max(1,2rs\,\operatorname{rd}(E)/\pi)\right]
               +\sum_\tau\log H_{rs}(rb_\tau)                  \notag\\
  &+\sum_\tau\sum_{i=0}^{r-1}\log\binom{b_\tau+i}{i}.
                                                       \label{eq:opt60}
\end{align}
The last two sums together are at most
\[
 D\big(1+\log(1+B/s)\big)+\tfrac12fr(r-1)\log(1+B+r).
\]
\end{theorem}
\begin{proof}
Fix an embedding \(\rho\) of \(J\). The metric supremum bound in
Theorem \ref{thm:transport59}, and the coefficient bound give, at each
target class \(h\),
\[
 \left|\sum_c m_c F_{ch}^\rho(Z)\right|^2
 \le M_{\rm in}^2 Q^{rs}
     [e\max(1,2rs\,\operatorname{rd}(E)/\pi)]^D
            \nu_\rho^D\, e^{(\pi/A_\rho)|Z|^2}.
\]
The field embeddings act on the indicated Taylor coefficients.
For the sum of squared moduli of \(\rho(v_{h,T})\), the exact
minor identity in Theorem \ref{thm:frame55}, gives the product
\(\prod_\tau C_{r,b_\tau}\) times the square of a Fischer
orthogonal projection of this Taylor component. Discard the
projection using contraction and apply Lemma~\ref{lem:sphere60}
with \(j_\tau=k_\tau\). Sum over \(h\) with weights \(m'_h\).
Thus the two elliptic factors in the embedding bound are now
\((\pi/A_\rho)^k\nu_\rho^D\).

If \(\iota\) is complex conjugation on \(J\), apply
Cauchy--Schwarz to \(\rho(\mathcal N)\) using the two vectors of
coordinates at \(\rho\) and \(\rho\iota\). On averaging logarithms
each set of coordinates occurs equally often. This bounds
\(\mathfrak h\) by the mean logarithm of the squared-modulus sum
just estimated. Theorems~\ref{thm:mean56} and~\ref{thm:norm58}
give the means \(L_d\) and \(V_d\) for \(\log(\pi/A_\rho)\) and
\(\log\nu_\rho\). Finally
\[
       C_{r,b}=(b!)^r\prod_{i=0}^{r-1}\binom{b+i}{i}.
\]
This proves \eqref{eq:opt60}. Use the bound on \(H_l\) with
\(k_\tau/l=b_\tau/s\), and
\(\log\binom{b+i}{i}\le i\log(1+B+r)\), to obtain the last
estimate.
\end{proof}

\begin{corollary}[The linear term after the factorials are expanded]\label{cor:linear60}
With \(0\log0=0\), the bound in Theorem~\ref{thm:opt60} implies
\begin{align*}
 \mathfrak h \le {}&
  \log(M_{\rm out}M_{\rm in}^2)+rs\log Q+D V_d
       +r\sum_\tau b_\tau\log b_\tau+k(L_d-1) \\
  &\quad +O\!\left(D\log\!\big( (2+rs\,\operatorname{rd}(E))
                                            (2+B+r)\big)\right).
\end{align*}
The implied constant is absolute, including when some weights
are zero.
\end{corollary}
\begin{proof}
Comparison of \(\sum_{j=1}^b\log j\) with the integral of \(\log t\)
gives \(\log(b!)=b\log b-b+O(\log(2+b))\) for every \(b\ge0\).
The sum of the resulting errors is \(O(fr\log(2+B))\).
Use \(r\le s\), \(fr\le D\), and \(fr(r-1)\le D\) in the last
estimate of Theorem~\ref{thm:opt60}.
\end{proof}

\begin{corollary}[Mass and level on the prescribed interval]\label{cor:mass60}
Use the ranks of Theorem~\ref{thm:rangenew50} with
\(r\le s\le x^{402}\). Take the input level to be the full
lattice stabilizer and the output level of index at most
\(Q^{s^2+4}\) in the full stabilizer, as in
\srcloc{Development/AbelianVarietyTheory/Definitions/Foundations/NormSumPairRankZeroOfCurveCoordNorm}{145}.
If \(B\le2\lceil\exp(x^\eps)\rceil\), then uniformly on the interval
\[
 \mathfrak h\le r\sum_\tau b_\tau\log b_\tau+k(L_d-1)
        +(rs+s^2+4)\log Q+O_{\psi,\eps}(x^{1630}) .
\]
For fixed \(\eps>0\) and \(A_0>0\), the polynomial error is
\(o(\exp(x^\eps)/x^{A_0})\).
\end{corollary}
\begin{proof}
For a full-level class with finite stabilizer \(\Gamma_c\),
the finer classes above it correspond to the orbits of
\(\Gamma_c\) on the cosets of the finer level in the full level.
The reciprocal stabilizer weights over those orbits sum to the
index divided by \(|\Gamma_c|\), by the orbit--stabilizer
formula. The input full mass is bounded above by the number of
full-level classes, and the output mass is its own full mass
times its level index.

Write \(e=[E:\QQ]\) and \(E\subset\QQ(\zeta_N)\) for the
conductor multiple chosen earlier in this section. The full-level
part of the proof of
Lemma \ref{lem:mass32}
counts genus classes of rank \(s\) and gives the bound
\(C(se)^4\log(2seN)\) for their logarithmic number. It applies
also in rank \(r\). The result follows there by counting Gram
matrices and matrices of one generator of \(E\) in a bounded
integral basis. Thus
\[
 \log(M_{\rm out}M_{\rm in}^2)
       \le (s^2+4)\log Q + O\!\left((se)^4\log(2seN)\right).
\]
The construction gives \(e\ll_\psi x^3\) and
\(\log N\ll_{\psi,\eps} x^{\eps/5}\log^2(2x)+\log(2x)\).
Moreover \(\operatorname{rd}(E)\le N\) by the
conductor--discriminant formula, and
\(|V_d|\ll_{\psi,\eps} x^{\eps/4}\) by Theorem~\ref{thm:twists}.
Therefore \(se\ll_\psi x^{405}\), \(D\ll_\psi x^{807}\),
and each logarithmic factor in the error of
Corollary~\ref{cor:linear60} is \(O_{\psi,\eps}(x)\).
Combining these estimates proves the indicated polynomial bound.
Finally \(x^\eps/(\log x)\to\infty\), which proves the last claim.
\end{proof}

This bound retains the signed term in \(k\) needed when comparing
upper bounds at infinity with finite divisibilities. To
derive \eqref{eq:gap} the comparison must reach the specific
integral class sections in the iterated lifts and their nonzero
row periods. The complementary odd-prime groups in
Proposition~\ref{prop:res23} also require a uniform estimate.

\subsection{The first inert-prime congruence in the Hermitian frame}\label{nav:V-23}
Let \(p\ge5\) be inert in \(K\). Work over the completed integer
ring \(R\) of a finite extension of \(\QQ_p\) containing the
coefficients and torsion under consideration. Normalize \(v_p(p)=1\).
The discriminant-one elliptic curve \(C\) has good reduction by
Theorem \ref{thm:germ57}.
Let \(\partial\) be the invariant vector field dual to \(dx/y\).
For any nonzero \(t\in C[2]\), put
\(\kappa_t=-(\partial^{p-1}x)(t)\). Here the coordinate \(x\)
in the Weierstrass equation is a function on \(C\).

\begin{lemma}[A nonzero constant on the supersingular fiber]\label{lem:kappa62}
On the good model over \(R/pR\), one has \(\partial^p=0\) on
functions. On the geometric special fiber \(\partial^{p-1}x\)
is a nonzero constant. In particular \(\kappa_t\in R^\times\).
\end{lemma}
\begin{proof}
In characteristic \(p\) the \(p\)-fold iterate of a derivation is
a derivation, by Leibniz's rule. The iterate of an invariant vector
field is invariant, and hence \(\partial^p=c\partial\) over
\(R/pR\), where \(c\in R/pR\). These assertions hold over this
ring even when it has nilpotents, because they follow from group
translation and the binomial coefficients. Choose a CM scalar
\(\alpha\in\mathcal O_K\) whose residue generates
\(\mathbb F_{p^2}\) over \(\mathbb F_p\), using inertness.
Its isogeny has degree prime to \(p\) and extends to the good
model after finite extension. Its tangent multiplier is the
chosen image \(\lambda\) of \(\alpha\). Pullback of functions
intertwines \(\partial\) with \(\lambda\partial\).
Iterating \(p\) times and using \(\partial^p=c\partial\)
gives \(c(\lambda^p-\lambda)=0\). One can evaluate the identity
on a parameter at the origin, whose first derivative is a unit.
Since \(\lambda^p-\lambda\) is a unit, \(c=0\).

Pass to the algebraic closure \(k_0\) of the residue field and
write \(L=k_0(C)\). The constants of \(\partial\) in \(L\)
are \(L^p\) and \([L:L^p]=p\). To see this directly, \(x\) is
separating, \(y=y^p/(4x^3-g_2x-g_3)^{(p-1)/2}\), and
\(L=L^p(x)\). The element \(x\) is not a \(p\)-th power since
\(\partial x=y\ne0\). In the basis \(1,x,\ldots,x^{p-1}\)
the ordinary derivative in \(x\) has kernel \(L^p\), as does
\(\partial=y\,d/dx\). Thus \(\partial\) on this \(p\)-dimensional
space is nilpotent with one Jordan block. In particular
\(\ker(\partial^{p-1})=\operatorname{im}(\partial)\).

Suppose \(\partial^{p-1}x=0\). Both \(1\) and \(x\) would then
be in \(\operatorname{im}(\partial)\), so the differentials
\(dx/y\) and \(x\,dx/y\) would both be exact in \(L\).
Use a linear change of \(x\) and rescale \(y\) to obtain
\(Y^2=X(X-1)(X-\lambda_0)=f(X)\) with \(\lambda_0\ne0,1\).
The two differentials span the same \(k_0\)-plane as
\(dX/Y\) and \(X\,dX/Y\). Write \(m=(p-1)/2\). Since
\(1/Y=f(X)^m/Y^p\), exactness for both requires the
coefficients of \(X^{p-1}\) and \(X^{p-2}\) in \(f(X)^m\)
to vanish. Indeed \(d/dX\) on the \(L^p\)-basis through
\(X^{p-1}\) misses precisely the last basis vector, and
\(\deg(X f(X)^m)<2p-1\).

Put \(P(X)=(X-1)^m(X-\lambda_0)^m=\sum_j a_jX^j\), so the two
vanishing coefficients give \(a_m=a_{m-1}=0\). Logarithmic
differentiation of this polynomial gives
\[
  \lambda_0(j+1)a_{j+1}+(m-j)(1+\lambda_0)a_j
                         +(j-1-2m)a_{j-1}=0 .
\]
For \(1\le j\le m-1\) the last coefficient is \(j\) in \(k_0\).
Starting at \(j=m-1\) and descending forces \(a_0=0\), whereas
\(a_0=\lambda_0^m\ne0\). This proves \(\partial^{p-1}x\ne0\).
It is a \(p\)-th power in \(L\) because its derivative is zero.
It is regular off the origin and has a pole there of order at
most \(p+1\), since \(x\) has order two and the invariant vector
field increases pole orders by at most one. Its \(p\)-th root
therefore has at most a simple pole at the origin and no other
poles. Riemann--Roch in genus one gives only constants with this
pole bound. We have proved that \(\partial^{p-1}x\) is a nonzero
constant. The mark specializes away from the origin, proving the
unit assertion.
\end{proof}

\begin{theorem}[A recurrence for the actual elliptic theta derivatives]\label{thm:hurwitz62}
Let \(F=U_w\Theta_t\), where \(w\) has order prime to \(p\)
on \(C\), and include \(w=0\). Write \(H_n(F)=F^{(n)}(0)\)
in the normalized Hermitian frame of Proposition~\ref{prop:ellnorm56}.
For \(n\ge0\), with the right side interpreted as zero at \(n=0\),
\begin{equation}
             H_{p+n}(F)\equiv
                \kappa_t\, n\,H_{n-1}(F) \pmod {pR}.       \label{eq:rec62}
\end{equation}
For the unshifted vacuum and all \(0\le j<p\),
\[
           H_{j(p+1)}(\Theta_t)\equiv j!\,\kappa_t^j\pmod{pR}.
\]
In particular \(H_{p+1}(\Theta_t)\) is a unit. For every indicated
translate, \(H_n(F)\in pR\) for every \(n\ge p^2\).
\end{theorem}
\begin{proof}
It is important here to distinguish a sequence of undivided
derivatives from coefficients of an ordinary power series.
We use the ring of sequences over \(R\) in which multiplication is
\[
        (ab)_n=\sum_{i=0}^n\binom ni a_i b_{n-i}.
\]
The shift \(D(a)_n=a_{n+1}\) is a derivation. Taking all
undivided derivatives of a germ at zero respects these operations.
Let \(Z=(0,1,0,\ldots)\) represent the linear coordinate.
Modulo \(pR\) the \(p\)-th power of a sequence is the constant
sequence with value \(a_0^p\). Its shift vanishes by Leibniz's rule,
and its degree-zero term has that value.

For \(F_0=\Theta_t\) let \(\ell=(\log F_0)'\).
All derivatives of both functions are integral at \(p\), by
the algebraic formulas of Theorem \ref{thm:germ57}.
In particular \(\ell(0)=0\),
\(\ell'=-u_\Lambda-x(z+t)\) with \(u_\Lambda\in R\), and
\(F_0\) is even. Hence \(\ell^{(p-1)}(0)=0\) and
\(\ell^{(p)}(0)=\kappa_t\). Every subsequent derivative of
\(\ell\) vanishes modulo \(pR\) at zero, by the first assertion
of Lemma~\ref{lem:kappa62}. In the sequence ring this means
\(D^{p-1}\ell=\kappa_t Z\) after reduction.

From \(DF_0=\ell F_0\) one obtains
\[
       D^p F_0=(D^{p-1}\ell+\ell^p)F_0\pmod {pR}.
\]
For an elementary proof, repeat the product rule. The terms in
\(D^p F_0/F_0\) correspond to set partitions of \(p\) positions,
with a block of size \(i\) contributing \(D^{i-1}\ell\).
This follows by adjoining one position, either in a new block
or in an existing block. The cyclic permutation of the \(p\)
positions groups the terms into orbits of size \(p\), except for
the partition into singletons and the partition with one block.
Indeed the blocks in any translation-invariant partition of
\(\ZZ/p\ZZ\) are cosets of a subgroup. This proves the formula.
Now \(\ell^p=0\) after reduction since \(\ell(0)=0\).
Thus \(D^p F_0=\kappa_t ZF_0\) in the reduced sequence ring.

For a translate, the proof of Theorem \ref{thm:trans57}, gives
\(F=R_wF_0\) with \(R_w\) integral in a formal parameter of the
smooth elliptic model. The invariant derivative \(\partial^p\)
annihilates this formal ring modulo \(pR\) by Lemma~\ref{lem:kappa62}
and continuity. Every derivative after the first \(p\) of
\(R_w\), starting at order \(p\), is consequently zero modulo
\(pR\). Since \(D^p\) is a derivation in characteristic \(p\),
we get \(D^p F=\kappa_t ZF\). The \(n\)-th entry of \(ZF\)
is \(nH_{n-1}(F)\), giving \eqref{eq:rec62}.
Apply it at \(n=j(p+1)-p\) and use
\(j(p+1)-p\equiv j\bmod p\), starting with \(F_0(0)=1\).
This gives the factorial congruence. Finally the derivation
\(D^p\) kills \(Z\), so iterating it \(p\) times gives
\(D^{p^2}F=\kappa_t^pZ^p F=0\) modulo \(pR\).
\end{proof}

\begin{corollary}[Failure of a scalar integral frame for this germ]\label{cor:obstacle62}
At every such inert prime the analytically normalized
\(\Theta_t(z)\), expressed in a good elliptic formal parameter,
does not belong to its integral power series ring. Multiplication
by any unit in that ring preserves this obstruction.
\end{corollary}
\begin{proof}
For an integral scalar series every derivative of order at least
\(p\) is zero modulo \(pR\), by the first part of
Lemma~\ref{lem:kappa62}. The order \(p+1\) of \(\Theta_t\)
has unit value by Theorem~\ref{thm:hurwitz62}. Finally a unit
and its inverse both preserve the integral power series ring.
\end{proof}

This calculation takes place on the elliptic curve underlying the
same class kernels as in Theorem~\ref{thm:opt60}. It explains why
the scalar integral-section estimates in this section cannot be
applied directly to the analytic germ of the prescribed vacuum.
Equation~\eqref{eq:rec62} gives a comparison in its actual
Hermitian frame.

\subsection{Higher powers at an inert prime}\label{nav:V-24}
Keep the elliptic curve, mark, invariant differential and Hermitian
normalization of Theorem~\ref{thm:hurwitz62}. Write \(Z\) for the
additive coordinate in which this differential is \(dZ\).
For \(n\ge0\) put
\[
    F_p(n)=v_p(n!)=\sum_{a\ge1}\left\lfloor n/p^a\right\rfloor,
       \qquad \omega_p=\frac{p}{p^2-1}.
\]
Throughout, \(v_p(0)=+\infty\). We first locate the coefficient of
the logarithm which determines the radius of the normalized theta
function. All valuations are taken with \(v_p(p)=1\).

\begin{lemma}[The lowest coefficients of the scaled logarithm]\label{lem:scale64}
Put \(L_j=(d/dZ)^j\log\Theta_t(0)\), meaning the derivative at zero
of the formal logarithm with constant term zero. The odd \(L_j\)
vanish. For each even \(j\ge2\) set \(h_j=\lfloor(j-2)/p\rfloor\).
Then \(L_j\in p^{h_j}R\), \(L_{p+1}\) is a unit, and
\begin{equation}
      h_j-F_p(j)+j\omega_p\ \ge\ \frac{1}{p-1}.
                                                \label{eq:lowest64}
\end{equation}
Equality in the numerical inequality occurs exactly at \(j=p+1\).
\end{lemma}
\begin{proof}
Near the marked point, \(g=-u_\Lambda-x\) is an integral regular
germ and \(L_j=(\partial^{j-2}g)(t)\). This follows by differentiating
\((\log\Theta_t)''(Z)=-u_\Lambda-x(Z+t)\).
The invariant derivation preserves the completed local ring at \(t\).
Lemma~\ref{lem:kappa62} gives \(\partial^p R[[T]]\subset pR[[T]]\)
in a parameter there. By \(R\)-linearity, each block of \(p\)
derivatives gains a factor of \(p\). The parity follows from
evenness, and \(L_{p+1}=\kappa_t\) is the unit in the same lemma.

If \(2\le j<p\), then \(F_p(j)=0\) and
\(2\omega_p>1/(p-1)\). For \(p+1\le j<2p\) we have
\(F_p(j)=1\) and \((p+1)\omega_p-1=1/(p-1)\).
For \(j\ge2p\), the base-\(p\) digit formula gives
\(F_p(j)\le(j-1)/(p-1)\). Also
\[
       h_j\ge (j-p-1)/p>j/(p^2-1).
\]
The last strict inequality holds at \(j=2p\) for \(p\ge5\),
and the left side increases faster thereafter. Using
\(1/(p-1)-\omega_p=1/(p^2-1)\) proves the strict inequality
in this range.
\end{proof}

\begin{theorem}[Valuations and the radius of the elliptic theta series]\label{thm:powers64}
For every prime-to-\(p\) torsion point \(w\), including zero, and
every \(n\ge0\), the undivided derivatives in the analytic frame
satisfy
\begin{equation}
 v_p\big(H_n(U_w\Theta_t)\big)\ \ge     \max\{0,F_p(n)-\lfloor\omega_p n\rfloor\}.             \label{eq:higher64}
\end{equation}
For the marked section itself, put \(n_e=(p+1)p^e\), \(e\ge0\).
Then
\begin{equation}
     v_p\big(H_{n_e}(\Theta_t)\big)=\frac{p^e-1}{p-1}.    \label{eq:exact64}
\end{equation}
In particular its ordinary Taylor series in \(Z\), over the completed
algebraic closure of \(\QQ_p\) with \(|p|_p=p^{-1}\), has radius
\(p^{-\omega_p}\).
\end{theorem}
\begin{proof}
Let \(A=\log\Theta_t\). Its constant term is zero and its first
possible degree is two. The coefficient of \(Z^n\) in \(\exp A\)
is a finite sum over multiplicities \(m_j\ge0\) with
\(\sum j m_j=n\), indexed by even \(j\ge2\). The valuation of
each summand is at least the integer
\[
     I=\sum_j m_j(h_j-F_p(j))-\sum_j F_p(m_j).
\]
By \eqref{eq:lowest64} and \(F_p(m)\le m/(p-1)\), we have
\(I\ge-\omega_p n\). Thus every summand belongs to
\(p^{-\lfloor\omega_p n\rfloor}R\), even if \(R\) is ramified.
Multiplying the coefficient by \(n!\) and using the integrality
of undivided derivatives in Theorem~\ref{thm:hurwitz62} gives
\eqref{eq:higher64} for \(w=0\).

For a translate write \(U_w\Theta_t=R_w\Theta_t\).
Theorem~\ref{thm:trans57}
constructs \(R_w\) as an integral series in a good elliptic parameter.
It follows from the same block of \(p\) derivatives that
\(v_p(R_w^{(j)}(0))\ge\lfloor j/p\rfloor\).
If \(j<p\), \(F_p(j)=0\). For \(j\ge p\), putting \(a=\lfloor j/p\rfloor\)
gives \(j\le pa+p-1\le (2p-1)a\le(p^2-1)a\). Therefore
\[
 F_p(j)-\lfloor j/p\rfloor
       \le j/(p-1)-j/(p^2-1)=\omega_p j .
\]
Its left side is an integer. Coefficient multiplication now gives the
same denominator bound for the product, since
\(\lfloor\omega_p j\rfloor+\lfloor\omega_p(n-j)\rfloor
\le\lfloor\omega_p n\rfloor\). The integrality of the undivided
derivatives of the translate completes this part of the proof.

For equality, make a finite extension containing elements \(a,b\)
of valuations \(\omega_p,1/(p-1)\), respectively. Set
\(U(Z)=b^{-1}A(aZ)\). By Lemma~\ref{lem:scale64} its coefficients
are integral and its reduction modulo the maximal ideal is exactly
\(cZ^{p+1}\) for a nonzero residue \(c\). The equality assertion
in that lemma and the unit value of \(L_{p+1}\) are both used here.
For \(m>0\) the valuation of \(b^m/m!\) is
\(s_p(m)/(p-1)\), where \(s_p(m)\) is the sum of the base-\(p\)
digits. This follows by summing the floors in \(F_p(m)\).
In degree \(n_e\), expand \(\Theta_t(aZ)=\sum_m b^m U(Z)^m/m!\).
Only finitely many \(m\) occur in this coefficient.
A term of valuation \(1/(p-1)\) requires \(m\) to be a power of
\(p\) and the coefficient of \(Z^{n_e}\) in \(U^m\) to be a unit.
Its reduction is determined by \(c^mZ^{(p+1)m}\), so precisely
\(m=p^e\) gives that valuation. All other terms have strictly larger
valuation. Undoing the scale gives
\[
 v_p(H_{n_e}(\Theta_t))=F_p(n_e)-\omega_p n_e+\frac1{p-1}.
\]
Here \(s_p(n_e)=2\), which gives \eqref{eq:exact64}. The lower
bound for all ordinary coefficients and equality on this unbounded
sequence give \(\liminf_{n\to\infty}n^{-1}v_p([Z^n]\Theta_t)
=-\omega_p\). The nonarchimedean root test proves the radius claim.
\end{proof}

For example \(H_{p+1}\) is a unit, \(H_{p(p+1)}\) has valuation
one, and \(H_{p^2(p+1)}\) has valuation \(p+1\). The units and
the higher powers occur in the same normalization that defines the
archimedean kernel in Theorem~\ref{thm:opt60}. In several variables
an integral change of coordinates preserves an ordinary coefficient
bound in the total degree. In contrast, the factorial recovered on
taking a derivative depends on the separate coordinate orders.

\subsection{Odd powers of Frobenius in the class-section derivatives}\label{nav:V-25}
Use the data of Sections \ref{nav:V-14}--\ref{nav:V-16}. Fix \(x\) in the common interval, one set \(S\)
of the indicated inert primes and the one elliptic product basis
integral at all of them. For \(p\in S\) put
\[
 F_{p^2}(n)=\sum_{j\ge1}\lfloor n/p^{2j}\rfloor,\qquad
 B_p(n)=F_p(n)-F_{p^2}(n)=\sum_{j\ge1,\ j\ {\rm odd}}
                                           \lfloor n/p^j\rfloor .
\]
Let \(\Sigma\) be the CM embeddings over \(K=\QQ(\sqrt{-d})\);
arithmetic Frobenius \(\sigma_p\) and complex conjugation \(c_E\)
give a permutation \(\pi=\sigma_p c_E\) of \(\Sigma\). Each
even cycle of \(\pi\) splits into two orbits under \(\pi^2\).
We call these its two halves. For example when \(K_q\) splits at
\(p\), \(\pi\) interchanges the two restrictions of embeddings
to \(K_q\), since \(c_E\) acts there as conjugation. In this case
every cycle is even.

\begin{lemma}[Elliptic parameters and a logarithmic derivative]\label{lem:multijet67}
Let \(c_{i\nu}\in R\) and \(u_i=\sum_\nu c_{i\nu}Z_\nu\).
For an integral germ \(H\) in the elliptic parameters \(T_i\)
the ordinary coefficient of \(Z^\beta\) in
\(H(T(u_1),\ldots,T(u_n))\) has valuation at least
\(-\sum_\nu F_{p^2}(\beta_\nu)\). Put
\(\ell=(\log\Theta_{t_*})'\). For one integral linear form \(u\)
the coefficient of \(Z^\beta\) in \(\ell(u)\), for \(\beta\ne0\),
has valuation at least
\begin{equation}
 -\sum_\nu F_{p^2}(\beta_\nu)-
    \min_{\nu:\,\beta_\nu>0}\#\{j\ge1:\ j\text{ odd},\ p^j\mid\beta_\nu\}.
                                                       \label{eq:loss67}
\end{equation}
Its constant coefficient vanishes.
\end{lemma}
\begin{proof}
The formal exponential in an integral elliptic parameter with
unit linear coefficient has \(v_p([Z^m]T(Z))\ge-F_{p^2}(m)\).
We recall the hypotheses of the factorial argument in
Lemma \ref{lem:factorial41}.
The formal CM endomorphisms extend from \(\mathcal O_K\) to
\(O=\mathcal O_{K_p}\) by coefficientwise approximation, as in
Corollary \ref{cor:class66}. Choose a parameter of linear coefficient
one and average \(\zeta^{-1}[\zeta](T)\) over
\(\zeta\in\mu_{p^2-1}\subset O^\times\), dividing by \(p^2-1\).
This gives an integral change with linear coefficient one in which
these roots of unity act linearly. Commutation with them forces
the first possible nonlinear degree in \([p](T)\) to be \(p^2\).
The cited lemma gives the coefficient bound in this coordinate.
The functions \(F_p,F_{p^2},B_p\) are superadditive on nonnegative integers, by the floor formulas;
in particular integral postcomposition preserves the bound for
\(T(Z)\). Translation of an integral regular germ at the mark
to the origin also gives an integral elliptic germ.

For a monomial of degree \(m=|\beta|\) obtained from \(u^m\), the
multinomial coefficient adds valuation
\(F_p(m)-\sum_\nu F_p(\beta_\nu)\). Superadditivity of \(B_p\)
gives the required lower bound for each \(T(u_i)\).
Products preserve the bound since each \(F_{p^2}\) is
superadditive. In any fixed degree only finitely many terms
of \(H\) enter.

Now \(\ell(0)=0\) and \(\ell'=-u_\Lambda-x(Z+t_*)\) is an
integral elliptic germ. Consequently \(v_p(\ell^{(m)}(0))\ge
B_p(m-1)\) for \(m\ge1\), by the first assertion for this
one-variable germ followed by multiplication by \((m-1)!\).
For each \(\beta_\nu>0\), write \(m-1\) as the sum of the
coordinates of \(\beta-e_\nu\). Superadditivity gives
\[
 B_p(m-1)\ge\sum_\mu B_p(\beta_\mu)-
     \big(B_p(\beta_\nu)-B_p(\beta_\nu-1)\big).
\]
The difference in parentheses counts the odd \(j\) with
\(p^j\mid\beta_\nu\). The coefficient after linear substitution is
\(\ell^{(m)}(0)c^\beta/\beta!\), which proves \eqref{eq:loss67}.
\end{proof}

Write \([y]_+=\max(0,y)\). A matrix coordinate at \(\tau\) has
indices \((\tau,i,j)\), with row \(i\) and column \(j\).
For every even cycle \(C\), and separately for every pair \(i,j\),
choose one half \(H\) of \(C\) and write \(L=C\setminus H\).
In each expression below, \(H,L\) refer to that choice in the
indicated block.

\begin{theorem}[Valuations on the two halves of a Frobenius cycle]\label{thm:ramp67}
Let \(F_{ch}\) be any of the prescribed single-lift class sections
in the matrix coordinates. At every place over the chosen \(p\)
and for every multi-index \(\alpha\), one has
\begin{align}
 v_p(\partial^\alpha F_{ch}(0))\ \ge \sum_{i,j}\sum_{C\ {\rm even}}\sum_{a\ge1,\ a\ {\rm odd}}
  \left[\sum_{\tau\in H}
               \left\lfloor\frac{\alpha_{\tau ij}}{p^a}\right\rfloor
                        -\sum_{\tau\in L}\alpha_{\tau ij}\right]_+ .
                                                     \label{eq:ramp67}
\end{align}
All choices of halves are allowed in this inequality.
\end{theorem}
\begin{proof}
First decompose the trace lattice at \(p\). The local lattice
of the standard vacuum is a copy of
\((\mathcal O_E\otimes\ZZ_p)^{rs}\), since \(E\) is unramified
and the relative different is a unit at \(p\). The chosen
orthonormal basis over \(K\) gives an integral basis with integral
inverse here. The factors of \(\mathcal O_E\otimes\ZZ_p\), as an
\(O\)-algebra, correspond to orbits of \(\sigma_p^2=\pi^2\) on
\(\Sigma\). Indeed \(K_p/\QQ_p\) is unramified quadratic and its
arithmetic Frobenius acts on these embeddings by \(\sigma_p^2\).
After unramified base change splitting \(E\), the homomorphism
given by the embeddings identifies each local factor with the
product of the coordinate rings in its orbit. This follows from
the factorization of a separable residue polynomial and Hensel
lifting, or from the integral splitting of a finite unramified
algebra. Thus a basis of each local factor gives an invertible
integral coordinate change between it and the corresponding
matrix variables after that extension.

For fixed \(i,j\), a cycle of \(\pi\) gives an orthogonal block
for the trace Hermitian form. On an even cycle its two orbit
lattices are isotropic and pair with each other. In fact the form
pairs a local field factor only with its conjugate factor.
Complex conjugation takes the \(\pi^2\)-orbit to the orbit
described on \(\Sigma\) by \(\sigma_pc_E=\pi\): the extra
Frobenius returns to the fixed embedding of \(K\) and fixes
the place of the normal extension \(E\). Equivalently for
\(y\in E\), one has
\(\tau(c_Ey)=\sigma_p^{-1}((\pi\tau)(y))\); substituting this
identity in the trace pairing verifies the claim. Distinct matrix
entries are orthogonal by the matrix trace.

Choose \(O\)-bases of all these orbit lattices, adjoining the
two bases within each even block. Write the full matrix of
columns in the original elliptic product basis as \(A\); it and
its inverse are integral over \(O\). If \(u\) denotes the new
variables, the product vacuum is \(G_A(u)=\prod_i\Theta_{t_*}((Au)_i)\).
For a cycle block \(b\) consisting of one cycle and one matrix
entry, let \(A_b\) denote its columns, and write \(u_b\) for
its variables. Orthogonality says that the Gram matrix of \(A\)
is the sum of those of the matrices \(A_b\) extended by zero
outside their own columns. By
Theorem \ref{thm:gram66} and Corollary \ref{cor:class66}, the actual section has the factorization
\[
             F_{ch}=U(T(u))\prod_b G_{A_b}(u_b),
             \qquad U\in R[[T_1,\ldots,T_D]] .
\]
Here \(D=rs[E:K]\). Each block factor, after the variables in
the unchosen half of an even block are set to zero, is an
integral elliptic unit. This follows from the zero Gram matrix
on the chosen half and the same theorem. The basis change from
\(u\) to the matrix variables and its inverse preserve each
\(\pi^2\)-orbit and are integral after the unramified extension
just described.

Differentiate first in all variables of the unchosen halves
and in all the variables of the odd cycles, to their prescribed
orders. On \(U\) any such derivative is an integral linear
combination of the invariant derivatives in the elliptic
parameters, so it preserves integrality. On an even block write
\(\ell=(\log\Theta_{t_*})'\). Differentiating its factor once
multiplies it by an integral linear combination of
\(\ell((A_bu_b)_i)\). Successive differentiations of these
logarithmic derivatives contribute integral elliptic germs,
because \(\ell'\) is such a germ and each row of \(A_b\) is over
\(O\). Iterating the product rule gives terms with at most
\(l_b=\sum_{\tau\in L}\alpha_{\tau ij}\) factors of \(\ell\)
from this block, times the original \(G_{A_b}\) and integral
elliptic germs. On an odd block every undivided derivative of
\(G_{A_b}\) at zero is integral, by Theorem~\ref{thm:hurwitz62}
and the integral linear change of variables.

Set all the differentiated variables to zero. Every remaining
term is now an integral elliptic germ in the chosen halves
times at most \(l_b\) factors of \(\ell\) in the variables of
each even block \(b\). The arguments of these factors are
integral linear forms in the matrix variables of that block.
Apply Lemma~\ref{lem:multijet67} and extract the coefficient
in the remaining coordinates. The ordinary factorial denominator
from all factors together is bounded by
\(\sum_{b,\tau\in H}F_{p^2}(\alpha_{\tau ij})\), by superadditivity.
For fixed odd \(a\), an \(\ell\)-factor contributes one additional
unit of denominator in \eqref{eq:loss67} at this level only if
all its positive monomial exponents are divisible by \(p^a\).
It then consumes at least \(p^a\) in one chosen coordinate. Since
the total in each coordinate is prescribed, the number of such
factors is at most
\[
       \min\left(l_b,\sum_{\tau\in H}
                             \lfloor\alpha_{\tau ij}/p^a\rfloor\right).
\]
No zero-degree \(\ell\)-factor contributes. Finally multiply by
the factorials of the remaining coordinate orders to recover
their undivided derivative. Subtraction of the \(F_{p^2}\) terms
leaves \(\sum B_p(\alpha_{\tau ij})\). Subtracting the last
display, level by level in each block, gives \eqref{eq:ramp67}.
This argument applies to every summand, hence to the full derivative.
\end{proof}

\begin{corollary}[Truncated prime-power terms in the minor projections]\label{cor:minor67}
Put \(d_P(b,s)=\max(0,\lceil(b-s(P-1))/P\rceil)\). For each
even cycle \(C\) write \(C^+,C^-\) for its halves. At every
place over \(p\), the valuation of every actual
\(P_T(\partial)F_{ch}(0)\) and of every \(A_r v_{h,T}\) is at least
\begin{equation}
 t_p=r\sum_{C\ {\rm even}}\max_{e\in\{+,-\}}
    \sum_{a\ge1,\ a\ {\rm odd}}
 \left[\sum_{\tau\in C^e}d_{p^a}(b_\tau,s)
                             -\sum_{\tau\in C^{-e}}b_\tau\right]_+ .
                                                       \label{eq:minor67}
\end{equation}
Here \(-e\) denotes the opposite half. The valuation of
\(Y=A_s A_r^2\mathcal N\) is at least \(2t_p\). The lower bound
in Theorem \ref{thm:joint65}, holds with \(m_p\) replaced at
each prime by \(\max(m_p,t_p)\) on the same interval.
\end{corollary}
\begin{proof}
In one column each minor monomial has row sum \(b_\tau\) at
\(\tau\). Its sum of quotients on division by \(P\) is at least
\(d_P(b_\tau,s)\), by bounding the \(s\) remainders by \(P-1\).
Choose in every row and column the same orientation of a given cycle
that maximizes \eqref{eq:minor67}. For each level sum the positive
parts of \eqref{eq:ramp67} over rows, using
\(\sum_i[y_i]_+\ge[\sum_i y_i]_+\). The other half has row
sum exactly \(\sum_{\tau\in C^{-e}}b_\tau\). Sum over the \(r\)
columns and the cycles, then take the integral linear combination
specified by the minors. Class summation with the cleared
weights and complex conjugation at all places over \(p\) give
the result for \(Y\), as in Theorem \ref{thm:joint65}.
\end{proof}

For a two-element cycle, one row and one column, the formula
contains \(\sum_{a\ {\rm odd}}[\lfloor b_+/p^a\rfloor-b_-]_+\).
The term at \(a=1\) can therefore have order \(b_+/p\), and
the higher odd powers retain the same cutoff. Formula~\eqref{eq:minor67}
depends on the actual weights in a single lift. For the full
period argument one must still sum the gains after the prescribed
maps, compare the resulting nonzero Bessel projection with its
archimedean norm, and establish the odd-primary length estimate on
these fields.

The supplied source uses a Gram identity to pass to an isotropic
family and to a pair of dual isotropic families in
\srcloc{Formalism/SharpRestriction/GlobalSchurBounds/GlobalRankLatticeProofs/UnitTraceSumBounds/PhIsotropicOf}{127}
and line~132 of that file. The reduction in a dual basis is
spelled out in
\srcloc{Formalism/PeriodIntegrals/UnitaryGroupMassFormula/DualGramSubstitution/PhDualGramMain}{11}.
The proof above uses the Gram comparison of Section~\ref{nav:V-15} and counts
the logarithmic derivatives directly.

\subsection{Exact coefficients at the odd powers of Frobenius}\label{nav:V-26}
Keep a good prime \(p\ge5\) from Theorem~\ref{thm:ramp67}.
Put \(Q_2=p^2\), \(F_{Q_2}(n)=\sum_{i\ge1}\lfloor n/Q_2^i\rfloor\)
and \(\lambda(z)=(\log\Theta_{t_*})'(z)\). Write
\(g_m=[z^m]\lambda(z)\). For the marked elliptic curve the
possible exceptions in
Lemma \ref{lem:add76}
can all be evaluated in the same additive coordinate.

\begin{theorem}[The exceptional coefficients of the marked theta]
\label{thm:sharp78}
For every \(i\ge0\), putting \(m=p^{2i+1}\), one has
\[
                 v_p(g_m)=-F_{Q_2}(m)-1.
\]
For every other positive \(m\), \(v_p(g_m)\ge-F_{Q_2}(m)\).
\end{theorem}
\begin{proof}
The second assertion was proved in Lemma \ref{lem:add76}, using
the integral elliptic germ \(\lambda(u+v)-\lambda(u)-\lambda(v)\).
Iterate this addition identity with \(u=jz,\ v=z\) for
\(1\le j<p\). Multiplication by any integer is an integral
endomorphism of the elliptic formal group. Thus
\[
                 \lambda(pz)-p\lambda(z)=S(T(z))
\]
with \(S\) an integral series in an elliptic parameter \(T\).
Use the parameter in
Lemma \ref{lem:unit41},
with linear coefficient one. Its inverse under composition is
integral through degree \(p\), because the ordinary coefficients
of \(T(z)\) through that degree are integral by Lemma~\ref{lem:multijet67}.
The coefficient in degree \(m\) of the left side is
\((p^m-p)g_m\). All coefficients in degrees \(m<p\) are in
\(pR\), including the zero constant term. In degree \(p\) the
coefficient is a unit, since \(p!g_p=\kappa_{t_*}\) by
Lemma~\ref{lem:kappa62} and \((\log\Theta_{t_*})''=-u_\Lambda-x(z+t_*)\).
Composition with the inverse parameter shows that the coefficients
of \(S\) below degree \(p\) lie in \(pR\) and \(S_p=[T^p]S\)
is a unit.

Adjoin \(\rho\) with \(\rho^{Q_2-1}=p\), and reduce modulo the
maximal ideal using a bar. Lemma \ref{lem:unit41} gives
\[
  H(Z)=T(\rho Z)/\rho\text{ integral},\qquad
       \overline H(Z)=\sum_{i\ge0}h_i Z^{Q_2^i},
                   \qquad h_i\ne0 .
\]
This follows from the multiplication-by-\(p\) series. In the
chosen parameter it begins with \(pT+b_{Q_2}T^{Q_2}\), where
\(b_{Q_2}\) is a unit. Rescaling its exponential identity and
reducing gives \(Z=\overline H+\overline b_{Q_2}\overline H^{Q_2}\),
whose successive coefficients give the displayed series.
On substituting \(z=\rho Z\), the reduction of
\(\rho^{-p}S(T(\rho Z))\) is \(\overline S_p\,\overline H^p\).
Indeed the coefficients of \(S\) below \(p\) are divisible by
\(p=\rho^{Q_2-1}\), and the terms above \(p\) are multiplied
by positive powers of \(\rho\).

All coefficients of \(\rho^{-p}\lambda(p\rho Z)\) have positive
valuation. For \(m=1\) this follows by integrality of \(g_1\).
For \(m\ge2\) use
\(F_{Q_2}(m)\le m/(Q_2-1)\) and the bound
\(v_p(g_m)\ge -F_{Q_2}(m)-1\) from Volume~XIII. After the
substitution the valuation is at least
\(m-1-p/(Q_2-1)>0\).
Consequently
\[
       \overline{p\rho^{-p}\lambda(\rho Z)}
              =-\overline S_p\sum_{i\ge0}h_i^p Z^{pQ_2^i}.
\]
Every displayed coefficient is nonzero. At \(m=pQ_2^i\),
\((m-p)/(Q_2-1)=F_{Q_2}(m)\). Undoing the scale proves
the first assertion.
\end{proof}

\begin{corollary}[A mixed derivative which detects the individual translate]
\label{cor:detect78}
In one matrix entry let \(\nu,\mu\) lie in opposite halves
of an even cycle and \(\pi=\sigma_pc_E\). For \(m=p^{2i+1}\) put
\(u_i=p^{F_{Q_2}(m)+1}g_m\), a unit. If \(F\) is any of the
class sections of Section \ref{nav:V-14}, then
\[
 \frac{p^{F_{Q_2}(m)+1}}{m!}\,
       (\partial_\nu^{\,m}\partial_\mu F)(0)
       \equiv u_i F(0)\,
          \mathbf1_{\{\mu=\pi^{2i+1}\nu\}}\pmod {pR}.
\]
In particular if \(F(0)\) is a unit, the indicated derivative
has valuation \(v_p(m!)-F_{Q_2}(m)-1\) at the matching
translate and at least \(v_p(m!)-F_{Q_2}(m)\) at every other
such \(\mu\). This holds in particular for the original product
vacuum in its orthonormal elliptic coordinates, expressed in
the matrix variables.
\end{corollary}
\begin{proof}
On the line \(Z_\nu=z\) with all other variables zero, write
\(A(z)\) for the restriction of \(F\). The factorization in
Theorem \ref{thm:direct76}, and the product rule give
\[
        (\partial_\mu F)|_{\nu}= B(z)+
             A(z)\sum_k\overline V_{\mu k}
                               \lambda(\overline V_{\nu k}z).
\]
The coefficients of both \(A\) and \(B\) in degree \(n\) have
valuation at least \(-F_{Q_2}(n)\). For \(A\) use isotropy
of the chosen half. For \(B\) the derivative falls on the
integral factor of the block decomposition, and the same
restriction applies. These statements also apply directly
to the product vacuum, whose value at zero is one.

Multiply the degree \(m\) coefficient by \(p^{F_{Q_2}(m)+1}\).
The term from \(B\) and every nonexceptional coefficient in
the product vanish modulo \(pR\) by superadditivity.
If an exceptional coefficient has degree \(n=pQ_2^a<m\), the
floor at \(Q_2^{a+1}\) in \(F_{Q_2}(m)\) is strictly larger
than the sum of the corresponding floors at \(n\) and \(m-n\).
Indeed \(0<n<Q_2^{a+1}\) and \(m\) is divisible by \(Q_2^{a+1}\).
That factor of \(p\) also makes the term vanish modulo \(pR\).
Only \(n=m\) remains. Its multiplier is \(A(0)\) times the row
pairing \(\sum_k\overline V_{\mu k}\overline V_{\nu k}^{\,m}\),
which is the stated Kronecker symbol modulo \(pR\) by Lemma \ref{lem:partner76}. Finally multiplication by \(m!\) converts the
coefficient into the undivided derivative.
\end{proof}

\subsection{Individual Frobenius translates in the theta derivatives}\label{nav:VI-00}
\label{sec:direct76}

Fix a good prime \(p\ge5\) inert in \(K=\QQ(\sqrt{-d})\) and
unramified in \(E\). Keep the marked even elliptic theta function
\(\vartheta=\Theta_{t_*}\), the single-step class sections \(F_{ch}\)
and their matrix coordinates \(Z_\nu\) from Section~\ref{sec:period63}.
Write \(T(z)\) for the good elliptic parameter in additive coordinate
\(z\), normalized to have linear coefficient one. Put
\[
       Q_2=p^2,\qquad F_{Q_2}(n)=\sum_{i\ge1}\lfloor n/Q_2^i\rfloor,
       \qquad \lambda(z)=\vartheta'(z)/\vartheta(z).
\]
All valuations assign one to \(p\). We first determine where a
coefficient of \(\lambda\) can have a larger denominator than an
integral function of \(T\). This will identify the low direction
responsible for each possible loss in a mixed derivative.

\begin{lemma}[The addition formula and the exceptional degrees]
\leavevmode\label{lem:add76}
The difference \(\lambda(u+v)-\lambda(u)-\lambda(v)\) is an
integral germ in \(T(u),T(v)\). If \(g_m=[z^m]\lambda(z)\) and
\(m>0\), then
\begin{equation}
 v_p(g_m)\ \ge\ -F_{Q_2}(m)
         -\mathbf 1_{\{m=p^j\ {\rm for\ an\ odd}\ j\ge1\}} .
 \label{eq:add76}
\end{equation}
Here \(g_0=0\).
\end{lemma}
\begin{proof}
Consider the germ
\[
 U(u,v,w)=
 \frac{\vartheta(u+v+w)\vartheta(u)\vartheta(v)\vartheta(w)}
      {\vartheta(u+v)\vartheta(u+w)\vartheta(v+w)} .
\]
The signed sum of the seven row Gram matrices is zero. Each
diagonal entry occurs twice with either sign and each off-diagonal
entry occurs once with either sign. Lemma~\ref{lem:rational66}
therefore makes \(U\) a unit in the integral elliptic parameters.
The invariant derivative in the third elliptic variable preserves
the ring of integral germs. Divide that derivative of \(U\) by
\(U\) and put \(w=0\). Evenness gives \(\lambda(0)=0\), and
logarithmic differentiation gives exactly the difference asserted.
This three-variable quotient is the local theta expression for
the theorem of the cube. \href{https://www.jmilne.org/math/CourseNotes/AVc.pdf}
{See J.~S. Milne, \emph{Abelian Varieties}, the section on the theorem of the cube.}

The coefficient of \(u^k v^{m-k}\) in this difference is
\(\binom{m}{k}g_m\) for \(0<k<m\). An integral elliptic germ has
coefficient valuation at least \(-F_{Q_2}(k)-F_{Q_2}(m-k)\) there.
Indeed the formal exponential in a good parameter has degree \(n\)
coefficient of valuation at least \(-F_{Q_2}(n)\), and the floor
sum is superadditive. The formal-exponential bound and its
invariance under integral change of parameter are proved in
Lemma \ref{lem:multijet67}.

If \(m>1\) is not a power of \(p\), choose \(k=p^a<m\) at a
nonzero base-\(p\) digit of \(m\). The subtraction \(m-k\) uses
no borrow. Legendre's factorial formula gives
\(v_p\binom{m}{k}=0\), and
\(F_{Q_2}(k)+F_{Q_2}(m-k)=F_{Q_2}(m)\). If \(m=p^j\), \(j\ge1\),
take \(k=p^{j-1}\). There is exactly one carry in the addition,
across the boundary \(p^j\). The binomial valuation is one,
and the difference
\(F_{Q_2}(m)-F_{Q_2}(k)-F_{Q_2}(m-k)\) equals one for even
\(j\) and zero for odd \(j\). This gives \eqref{eq:add76}.
For \(m=1\), \(\lambda'(0)=-u_\Lambda-x(t_*)\) is integral
by good reduction at the mark.
\end{proof}

Call a series in a finite set of additive variables \(Z_\nu\)
bounded by the elliptic factorials if its coefficient of \(Z^\beta\)
has valuation at least \(-\sum_\nu F_{Q_2}(\beta_\nu)\).
This condition is preserved by sums and products. It holds for
every integral elliptic germ evaluated at \(T(L_i(Z))\) for
integral linear forms \(L_i\), even with coefficients in an
extension of \(\mathcal O_{K_p}\). In fact the multinomial
valuation in \(L_i^m\) is \(v_p(m!)-\sum_\nu v_p(\beta_\nu!)\).
The difference \(v_p(n!)-F_{Q_2}(n)\) is superadditive, which
gives the stated bound after substitution and multiplication.

Let \(\pi=\sigma_p c_E\) act on the embedding index and fix
each matrix entry. On an even cycle of this permutation the two
orbits of \(\pi^2\) are denoted by \(H,L\), in either order.
Fix one matrix entry in that cycle. Write \(c_{\nu i}=\overline
V_{\nu i}\) for the entries of the change of coordinates
\(z_i=\sum_\nu c_{\nu i}Z_\nu\) of Section~\ref{sec:period63};
the embedding and matrix-entry labels are combined in \(\nu\).
For \(\mu\in L\) define
\[
       E_\mu(Z_H)=\sum_i c_{\mu i}
                     \lambda\left(\sum_{\nu\in H}c_{\nu i} Z_\nu\right).
\]

\begin{lemma}[The coefficient attached to a specified translate]
\label{lem:partner76}
The coefficients of \(E_\mu\) satisfy the elliptic factorial
bound, with a possible further loss of one exactly in the
following list of positions:
\[
          \beta=p^j e_\nu,\qquad j\ge1\ {\rm odd},\quad
                       \nu\in H,\quad\pi^j\nu=\mu .
\]
At each listed position the valuation is at least
\(-F_{Q_2}(p^j)-1\).
\end{lemma}
\begin{proof}
Apply the addition formula in Lemma~\ref{lem:add76} to each
linear form one summand at a time. The difference between
\(E_\mu\) and \(\sum_{\nu,i} c_{\mu i}\lambda(c_{\nu i}Z_\nu)\)
satisfies the elliptic factorial bound by integral substitution
as above. This treats every coefficient involving two different
variables. By \eqref{eq:add76}, the remaining possible exception
is the pure coefficient at \(m=p^j\), \(j\) odd. Its multiplier
of \(g_m\) is
\[
            \sum_i c_{\mu i}c_{\nu i}^{p^j}
                 \equiv \sum_i\overline V_{\mu i} V_{\pi^j\nu,i}
                 =\mathbf1_{\{\mu=\pi^j\nu\}}\pmod{p}.
\]
The elements of \(E\) in this formula are integral in an
unramified local ring at \(p\). Raising to \(p^j\) modulo \(p\)
acts there by \(\sigma_p^j\). Since \(j\) is odd,
\(\sigma_p^j c_E=\pi^j\). The last equality is row
orthogonality \(VV^*=1\). In particular the congruence takes
place modulo \(p\) itself after any extension of coefficients.
When the indices differ it supplies the full missing factor.
\end{proof}

\begin{theorem}[Divisibility at each odd translate of Frobenius]
\label{thm:direct76}
In each matrix entry and each even cycle choose a half \(H\).
For any multi-index \(\alpha\), at every place over \(p\) one has
\begin{equation}
 v_p(\partial^\alpha F_{ch}(0))\ \ge\
 \sum_{\text{\rm entries}}\ \sum_{C\ {\rm even}}\ 
 \sum_{\nu\in H}\ \sum_{j\ge1,\ j\ {\rm odd}}
 \left[\left\lfloor\frac{\alpha_\nu}{p^j}\right\rfloor
                                 -\alpha_{\pi^j\nu}\right]_+ .
 \label{eq:direct76}
\end{equation}
Here \(\nu\) includes the matrix entry in the outer sum and
\([u]_+=\max(u,0)\). The same choice of basis and vacuum
works at any finite set of the indicated primes.
\end{theorem}
\begin{proof}
Use the integral splitting of the trace lattice in the proof of
Theorem \ref{thm:ramp67}.
For clarity, at \(p\) the local field factors correspond to the
orbits of \(\pi^2\) in the embeddings. Hermitian duality pairs
the two factors in an even cycle, each of which is isotropic.
Different cycles and different matrix entries are orthogonal.
Choosing integral bases in the local factors, write \(A_b\)
for the columns of the resulting change of basis in block \(b\).
The block here consists of one matrix entry and one cycle.
Corollary~\ref{cor:class66} and Theorem~\ref{thm:gram66} express
\(F_{ch}\) as an integral elliptic germ times
\(\prod_b G_{A_b}(u_b)\). Indeed the full Gram matrix is the
sum of its orthogonal block Gram matrices. Each block vacuum
on the chosen isotropic half alone is an integral elliptic
unit by the same theorem. Passing from a factor basis to its
embedding variables is an integral linear substitution in
the additive variables. A matrix vector supported on one
block has original product coordinates
\(\sum_{\nu\in b}c_{\nu i} Z_\nu\) in position \(i\).

Apply all derivatives in \(L\) and in the odd cycles before
extracting coefficients in the chosen halves. For an even
block the first logarithmic derivative in direction \(\mu\in L\),
restricted to \(H\), is precisely \(E_\mu\). A logarithmic
derivative of order at least two is integral in the local
elliptic parameters before passing to embedding coordinates.
In fact \(\lambda'(z)=-u_\Lambda-x(z+t_*)\) is an integral
elliptic germ. Its further invariant derivatives are integral,
and the rows of \(A_b\) have entries in \(\mathcal O_{K_p}\),
which act integrally on the elliptic formal group. The scalar
coefficients introduced by differentiation in embedding
directions are integral. The same differentiation of the
initial integral factor also preserves integrality. Derivatives
of odd block vacua evaluated at zero are integral by
Theorem~\ref{thm:jet63}'s undivided-derivative argument for
products of marked elliptic thetas.

By the product rule, after setting the differentiated variables
to zero, every term is a series satisfying the elliptic factorial
bound times a product of \(E_\mu\)'s. There are at most
\(\alpha_\mu\) such factors bearing a given label \(\mu\).
Indeed each requires one of the derivatives in that direction
to act on a block vacuum as a singleton, and an additional
derivative on that singleton makes it an integral factor.
Any multiplicities from the product rule are integers.

Extract the coefficient of the prescribed monomial in the
remaining variables. Superadditivity bounds the common
factorial denominator by \(\sum_{b,\nu\in H}F_{Q_2}(\alpha_\nu)\).
For fixed \((\nu,j)\) an exceptional factor must use exactly
\(p^j\) powers of \(Z_\nu\) and bear the label \(\pi^j\nu\),
by Lemma~\ref{lem:partner76}. Thus it can occur at most
\(\min(\lfloor\alpha_\nu/p^j\rfloor,\alpha_{\pi^j\nu})\) times.
Each such occurrence costs at most one further power of \(p\).
On multiplying by the factorials of the remaining orders,
the difference \(v_p(\alpha_\nu!)-F_{Q_2}(\alpha_\nu)\)
is the sum of \(\lfloor\alpha_\nu/p^j\rfloor\) over odd \(j\).
Subtracting the exceptional losses gives \eqref{eq:direct76}
in each term, hence in their sum. Lemma~\ref{lem:basis63}
gives the common basis.
\end{proof}

The source uses two-variable series in a height-two formal group and
an odd-power formula in
\srcloc{Formalism/SharpRestriction/GlobalSchurBounds/GL3Theta}{382}.
The calculation above derives the bound directly in the additive
coordinates of the marked analytic theta function.

\subsection{One minor projection for the translates at primes split in the second field}\label{nav:VI-01}
\label{sec:signal77}

We apply Section~\ref{sec:direct76} on the common intervals
of Section \ref{nav:I-28}.
Write \(K_d=\QQ(\sqrt{-d})\), \(K_q=\QQ(\sqrt{-q})\), and
\(\Phi=\Gal(E/K_d)\). In the coordinates of Theorem
\ref{thm:orientation46}, with the correction for an odd \(\psi\) specified there,
complex conjugation \(c_E\) reverses both quadratic signs.
We use only the signs themselves in the proof. Decompose
\(\Phi=\Phi_0\sqcup\Phi_1\) according as an embedding fixes
\(K_q\) or acts on it by conjugation. Their sizes are equal
to \(n\). Put the high orders \(a_t\) at the embeddings
\(t\in T\subset\Phi_0\), where \(h=|T|\) and
\(|a_t-xb|\le b\). They are positive integers. The sparse
construction uses \(a_t=\lceil xb\rceil\). The low thresholds
\(w_\mu\) equal \(B'=\lceil b\rceil\) on the chosen translate
of \(T_2\) in \(\Phi_1\), of size \(h_2\), and \(W=\lceil\sqrt
b\rceil\) everywhere else. As before \(b=\exp(x^\eps)\).

For a list \(\mathbf b\) of minor degrees let
\(P_{\mathcal I}\) range over products of \(r\)-rowed determinants
at the embeddings, in the \(s\)-by-\(r\) matrix variables of
\eqref{eq:pair55}. Write \(d_M(b,s)=
\max(0,\lceil(b-s(M-1))/M\rceil)\). A place over \(p\)
has \(v_p(p)=1\).

\begin{corollary}[Frobenius translates after taking minors]
\label{cor:minor77}
At each of the good inert primes of Section~\ref{sec:period63},
choose a half \(H_C\) in each even cycle \(C\) on the embeddings.
At every place over \(p\) the valuations of
\(P_{\mathcal I}(\partial)F_{ch}(0)\) and of its corresponding
cleared class coordinates \(A_r v_{h',\mathcal I}\) are at least
\begin{equation}
 r\sum_{C\ {\rm even}}\sum_{\nu\in H_C}\sum_{j\ge1,\ j\ {\rm odd}}
          [d_{p^j}(b_\nu,s)-b_{\pi^j\nu}]_+ .
 \label{eq:minor77}
\end{equation}
\end{corollary}
\begin{proof}
Use the same half in all entries of a matrix block. In any
monomial of a determinant product the sum of the row exponents
within each column at \(\nu\) equals \(b_\nu\). The sum of their
quotients on division by \(p^j\), rounded down, is at least
\(d_{p^j}(b_\nu,s)\), since each of the \(s\) remainders is at
most \(p^j-1\). At \(\pi^j\nu\) the row exponents in the same
column sum to \(b_{\pi^j\nu}\). Sum \eqref{eq:direct76} over
these rows using \(\sum_i[y_i]_+\ge[\sum_i y_i]_+\).
Each of the \(r\) columns gives the asserted bound. The
determinant polynomials have integer coefficients. Taking
their sum and taking a class sum with its integral coefficients
\(A_r m_c\) both preserve the bound.
\end{proof}

For each translate of the low thresholds use a single list
\begin{equation}
 b_\nu=
 \begin{cases}
      a_\nu,&\nu\in T,\\
      0,&\nu\in\Phi_0\setminus T,\\
      w_\nu,&\nu\in\Phi_1 .
 \end{cases}                                                   \label{eq:degrees77}
\end{equation}
Let \(C_{\rm sp}(t,a_t,w)\) denote the contribution to
\(C_d(t,a_t,w)\) from Galois elements which fix \(K_q\) and
act nontrivially on \(K_d\). At ramified primes use the same
inertia average as in \eqref{eq:cp73} and the definition
preceding it. Let \(S_{\rm sp}\) consist of the unramified
primes \(p\ge5\) with these two signs up to
\(X=\max_t a_t/W\), excluding the primes dividing the level.
For a nonzero algebraic number \(Y\) in a finite Galois field \(J\) put
\[
 \mathcal H_S(Y)=\frac{\log|N_{J/\QQ}Y|}{[J:\QQ]}
    +\sum_{p\notin S}\sum_{v\mid p}
        \frac{[J_v:\QQ_p]}{[J:\QQ]}
                 \max(0,-v_p(Y))\log p .
\]

\begin{theorem}[Simultaneous divisibility for all high embeddings]
\label{thm:signal77}
Fix the positive split intertwiners and representatives of
Theorem~\ref{thm:nonzero74}. Fix \(r\le s\) with \(rs\) even
and \(s\) bounded by a fixed power of \(x\). There is one
rational torsion input, supported at the fixed split level prime,
giving nonzero squared class sums
\(Y_u=A_s A_r^2\mathcal N_u\) simultaneously for all the lists
\eqref{eq:degrees77} indexed by \(u\). On this field and interval,
\begin{equation}
       \sum_{t\in T} C_{\rm sp}(t,a_t,w)
       \ \le\ \frac{\mathcal H_{S_{\rm sp}}(Y_u)}{r}
                         +O_{\psi,\eps}(hb/x).
       \label{eq:signal77}
\end{equation}
The left side uses the list of low thresholds for this \(u\).
If a finite family of such lists occurs at the same definite
step, it can be included in the same torsion choice.
\end{theorem}
\begin{proof}
For \(p\in S_{\rm sp}\), put \(\pi=\sigma_p c_E\).
Its action interchanges \(\Phi_0\) and \(\Phi_1\). Therefore
every cycle has even length and its intersection with \(\Phi_0\)
is one entire half. Choose these halves in Corollary~\ref{cor:minor77}.
For every \(t\in T\) and odd \(j\) the order at \(\pi^jt\)
is precisely \(w_{\pi^jt}\). At an unramified prime the count
in these signs has the exact term
\(2\log p\sum_{j\ {\rm odd}}[a_t/p^j-w_{\pi^jt}]_+\).
Use \(d_{p^j}(a_t,s)\ge a_t/p^j-s\) in \eqref{eq:minor77}.
Summing positive parts in that formula keeps every \(t\)
simultaneously. At most \(s\) is lost per \(t\) and per active
odd exponent. This proves divisibility before any squaring
or archimedean estimate.

The cleared squared norm is a sum of products of cleared
coordinates and their conjugates with integral outer weights.
The lower bound just proved holds at every coefficient place
over \(p\), so it doubles on those products and holds on their
sum. Theorem~\ref{thm:nonzero74} applies to any finite family
of minor lists, in particular \eqref{eq:degrees77}. Its positive
translation mean and the density of the common split torsion
grid supply the nonzero numbers asserted here. At the level
prime the splitting in \(K_d\) rules out a contribution.
For any other unramified prime, an odd power has the two
required signs precisely when the prime itself does. An even
power has trivial quadratic signs.

The primes dividing \(2\operatorname{cond}(\psi)\), and the prime
3, split in \(K_d\) on sufficiently distant intervals. At the
remaining ramified primes \(q,d\) the normalized count per \(t\)
is at most \(2a_t\log p/(p-1)\), by the inertia average.
Their combined contribution is \(o(hb/x)\) since both primes
exceed every fixed power of \(x\) on these intervals. The total
rounding loss, after accounting for the square and dividing
by \(r\), is at most \(2hsX\log(2X)=o(hb/x)\), since
\(X\asymp x\sqrt b\). Finally
\(\mathcal H_S(Y)\) bounds from above the normalized sum of the
valuations over \(S\), each weighted by \(\log p\), by the
product formula. This gives \eqref{eq:signal77}.
\end{proof}

\begin{proposition}[The other sign and the cost of the lower degrees]
\label{prop:others77}
Let \(D_\rho=L'(1,\rho)/L(1,\rho)\) for a nonprincipal primitive
character, and write \(\xi_{qd}\) for the primitive product of
the two quadratic characters. Uniformly for \(t\in T\) and \(u\),
\begin{align}
 C_d(t,a_t,w)-C_{\rm sp}(t,a_t,w)
  &= M_{\rm both}(a_t,W)+O_{\psi,\eps}(b/x),\notag\\
 M_{\rm both}(a,W)
  &=\frac{a}{2}
     \{\log(a/W)-1-\gamma+D_{\xi_q}+D_{\xi_d}-D_{\xi_{qd}}\}.
                                                        \label{eq:both77}
\end{align}
The lower orders in \eqref{eq:degrees77} have the following
contribution to the exact Fischer factorials of
Theorem~\ref{thm:frame55}.
\begin{equation}
  0\le\frac{1}{rh}\sum_{\mu\in\Phi_1}\log C_{r,w_\mu}
           \ll_{\psi,\eps} b x^{9\eps/8}=o(b\sqrt x).
                                                        \label{eq:lowcost77}
\end{equation}
In \eqref{eq:lowcost77} the field and the sizes of \(T,T_2\)
are those of Theorem \ref{thm:rangenew50}, with the optional field
shortening of Proposition~\ref{prop:desc50}.
\end{proposition}
\begin{proof}
In the omitted part the Galois element \(g\) reverses both
quadratic signs. Thus \(t+c_E+g\) fixes \(K_q\) and its threshold
is \(W\). The sign indicator with the factor two from the count is
\((1-\xi_q(g))(1-\xi_d(g))/2\). Average this identity on each
inertia coset. A character trivial on inertia takes the value
of its primitive Dirichlet character on the prime power. One
nontrivial on inertia has average zero and the primitive
character also vanishes there. Writing \(Y=a_t/W\), the
omitted part is consequently
\[
       \frac W2\{R_{\mathbf1}(Y)-R_{\xi_q}(Y)
                         -R_{\xi_d}(Y)+R_{\xi_{qd}}(Y)\}.
\]
Both quadratic fields are distinct and their product character
is nonprincipal. Insert the linear and logarithmic main terms
of equation \ref{eq:Gchi} in the notation \(R_\rho,G_\rho\)
of its explicit formula. Equivalently these main terms are
\(Y(\log Y-1-\gamma)\) for the principal character and
\(-YD_\rho\) for each of the other three. The selected-family
estimate at \(Y\asymp x\exp(x^\eps/2)\) bounds each of the four
remaining terms after multiplication by \(W\) by
\(O_{\psi,\eps}(a_t\exp(-c x^{\eps/8}))=o(b/x)\).
See Lemma \ref{lem:far14}.
This proves \eqref{eq:both77}, including the ramified primes.

At \(\mu\) the product \(C_{r,w_\mu}\) contains \(r w_\mu\)
positive integers at most \(w_\mu+r\). The sum of the \(w_\mu\)
over \(\Phi_1\) is \(h_2 B'+(n-h_2)W\). Since \(h_2\le n\)
and \(h\ge n x^{-\eps/8}/2\), division by \(h\) bounds it
by \(2B'x^{\eps/8}\). Finally \(\log(B'+r)=O_\eps(x^\eps)\)
for the stated polynomial ranks. This proves \eqref{eq:lowcost77}.
\end{proof}

In particular the centered count at this step is bounded above by
\[
 \frac{\mathcal H_{S_{\rm sp}}(Y_u)}r-
 \sum_{t\in T}\bigl(M_d(a_t,w)-M_{\rm both}(a_t,W)\bigr)
       +O_{\psi,\eps}(hb/x).
\]
The bound \eqref{eq:lowcost77} accounts for one explicit factor
of this arithmetic norm. The remaining upper estimate concerns
the high orders, the projection, and the finite valuations outside
\(S_{\rm sp}\). Its required precision is that of equation \ref{eq:gap}, after the sum over all steps.

\subsection{Composing the polynomial maps through the class sums}\label{nav:VI-02}
\label{sec:chain81}

Fix the field, CM type \(\Phi\), additive matrix coordinates, positive
split intertwiners and definite class representatives of
Sections~\ref{sec:period63} and~\ref{sec:mean74}. We first describe
the map on polynomials in the square frame variables. Its coefficients
will let us apply Theorem~\ref{thm:direct76} to an entire composition.

At rank \(r\) put one \(r\)-by-\(r\) matrix \(X_\tau\) at each
\(\tau\in\Phi\). Over an integral domain \(R\), let \(V_\lambda(R)\)
consist of the polynomials homogeneous of degree \(\lambda_{\tau,i}\)
in column \(i\) of \(X_\tau\), invariant under adding an indeterminate
multiple of any earlier column to any later one in the same matrix.
At each \(\tau\) the list \(\lambda_\tau\) is a decreasing list in the
weak sense of \(r\) nonnegative integers. Write
\(|\lambda_\tau|=\sum_i\lambda_{\tau,i}\). These are identities of
polynomials in the matrix entries and in the indeterminate.
Consequently \(V_\lambda(R)=V_\lambda(K)\cap R[X]\) for the
fraction field \(K\).

For \(r\le s\), introduce \(s\)-by-\(r\) matrices \(D_\tau\) and
\(s\)-by-\(s\) matrices \(Y_\tau\). Choose nonnegative integers
\(n_\tau\). If \(F\) is a holomorphic theta section in the matrix
coordinates, define a linear functional on the monomials in \(D\) by
\(J_F(D^\alpha)=\partial^\alpha F(0)\), with all derivatives
undivided. It acts coefficientwise in \(Y\). Let \(Y_{\tau,[r]}\)
denote the first \(r\) columns and put
\begin{equation}
 {\cal T}_{F,n}(Q)(Y)=J_F\left(
   \left.Q(X)\prod_{\tau\in\Phi}\det(X_\tau)^{n_\tau}
             \right|_{X_\tau=D_\tau^{\mathsf T}Y_{\tau,[r]}}\right).
 \label{eq:sub81}
\end{equation}
The transpose in this substitution is the ordinary transpose.

\begin{proposition}[Column weights and divisibility of every coefficient]
\label{prop:op81}
The output of \eqref{eq:sub81} has column weights
\(\lambda'_\tau=(\lambda_{\tau,1}+n_\tau,\ldots,
\lambda_{\tau,r}+n_\tau,0,\ldots,0)\) and belongs to \(V_{\lambda'}\).
Set \(B_\tau=|\lambda_\tau|+rn_\tau\), and take \(F=F_{ch}\) to
be one of the translated class sections at a good prime \(p\) of
Theorem~\ref{thm:direct76}. Choose one half \(H_C\) in each even
cycle \(C\) of \(\pi_p\), the same half for all matrix entries.
If the coefficients of \(Q\) are integral at a place over \(p\),
every output coefficient has valuation at least
\begin{align}
 g_p(\lambda,n;r,s)&=\sum_{C\ {\rm even}}\ \sum_{\nu\in H_C}\
               \sum_{k\ge1,\ k\ {\rm odd}}\max(0,L_{\nu,k},M_{\nu,k}),
                                                        \label{eq:op81}\\
 L_{\nu,k}&=r d_{p^k}(n_\nu,s)-B_{\pi_p^k\nu},\qquad
 M_{\nu,k}=d_{p^k}(B_\nu,rs)-B_{\pi_p^k\nu}.             \notag
\end{align}
Here \(d_l(b,s)=\max(0,\lceil(b-s(l-1))/l\rceil)\).
\end{proposition}
\begin{proof}
Multiplication by the determinant power adds \(n_\tau\) to
each column degree and preserves all the column-addition
identities. A column operation on the first \(r\) columns of \(Y_\tau\)
induces the identical one on \(D_\tau^{\mathsf T}Y_{\tau,[r]}\).
The later columns of \(Y_\tau\) do not enter the expression.
Taking \(J_F\) commutes with these operations, including with
scaling a column by an indeterminate. This proves the asserted
column conditions over the coefficient field.

Fix a monomial from \(Q\), with row degrees \(e_{\tau,i}\) at
\(\tau\). Each determinant monomial has degree one in each row.
Every resulting monomial in \(D,Y\) thus has degree
\(n_\tau+e_{\tau,i}\) in the \(i\)-th column of \(D_\tau\).
Its total degree there, over all columns, is \(B_\tau\).
The sum of the floors of \(s\) nonnegative integers divided by
\(p^k\) is at least \(d_{p^k}(b,s)\) when their sum is \(b\), by
the bound \(p^k-1\) on each remainder. For fixed \(\nu,k\) in
Theorem~\ref{thm:direct76}, first sum the positive parts over the
\(rs\) entries. Subtracting all exponents at \(\pi_p^k\nu\) costs
exactly \(B_{\pi_p^k\nu}\). Applying the floor bound in each column
gives at least \(L_{\nu,k}\), since \(e_{\nu,i}\ge0\). Applying it
to all \(rs\) entries at once gives at least \(M_{\nu,k}\).
The sum of positive parts is nonnegative and bounds both quantities
for this same monomial. Summing their maximum with zero proves
\eqref{eq:op81}. Expansion of the substitution introduces integer
coefficients, so the bound persists in each resulting coefficient.
After dividing by \(p^{g_p}\), the polynomial is still of the
asserted weights over the valuation ring by the intersection identity.
\end{proof}

Let the successive ranks be \(N_0\le\cdots\le N_m\), with
\(N_{j-1}N_j\) even. Write \(\mathcal C_j\) for the finite set
of representatives at rank \(N_j\) and \(m_c>0\) for the inverse
stabilizer order. Choose a basis as in Lemma~\ref{lem:basis63} at
each step, integral at a prescribed finite set of good primes.
For every pair at step \(j\) choose the positively normalized
section \(G^{(j)}_{ch}\) of Section~\ref{sec:mean74}. Apply the
same translation \(U_{w_j}\) to every such pair at this step.
Given an array \(Q_0\) on \(\mathcal C_0\), define recursively
\begin{equation}
  Q_{j,h}(Y)=\sum_{c\in\mathcal C_{j-1}}m_c\,
       {\cal T}_{\,U_{w_j}G^{(j)}_{ch},\,n_j}(Q_{j-1,c})(Y).
                                                        \label{eq:array81}
\end{equation}
For \(P_0\)-power torsion inputs use the corrected rational
translations of Theorem~\ref{thm:nonzero74}. The correction multiplies
the entire \(j\)-th map by one common unit complex scalar of finite order.

\begin{theorem}[Nonvanishing through the intermediate class sums]
\label{thm:compose81}
Use one nonzero polynomial \(Q_0\in V_{\lambda_0}(\overline{\QQ})\)
at every initial class. For any finite set of such polynomials,
degree lists and target classes, there is one tuple of rational
translations, supported at the fixed split prime \(P_0\) as in
Section~\ref{sec:mean74}, for which all the values \(Q_{m,h}(I)\)
in \eqref{eq:array81} are nonzero and algebraic. Here \(I\)
is the identity frame at every embedding. Put
\(A_N=|\mathrm{GL}_{N[E:\QQ]}(\mathbb F_3)|\) and
\(a_{\rm cl}=\prod_{j=1}^m A_{N_{j-1}}\).
If \(Q_0\) is integral at the chosen coefficient places, then
\[
  v_p\big(a_{\rm cl} Q_{m,h}(I)\big)\ \ge\
     \sum_{j=1}^m g_p(\lambda_{j-1},n_j;N_{j-1},N_j)
\]
at each such place. The successive \(\lambda_j\) are obtained by
the column-weight rule of Proposition~\ref{prop:op81}.
\end{theorem}
\begin{proof}
Fix one of the data sets and a terminal class \(h\). At each step
use a common period lattice \(\Gamma_j\) for its finitely many
sections, as in Section~\ref{sec:mean74}. On the product of the
resulting tori form the section
\[
 {\mathcal G}_h=
   \sum_{c_0,\ldots,c_{m-1}}
   \left(\prod_{j=1}^m m_{c_{j-1}}\right)
   \bigotimes_{j=1}^m G^{(j)}_{c_{j-1}c_j},\qquad c_m=h.
\]
In the proof of Proposition~\ref{prop:overlap74} the overlap of
each pair section with its step's unshifted vacuum \(G_{0,j}\)
is \(\|G_{0,j}\|^2 n_{ch}^{-1/2}>0\). Pairing \({\mathcal G}_h\)
with the tensor of these vacua gives a sum of products of positive
numbers, by product integration. In particular \({\mathcal G}_h\ne0\).

Set \(X_{m,\tau}=I\) and work backwards with symbolic matrices
\(X_{j-1,\tau}=D_{j,\tau}^{\mathsf T}X_{j,\tau,[N_{j-1}]}\).
Define a single polynomial in the independent matrices \(D_j\) by
\[
 P(D)=Q_0(X_0)\prod_{j=1}^m\prod_{\tau\in\Phi}
                         \det(X_{j-1,\tau})^{n_{j,\tau}} .
\]
To see that \(P\ne0\), choose \(D_j\) for \(j\ge2\) to be the
standard rectangular inclusions at every embedding. Then all
\(X_j\), \(j\ge1\), are identity matrices. The top square block
of each \(D_{1,\tau}\) can give any \(X_{0,\tau}\) after transposition.
Invertible tuples are Zariski dense in the product of the square
matrix spaces, so choose such a tuple with \(Q_0(X_0)\ne0\).

Applying \(P(\partial)\) at zero to the product translation of
\({\mathcal G}_h\) gives \(Q_{m,h}(I)\) with the analytic translations.
Indeed, on each path the rule \eqref{eq:sub81} first replaces
monomials in \(D_1\) by the derivatives of the first section.
The next substitution acts on the surviving frame variables and
commutes with this replacement. Iterating and evaluating at \(I\)
gives the displayed \(P\) followed by the independent derivative
evaluations. Summation over paths gives \eqref{eq:array81}.
The degree of \(P\) in the \(j,\tau\) block is
\(B_{j,\tau}=|\lambda_{j-1,\tau}|+N_{j-1}n_{j,\tau}\). Indeed,
before applying the \(j\)-th substitution the degree in that frame
is \(B_{j,\tau}\), and all later substitutions leave \(D_j\)
untouched. Put \(b_j=\sum_\tau B_{j,\tau}\).

If \(c_j^*=\pi/A_j\) is the metric constant at the \(j\)-th step,
the product-torus form of Lemma~\ref{lem:mean74} is
\begin{equation}
 \int\left|(P(\partial) U_{\boldsymbol w}{\mathcal G}_h)(0)\right|^2
                 d\boldsymbol w
     =\left(\prod_{j=1}^m(c_j^*)^{b_j}\right)
                    \|P\|_{\rm F}^2\|{\mathcal G}_h\|^2>0. \label{eq:avg81}
\end{equation}
For details, expand the polynomial into monomials and the section
into its finite sum of tensors. For two sections \(F,F'\) on one
torus and two multi-indices of the same total degree, the adjoint
calculation in Lemma~\ref{lem:mean74} gives the inner product
\(\mathbf1_{\alpha=\beta}(c_j^*)^{|\alpha|}\alpha!
\langle F,F'\rangle\) for their translated derivatives. Integrate
separately over the tori and sum. The diagonal monomial factors
give exactly the Fischer norm \(\sum_\alpha\alpha!|P_\alpha|^2\).

The squared absolute value on the left is a real analytic
function on a connected product torus. Its nonvanishing set is
open and dense by \eqref{eq:avg81}. A finite intersection of
these sets is open and dense. The product of the \(P_0\)-power
torsion grids for the common lattices is dense there, so one
tuple belongs to that intersection. The translation identity
in Theorem~\ref{thm:nonzero74} moves each selected translation
to one finite input before all the representatives at that step.
Its common phase multiplies the whole path sum by the product
of the \(m\) phases. The same theorem, using the algebraic
addition formulas for rational torsion, gives algebraicity of
the corrected derivatives. Finite sums, products and evaluation
at \(I\) preserve it. Work in one finite normal coefficient field
containing the finitely many numbers in question.

For the valuation claim replace every weight \(m_c\) at step
\(j\) by \(A_{N_{j-1}}m_c\), an integer by the stabilizer argument
in Theorem~\ref{thm:joint63}. By Proposition~\ref{prop:op81} the
cleared map at that step takes an integral input into polynomials
divisible coefficientwise by its \(p^{g_p}\). After dividing by
that power the inputs for the next step are integral of the new
column weights. Induction proves the stated sum in every final
coefficient. Evaluation at identity frames preserves the bound.
\end{proof}

The source defines its polynomial substitution and differentiated
table in \srcloc{Formalism/LocalToGlobal/KernelLattices/ThDistC}{23}
and \srcloc{Formalism/LocalToGlobal/KernelLattices/ThDistC}{60}.
Multiplying each derivative in direction \(\tau\) by a scalar
\(c_{D,\tau}\) multiplies its entire step by
\(\prod_\tau c_{D,\tau}^{B_{j,\tau}}\) by homogeneity.
The distribution in
\srcloc{Formalism/LocalToGlobal/KernelLattices/ThDistC}{137} also
uses character and frame factors. Theorem~\ref{thm:compose81}
uses positive split representatives, the additive coordinates
of Section~\ref{sec:period63}, and evaluation at the identity
frame. Comparison at the prescribed levels with the terminal
Bessel functional, including its normalization, remains open.

\subsection{The factorial and field norm of the composed value}\label{nav:VI-03}
\label{sec:fischer83}

We compute the archimedean polynomial factor for the same scalar
as in Section~\ref{sec:chain81}. Use \(Q_0=1\), the identity
terminal frame, and any nonnegative determinant powers \(n_{j,\tau}\).
Write \(R_j=N_{j-1}\), \(s_j=N_j\). Extend each preceding column
list by zeros after its last column. The partition on the first
\(R_j\) columns after step \(j\) is
\[
 v_{j,\tau,i}=\lambda_{j-1,\tau,i}+n_{j,\tau}
       \quad(1\le i\le R_j).
\]
Let \(b_j=\sum_{\tau,i}v_{j,\tau,i}\), \(d_j=|\Phi|R_js_j\), and
let \(P\) be the polynomial in the proof of Theorem~\ref{thm:compose81}.
Use \((k)_v=k(k+1)\cdots(k+v-1)\) with \((k)_0=1\).

\begin{theorem}[Exact factorials through successive substitutions]
\label{thm:factor83}
In the Hermitian additive coordinates the squared Fischer norm is
\begin{equation}
 \|P\|_{\rm F}^2=
     \prod_{j=1}^m\prod_{\tau\in\Phi}
                   \prod_{i=1}^{R_j}(R_j-i+1)_{v_{j,\tau,i}} .
                                                        \label{eq:factor83}
\end{equation}
\end{theorem}
\begin{proof}
Give every entry of the independent matrices \(D_{j,\tau}\) the
standard complex Gaussian density \(\pi^{-1}e^{-|z|^2}\).
Polar integration gives
\(\mathbb E[z^k\bar z^l]=\mathbf1_{k=l}k!\), so the squared
Fischer norm is \(\mathbb E|P(D)|^2\). The embeddings separate
because \(Q_0=1\), and we prove the formula for one embedding.

For a full-rank matrix \(A\) with \(r\) columns put
\(H_v(A)=\prod_{i=1}^r t_i(A)^{2v_i}\), where \(t_i(A)\)
is the length left after removing from column \(i\) its orthogonal
projection on the previous columns. Gram--Schmidt writes \(A=UT\)
with \(U^*U=I_r\) and \(T\) upper triangular with positive
diagonal \(t_i(A)\). For any full-rank \(B\), right multiplication
by such a \(T\) multiplies its \(i\)-th residual length by \(T_{ii}\).
Indeed the earlier column spans agree, and the new \(i\)-th column
modulo this span is \(T_{ii}\) times the old one. For square \(B\),
\(|\det B|^2=\prod_i t_i(B)^2\).

Hold \(X_j\) fixed in the backwards substitutions and put
\(A=X_{j,[R_j]}\). Conditional on this \(A\), each row of
\(D_j^{\mathsf T}U\) is a standard complex Gaussian row of length
\(R_j\). Indeed \(U\) extends to a unitary matrix of size \(s_j\),
and the Gaussian row in \(D_j^{\mathsf T}\) is invariant under
right multiplication by that unitary matrix. The rows are independent.
Thus \(X_{j-1}=D_j^{\mathsf T}A\) has the same conditional
distribution as \(GT\) for a square standard Gaussian \(G\).
These matrices are of full rank almost surely, starting at \(X_m=I\)
and working backwards.

Given the first \(i-1\) columns of \(G\), the random variable
\(t_i(G)^2\) is the sum of \(R_j-i+1\) independent exponential
variables of mean one. Its conditional \(v_i\)-th moment is
\((R_j-i+1)_{v_i}\), by integration of the gamma density.
The moment is constant in the earlier columns. Taking conditional
expectations starting with the last column gives
\[
 \mathbb E_{D_j}\bigl[H_v(X_{j-1})\,\big|\,X_j\bigr]
        = H_v(X_{j,[R_j]})\prod_{i=1}^{R_j}(R_j-i+1)_{v_i}.
\]
When \(H_v\) is put on the whole next frame, adjoining zeros to
\(v\) has exactly the effect of restricting to the first \(R_j\)
columns. Begin with \(|\det X_0|^{2n_1}\), integrate \(D_1\)
by this formula, and multiply the remaining function by
\(|\det X_1|^{2n_2}\). This adds \(n_2\) to each current column
weight, including the new ones. Repeat to \(j=m\). The remaining
factor at \(X_m=I\) is one, giving precisely \eqref{eq:factor83}.
\end{proof}

For \(d>0\), \(b>0\) put
\[
 e(b,d)=b\log(1+d/b)+d\log(1+b/d),\qquad e(0,d)=0.
\]
The next estimate allows \(b\) to be much larger than the
dimension of the derivative variables.

\begin{lemma}[Gaussian estimate for a multihomogeneous derivative]
\label{lem:opt83}
Let \(H(z)\) be entire on blocks of dimensions \(d_j\). Suppose
\(|H(z)|^2\le S\exp(\sum_j c_j|z_j|^2)\) with \(S,c_j>0\).
For any polynomial \(P\) of block degrees \(b_j\), homogeneous in
each block,
\[
 |P(\partial)H(0)|^2
   \le S\|P\|_{\rm F}^2
             \prod_j c_j^{\,b_j}\exp(e(b_j,d_j)).
\]
Moreover \(e(b,d)\le d[1+\log(1+b/d)]\).
\end{lemma}
\begin{proof}
Integrate against independent complex Gaussians of density
\((t_j/\pi)^{d_j}e^{-t_j|z_j|^2}\), with \(t_j>c_j\).
The squared integral norm is at most
\(S\prod_j(1-c_j/t_j)^{-d_j}\). Rotation of independent blocks
is unitary for this measure. Its Fourier projection onto the block
degrees \(b_j\) is the corresponding Taylor part, by integration
of the Taylor series on each compact rotation orbit. It is a
contraction in the integral norm. Orthogonality of Gaussian
monomials identifies the squared Fischer norm of that Taylor part
with its squared integral norm times \(\prod_j t_j^{b_j}\).
Pair its coefficients with the coefficients of \(P\), using
Cauchy--Schwarz with weights \(\alpha!\). The resulting product
is minimized in each block with \(b_j>0\) at
\(t_j=c_j(1+d_j/b_j)\), by logarithmic differentiation.
With \(b_j=0\) take \(t_j\) to infinity. These choices give
the asserted factor. The last inequality follows from
\(\log(1+d/b)\le d/b\).
\end{proof}

Return to one of the torsion tuples supplied by
Theorem~\ref{thm:compose81}. Let \(y=a_{\rm cl}Q_{m,h}(I)\) for
a fixed terminal class and put
\[
 M=\prod_{j=1}^m\left(\sum_{c\in\mathcal C_{j-1}}m_c\right),
 \qquad E_j=e\max(1,2R_js_j\,\operatorname{rd}(E)/\pi).
\]
Here \(e\) in \(E_j\) is the base of natural logarithms. In the
discriminant-one elliptic normalization for \(K_d\) set
\[
 L_d=\log(\sqrt d/(2\pi))-\gamma+D_{\xi_d},
 \qquad V_d=\tfrac14\log d-\tfrac12\gamma+\tfrac12D_{\xi_d}.
\]

\begin{theorem}[An upper estimate at infinity for the iterated scalar]
\label{thm:norm83}
For the nonzero \(y\) just defined, in any common finite Galois
coefficient field \(J\),
\begin{align}
 \frac{\log|N_{J/\QQ}(y\bar y)|}{[J:\QQ]}
 \le{}&2\log(a_{\rm cl}M)+\log\|P\|_{\rm F}^2 \notag\\
 &+\sum_{j=1}^m
      \bigl(b_j L_d+d_j(V_d+\log E_j)+e(b_j,d_j)\bigr).
                                                        \label{eq:norm83}
\end{align}
For the balanced degrees of
Lemma \ref{lem:balance82},
the contribution of the \(e(b_j,d_j)\) is \(o_{\psi,\eps}(b/x)\)
on the common intervals, where \(b=\exp(x^\eps)\).
\end{theorem}
\begin{proof}
At any embedding \(\rho\) use the conjugate sections of
Theorem~\ref{thm:transport59}. Each corrected torsion translate
has the same metric supremum as its vacuum there. Its squared
absolute value in block \(j\) is bounded by
\((E_j\nu_\rho)^{d_j}\exp(\pi |z_j|^2/A_\rho)\), by
\eqref{eq:vac59}. There is one translate with coefficient of
magnitude one per pair. The coefficients of \(P\) are integers
since the initial polynomial is one. Thus its conjugated
derivative is computed by the same \(P\) on the sum over all
paths of the products of these conjugated sections. Roots of
unity have absolute value one at every embedding. The triangle
inequality gives a squared envelope for this sum with
\(S=M^2\prod_j(E_j\nu_\rho)^{d_j}\) and \(c_j=\pi/A_\rho\).
Apply Lemma~\ref{lem:opt83} and include the factor \(a_{\rm cl}^2\).

For complex conjugation \(\iota\) at the original embedding,
the map \(\rho\mapsto\rho\iota\) permutes the embeddings. The
normalized logarithm on the left of \eqref{eq:norm83} therefore
equals the mean of \(\log|\rho(y)|^2\). The means of
\(\log(\pi/A_\rho)\) and \(\log\nu_\rho\) are \(L_d,V_d\),
respectively, by the elliptic area formula and the formula for
the marked two-torsion point in
Theorems \ref{thm:mean56} and \ref{thm:norm58}.
Taking these means and using the exact factorial in \eqref{eq:factor83}
gives \eqref{eq:norm83}. On the given intervals \(d_j,m,R_j,s_j,|\Phi|\)
are bounded by fixed powers of \(x\). The balanced degrees give
\(b_j\le R_j|\Phi|(a+B'+1)\), so \(\log(1+b_j)=O_\eps(x^\eps+\log x)\).
The final assertion now follows from the last bound of
Lemma~\ref{lem:opt83} and the growth of \(b\).
\end{proof}

For \(Z_u=y_u\overline{y_u}/D_x\) in Theorem \ref{thm:single82},
subtract \(\log D_x\) from the right of \eqref{eq:norm83} to
bound its normalized logarithmic field norm. Its \(\mathcal H_S\)
also includes the negative valuations outside \(S\). Equation
\eqref{eq:norm83} evaluates the archimedean factorial for the
same composed and translated scalar as the finite divisibilities.
The comparison with the centered main term and the control of those
remaining finite valuations are open.

\subsection{A finite torsion level for simultaneous nonvanishing}\label{nav:VI-04}
\label{sec:grid85}

The torsion choice in Theorem~\ref{thm:compose81} used density.
We now give a sufficient level in terms of the derivative orders
and a common period lattice for the fixed representatives. The
estimate will apply to the product of all the tests for which
nonvanishing is required.

Let \(T=\CC^n/\Gamma\), \(d_0=2n\), with Euclidean real inner
product. A smooth section \(f\) has the multiplier and metric
of Section~\ref{sec:mean74} with parameter \(c>0\).
Its squared metric magnitude \(g=|f|^2e^{-c|z|^2}\) is periodic.
Use normalized Haar measure and let \(\mu=\int_T g\).
For a unit vector \(v\in\CC^n\) put
\[
 A_v=\sum_i\bar v_i\partial_{\bar z_i},\qquad
 B_v=-\sum_i v_i(\partial_{z_i}-c\bar z_i).
\]
Assume \(A_v^{K+1}f=0\) for every such \(v\). Let \(\lambda\)
denote the shortest nonzero length in the real dual lattice
\(\Gamma^\vee=\{u:\langle u,\gamma\rangle_{\RR}\in\ZZ
\text{ for every }\gamma\in\Gamma\}\).

\begin{lemma}[Sampling a squared section at torsion points]\leavevmode\label{lem:grid85}
For any positive integer \(N\) with
\(N\lambda\ge8\sqrt{c(K+2d_0+2)}\), one has
\[
 \left|\frac1{N^{d_0}}\sum_{w\in N^{-1}\Gamma/\Gamma}g(w)-\mu\right|
     \le \frac{\mu}{2}.
\]
\end{lemma}
\begin{proof}
If \(\mu=0\), the continuous nonnegative function \(g\) is zero
everywhere. Assume henceforth \(\mu>0\).
Fix \(v\) and abbreviate \(A=A_v,B=B_v\).
Integration by parts on \(T\) gives \(A^*=B\) and direct
differentiation gives \([A,B]=c\), as in Lemma~\ref{lem:mean74}.
On smooth sections killed by \(A^{k+1}\) one has an orthogonal
expansion into \(B^l f_l\), \(0\le l\le k\), with \(Af_l=0\).
Indeed \(A^k f\) is killed by \(A\) and
\(A^kB^k A^kf=c^k k!A^kf\). Subtract
\(B^k A^k f/(c^k k!)\) and repeat with \(k-1\).
Distinct levels are orthogonal by moving the larger power to
the other argument and commuting the resulting \(A\)'s to the right.
The same computation gives
\(\|B^l f_l\|^2=c^l l!\|f_l\|^2\). It follows by orthogonality that
\[
       \|Af\|\le\sqrt{ck}\|f\|,\qquad
       \|Bf\|\le\sqrt{c(k+1)}\|f\|.
\]
The real covariant derivative \(D_v=A-B\) sends these sections
into the kernel of \(A^{k+2}\), by the commutator. Thus for \(s\ge1\)
\(\|D_v^s f\|\le(2\sqrt{c(K+s)})^s\|f\|\).

In the real direction \(v\), differentiation of a pointwise
metric pairing differentiates each argument by \(D_v\). The
two connection terms sum to the derivative of the weight
\(e^{-c|z|^2}\). Apply this product rule \(s\) times to \(g\).
Cauchy--Schwarz for each integrated absolute product yields
\(\|\partial_v^s g\|_{L^1}\le(4\sqrt{c(K+s)})^s\mu\).
For \(u\ne0\) in the dual lattice, integrate by parts in the
direction \(u/|u|\) against the character
\(\exp(-2\pi i\langle u,z\rangle_{\RR})\). Then
\[
 |\widehat g(u)|\le
       \mu\left(\frac{2\sqrt{c(K+s)}}{\pi |u|}\right)^s .
\]
The dual vectors in \(2^l\lambda\le|u|<2^{l+1}\lambda\)
number at most \(5^{d_0}2^{ld_0}\) for \(l\ge0\). Disjoint
balls of radius \(\lambda/2\) around them fit in the ball of
radius \(2^{l+1}\lambda+\lambda/2\), which gives the count by volume.
Take \(s=2d_0+2\). The Fourier series now converges absolutely,
and its finite grid average retains exactly the frequencies
in \(N\Gamma^\vee\) by finite character orthogonality. Hence
the discrepancy is at most
\[
 2\,5^{d_0}\mu\left(\frac{2\sqrt{c(K+s)}}{\pi N\lambda}\right)^s
 \le 2\,5^{d_0}(4\pi)^{-2d_0-2}\mu<\mu/2 .
\]
\end{proof}

\begin{theorem}[One bounded level for the composed values]\label{thm:level85}
Fix the positive definite class data of Theorem~\ref{thm:compose81}
and \(H\ge1\) tests on this chain, allowing different initial
polynomials, terminal classes and degree lists. For test \(i\)
let \(b_{ij}\) be its degree in the \(j\)-th derivative block.
Set \(K=\sum_{i,j} b_{ij}\), \(d_0=2\sum_j|\Phi|R_js_j\).
At step \(j\) write \(\Lambda_j\) for the period lattice of the
standard trace vacuum. Fix a positive integer \(L_j\) such that
\(\Gamma_j=L_j\Lambda_j\) is a common lattice as in
Section~\ref{sec:mean74}, and set \(L_*=\max_j L_j\).
Let \(f_0=[E:K_d]\). One may choose all the rational translations
of Theorem~\ref{thm:compose81} with denominator dividing the
same \(P_0\)-power \(N\), where
\begin{equation}
 N\ \le\ P_0\max\left(1,\
       8L_*\sqrt{\frac{\pi\operatorname{rd}(E)}{2f_0}}\,
                         \sqrt{H(K+2d_0+2)}\right).            \label{eq:level85}
\end{equation}
For all the translates of
Theorem \ref{thm:single82}
with one terminal class this gives
\(\log N\le\log L_*+O_{\psi,\eps,P_0}(x^\eps)\).
\end{theorem}
\begin{proof}
Rescale the \(j\)-th block of the product torus by
\(\sqrt{\pi/A_j}\), so its Gaussian metric becomes the one with
parameter one. The derivative evaluation of test \(i\), as a
function of the translation, has squared magnitude equal to the
squared metric magnitude of the smooth section obtained by applying
its derivative polynomial in the raising operators to its
holomorphic path sum. This is the differentiated translation
formula in Lemma~\ref{lem:mean74} applied in every block.
The rescaling only multiplies homogeneous derivatives by fixed
nonzero constants. Each antiholomorphic directional derivative
commutes through the raising operators at the cost of removing
one of them or else kills the holomorphic factor. Consequently
every direction annihilates this smooth section after
\(1+\sum_j b_{ij}\) successive antiholomorphic derivatives.

Take the product of the \(H\) sections. Its squared metric
magnitude \(g\) is the product of the \(H\) squared tests, and
its metric has parameter \(H\). The ordinary product rule
gives \(A_v^{K+1}f=0\) for the product. Each test has positive
mean and real analytic squared magnitude by \eqref{eq:avg81}.
Their nonzero sets are open and dense on the same connected
torus. The product is positive on a nonempty open set, hence
\(\int g>0\).

It remains to bound the real dual minimum after rescaling.
In a rescaled standard trace lattice \(\Lambda_j'\), the
alternating form \(\pi^{-1}\operatorname{Im}\langle z,w\rangle\)
is unimodular by Theorem~\ref{thm:lattice55}. Hence its real
Euclidean dual lattice equals \(i\Lambda_j'/\pi\), up to sign.
Indeed pairing with \(i z/\pi\) converts the real pairing into
the given alternating one, and unimodularity gives both
inclusions. The squared minimum of \(\Lambda_j'\) is at least
\(2\pi f_0/\operatorname{rd}(E)\) by the same theorem. Passing
to the dual of \(L_j\Lambda_j'\) divides its minimum length
by \(L_j\). The dual minimum of the orthogonal product therefore
satisfies
\[
        \lambda\ \ge\ L_*^{-1}
                  \sqrt{2f_0/(\pi\operatorname{rd}(E))}.
\]
Choose the least nonnegative power of \(P_0\) at least the
quantity inside the maximum in \eqref{eq:level85}, including
the lower bound one. Lemma~\ref{lem:grid85} makes the average of
\(g\) on that grid positive, so one grid point makes every test
nonzero. Before scaling its \(j\)-th component is in
\(N^{-1}\Gamma_j/\Gamma_j\). This is exactly the grid of
rational translations used in Section~\ref{sec:mean74}, with
order dividing \(N\) modulo the original lattice. The same
commutation with the split representatives gives one input per
step and preserves all the previously obtained divisibilities.

For the indicated translates \(H\le|\Gal(E/\QQ)|\), and
\(m,d_0,|\Phi|,R_j,s_j\) are bounded by fixed powers of \(x\).
By Lemma \ref{lem:balance82}, the total derivative order of one test is
at most \(\sum_j R_j|\Phi|(a+B'+1)\). Thus
\(\log(H(K+2d_0+2))=O_\eps(x^\eps+\log x)\).
The conductor--discriminant formula gives
\(\log\operatorname{rd}(E)\le\log\operatorname{cond}(E)
=O_{\psi,\eps}(x^{\eps/4})\) on this family. This proves the
logarithmic estimate.
\end{proof}

For lattice sampling of polyanalytic functions in the complex
plane see \href{https://arxiv.org/abs/0901.4386}{L.~D.~Abreu,
\emph{Sampling and interpolation in Bargmann--Fock spaces of
polyanalytic functions}}. Lemma~\ref{lem:grid85} proves directly
the estimate for the compact product torus and the lattice
used here. Formula \eqref{eq:level85} retains the dependence on
the common period indices of the actual representatives. The
height of the divided periods in Section \ref{nav:I-28}, still
requires a sharper estimate.

\subsection{Differentiating torsion translates at primes dividing the order}\label{nav:VII-62}
Let \(P=[w]\in C\) be a torsion point on the marked elliptic curve
of Theorem~\ref{thm:elliptic312}. The function \(U_w\Theta_{t_*}\)
uses the Hermitian translation of Theorem~\ref{thm:trans57}.
Use the complex embedding of \(K\) giving its action on the tangent
line, and extend its image to the coefficient field. Let
\(\mathfrak a(P)=\{\beta\in\OO_K:\beta P=0\}\) and define
\[
 \delta_{\mathfrak p}(P)=
   \min_{0\ne\alpha\in\overline{\mathfrak a(P)}}v_{\mathfrak p}(\alpha).
\]
The minimum is attained since this is a nonzero integral ideal and
the local value group of \(K^\times\) is discrete. If \(mP=0\), then
\(0\le\delta_{\mathfrak p}(P)\le v_{\mathfrak p}(m)\).

\begin{theorem}[Torsion translates in the original differential]\leavevmode
\label{thm:torsion312}
At every finite place and every \(j\ge0\), with
\(\delta=\delta_{\mathfrak p}(P)\) and \(a_+=\max(a,0)\), one has
\[
 v_{\mathfrak p}\big((U_w\Theta_{t_*})^{(j)}(0)\big)
 \ge -j\delta-\lfloor j/2\rfloor(\kappa_p-2\delta)_+ .
\]
In particular the value at zero is integral for torsion of arbitrary
order. At a split prime the derivative loss due to the translation
uses only the component of \(\overline{\mathfrak a(P)}\) at the
chosen place.
\end{theorem}
\begin{proof}
Set \(R=U_w\Theta_{t_*}/\Theta_{t_*}\) and fix \(m\ge1\) with
\(mP=0\). The transformation law in Theorem~\ref{thm:trans57}
gives \(R^m=ch\) for a rational function \(h\) with divisor
\(m([t_*-P]-[t_*])\). Its proof also gives
\[
 c^m\prod_{i=0}^{m-1}h(z+iP)=\epsilon_{t_*}(mw)^m,\qquad
                         \epsilon_{t_*}(mw)\in\{1,-1\}.
\]
Here the constants and the rational function are algebraic.
Over a local field of good reduction, scale \(h\) so its order at the
generic point of the special fiber is zero and change \(c\) inversely.
This scaling is possible since the smooth special fiber is reduced
and irreducible, so a uniformizer has order one there. Translation
by \(iP\) extends to an automorphism of the good model. Each factor
in the product has vertical order zero. The displayed equation
then proves that \(c\) is a unit.

Choose \(0\ne\alpha\in\overline{\mathfrak a(P)}\) attaining the
minimum. For \(\lambda\in\Lambda\), the multiplier of \(R\) under
\(\alpha\lambda\) is the exponential of \(2\pi i\) times, up to
sign, the alternating pairing of \(\lambda\) with \(\bar\alpha w\).
This follows from the adjoint identity for the Hermitian form on
the CM line. Since \(\bar\alpha w\in\Lambda\), the pairing is integral.
Consequently \(F(z)=R(\alpha z)\) is a meromorphic function on \(C\)
and \(F^m=c\,h\circ[\alpha]\). It is a rational function over a
finite number field. Indeed its divisor is the algebraic divisor
\([\alpha]^*([t_*-P]-[t_*])\). A degree-zero algebraic divisor
which is principal over \(\CC\) is already principal over
\(\overline{\QQ}\), by base change for global sections of its
line bundle. Divide \(F\) by an algebraic function with that divisor.
The quotient is constant and its \(m\)-th power is algebraic by
the displayed identity for \(F^m\). Enlarge the local field to
contain these data.

The endomorphisms \([\alpha]\) and \([\bar\alpha]\) extend over
good reduction by the N\'eron property of the abelian scheme.
Their composite on the special fiber is multiplication by the
nonzero integer \(N_{K/\QQ}(\alpha)\), which is finite and
surjective in every characteristic. Thus \([\alpha]\) is dominant
on the special fiber. At its generic point \(h\) is a unit, so
its pullback is a unit there as well. The identity for \(F^m\)
gives vertical order zero for \(F\). The possible horizontal poles
of \(F\) are supported in the inverse image of the section \(t_*\).
This closed inverse image misses the origin in the special fiber,
since \([\alpha](0)=0\) and \(t_*\) has nonzero reduction by
Theorem~\ref{thm:elliptic312}. Thus \(F\) has no pole along any
height-one prime of the local ring at the closed origin. The ring
is regular and equals the intersection of its localizations at
height one. Hence \(F\) is regular at that origin. These standard
facts are stated in Milne, \href{https://www.jmilne.org/math/CourseNotes/CM.pdf}
{\emph{Complex Multiplication}}, Proposition 6.4, and in the
Stacks Project, \href{https://stacks.math.columbia.edu/tag/0BFG}{Tag 0BFG}
and \href{https://stacks.math.columbia.edu/tag/031T}{Tag 031T}.

The integral invariant derivation in
Theorem~\ref{thm:elliptic312} therefore gives
\(\partial^iF(0)=\alpha^i R^{(i)}(0)\) integral for every \(i\).
In the Leibniz formula for the \(j\)-th derivative of
\(R\Theta_{t_*}\), the term with \(i\) derivatives on \(R\)
has valuation at least \(-i\delta-\lfloor(j-i)/2\rfloor\kappa_p\).
The largest possible loss is obtained by maximizing
\((j-2a)\delta+a\kappa_p\) over the integers
\(0\le a\le\lfloor j/2\rfloor\). Indeed among the orders with
\(\lfloor(j-i)/2\rfloor=a\), the largest \(i\) is \(j-2a\).
The maximum equals the stated expression. This calculation uses
ordinary derivatives in \(dz\) throughout.
\end{proof}

\section{Signed theta periods and CM derivatives}\label{nav:g-V-3}


\subsection{Congruences for a prescribed signed theta period}\label{nav:V-27}
There is a second way to choose the two narrow-class representatives
of Lemma~\ref{lem:good175}. It starts with the original
test levels but chooses the representatives jointly with the norm
element. This permits simultaneous congruences for a signed sum of
columns. We use the fixed splittings and finite-order determinant
tuples of Theorem~\ref{thm:even233}.

\begin{proposition}[Two auxiliary prime ideals and prescribed norm congruences]
\label{prop:ray233}
Fix a period field \(E/F\) in the interval of Theorem~\ref{thm:conductor232}, and a finite list of the prescribed
split tests and determinant characters. Put
\(\mathfrak b=\mathfrak d_F\mathfrak a_0^{-1}\). Let \(\mathcal L\)
be a finite set of odd rational primes disjoint from the rational
support of \(\mathfrak b\), with positive integers \(h_l\) for
\(l\in\mathcal L\). One can choose the representatives
\(\mathfrak p,\mathfrak q\), totally positive \(\beta,\gamma\) with
\((\beta)\mathfrak d_F=\mathfrak p\),
\((\gamma)\mathfrak a_0=\mathfrak q\), and \(\xi\in E^\times\)
so that
\[
   u=\xi\bar\xi=(d_0\beta\gamma)^{-1},\qquad
   \xi\equiv 1\pmod{l^{h_l}\OO_{E,w}}
                   \quad(l\in\mathcal L,\ w\mid l).
\]
At every split place choose a position component, keeping all the
original choices. In that component \(\xi\) fixes every prescribed
split position test away from \(\mathfrak p\) and all other split
lattice indicators. At \(\mathfrak p\) it acts by a scalar unit
in the position space. The determinant characters are unramified
there. The only rational primes supporting \((\xi)\) are those
under \(\mathfrak b\), \(\mathfrak p\), and \(\mathfrak q\).
\end{proposition}
\begin{proof}
Use the prime selection of Theorem~\ref{thm:common233},
to represent the narrow class of \(\mathfrak a_0\) by a split
\(\mathfrak q\) of norm larger than every prescribed rational prime,
including every prime in \(\mathcal L\). This gives
a positive \(\gamma\). At every \(v\) supporting \(\mathfrak b\)
let \(w_v\) denote the position prime and set
\(J=\prod_{v\mid\mathfrak b}\bar w_v^{v(\mathfrak b)}\).
All these places split by Theorem~\ref{thm:scale232}.
At \(\mathfrak q\) choose the prime \(Q\) opposite its position
component. The fractional ideal \(JQ^{-1}\) has norm
\(\mathfrak b/\mathfrak q\) in \(F\).

At each prescribed nonspherical split place, including
\(\mathfrak q\), and each split divisor of \(\mathfrak b\), put
only its position prime in a ray modulus with a positive exponent
large enough that the corresponding scalar units fix the tests.
Also include the full ideals \(l^{h_l}\OO_E\) for \(l\in\mathcal L\),
taking the larger exponent where requirements coincide. This
modulus is coprime to \(JQ^{-1}\). Represent the ray class of
\(JQ^{-1}\) by a prime \(P\) of absolute prime norm greater than
all the specified rational primes and the rational primes ramified
in \(E\) or in the determinant characters. The construction over
\(E\) in Theorem~\ref{thm:common233}, gives such a
prime. Its rational prime splits completely. Set
\(\mathfrak p=P\cap F\). Equality in the ray class group gives
an element \(\xi\) congruent to one modulo the ray modulus with
\((\xi)=JQ^{-1}P^{-1}\). Thus
\((u)=\mathfrak b/(\mathfrak q\mathfrak p)\), and \(u\) is
totally positive. Defining \(\beta=(d_0\gamma u)^{-1}\) gives
\((\beta)\mathfrak d_F=\mathfrak p\) by \(\mathfrak a_0^2=(d_0)\).
Take the position component at the new \(\mathfrak p\) to be
\(\bar P\), where \(\xi\) is a unit. At other split places outside
the modulus both components are units, except at primes opposite
those already protected. This proves all the claims.
\end{proof}

\begin{lemma}[Parity forced by the internal unitary centers]
\label{lem:parity233}
In an even chain with the splittings fixed, all columns giving
nonzero operators have the same \(D_{\tau,c}\bmod 2\) at each real
place \(\tau\). This holds for every choice of the finite-order
intermediate determinant characters.
\end{lemma}
\begin{proof}
Take the norm-one idele equal to \(-1\) at \(\tau\) and to one
elsewhere. Act on every test by this source center. A homogeneous
real polynomial on edge \(i\) is multiplied by
\((-1)^{d_{i,\tau,c}}\) times the vacuum character at \(-1\).
Multiplication through the chain gives a fixed phase times
\((-1)^{D_{\tau,c}}\). On the other hand, transferring the centers
on the odd edges to their targets and changing the odd integration
variables, exactly as in Theorem~\ref{thm:even233}, gives
a phase depending only on the splittings. Each finite-order character
is trivial on the connected real scalar circle. Neither endpoint
is changed. Comparing the two multiples of any nonzero column
proves the assertion at \(\tau\).
\end{proof}

\begin{theorem}[Simultaneous signed Gram congruences at cohomological depth]
\label{thm:signed233}
Use Proposition~\ref{prop:ray233} and retain the original terminal
Bessel character and any fixed complex-linear projection. In the
fixed-coordinate comparison let \(G,G'\) be the Gram matrices before
and after the change from \(\beta\) to \(\beta'=(d_0\gamma)^{-1}\).
For every pair of nonzero operator columns there are totally positive
\(b_{cc'}\in F^\times\) with
\[
  G_{cc'}=b_{cc'}G'_{cc'},\qquad
  b_{cc'}\equiv1\pmod{l^{h_l}\OO_{F,v}}\quad
                  (l\in\mathcal L,\ v\mid l).
\]
Zero operator columns give zero on both sides. In particular, let
\(s\) be any fixed vector of algebraic coefficients, including a
prescribed choice of signs. If \(G'\) has algebraic entries, take
a number field containing the entries, coefficients, their conjugates
and \(F\). At each prime \(\mathfrak l\mid l\) of this field one has
\[
    s^*G s-s^*G's\ \in\ l^{h_l}
       \sum_{c,c'}\OO_{\mathfrak l}\,\overline{s_c}s_{c'}G'_{cc'}
                     \qquad(l\in\mathcal L).
\]
On the interval of Theorem~\ref{thm:cutoff233}, one may
take every odd \(l\le Y\) prime to \(2F_\psi\) simultaneously and
\(h_l=\lceil V^{203}\rceil\). For the original finite levels and
polynomial test depths, both representative primes still have
logarithmic norms polynomial in \(x\). This depth kills the full
groups covered by that theorem on the same fields and intervals.
\end{theorem}
\begin{proof}
The argument of Theorem~\ref{thm:even233}, applies to
any norm element with the stated local test actions. Here the only
possible finite translation is at \(\mathfrak p\), and it belongs
to the central integral units. The intermediate characters are
unramified there, so every one has value one on this translation.
Their values on the archimedean rotations are also one. All inert
lattice vacua are carried to the new vacua since \(\xi\) is a unit
at every inert place. Thus each operator satisfies \(A_c=C a_c A'_c\)
with the same unitary phase and with
\(a_c=\prod_\tau(\tau u)^{D_{\tau,c}/2}\).

For nonzero columns the lemma gives
\((D_{\tau,c}+D_{\tau,c'})/2\in\ZZ\) for every \(\tau\). The phase
cancels in the Gram pairing. Its multiplier is therefore
\[
  b_{cc'}=\prod_\tau(\tau u)^{(D_{\tau,c}+D_{\tau,c'})/2}\in F^\times.
\]
The full congruence modulus at \(l\) is Galois stable. Both \(\xi\)
and \(\bar\xi\) are one modulo \(l^{h_l}\) at all places above
\(l\), so every \(\tau u\) is one there to the same depth. Taking
integer powers and multiplying proves the first congruence. Expanding
the two quadratic forms and subtracting proves the assertion for
the given coefficients, including at places where cancellation
increases the valuation of the sum. Explicitly, if any summand is
nonzero, put
\[
 t=\min_{c,c'}v_{\mathfrak l}(\overline{s_c}s_{c'}G'_{cc'})
              +h_l v_{\mathfrak l}(l),
 \qquad v_{\mathfrak l}(0)=+\infty .
\]
The difference has valuation at least \(t\). The ultrametric
inequality then gives
\(\min\{v_{\mathfrak l}(s^*G s),t\}
  =\min\{v_{\mathfrak l}(s^*G's),t\}\).
In particular, a computed valuation strictly below \(t\) on either
side establishes nonvanishing of the prescribed projection on both.

For the quantitative choice, \(q>V^{401}>Y\) and the support of
\(\mathfrak b\) lies above \(2F_\psi q\). Thus all the stated primes
are available. The added full rational ideals have log norm at most
\([E:\QQ]\lceil V^{203}\rceil Y\log(2Y)\), a fixed power of \(V\).
The remaining ray modulus also has polynomial logarithmic norm by
the finite-level bounds in Theorem~\ref{thm:common233}.
In its prime-selection argument choose the auxiliary cyclotomic
prime above \(Y\) and outside the bounded number of other forbidden
rational primes. Repeated application of Bertrand's postulate gives
such a prime of size \(O_\psi(Y)\). The log norm estimate for the
ray prime is therefore polynomial in \(V\) and hence in \(x\).
This proves the simultaneous assertion.
\end{proof}

The modulus in the last theorem controls the difference of the two
signed squared norms relative to the indicated fractional ideal.
It allows both norms to vanish. A lower bound for the designated
projected vector, and a uniform archimedean size estimate for the
generators subject to these ray congruences, are separate questions. The next theorem proves the uniform estimate for those generators.

\subsection{Real sizes and prime-power congruences for the same theta periods}\label{nav:V-28}
\begin{theorem}[Nearly equal conjugates at the full prescribed ray depth]
\label{thm:balanced234}
Keep the fields, intervals, original finite test levels and
coefficient primes of Theorem~\ref{thm:signed233}. Fix \(A>0\)
and put \(\delta_x=e^{-x^A}\). In Proposition~\ref{prop:ray233}
one can choose the two auxiliary prime ideals so that
the same \(\beta,\gamma,\xi\) satisfy all its finite congruences
and, for \(u=\xi\bar\xi=(d_0\beta\gamma)^{-1}\),
\[
 \begin{split}
 \max_\tau\left|\log\tau\gamma-f^{-1}\log N_{F/\QQ}\gamma\right|
             &\le\delta_x,\\
 \max_\tau\left|\log\tau u-f^{-1}\log N_{F/\QQ}u\right|
             &\le2\delta_x,\\
 \max_\tau\left|\log\tau\beta-f^{-1}\log N_{F/\QQ}\beta\right|
             &\le3\delta_x.
 \end{split}
\]
Here \(F=E_x^+\), \(f=[F:\QQ]\), and \(\tau\) runs through its
real embeddings. For fixed \(A\) and polynomial test depths the
logarithms of the two prime norms and the absolute logarithms
of all conjugates of \(\beta,\gamma,u\) are bounded by a fixed
power of \(x\), uniformly on the interval. In particular this
holds with every \(l\le Y,\ l\nmid2F_\psi\), at once, to the
depth \(h_l=\lceil V^{203}\rceil\) of Theorem~\ref{thm:signed233}.
\end{theorem}
\begin{proof}
Use the polynomial conductor \(N_x\) of Theorem~\ref{thm:conductor232}. Over \(F\) there are \(f-1\)
independent log vectors of totally positive units with Euclidean
norm at most \(2\sqrt f\,\varphi(N_x)\log(2N_x)\), by the proof
of Lemma~\ref{lem:unit179}. Their images in the CM
field \(E=E_x\) have the same absolute values. Given any finite
ray modulus \(\mathfrak m\), raising all these units to the order
of \((\OO_E/\mathfrak m)^\times\) puts them in the ray unit
group. This order is at most \(N\mathfrak m\). Thus the logarithm
of the length bound for independent ray-unit vectors is polynomial
in \(x\) whenever \(\log N\mathfrak m\) is. The same conclusion
holds over \(F\).

Apply Theorem~\ref{thm:sector234}, to \(F\), with
the narrow class of \(\mathfrak a_0\), finite modulus one,
\(y=0\) and \(\delta=\delta_x\). Choose its parameter \(B\)
polynomial in \(x\) and large enough that \(X\) exceeds all
the prescribed primes, the ramified rational primes and any
desired \(\exp(x^{A_0})\) with \(A_0\) fixed. This is possible
by the original level bounds and \(\log d_F\le f\log N_x\),
where here \(d_F\) denotes the absolute discriminant.
We obtain \(\mathfrak q=(\gamma)\mathfrak a_0\) with positive
\(\gamma\) and the first asserted inequality. The quadratic
character of this narrow class is trivial by Lemma~\ref{lem:good175}. Hence \(\mathfrak q\) splits in \(E/F\).

Form \(J,Q,\mathfrak m\) exactly as in Proposition~\ref{prop:ray233},
where \(Q\) is the prime over \(\mathfrak q\) opposite the chosen
position component. The log norm of the ray modulus is polynomial
in \(x\). The new prime \(\mathfrak q\) has polynomial log norm,
each ordinary test depth is polynomial and the full rational part
contributes at most \([E:\QQ]\lceil V^{203}\rceil Y\log(2Y)\).
Apply the same theorem over \(E\) to the fractional ideal
\(JQ^{-1}\), with \(y=0,\ \delta=\delta_x\). Increase its
polynomial parameter to exclude all primes already specified,
including the rational prime under \(\mathfrak q\).
We obtain a prime \(P=(s)JQ^{-1}\) of absolute degree one
and \(s\equiv1\bmod\mathfrak m\) with \(\|\lambda(s)\|<\delta_x\).
Its rational prime is unramified and \(E/\QQ\) is Galois, so that
prime splits completely. Set \(\mathfrak p=P\cap F\),
\(\xi=s^{-1}\) and \(\beta=(d_0\gamma \xi\bar\xi)^{-1}\).
The ideal and ray verifications in Proposition~\ref{prop:ray233}
now apply to these very elements. At infinity \(\log\tau u=
-2\log|\sigma s|\) for \(\sigma\mid\tau\). Center the identity
for \(\log\tau\beta\) by subtracting its mean. This proves the
other two inequalities. Polynomial bounds for the means follow
from \(N\mathfrak a_0=d_0^{f/2}\),
\(N_{F/\QQ}u=N\mathfrak b/(N\mathfrak p\,N\mathfrak q)\),
and \(N\mathfrak b\le d_F\). The same estimates and prime choices
apply for each \(x\) on the common cohomology interval.
\end{proof}

\begin{corollary}[Real comparison for the specified homogeneous tests]
\label{cor:parallel234}
In the fixed-coordinate comparison of Theorem~\ref{thm:signed233}
put \(d_c=\sum_\tau D_{\tau,c}\) and \(d_{\max}=\max_c d_c\).
For \(a_c=\prod_\tau(\tau u)^{D_{\tau,c}/2}\), set
\(a_c^0=(N_{F/\QQ}u)^{d_c/(2f)}\), taking the positive real root.
Then \(|\log(a_c/a_c^0)|\le d_c\delta_x\). With the same splittings
and the same terminal Bessel character as in that theorem, let
\(v_c,v'_c\) be corresponding vectors after any common fixed
linear projection. For arbitrary complex signs or coefficients
\(s_c\), the common phase \(C\) of Theorem~\ref{thm:even233},
satisfies
\[
 \left\|\sum_c s_c v_c-C\sum_c s_c a_c^0 v'_c\right\|
  \le (e^{d_{\max}\delta_x}-1)\sum_c |s_c|a_c^0\|v'_c\|.
\]
All the finite signed Gram congruences of Theorem~\ref{thm:signed233}
hold for these same data whenever its algebraicity hypothesis holds.
For the terminal scalar translates of one fixed homogeneous theta
operator the factor is exactly \(C a_c\) on their whole signed sum.
\end{corollary}
\begin{proof}
The second bound of the theorem, multiplied by each nonnegative
\(D_{\tau,c}/2\) and summed over \(\tau\), bounds the logarithm
of the ratio. For these ray primes the phase contributed by the
finite-order intermediate determinant characters is one, by the
proof of Theorem~\ref{thm:signed233}. Thus \(v_c=C a_c v'_c\).
The triangle inequality gives the estimate. For one fixed operator
the same identity holds on every input, so it applies directly to
any linear combination of its terminal input translates. The same
reasoning applies to a fixed linear projection on its output.

For the explicit real tests in Theorem~\ref{thm:harmonic188},
the factor \(d_{\max}\) is at most \(x^{C_0}\exp(x^\varepsilon)\)
for a fixed \(C_0\). Indeed the degree on an edge is the sum of
the absolute partition entries in that theorem. Their number, the
ranks, \(t_{\rm r}\), the number of edges and \(f\) are polynomial
in \(x\). Each entry is an original row parameter plus splitting
and rank shifts. The lists in equation~\eqref{eq:letters111},
and Theorem~\ref{thm:tower170}, use the \(B_\tau\)
of size \(O(x\exp(x^\varepsilon))\) and polynomial spacings.
At most a polynomial number of these terms occurs in an entry.
This proves the bound. For these tests, taking \(A>\varepsilon\)
makes the relative error factor at most \(e^{-x^A/2}\) for large
\(x\). This controls the change in real scale. The arithmetic size
of a nonzero reference Bessel vector still enters the lower bound
in Section~\ref{sec:missing}.
\end{proof}

\begin{theorem}[Cancellation among different degree profiles]
\label{thm:sampling234}
Fix the reference character \(\beta'=(d_0\gamma)^{-1}\), one
finite family of homogeneous columns, their inputs, a terminal
linear projection and complex coefficients \(s_c\). Fix the
reference tests before varying the second prime, using spherical
position tests away from a prescribed finite set. Write \(v'_c\)
for the projected reference vectors, \(\mathbf d_c=(D_{\tau,c})_\tau\),
and form the vectors
\[
 U_{\mathbf d}=\sum_{\mathbf d_c=\mathbf d}s_c v'_c .
\]
There exists a choice of the second auxiliary prime with nonzero
signed projected vector if and only if some \(U_{\mathbf d}\ne0\).
Here one varies only the norm comparison of
Theorem~\ref{thm:balanced234}, using the transported splittings
from the common reference, the original Bessel character and
the finite congruences of Theorem~\ref{thm:signed233}.
One can retain any prescribed accuracy \(0<\eta\le1/10\) in the
centered logarithms of \(u\). If the number of distinct profiles,
\(\log(2+d_{\max})\), and \(\log(1/\eta)\) are polynomial in \(x\),
the logarithmic prime norms are still polynomial in \(x\).
\end{theorem}
\begin{proof}
Use \(K=E,\ n=2f\), and fix \(A_0=JQ^{-1}\) and the ray modulus
from the proof of Theorem~\ref{thm:balanced234}. For a selected
prime \(P=(s)A_0\) the signed vector, up to a common unit scalar,
is the value of the vector polynomial
\[
   \sum_{\mathbf d}U_{\mathbf d}\prod_\tau t_\tau^{d_\tau},
              \qquad t_\tau=|\sigma s|^{-1},\quad \sigma\mid\tau .
\]
Indeed the real multiplier is exactly \(a_c\). At all protected
split tests the chosen component of \(s^{-1}\) is one to the
prescribed depth. At \(P\cap F\) take the unit position component
as in Proposition~\ref{prop:ray233}. Other fresh places carry the
reference spherical tests. This proves necessity.

Let \(h\) be the number of profiles and set \(L=h(h-1)/2\).
Choose integers \(0\le b_\tau\le L\) such that the numbers
\(e_{\mathbf d}=\sum_\tau b_\tau d_\tau\) are distinct. Each pair
of distinct profiles rules out at most a fraction \(1/(L+1)\)
of the integer box, by fixing all but a coordinate where they
differ. The union of the \(L\) forbidden sets leaves a choice.
Put \(D=d_{\max}\), \(M=\lceil4h(1+L)(1+D)/\eta\rceil\), and
\(z_{\mathbf d}=\exp(e_{\mathbf d}/M)\). These numbers have
mutual spacing at least \(1/M\), and \(1\le z_{\mathbf d}^{\,j}
\le2\) for \(0\le j<h\). For any vectors \(W_{\mathbf d}\) we have
\[
   \max_{0\le j<h}\left\|\sum_{\mathbf d}z_{\mathbf d}^{\,j}
            W_{\mathbf d}\right\|
    \ge (3M)^{1-h}\max_{\mathbf d}\|W_{\mathbf d}\| .
\]
To extract a given vector, take the linear combination of the
\(h\) sums with coefficients of its Lagrange polynomial on the
nodes \(z_{\mathbf d}\). The sum of the absolute coefficients is
at most \((3M)^{h-1}\) by multiplying the linear factors.

Choose \(\delta=\min(\eta/100,[100h(1+D)(3M)^h]^{-1})\).
Select \(B\) meeting all the hypotheses of Theorem~\ref{thm:sector234}, for this accuracy, large enough to
exclude the given places and to have \(n+h\le B\).
Let \(X_0=\exp((2B)^{20})\), \(t_0=(NA_0/X_0)^{1/n}\), and
\(\bar b=f^{-1}\sum_\tau b_\tau\). For each \(j=0,\ldots,h-1\)
prescribe
\[
   X_j=X_0\exp(-n j\bar b/M),\qquad
   y_{\tau,j}=-j(b_\tau-\bar b)/M .
\]
Use the theorem with parameter \((\log X_j)^{1/20}\ge B\).
It gives \(|\log(t_\tau/t_0)-j b_\tau/M|\le2\delta\).
The weighted lower bound \(X_j e^{-O(B^2)}\) in its proof dominates
the weight \(O(hn\log X_j)\) at all earlier choices. Thus the
primes can be taken above distinct rational primes, and the
position halves at these spherical places can be fixed simultaneously.
The centered logarithms of \(u\) have modulus at most
\(2j\max_\tau|b_\tau-\bar b|/M+2\delta<\eta\).

Now take \(W_{\mathbf d}=t_0^{|\mathbf d|}U_{\mathbf d}\), where
\(|\mathbf d|=\sum_\tau d_\tau\). Replacing an ideal sample
\(z_{\mathbf d}^{\,j}\) by the actual product of the \(t_\tau/t_0\)
changes each sum in norm by at most
\(2h(e^{2D\delta}-1)\max_{\mathbf d}\|W_{\mathbf d}\|\).
This is at most half the lower bound just obtained. Thus one
of the actual signed vectors has norm at least
\(\max_{\mathbf d}t_0^{|\mathbf d|}\|U_{\mathbf d}\|/
[2(3M)^{h-1}]\). Finally \(\log(1/\delta)\ll
h\log(2h(1+D)/\eta)\), so the choices of \(B\) and
\(\log X_j\) have the asserted polynomial bounds.
\end{proof}

All scalar translates of one fixed word have the same degree
profile. In that case the vector \(U_{\mathbf d}\) is the signed
reference projection itself, and a change of scale preserves its
vanishing exactly. For the normalized Bessel periodization,
Theorem~\ref{thm:dyad234}, gives a test at the
specified terminal character.

\subsection{From the adelic theta distribution to CM derivatives}\label{nav:V-29}
Fix one of the fields \(E=E_x\) of Section~\ref{sec:ranks47} on the
interval of Section~\ref{sec:missing}. Put \(F=E^+\),
\(K=\QQ(\sqrt{-d_K})\), where \(d_K\) is the ramified prime denoted
by \(d_0\) in Theorem~\ref{thm:scale232}, and \(f=[E:K]\).
Consider two of the standard definite spaces of ranks \(r,s\) with
\(rs\) even. Use all extensions \(\Sigma\) of the chosen embedding
of \(K\), the matrix space \(W=\operatorname{Mat}_{s\times r}(E)\)
over \(\QQ\), and \(D=frs\). Write \(h(a,b)=\operatorname{tr}(b^*a)\)
and \(h_K=\operatorname{Tr}_{E/K}h\). In the marked elliptic
differential of Lemma~\ref{lem:basis59}, put
\[
 v(a)=\delta\Phi(a),\quad A=|\delta|^2\sqrt{d_K}/2,\quad
 c=\pi/A,\qquad B(a,b)=\operatorname{Tr}_{E/\QQ}
                         \bigl(-h(a,b)/\sqrt{-d_K}\bigr).
\]
Thus \(B(a,b)=-\imag(v(a)\cdot\overline{v(b)})/A\).
This is the alternating form for the character of
Theorem~\ref{thm:scale232}. Indeed, multiplying the skew scalar
\(\gamma\sqrt{-d_K}\) by \((d_K\gamma)^{-1}\) gives
\(-1/\sqrt{-d_K}\), and composing the two traces gives \(B\).
Use the symmetric Weyl law with central term \(B/2\), and the
rational adelic character \(\psi\) with \(\psi_\infty(t)=e^{-2\pi it}\).
The oriented matrix action and the half factor can also be found in
\srcloc{Formalism/LocalToGlobal/IndefiniteReal/LocalNewformFunction}{25}.
For the original character one transports the actual functions and
splittings by Theorem~\ref{thm:transport233}.

\begin{lemma}[Mean square of the adelic theta distribution]
\label{lem:zak242}
Choose a rational polarization of \(W\) into a position space \(X\)
and its dual \(X^\vee\). Use the dual adelic measures with both
rational quotients of volume one. For a Schwartz--Bruhat vector
\(l\), put \(\vartheta(l)=\sum_{q\in X(\QQ)}l(q)\).
If \(R(w)\) denotes symmetric Weyl translation, then
\[
 \int_{W(\QQ)\backslash W(\mathbb A)}
                  |\vartheta(R(w)l)|^2\,dw=\|l\|_2^2 .
\]
If the finite factor of \(l\) is an eigenvector for the Weyl
translations of a full finite lattice \(\widehat M\), the left
side is also the average on \(W_{\RR}/M\) of the same expression
with finite argument zero, where \(M=W(\QQ)\cap\widehat M\).
\end{lemma}
\begin{proof}
For \(w=(x,y)\) the modulus of the summand after periodizing is
the modulus of
\(\sum_q l(q+x)\psi(\langle q,y\rangle)\). The factors removed
are common to the entire sum and have modulus one. Replacing \(x\)
by a rational translate reindexes \(q\) and replacing \(y\) by a
rational frequency leaves every character on \(q\) fixed. These
characters of \(X^\vee(\mathbb A)/X^\vee(\QQ)\) are mutually
orthogonal for probability Haar measure. The sum converges
absolutely for every \(x\), since the finite support condition
puts \(q\) in finitely many translates of a real lattice and
the real factor decreases rapidly. Integration in \(y\) gives
\(\sum_q |l(q+x)|^2\). Tonelli's theorem and integration over a
fundamental domain in \(x\) give precisely \(\int |l|^2\).
For the last assertion the modulus is unchanged by finite
translations in \(\widehat M\), by the eigenvector hypothesis and
the Weyl law. Every finite adele is rational modulo \(\widehat M\),
by the Chinese remainder theorem in lattice coordinates. The
kernel of the resulting map from \(W_{\RR}\) to the double
quotient is \(M\). The pushforward of Haar probability to this
torus is Haar probability.
\end{proof}

Let \(M=\mathfrak b^{-1}\operatorname{Mat}_{s\times r}(\OO_E)\)
be the standard module in Lemma~\ref{lem:basis59}, and
\(\phi_0\) its unit finite vacuum with the corrected period character.
Write \(G_\phi\) for the linear realization of a finite vector
\(\phi\) in the CM sections of Theorem~\ref{thm:one61}. In
particular \(G_{\phi_0}=G_0\), the image of the product of the
marked elliptic theta functions under the normalized lattice
intertwiner. Here \(\nu\) denotes the averaged squared norm of
\emph{one} marked elliptic theta in Proposition~\ref{prop:ellnorm56}.
Use the antiholomorphic Fock realization at infinity. Its vacuum
\(1\) has norm one for Gaussian probability of density
\((c/\pi)^D e^{-c|w|^2}\) in the coordinate \(w=v(a)\).
A polynomial is written \(P(\sqrt c\,\bar w)\), so its squared
norm is the Fischer norm squared of \(P\).

\begin{theorem}[The theta functional in the marked elliptic differential]
\label{thm:vac242}
Fix the Schr\"odinger--Fock identification once. There is a common
scalar \(C\), of absolute value \(\nu^{-D/2}\), such that for every
smooth finite Schwartz vector \(\phi\) and every homogeneous
polynomial \(P\) of degree \(d\),
\[
 \vartheta\bigl(P(\sqrt c\,\bar w)\otimes\phi\bigr)
        =C\,c^{-d/2}\,P(\partial)G_\phi(0).
\]
The same \(C\) works for every translate, class representative and
degree. The differential operator uses the coefficients of \(P\)
as written. The phase of \(C\) is fixed by the ordinary theta sum on the marked product vacuum.
\end{theorem}
\begin{proof}
Set \(K_z(\bar w)=\exp(c\,\bar w\cdot z)\) and
\(F_\phi(z)=\vartheta(K_z\otimes\phi)\). This is an entire
function. In real position coordinates the inverse Bargmann
transform of \(K_z\) is a Gaussian times the exponential of an
affine function of position with coefficients holomorphic in \(z\).
A fixed real symplectic change of polarization acts continuously on
Schwartz space. On compact sets of \(z\) every derivative is
therefore bounded in each Schwartz seminorm. This gives local
uniform convergence of the theta sum and its derivatives. These
are the Gaussian kernel and intertwining formulas of
\href{https://doi.org/10.1002/cpa.3160140303}{V.~Bargmann, Part~I}.

Let \(U_vH(z)=e^{-cz\cdot\bar v-c|v|^2/2}H(z+v)\), and let
\(\bar U_v\) act by the conjugate formula on antiholomorphic
functions. The real symmetric Weyl action is \(\bar U_v\).
Indeed \(U_vU_{v'}\) has multiplier \(e^{\pi iB(v,v')}\)
in the real coordinates just chosen, while \(\psi_\infty(B/2)\)
is its conjugate. The unitary Schr\"odinger--Fock intertwiner is
chosen for these translation operators. For \(a\in W(\QQ)\)
the full rational Weyl translation fixes \(\vartheta\), by
reindexing positions and using the triviality of \(\psi\) on
\(\QQ\). The finite corrected translate is
\(W_f(a)=e^{\pi i h_K(a,a)}R_f(a)\). On rational elements this
correction gives exactly the cocycle
\(e^{-2\pi i(\real h_K(a,b)+\imag h_K(a,b)/\sqrt{d_K})}\)
used in Theorem~\ref{thm:one61}. Transferring the real operator
across the rational invariance and applying its formula to \(K_z\)
gives
\[
 F_{W_f(a)\phi}(z)=e^{\pi i h_K(a,a)}U_{v(a)} F_\phi(z),
 \qquad
 \bar U_{-v}K_z=e^{-cz\cdot\bar v-c|v|^2/2}K_{z+v}.
\]
Thus \(F_{\phi_0}\) has the same theta multipliers as \(G_0\) on
\(v(M)\). This is a principal polarization by
Theorem~\ref{thm:lattice55}. Its holomorphic section space has
dimension one by Riemann--Roch and vanishing of higher cohomology
for an ample line bundle on a complex abelian variety. Hence
\(F_{\phi_0}=C G_0\). Apply Lemma~\ref{lem:zak242} to
\(1\otimes\phi_0\). The real translate of \(1\) at \(-z\) is
\(e^{-c|z|^2/2}K_z\). Consequently
\(\int_{\,\CC^D/v(M)} e^{-c|z|^2}|F_{\phi_0}(z)|^2\,dz_{\rm prob}=1\).
The corresponding integral for \(G_0\) is \(\nu^D\) by
Theorem~\ref{thm:transport59}. This proves the asserted modulus.

In the finite lattice model the translates \(W_f(a)\phi_0\), one
per coset in \(W(\mathbb A_f)/\widehat M\), form its standard
orthonormal basis. One can see finiteness of the expansion of a
smooth vector by restricting its period character to a sufficiently
small open subgroup of \(\widehat M\) fixing it. Only the cosets in
the annihilator of this subgroup can occur, by the Weyl commutator.
There are finitely many, since \(\widehat M\) is self-dual. Rational
representatives exist by the lattice calculation in the lemma.
The translation identity and linearity give \(F_\phi=C G_\phi\)
for every such vector. To recover the phase, apply the reverse
normalized lattice sum of Lemma~\ref{lem:basis59} to \(\phi_0\),
obtaining \(\phi_{\rm prod}\) on the product lattice. On sections
the reverse sum recovers the original section. Averaging over the source periods
kills every nonzero coset of the target modulo the intersection,
by the commutator. The two indices agree by self-duality, so the
zero coset has coefficient one after the two normalizations.
Thus \(G_{\phi_{\rm prod}}\) is the marked product section, of
value one at zero, and \(C=\vartheta(1\otimes\phi_{\rm prod})\).
Finally \(P(\partial_z)K_z|_{z=0}
=c^{d/2}P(\sqrt c\,\bar w)\). Applying this identity inside the
theta sum proves the formula.
\end{proof}

\begin{theorem}[A compact theta period with its finite projections]
\label{thm:edge242}
Write \(H=U_r,\ H'=U_s\) for the two definite groups and
\([H]=H(F)\backslash H(\mathbb A_F)\), with analogous notation
for \(H'\). Let \(a,b\) be continuous forms of right finite
level characters \(\chi,\chi'\) on \(J,J'\), respectively.
Use their given values, the chosen splitting \(\omega\), the
polynomial \(P\) and the finite test \(\phi\).
For finite class representatives \(u,v\), set
\[
 m_u=\frac{\operatorname{vol}(H_\infty\times J)}
                {|H(F)\cap uJu^{-1}|},\qquad
 m'_v=\frac{\operatorname{vol}(H'_\infty\times J')}
                {|H'(F)\cap vJ'v^{-1}|}.
\]
In the denominator only the finite components of a rational
element are tested. Use probability Haar measure for the following
averages, with \(k,l\) denoting the real variables.
\begin{align*}
 \phi_{uv}&=\omega_f(u,v)\int_{J\times J'}
       \chi(j)\overline{\chi'(j')}\,\omega_f(j,j')\phi\,dj\,dj',\\
 P_{uv}&=\int_{H_\infty\times H'_\infty}
            a(k,u)\overline{b(l,v)}\,\omega_\infty(k,l)P\,dk\,dl .
\end{align*}
Here the real action is written on polynomials in \(\sqrt c\,\bar w\),
including its vacuum character. If \(\Theta_P\) denotes the adelic
theta kernel of \(P\otimes\phi\) in these coordinates, then
\[
 \int_{[H]\times[H']}a(g)\overline{b(h)}\,\Theta_P(g,h)\,dg\,dh
  = Cc^{-d/2}\sum_{u,v}m_um'_v\,P_{uv}(\partial)G_{\phi_{uv}}(0).
\]
The formula includes, by linearity, a finite signed sum of tests.
For a lift pairing the kernel with a conjugate input, take \(a\)
to be that conjugate, with its resulting level character.
\end{theorem}
\begin{proof}
The quotients are compact and have finitely many finite classes.
Above \(u\) the group \(H(F)\cap uJu^{-1}\) acts on
\(H_\infty\times J\) by left translations, freely, through its
real component and its finite component conjugated by \(u\). It
is finite since rational points are discrete in the adeles.
This gives exactly the mass \(m_u\) for normalized averages on
this fiber. The same argument gives \(m'_v\). The Weil action
factors into its real and finite actions in the chosen splitting.
Transforming \(a\) and \(b\) under \(j,j'\) at the finite level
gives precisely the two character factors displayed above.
This finite average is a finite sum, by smoothness and compactness.
At infinity the unitary action preserves the Gaussian and rotates
the coordinates, up to its vacuum character. It preserves degree
and acts on a finite-dimensional space. Compact integration
therefore gives exactly \(P_{uv}\). The theta distributions of
these polynomials converge uniformly on the compact groups by the
Schwartz estimates in Theorem~\ref{thm:vac242}. Fubini's theorem
and that theorem give the identity with the original measures.
\end{proof}

With \(\operatorname{vol}(H_\infty\times J)
=\operatorname{vol}(H'_\infty\times J')=1\), the masses are exactly
the rational weights of \eqref{eq:pair55}. The formula identifies
each ordinary derivative and its measure factor before any square or
class sum is taken. For a general compact type the polynomial to be
differentiated is \(P_{uv}\). A valuation estimate proved for powers
of minors applies when those are the polynomials occurring in the sum.

\subsection{Local divisibility of the prescribed compact coefficient}\label{nav:VIII-34}
Use the rectangular compact period \(\mathcal P=\langle L(Q),b_*\rangle\)
of Theorem~\ref{thm:fixed243}, with \(b_*\) the specified terminal
test and \(B_v\) its covectors there. Take the degrees
\eqref{eq:degrees73} and restrict the target to its standard
hyperplane. Let \(S\) be the set of good inert primes used in
Theorem~\ref{thm:count73}. At their finite places the original
levels lie in the standard integral groups. Choose representatives
away from these primes by Theorem~\ref{thm:class61}. Use a common
number field containing the arithmetic coefficients and their
complex conjugates, with valuations normalized by \(v(p)=1\).

\begin{theorem}[A local bound for the integral against the fixed test]
\label{thm:div243}
Suppose the unaveraged finite tests and their class translates
satisfy the CM hypotheses of Theorem~\ref{thm:count73}. At every
place above \(S\) suppose the class values \(a_u\) are integral
and each \(B_v\) is integral on the module generated by the minor
products \(P_T\) of its branch. This hypothesis refers to algebraic
coefficients at the chosen representatives at the original levels.
It suffices for the specified terminal compact type to have an
integral matrix coefficient presentation as at the end of
Theorem~\ref{thm:fixed243}. The finite level characters have
finite order.

Choose open subgroups \(J_0\subset J,\ J_{*,0}\subset J_*\)
fixing the incident finite tests and killing the finite level
characters. Set \(I_0=[J:J_0][J_*:J_{*,0}]\). Let \(A_r,A_*\)
be common multiples of the respective stabilizer orders at the
original levels. For every \(p\in S\) for which the 2-primary
component of \(\delta_p\) is nonzero and every valuation \(v\)
above \(p\), the algebraic number
\[
 Y=I_0 A_r A_*\, d_V H\,
                \frac{c^{k/2}}{\kappa\kappa_* C}\,\mathcal P
\]
satisfies
\[
 v(Y)\ \ge\
 r\sum_{\substack{j\ge1\\j\ {\rm odd}}}
       [\,d_{p^j}(a,s)-w_{t+j\delta_p}\,]_+ .
\]
If the two finite character averages together fix the given
test already and the projected sections satisfy the stated CM
hypotheses, one can use \(I_0=1\). The inequality also holds
for integral signed combinations at the same degrees, branch
and normalization, using common subgroups and stabilizer multiples.

For the ranks of Proposition~\ref{prop:ranks47}, when \(r\ge xm\),
the entire contribution of the compact representation factor is
bounded uniformly by
\[
                0\le \frac{\log(d_V H)}r
                            \ll_{\psi,\eps}\frac{b}{\sqrt x},
                    \qquad b=\exp(x^\eps).
\]
All the assertions use the original fields, intervals and terminal
hyperplane.
\end{theorem}
\begin{proof}
The target test \(b_*\) is already in the range of its finite
orthogonal projector. In its pairing one can therefore integrate
the unprojected lift against \(b_*\). Cut up these same source and
target quotients at \(J_0,J_{*,0}\). The integrals and their
original measures do not change. At the refined levels the finite
averages of the test are identity operators. The two group-volume
factors at those levels multiply to \(\kappa\kappa_*/I_0\).
Apply the first identity of Theorem~\ref{thm:fixed243},
still using the originally specified \(a,b_*,Q\).

Choose the refined representatives also by
Theorem~\ref{thm:class61}, supported at primes completely split
in \(E\) away from the bad places. To compare representatives
at the two levels, write a refined representative as
\(\gamma u j\) with \(\gamma\) rational and \(j\) in the old
level, and likewise at the target. At every place over \(p\in S\)
both representatives have component one. Hence \(\gamma\) and
its inverse have integral matrices at every such place. This
remains true for all their algebraic conjugates at valuations
above \(p\), since integrality holds at all places in question.
On the source the change of the source coefficient is a
value of a determinant character, a unit at those valuations.
Indeed the scalar source action on the full-column minor is a
determinant power at each real place. Changing variables in the
compact definition of \(a_u\) gives that very character, together
with the finite level factor. On the target the rational change
of representative acts on \(B_v\) by composition with its
compact matrix action or its inverse. Cauchy--Binet expresses
each minor of this substitution as an integral linear combination
over the local ring of the minors of the same size. The matrices
act in the first \(s-1\) rows and preserve the last row, so they
preserve the submodule with \(l_\tau\) occurrences of that row.
Products and the determinant character of the vacuum preserve
the resulting module. The finite level multipliers are roots
of unity. This proves the asserted integrality at every refined
representative, also for complex conjugates of the covector
values since all places above \(p\) were included.

Every refined stabilizer embeds into its old stabilizer after
rational conjugation, so the original \(A_r,A_*\) clear its
denominator. For each refined pair the full target representative
is the block embedding \(v\oplus1\) of the chosen hyperplane
representative. It acts at the same good split places. The
derivative estimate in Theorem~\ref{thm:count73} is an estimate
for each individual minor product on the two moved lattices and
therefore also holds on this subset of pairs and minor indices.
In the first identity of Theorem~\ref{thm:fixed243} multiply by
\(A_rA_*\), and use the refined volume factors. Each derivative
is now multiplied by an element integral at \(v\). The
nonarchimedean triangle inequality gives the displayed bound
even when the signed sum vanishes. If the averaged finite data
satisfy the hypothesis at the original levels the same calculation
can be made there directly.

For the last assertion \(H=\prod_\tau\binom{b_\tau}{l_\tau}\)
satisfies \(\log H\le(\log2)\sum_\tau b_\tau\). At the high
embedding \(b_t=a=\lfloor bx\rfloor\), and at the other
embeddings the sum is \(O_\psi(b(1+\log x))\) by
Proposition~\ref{prop:cost73}. For a highest weight
\(\mu=(b_\tau^{r-1},b_\tau-l_\tau,0,\ldots,0)\) of length
\(n=s-1\), the Weyl dimension formula gives
\(\dim V_\mu=\prod_{i<j}(\mu_i-\mu_j+j-i)/(j-i)\).
It follows that
\(\log d_V\le [E:K]\,n(n-1)\log(1+\max_\tau b_\tau)/2\).
The degree and ranks are bounded by fixed powers of \(x\) and
\(\log(1+\max b_\tau)=O_\eps(x^\eps+\log x)\). This logarithm
of the dimension is consequently bounded by a fixed power of
\(x\) and is \(o(b/x)\). Finally \(r\ge xm\) and
\(m\asymp\sqrt x\) give the claimed estimate.
\end{proof}

The estimate applies to one prescribed integral compact coefficient.
It records separately the two finite-level indices and applies with
\(v(0)=+\infty\). Using a product formula for this value still
requires nonvanishing. For the full row of
Corollary~\ref{cor:original200}, one must also identify the
incoming row covectors and control their arithmetic size on this
interval.

\subsection{CM derivatives on the prescribed final row}\label{nav:VIII-35}
Take the last forward pair \(U(Y_{m-1}),U(Y_m)\) on the period
field of Proposition~\ref{prop:short107}. Write \(n=n_{m-1}\),
\(s=n_m\), and \(f=[E:K]\), where \(K=\QQ(\sqrt{-d_K})\) is
the marked imaginary quadratic field. Both groups are definite
with the standard forms. Their rank difference is odd, so \(ns\)
is even. At each real place use the actual partition
\(\lambda_\tau\) of Lemma~\ref{lem:degree181}, with
\(L_\tau=|\lambda_\tau|\) and \(d=\sum_\tau L_\tau\).
Write \(M\) for the tensor over these places of the maps in
Proposition~\ref{prop:cartan244}, permuting factors into the
indicated source and target order. Put
\[
 P_*=\prod_\tau p_{\lambda_\tau},\quad
 \mathfrak d=\prod_\tau\dim S_{\lambda_\tau}(\CC^n)\,
                             \dim S_{\lambda_\tau}(\CC^s).
\]
Thus \(P_*\) is the prescribed corner polynomial of
Proposition~\ref{prop:minor181}. We use the Gaussian coordinates
of Theorem~\ref{thm:vac242}, with \(c=\pi/A\). The original
character on the interval and its actual tests are carried to these data
by the pointwise transport in Theorem~\ref{thm:transport233}.

\begin{theorem}[The period at the full highest weights of the last row]
\label{thm:last244}
Use finite class representatives \(u,v\) at fixing levels of the
source and target. Let \(\gamma_u,\gamma'_v\) be their stabilizer
orders, and \(\kappa,\kappa'\) the respective volumes of the
real groups times these levels. Write \(G_{uv}\) for the marked
CM section of the actual finite translate at this pair. For a
forward lift the source function in the integral is conjugated.
Suppose the source and target functions before conjugation, in
their matching compact types, have class functions
\[
 \sum_a\langle\rho_n(g)x_{ua},y_{ua}\rangle,\qquad
 \sum_b\langle\rho_s(h)x'_{vb},y'_{vb}\rangle .
\]
Use the compact actions with their actual splittings. Denote the
two unit highest tensors by \(t_n,t_s\), now products over the
real places, and set
\[
   Y_u=\sum_a\langle t_n,x_{ua}\rangle y_{ua},\qquad
   Z_v=\sum_b\langle t_s,x'_{vb}\rangle y'_{vb}.
\]
For the period \(\mathcal P_*\) of the forward lift against the
target function with test \(P_*\), one has exactly
\[
 \frac{\mathfrak d\,c^{d/2}}{\kappa\kappa' C}\,\mathcal P_*
   =\sum_{u,v}\frac{
         M(Y_u\otimes Z_v)(\partial)G_{uv}(0)}{\gamma_u\gamma'_v},
           \qquad |C|=\nu^{-fns/2}.
\]
The complex identity allows the actual incoming class coefficients
to have arbitrary size. The right side is an algebraic bilinear
form in the tensor entries when the finite test has algebraic
coefficients in the lattice model of Theorem~\ref{thm:one61} with
its indicated splittings. It applies to signed inputs by linearity
in the displayed tensor coordinates.
\end{theorem}
\begin{proof}
At a fixed finite class pair, integrate the polynomial in the
archimedean oscillator against each pair of compact matrix
coefficients. The source function enters conjugated by the
forward convention, and the target function enters conjugated
by the Petersson pairing. Proposition~\ref{prop:cartan244}
and its tensor product give exactly
\(M(Y_u\otimes Z_v)/\mathfrak d\). Components of either input
of nonmatching type make zero contribution by compact
orthogonality. Every compact-finite function in the matching
type has such a finite matrix coefficient expansion. The
quantity extracted from a class is independent of the expansion,
as the one-group matrix-unit integral applied to the highest
tensor computes it, after multiplication by the representation
dimension.

Apply Theorem~\ref{thm:vac242} to the resulting polynomial at
each pair. The class measures are \(\kappa/\gamma_u\) and
\(\kappa'/\gamma'_v\) by Theorem~\ref{thm:edge242}. This proves
the identity. For a prescribed finite character projection one
can use the averaged test of that theorem at the original levels.
Alternatively, in a pairing with a test in the projected space,
refine inside those levels and use the fixing-level formula above.
The original integral is the same and the ratio of the two volume
products is the index product of Theorem~\ref{thm:div243}.
For algebraicity choose the fixing-level representatives as in
Theorem~\ref{thm:class61}. Their finite sections have algebraic
derivatives by Theorems~\ref{thm:one61} and \ref{thm:transport59}.
The determinant expansion and the rational stabilizer weights
then give algebraic coefficients of the bilinear form.

For a prescribed homogeneous polynomial \(P\) of the same block
degrees in place of \(P_*\), apply
Proposition~\ref{prop:cartan244} to each term \((a,b)\) before
summing. Its multiplier in the normalized formula becomes
\(\langle P,M(x_{ua}\otimes x'_{vb})\rangle_{\rm F}/
\prod_\tau h_{\lambda_\tau}\). In particular the formula keeps
the projections and phases of a fixed signed polynomial. When
the two right lines are highest and \(P\) satisfies the column
conditions of Theorem~\ref{thm:integral181}, the multiplier is
its value on the top identity frames. Indeed projecting to the
joint highest line retains the one diagonal monomial visible on
these frames, and \(P_*\) has value one there.
\end{proof}

\subsection{The marked elliptic derivatives at every prime}\label{nav:V-59}
Write \(d_0\) for the prime discriminant parameter of the interval
in Section~\ref{sec:missing}. Thus \(d_0\equiv7\bmod 8\) and
\(K=\QQ(\sqrt{-d_0})\). We use the discriminant-one elliptic curve
\(C\), its differential \(dz=dx/y\), and the mark \(t_*\) of
Lemma~\ref{lem:weber58}. At a place \(\mathfrak p\) above \(p\) let
\(v_{\mathfrak p}(p)=1\), and put
\[
 \kappa_p=
 \begin{cases}
  1 & p=3\text{ and }3\text{ is inert in }K,\\
  1/2 & p=d_0,\\
  0 & \text{otherwise}.
 \end{cases}
\]

\begin{theorem}[All-prime bound for the marked theta derivatives]\leavevmode
\label{thm:elliptic312}
For every \(n\ge0\) and every finite place of a coefficient field,
\[
 v_{\mathfrak p}\!\left(\Theta_{t_*}^{(n)}(0)\right)
                 \ge -\kappa_p\lfloor n/2\rfloor .
\]
In particular the ordinary derivatives are integral at \(2\),
including in the fixed discriminant-one differential.
\end{theorem}
\begin{proof}
We may pass to a finite local extension. Normalizing \(v_{\mathfrak p}(p)=1\)
makes the estimate descend. The CM \(j\)-invariant is integral.
The potential good reduction criterion gives a good integral
Weierstrass equation with affine coordinates \(X,Y\) and coefficients
\(a_i\). We can also take all the finitely many torsion points used
below to be rational. Write \(b_2=a_1^2+4a_2\). Completing the square
in this equation gives a cubic with quadratic coefficient \(b_2\)
and leading coefficient \(4\). Removing the quadratic term gives
\[
 x=u^2X+r,\qquad r=u^2 b_2/12,\qquad
 y=u^3(2Y+a_1X+a_3)
\]
for some \(u\) matching the specified differential. The two
discriminants are units and their ratio is \(u^{12}\), so \(u\)
is a unit. Thus \(dz=u^{-1}dX/(2Y+a_1X+a_3)\) generates the
invariant differentials over the local integer ring. Its dual
derivation \(\partial\) acts integrally on the structure sheaf.
For the criterion see Silverman, \emph{The Arithmetic of Elliptic
Curves}, Proposition VII.5.5.

The section \(t_*\) misses the origin in every special fiber.
For odd \(p\) this follows from injectivity of specialization
on two-torsion. At \(2\) the two idempotents in
\(\OO_K\otimes\ZZ_2=\ZZ_2\times\ZZ_2\) split the
two-divisible group into two height-one factors, one connected
in the special fiber and one \'{e}tale, as proved in
Lemma~\ref{lem:weber58}. The mark is the sum of the nonzero
points from the two factors. Its image in the \'{e}tale factor
remains nonzero. In particular \(X\) is regular at its reduction.
The identity
\((\log\Theta_{t_*})''(z)=-u_\Lambda-x(z+t_*)\) of
Theorem~\ref{thm:germ57} shows that every logarithmic derivative
at zero of order \(i\ge3\) is integral. Indeed \(\partial^{i-2}\)
kills the constant \(r\) and preserves regularity of \(u^2X\)
near the section.

It remains to estimate \(u_\Lambda+x(t_*)\). At \(p=2,3\) use the
CM expression from the proof of Theorem~\ref{thm:germ57}. Substitution
of the coordinate change, including its constant term, gives
\[
 u_\Lambda+x(t_*)=
 u^2\left(X(t_*)-\frac{1}{2d_0}
       \sum_{0\ne P\in C[\sqrt{-d_0}]}X(P)\right)
      +\frac{d_0+1}{2d_0}\,r .
\]
All points in the sum have order prime to \(p\) and their affine
coordinates are integral. Pairing \(P,-P\) in this odd-order kernel
makes the sum twice an integral element, since \(X(P)=X(-P)\).
The expression in parentheses is integral. At \(2\) the final
coefficient is divisible by \(4\), whereas \(v_{\mathfrak p}(r)\ge-2\).
At \(3\) we have \(v_{\mathfrak p}(r)\ge-1\), and if \(3\) splits
then \(-d_0\equiv1\bmod3\), so the coefficient is divisible by \(3\).
For \(p\ge5\) the logarithmic derivative of order two has valuation
at least \(-\kappa_p\) by the CM computation in
Theorem~\ref{thm:germ57}. The order-one logarithmic derivative is
zero. In the integral Leibniz recurrence for
\(\Theta_{t_*}^{(n)}(0)\) in that theorem, each occurrence of the
order-two derivative uses two of the \(n\) differentiations. This
proves the bound, including \(n=0\) by the normalization at zero.
\end{proof}

\subsection{Finite valuations of the CM splitting characters}\label{sec:cmplaces335}

The scalar in Theorem~\ref{thm:mixedfourier334} contains the value of
a unitary character at a uniformizer. For the splitting characters of a
period square one can determine its valuations from the CM infinity
type. This also specifies the coefficient primes at which the powers
of the uniformizer in the finite determinant average can affect an
integral period.

\Needspace{5\baselineskip}
\begin{proposition}[Valuations from the infinity type]\label{prop:cmplaces335}
Let \(E\) be a CM field, let \(\Sigma\) contain one complex embedding
from each conjugate pair, and let \(\nu\) be a unitary Hecke character
of \(E\) whose archimedean component is
\(\prod_{\tau\in\Sigma}(z_\tau/|z_\tau|)^{m_\tau}\) with integral
\(m_\tau\). Fix a finite prime \(\mathfrak p\) of \(E\) above \(p\)
and a local uniformizer \(\pi\). Put \(e_{\mathfrak p}=e(\mathfrak p/p)\)
and \(f_{\mathfrak p}=f(\mathfrak p/p)\). Then
\(u_{\mathfrak p}=\nu_{\mathfrak p}(\pi)\) is algebraic and every
algebraic conjugate in \(\CC\) has absolute value one. For an extension
\(v_\lambda\) to \(\overline{\QQ}\) of the valuation with
\(v_\lambda(\ell)=1\), it has valuation zero if \(\ell\ne p\).
At \(\ell=p\), let \(I_\tau=1\) if the prime on \(E\) induced by
\(v_\lambda\circ\tau\) is \(\mathfrak p\), and zero otherwise. Then
\[
 v_\lambda(u_{\mathfrak p})
   =-\frac{1}{2e_{\mathfrak p}}\sum_{\tau\in\Sigma}
               m_\tau (I_\tau-I_{\bar\tau}),\qquad
 |v_\lambda(u_{\mathfrak p})|\le
         \frac{f_{\mathfrak p}}2\max_{\tau\in\Sigma}|m_\tau|.
\]
\end{proposition}
\begin{proof}
Take \(h\ge1\) with \(\mathfrak p^h=(y)\), using finiteness of the
ideal class group. In the principal idele defined by \(y\), the finite
component is a unit at every place except \(\mathfrak p\), where it
is \(\pi^h\) times a unit. The values of a smooth character on the
units are roots of unity and only finitely many local restrictions to
units are nontrivial. Applying the character to this principal idele
and squaring gives
\[
       u_{\mathfrak p}^{\,2h}
       =\zeta\prod_{\tau\in\Sigma}
              \left(\frac{\tau y}{\bar\tau y}\right)^{-m_\tau}
\]
for a root of unity \(\zeta\). The right side is algebraic. Complex
conjugation on the source CM field pairs \(\tau\) with \(\bar\tau\),
and continues to pair them after any embedding of their images.
Every algebraic conjugate of the right side therefore has absolute
value one. At primes over \(\ell\ne p\) every conjugate of \(y\)
is a unit. At \(\ell=p\) one has
\(v_\lambda(\tau y)=h I_\tau/e_{\mathfrak p}\). Taking valuations
and dividing by \(2h\) gives the formula. Among all embeddings of
\(E\), exactly \(e_{\mathfrak p}f_{\mathfrak p}\) induce the prime
\(\mathfrak p\), since they correspond to the embeddings of \(E_{\mathfrak p}\)
in an algebraic closure of \(\QQ_p\). Bounding the sum by the total
number of those embeddings proves the inequality.
\end{proof}

On \(E_x\) use \(m=m_\alpha-m_\beta\) and the prime over \(v\)
specified by the split chart. Its residue field has size \(q_v\).
Thus the determinant and Fourier scalars of Theorem~\ref{thm:mixedfourier334}
are units away from the residue characteristic of \(v\), throughout
each interval of Section~\ref{sec:missing}.

\section{Marked elliptic jets and their determinants}\label{nav:chapter-split-elliptic-jets}

\subsection{Joint row degrees on the even Frobenius cycles}
\label{sec:inertrowdegrees461}

Fix a good \(p\ge5\), inert in \(K_d\) and unramified in \(E_x\),
and a definite edge with ranks \(r\le s\) and partitions
\(\nu_\tau\) of length at most \(r\). Use the permutation
\(\pi=\sigma_pc_E\) on the embedding labels of
Theorem~\ref{thm:direct76}. On each even cycle choose one
alternating half. Let \(H\) be the union of the chosen halves,
and \(L\) the union of their partners. Odd cycles belong to
neither set. We give the operator a lower bound for every such
choice and then maximize it.

For each \(\mu\in L\) form the list indexed by triples
\(\iota=(h,l,j)\), where \(h\in H,\ 1\le l\le r,\ j\ge1\) is
odd and \(\pi^j h=\mu\). The associated nonnegative integer is
\[
        D_\iota=d_{p^j}(\nu_{h,l},s),\qquad
        d_M(t,s)=\max\{0,\lceil(t-s(M-1))/M\rceil\}.
\]
Only indices with \(p^j\) at most the maximum block degree
need be used. Order this finite list decreasingly, writing its
members as \(D_{\mu,1},D_{\mu,2},\ldots\). Set \(\nu_{\mu,t}=0\)
for \(t>r\) and put
\[
 V_p(\nu;H)=\sum_{\mu\in L}\ 
     \max_{0\le k\le \#\{\iota:\pi^j h=\mu\}}
             \sum_{t=1}^k (D_{\mu,t}-\nu_{\mu,t}),
 \qquad V_p(\nu)=\max_H V_p(\nu;H).
\]
Both maxima include the zero contribution when the lists are
empty. These integers use only the permutation and the partitions.

\begin{lemma}[Simultaneous degrees of the exceptional coefficients]\leavevmode
\label{lem:jointdirections461}
In a monomial of the exceptional factors in the low symbol at
\(\mu\), let \(b_\iota\) count occurrences of the direction
indexed by \(\iota\). For any sublist \(I\) of size \(k\),
\[
             \sum_{\iota\in I} b_\iota
                            \le\sum_{t=1}^k\nu_{\mu,t}.
\]
\end{lemma}
\begin{proof}
After the high row substitution, the low symbol is
\(K_\mu(D+\sum_\iota y_\iota U_\iota)\), with a separate
indeterminate for each direction before its scalar Taylor
coefficient is restored. The matrices \(U_\iota\) are outer
products by Lemma~\ref{lem:rankoneshift449}. Treat the variables
outside \(I\) as coefficients. The columns supplied by the
perturbation from \(I\) lie in a space of dimension at most
\(k\), also after the right substitution and padding of the
flag symbol. A determinant of height \(v\) has joint degree at
most \(\min(k,v)\) in these \(k\) variables by its multilinear
expansion. Summing this bound over the diagram columns counts
the boxes in its first \(k\) rows. Tableau products span the
integral test lattice, so this proves the assertion for every
test, as in Lemma~\ref{lem:flagdegree456}. The statement holds
before specializing any of the formal variables.
\end{proof}

\begin{theorem}[The Weyl operator on independently oriented cycles]\leavevmode
\label{thm:cycleweyl461}
In the local integral frames of Theorem~\ref{thm:direct76},
the definite class operator sends its integral source Weyl
lattice into \(p^{V_p(\nu)}\) times its target Weyl lattice.
For any fixed \(H\), the summand at \(\mu\) is at least
\(\sum_\iota[D_\iota-\nu_{\mu,1}]_+\) over its list.
\end{theorem}
\begin{proof}
Fix the halves. Apply the operators in odd-cycle variables and
then set those variables equal to zero. In the factorization of
Theorem~\ref{thm:direct76}, blocks from different cycles and
matrix entries are orthogonal. An odd block vacuum depends only
on its own variables and its undivided derivatives at zero are
integral by Theorem~\ref{thm:jet63}. Derivatives of the initial
integral elliptic factor preserve integrality in the elliptic
parameters. Hence the product rule leaves on the even cycles
a sum of integral elliptic factors times the even block vacua.
The coefficients of the odd Weyl symbols and their row
parameters are integral. The argument of
Lemma~\ref{lem:rankoneshift449} now applies on each chosen
half: after removing its exponential of exceptional linear
terms in the low variables, all undivided low derivatives
satisfy the elliptic factorial bound in the high variables.
The Frobenius partner of \(h\) at exponent \(p^j\) is still
\(\pi^jh\) in its own cycle. No exceptional linear term joins
two cycles.

On a high derivative monomial let \(\alpha_{h,bl}\) be the
order in row \(b\) of column \(l\). The column sum over
\(1\le b\le s\) is \(\nu_{h,l}\). Put
\[
           C_\iota=\sum_b\left\lfloor
                              \alpha_{h,bl}/p^j\right\rfloor .
\]
In an individual term, restoration of the exceptional scalar
and extraction of high coefficients have valuation at least
\(\sum_\iota (C_\iota-b_\iota)\) after the undivided
differentiation. Indeed the factorials pay
\(\sum_\iota C_\iota\) above the even-power factorial bound,
and the coefficient \(g_{p^j}\) costs at most one more power
of \(p\) per occurrence, by the proof of
Theorem~\ref{thm:weylrowjet449}. Moreover \(b_\iota\le C_\iota\):
each occurrence consumes \(p^j\) powers of a single high
variable in that column. Division with remainder gives
\(C_\iota\ge D_\iota\).

For a fixed \(\mu\), retain just the \(k\) indices with largest
\(D_\iota\). The discarded differences \(C_\iota-b_\iota\)
are nonnegative. The lemma therefore bounds the retained sum
below by \(\sum_{t\le k}(D_{\mu,t}-\nu_{\mu,t})\).
Maximize independently at each \(\mu\). Every coefficient in
the independent row parameters satisfies the bound, so extracting
the source generators and pairing with the integral dual output
lattice proves the asserted inclusion for this \(H\). The class
operator is the same for each allowed \(H\). The block
factorizations used to prove the estimates all evaluate that
operator in the fixed physical coordinates of
Theorem~\ref{thm:direct76}. Taking the maximum is therefore valid.
Finally select the indices with \(D_\iota>\nu_{\mu,1}\); the
first \(k\) row lengths sum to at most \(k\nu_{\mu,1}\).
This gives the last assertion.
\end{proof}

\begin{corollary}[The extra divisor on the projected signed columns]\leavevmode
\label{cor:inertcorrection461}
Use \(q_c\) independent signed columns on the projected small
tests of Corollary~\ref{cor:shortcenter460}; here \(q_c\)
denotes their number. Let \(\nu_{ij}\) be the physical
partition tuple of column \(i\) at edge \(j\). At all
unramified \(p\ge5\) inert in \(K_d\), extend \(U_p\) by zero
when \(p\) is inert in \(K_q\). In the centered upper bound of
that corollary one can subtract the nonnegative quantity
\[
  \mathcal D_u=\frac{2}{q_c}
   \sum_{\substack{p\ge5\ {\rm unramified\ in}\ E_x\\
                         p\ {\rm inert\ in}\ K_d}}
       \left(\sum_{i=1}^{q_c}\sum_{j=1}^m V_p(\nu_{ij})
                           -q_cU_p\right)\log p .
\]
The sum is finite without a restriction on prime size.
\end{corollary}
\begin{proof}
In the signal signs, take \(H=\Phi_0\). The last assertion of
the theorem gives at least the row contribution of
Theorem~\ref{thm:weylrowjet449}, and Lemma~\ref{lem:wordfirstrow450}
bounds its sum over a word below by \(U_p\). In the other signs,
each \(V_p\) is nonnegative. For \(p\) larger than every block
degree and every relevant word cutoff the summand is zero,
by the definition of \(d_M\).

All primes of this finite support miss the marked levels.
Indeed every elementary test modification is at a rational
prime split completely in \(E_x\), as proved in
Lemma~\ref{lem:signalsupport452}. The simultaneous choice of
bases in that lemma works on this enlarged finite set: at
each prime it uses only that \(p\ge5\) is unramified in the
period field and inert in \(K_d\). Thus the translated elliptic
products have prime-to-\(p\) shifts with integral coefficients
and the class lattices match at every intermediate vertex.
Each path in column \(i\) gains \(\sum_j V_p(\nu_{ij})\) in
valuation before the masses and external scalars. Conjugating
a physical edge transports its partition tuple and its
Frobenius cycles together. They have the same sorted lists
and maximum because the field is abelian. Taking the signed
sums, with the same common coset on consecutive edges, preserves
the bound. Expansion of a minor adds the gains of its columns.

The two test lattices still have their original scalar, mass and
projector costs, as in Theorem~\ref{thm:fullsignal452}. The change
back to the fixed product bases preserves the local-degree
average by its phase calculation. The unramified inert primes
outside the signal signs previously had zero local field and
differential loss in Theorem~\ref{thm:outsidedifferent453}.
The terminal rescaling is a unit at all primes under consideration.
Summing the two ideal inequalities in
Theorem~\ref{thm:globalperiodcontent448} thus contributes twice
the indicated additional gains. Substitution in the centered
identity of Theorem~\ref{thm:polycenter460} gives the asserted
subtraction, with the same ramified term.
\end{proof}


\subsection{Elliptic theta derivatives at a split prime}
\label{sec:splitjets480}

Let \(C=\mathbb C/\Lambda\) be the marked elliptic curve with
complex multiplication by \(\mathcal O_{K_d}\) used in
Section~\ref{sec:ell57}. Write \(x=\wp,\ y=\wp'\),
\(\omega=dx/y\) and \(\partial=d/dz\).
Let \(t\) be its nonzero two-torsion mark and \(\Theta_t(0)=1\)
its even theta function in the Hermitian frame. Fix a prime
\(p\ge5\) split in \(K_d=\mathbb Q(\sqrt{-d_0})\), and a
coefficient place with good integral Weierstrass model and
nonvanishing reduction of \(\omega\). We may extend the
coefficient field to contain the indicated torsion points.
Write \(k\) for the algebraic closure of the residue field,
and use the same notation for reduced functions when working
on \(C_k\). Put
\[
 q_0(z)=-u_\Lambda-x(z),\qquad q_t(z)=q_0(z+t).
\]
Here \(q_t=(\log\Theta_t)''\) in the complex uniformization.

\begin{lemma}[A differential identity on the split reduction]\leavevmode
\label{lem:splitcmder480}
There is a nonzero scalar \(H_p\in k\), in the field of definition
of the reduced curve and differential, such that
\[
              \partial^p=H_p\partial,\qquad
              \partial^{p-1}q_t=H_pq_t
\]
on the reduced elliptic curve. In the first equality the powers
denote iterates on the rational function field.
\end{lemma}
\begin{proof}
The \(p\)-th power of a derivation is a derivation by the binomial
formula. Invariant vector fields on an elliptic curve form a line,
so invariance of \(\partial^p\) gives \(H_p\) over the field of
definition. For the moment it may be zero. Set
\(\Xi(f)=\partial^{p-1}f-H_pf\). Its derivative vanishes and
\(\Xi(\partial g)=0\) for every rational function \(g\).
The kernel of \(\partial\) on \(k(C)\) is \(k(C)^p\). Indeed
\(x\) is a separating coordinate and \(1,x,\ldots,x^{p-1}\)
are a basis over \(k(C)^p\), as in Lemma~\ref{lem:kappa62}.
Differentiating in this basis proves the assertion.
Consequently \(\Xi(q_0)\) is a \(p\)-th power. Its only possible
pole is the origin, of order at most \(p+1\), since \(x\) has
a double pole and the vector field is regular. A \(p\)-th root
would have at most a simple pole there. Riemann--Roch in genus
one makes that root constant.

We determine the constant using the original complex multiplication.
Take \(\alpha=\sqrt{-d_0}\). In characteristic zero put
\[
 B(z)=\frac1{\alpha}\sum_{0\ne P\in C[\alpha]}
              \frac{y(z)-y(P)}{2(x(z)-x(P))}.
\]
The zeta addition formula gives derivative \(x(z)-x(z+P)\)
for each fraction in the sum. The isogeny identity of
Section~\ref{sec:ell57} reads
\(\sum_{P\in C[\alpha]}x(z+P)=\alpha^2 x(\alpha z)+c_\alpha\),
where \(c_\alpha=\sum_{P\ne0}x(P)\) and
\(u_\Lambda=c_\alpha/(\alpha(\alpha-\bar\alpha))\).
Thus
\[
          \partial B(z)=\alpha q_0(\alpha z)-\bar\alpha q_0(z).
\]
This is an identity of rational functions with good reduction at
the chosen place. In fact \(p\nmid 2d_0\), every nonzero kernel
point specializes away from the origin, and all denominators
\(x-x(P)\) are nonzero rational functions on the special fiber.
The CM isogeny has degree prime to \(p\) and extends on the good
model. These facts also give integrality of \(u_\Lambda\).

Write \(a,a'\) for the reductions of \(\alpha,\bar\alpha\).
They are distinct elements of \(\mathbb F_p^\times\) by splitting.
The pullback of \(\omega\) by \([\alpha]\) is \(a\omega\).
It follows that
\(\Xi(a f\circ[\alpha])=a(\Xi f)\circ[\alpha]\), since \(a^p=a\).
Apply \(\Xi\) to the differential identity for \(B\).
Its left side vanishes and its right side is
\((a-a')\Xi(q_0)\), since \(\Xi(q_0)\) is constant.
Thus \(\Xi(q_0)=0\). Translation by \(t\) commutes with \(\partial\)
and gives the second identity of the lemma.

Finally, if \(H_p=0\), this argument would give
\(\partial^{p-1}x=0\). The function-field computation in the
proof of Lemma~\ref{lem:kappa62} shows that for a smooth cubic in
characteristic \(p\ge5\), \(\partial^p=0\) implies
\(\partial^{p-1}x\ne0\). In that computation nilpotence would
otherwise make both \(dx/y\) and \(x\,dx/y\) exact; their two
missing coefficients in the Legendre model cannot vanish
together. This portion of the argument uses only the smooth
cubic and nilpotence. We conclude \(H_p\ne0\).
\end{proof}

\begin{theorem}[All derivatives of a marked torsion translate at a split prime]\leavevmode
\label{thm:splitjets480}
Let \(F=U_w\Theta_t\), where \(w\) has order prime to \(p\) modulo
\(\Lambda\), and put \(J_n(F)=F^{(n)}(0)\) in the marked frame.
The undivided derivatives are integral at the chosen place and
satisfy, in its residue field,
\[
                    J_{n+p}(F)=H_pJ_{n+1}(F)
                           \qquad(n\ge0).
\]
This includes a translate vanishing at the origin. In particular,
for \(n\ge1\) write \(n=b(p-1)+r\), \(1\le r\le p-1\).
Then \(J_n(F)=H_p^bJ_r(F)\) in that field.
\end{theorem}
\begin{proof}
Use the ring of undivided derivative sequences from
Theorem~\ref{thm:hurwitz62}. Its multiplication uses binomial
coefficients and its derivation \(D\) shifts a sequence by one
place. The reference logarithmic derivative \(l=(\log\Theta_t)'\)
vanishes at zero and satisfies \(Dl=q_t\). Its higher entries are
integral since \(q_t\) is regular at the reduced origin and the
invariant vector field is integral. From \(D\Theta_t=l\Theta_t\)
and \(\Theta_t(0)=1\) all entries of the reference theta are
integral by induction. Theorem~\ref{thm:trans57} also gives
integrality of \(R_w=F/\Theta_t\) in an elliptic formal parameter,
including when its constant term vanishes.

Reduce these sequences. By the lemma,
\(D(D^{p-1}l-H_pl)=0\). A sequence killed by the shift is
constant. Its constant term here is zero, since \(l\) is odd
and \(p-1\) is even. Also \(l^p\) is constant by the derivation
rule and is zero at the origin. Repeated differentiation gives
\[
           D^p\Theta_t=(D^{p-1}l+l^p)\Theta_t
                       = H_pD\Theta_t .
\]
For the first equality, the product rule gives one contribution
per partition of \(p\) labeled positions, with a block of size
\(i\) contributing \(D^{i-1}l\). Cyclic permutation kills all
orbits of size \(p\) in the residue field. The only invariant
partitions are the single block and the singletons, since the
blocks must then be the cosets of a subgroup of \(\mathbb Z/p\).

On the completed elliptic local ring in characteristic \(p\),
\(\partial^p=H_p\partial\) as well. The two continuous
derivations agree on the local parameter and on constants, hence
on power series. Taking all further invariant derivatives of the
integral germ \(R_w\) consequently gives
\(D^pR_w=H_pDR_w\) as sequences. Now \(D^p\) is a derivation in
the reduced sequence ring. The product \(R_w\Theta_t\) satisfies
the same identity. Reading entry \(n\) proves the recurrence, and
iteration gives the formula for arbitrary positive orders.
\end{proof}

\begin{corollary}[Finite matrix evaluation for the prescribed period]\leavevmode
\label{cor:splitrank480}
Fix the fields, signed data and projected tests of
Theorem~\ref{thm:signedcovariance478}. Choose a good split prime
as above at which their scalar normalizations, changes of frame
and finite coefficients are integral and which is prime to the
orders of the torsion shifts. Every entry of the period matrix
at the original tests can be reduced by the following finite rule.
Use the elliptic product expansion of
Proposition~\ref{thm:classden312}. For a product coordinate with
derivative variable \(Z\), reduce its differential symbol in
\[
                        k[Z]/(Z^p-H_pZ).
\]
On its basis \(1,Z,\ldots,Z^{p-1}\) use the first \(p\)
derivatives of the indicated translate. Take tensor products
over all the coordinates, with the corresponding reduced
differentials on conjugated edges. Sum with the original signed
coefficients and the shared coupling index. A nonzero minor of
the resulting matrix proves independence of those columns in
characteristic zero.
\end{corollary}
\begin{proof}
The theorem states that the derivative functional on a monomial
kills every multiple of \(Z^p-H_pZ\). The ordinary chain rule
first expresses a matrix differential as a polynomial in the
commuting product-coordinate derivatives. On each tensor product
section these act on different factors independently. Thus the
calculation gives exactly the reduction of every original entry,
by the linear terminal and path formulas. The integers
\(b_i^{\rm per}\) of Lemma~\ref{lem:periodinteger463} are units
at these split primes, so the same rule applies to the divided
columns. Determinants commute with reduction.

For any finite set of such data, a single rational prime may
be taken split completely in their common normal coefficient
field and their period fields, outside all the finite exclusions,
by Chebotarev's theorem. Moreover, the characteristic-zero ranks
are preserved at all but finitely many of these primes. To see
this, in each positive-rank matrix fix one nonzero minor of its
rank and exclude the primes dividing its numerator or denominator.
All larger minors are already zero. This last existence assertion
gives no lower bound on the ranks. A lower bound for the family
in Section~\ref{sec:missing} requires a surviving minor for its
actual signed tests.
\end{proof}


\subsection{Simultaneous interpolation on the marked elliptic curve}
\label{sec:shiftedjets481}

At a good split prime \(p\), Theorem~\ref{thm:splitjets480}
reduces the derivatives on each elliptic coordinate of an actual
period to orders \(0,\ldots,p-1\). We determine which such
values can be prescribed simultaneously at a finite set \(B\)
of torsion shifts of order prime to \(p\). Throughout this
section \(B\) is nonempty. Use the same mark, differential and
good model, extending the coefficient valuation ring to contain
the finitely many data below, and take ranks over the algebraic
closure \(k\) of its residue field.

Fix \(M\ge2\) prime to \(p\). Let \(\mathcal L=\mathcal O_C([t])\),
with its section of divisor \([t]\), rigidified to have value one
at the origin in the marked fiber. On the elliptic curve itself
use the line bundle \(\mathcal L_M=[M]^*\mathcal L\).
In the analytic coordinate \(v\), write
\[
             s_a(v)=(U_a\Theta_t)(Mv),\qquad a\in C[M].
\]
Choose lifts of the torsion points consistently in all appearances
of this formula. Its meaning on the algebraic line bundle is
\(s_a=s_0 R_a(Mv)\), where \(R_a=(U_a\Theta_t)/\Theta_t\).

\begin{lemma}[A basis after reduction]\leavevmode
\label{lem:torsionbasis481}
The \(M^2\) sections \(s_a\) are integral on the good model
of \(\mathcal L_M\), and their reductions are a basis of
\(H^0(C_k,\mathcal L_{M,k})\).
\end{lemma}
\begin{proof}
Theorem~\ref{thm:trans57} makes \(R_a(Mv)\) a rational function
over a finite extension, with horizontal divisor
\([M]^*([t-a]-[t])\). Its possible vertical divisor is an integer
multiple of the special fiber, which is geometrically integral.
At the origin its leading coefficient in a local parameter is a
unit. To recall the reason, the \(M\)-th power of \(R_a\) is a
constant times a rational function with divisor
\(M([t-a]-[t])\). After giving the function vertical order
zero, the product of its translates by \(0,a,\ldots,(M-1)a\)
is a unit constant. The corresponding product for \(R_a\) is
a unit multiplier of the original theta. Hence the first constant
is a unit. Near the origin the rational function is a unit times
the \(M\)-th power of a parameter if \(a=t\), and a unit otherwise.
Distinct prime-to-\(p\) torsion specializes distinctly. Taking the
chosen \(M\)-th root preserves the unit leading coefficient,
also after pullback by \([M]\), whose tangent is a unit. Thus
there is no vertical contribution. Multiplication by \(s_0\)
leaves the effective divisor \([M]^*[t-a]\) and gives a nonzero
section after reduction.

Translation by \(b\in C[M]\) is a deck transformation of \([M]\).
It fixes \(s_0\) with the pullback action. On the ratio \(s_a/s_0\)
it acts by the commutator of \(a\) with the period \(Mb\).
In lattice coordinates this is
\(\exp(\pm2\pi i\,e(Mb,a))\), where \(e\) is the unimodular
alternating form of the polarization and the sign is fixed by
the convention for \(U\). As \(a\) varies modulo the period
lattice these characters are distinct. In fact a vector pairing
integrally with every period is itself a period. All character
values have order dividing \(M\); they reduce injectively on the
roots of unity of that order. The character eigenspaces on the
special fiber are therefore distinct and the nonzero sections
are independent. Their number is the dimension, since
\(\deg\mathcal L_M=M^2\) and Riemann--Roch on a genus-one curve
gives this dimension for a line bundle of positive degree.
\end{proof}

\begin{theorem}[The exact rank of the simultaneous derivative table]\leavevmode
\label{thm:shiftedjets481}
Consider the matrix with columns \(a\in C[M]\) and rows
\((b,r)\), \(b\in B\), \(0\le r<p\), whose entry is
\[
                J_r(U_b U_a\Theta_t)\pmod{\mathfrak p}.
\]
Its rank is \(\min(M^2,p|B|)\). The operator \(U_b\) is the
fixed outer translation in each row. The result includes rows
on which an individual translate vanishes at the origin.
\end{theorem}
\begin{proof}
For each chosen lift of \(b\) take the point \(v_b=b/M\)
on the cover, so \([M]v_b=b\). These points have order prime
to \(p\), and their reductions are distinct because their
images are distinct. At \(v_b\) use a local trivialization of
\(\mathcal L_M\) by some \(s_c\) with \(c\in C[M]\),
\(c+b\ne t\). Such a \(c\) exists since \(M\ge2\).
The horizontal divisor above shows \(s_c\) is a unit in this
trivialization. On a neighborhood of zero we have
\[
 \frac{(U_bU_a\Theta_t)(z)}{(U_bU_c\Theta_t)(z)}
                =(s_a/s_c)(v_b+z/M).
\]
The common translation factor cancels before any differentiation.
The denominator on the left has unit value at zero and has
integral undivided derivatives. Indeed \(U_bU_c\) is \(U_{b+c}\)
times a root of unity, and \(b+c\ne t\). The local calculation
in the lemma and Theorem~\ref{thm:splitjets480} give the claims.

Multiplication by the germ of this denominator acts triangularly
on the list of the first \(p\) undivided derivatives with unit
diagonal. Comparison with coefficients in an algebraic parameter
at \(v_b\) is triangular as well. In order \(r\) the leading
coefficient is \(r!\) times the \(r\)-th power of the nonzero
tangent conversion, including \(M^{-1}\). All are units for
\(0\le r<p\). Thus reduction gives an invertible row change,
independent of \(a\), from these \(p\) entries to restriction of
the section to the length-\(p\) subscheme supported at \(v_b\).
This argument uses only finite jets of the marked frame.

By the lemma the matrix therefore has the rank of restriction of
all sections of \(\mathcal L_{M,k}\) to
\(p\sum_{b\in B}[v_b]\). Its kernel is the space of sections
of \(\mathcal L_{M,k}(-p\sum_b[v_b])\). That bundle has degree
\(M^2-p|B|\), which cannot equal zero since \(p\nmid M\).
Its dimension is its degree when positive, and zero when
negative, by Riemann--Roch or the degree of the zero divisor.
Subtracting from \(M^2\) gives the formula.
\end{proof}

\begin{corollary}[One choice of additional translates for an entire edge]\leavevmode
\label{cor:productjets481}
Use an elliptic product presentation of an edge of the signed
column as in Proposition~\ref{thm:classden312}. In coordinate
\(l\), let \(B_l\) include every distinct original shift in
that presentation and suppose their orders are prime to \(p\).
Fix one prime \(P_0\ne p\) and put \(b_*=\max_l|B_l|\).
One may take \(M\) a power of \(P_0\) with
\[
                  \sqrt{p b_*}<M\le P_0\sqrt{p b_*}.
\]
There is one linear combination of the additional product
translates \(\bigotimes_lU_{a_l}\Theta_{t,l}\), \(a_l\in C_l[M]\),
giving any prescribed residue table of derivatives of coordinate
orders \(0,\ldots,p-1\) at all shift tuples \((b_l)_l\).
The coefficients can be integral at the testing place. All the
original translations are applied outside the added translations.
\end{corollary}
\begin{proof}
Take the least power strictly above \(\sqrt{p b_*}\).
The theorem gives full row rank in each coordinate. Its tensor
product is onto the space of arrays on
\(\prod_l(B_l\times\{0,\ldots,p-1\})\), since elementary tensors
span and each factor map is onto. At each coordinate a maximal
nonzero minor is a unit before reduction. Solving on those
columns, or lifting a residue solution, gives the asserted integral
coefficients after a finite extension if needed. A single product
input constructed with those coefficients serves every original
shift tuple simultaneously. Theorem~\ref{thm:splitjets480}
determines its higher orders from the prescribed table.
\end{proof}

In particular there is full local freedom in these low derivatives
when the product input can be varied at this auxiliary torsion
order. For the period at the specified input, the coefficients of
the additional translates are fixed, with only the original
vacuum present. Proving independence there still requires the
signed path sum to survive evaluation on that particular table.


\subsection{Positive derivatives of a fixed signed theta section}
\label{sec:positivejets482}

Retain the elliptic curve, differential, two-torsion mark \(t\)
and split prime from Theorem~\ref{thm:splitjets480}. Fix \(M\ge2\)
prime to \(p\), write \(d=M^2\), and suppose \(p>2d+2\).
Use the sections \(s_a(v)=(U_a\Theta_t)(Mv)\) of
\(\mathcal L_M=[M]^*\mathcal O([t])\) from
Lemma~\ref{lem:torsionbasis481}. Write \(k\) for the algebraic
closure of the residue field. The next argument works with the
given coefficients of a signed sum at a single prescribed
shift. In particular it will also apply when its value at that
shift is zero.

\begin{lemma}[The second logarithmic derivative on the torsion cover]\leavevmode
\label{lem:curvature482}
Let \(S\) be a nonzero section of \(\mathcal L_{M,k}\).
There is a rational function \(Q_S\) on the special fiber whose
poles are precisely the distinct zeros of \(S\), each of order two.
On the open set where \(s_c\ne0\) its formula, with \(h=S/s_c\),
is
\[
 Q_S(v)=M^2q_t(Mv+c)+
              \partial_v\!\left(\frac{\partial_v h}{h}\right),
 \qquad q_t(z)=-u_\Lambda-x(z+t).
\]
In particular \(Q_S\ne0\) and its pole divisor has degree at most
\(2d\).
\end{lemma}
\begin{proof}
The open sets in the statement cover the curve. The zero divisor
of \(s_c\) is \([M]^*[t-c]\); two different \(c\)'s have disjoint
zero divisors. On an overlap the ratio \(s_{c'}/s_c\) is a unit.
In characteristic zero its second logarithmic derivative is the
difference \(M^2(q_t(Mv+c')-q_t(Mv+c))\). Indeed the linear
exponential in each Hermitian translate has zero second
logarithmic derivative. Both sides are rational functions and
can be reduced on the overlap: the indicated sections,
the functions \(q_t\), and the invariant vector field are
integral there. The quotient rule now gives agreement of the
two formulas for any reduced section \(S\).

The first summand is regular wherever \(s_c\ne0\), because its
only possible poles there would have \(Mv+c=-t=t\).
At a zero \(P\) of \(S\) of multiplicity \(j\), write
\(h=w^j u\) for a local parameter \(w\) and a unit \(u\).
The coefficient of \(w^{-2}\) in the second summand equals
\(-j(\partial_v w(P))^2\). It is nonzero since the invariant
vector field never vanishes and \(1\le j\le d<p\). Away from
these zeros \(h\) is a unit in a suitable chart and both terms
are regular. The zero divisor of \(S\) has degree \(d>0\);
there is a pole and there are at most \(d\) distinct poles.
\end{proof}

\begin{theorem}[Positive derivatives detect every fixed signed section]\leavevmode
\label{thm:positivejets482}
Let \(b\) be any fixed torsion shift of order prime to \(p\).
Over \(k\) the matrix with columns \(a\in C[M]\), rows
\(1\le r\le 2d+2\), and entries
\[
                   J_r(U_bU_a\Theta_t)
\]
has column rank \(d\). The same holds upon replacing every \(r\)
independently by \(r+m_r(p-1)\), with \(m_r\) a nonnegative
integer. In particular the positive orders in any \(p-1\)
consecutive positive integers suffice.
\end{theorem}
\begin{proof}
Take a nonzero vector \((u_a)_a\) of residue coefficients and set
\(S=\sum_a u_as_a\). The torsion basis makes \(S\ne0\).
Write \(F_b\) for the marked derivative sequence
\(\sum_a u_a J_\bullet(U_bU_a\Theta_t)\) and \(P=b/M\) for the
chosen point of the cover. Let \(R=2d+2\). Through degree \(R\)
use the ordinary truncated series with coefficient
\(F_b^{(r)}(0)/r!\) in degree \(r\); all these factorials are
units. To relate this frame to the algebraic curve, take
\(s_c(P)\ne0\). The quotient formula in the proof of
Theorem~\ref{thm:shiftedjets481} expresses the series as the
product of \(U_bU_c\Theta_t\) and
\((S/s_c)(P+z/M)\) through every degree under consideration.
The first factor has unit constant term and the map from a local
parameter at \(P\) to \(z\) has unit tangent. These comparisons
use only the indicated finite Taylor coefficients and commute with
reduction.

If \(S(P)=0\) with multiplicity \(j\), then \(1\le j\le d\).
The coefficient of \(z^j\) in this product is nonzero, so its
\(j\)-th derivative is nonzero. If \(S(P)\ne0\), suppose that
all the derivatives of orders \(1,\ldots,R\) vanish. The truncated
series is then a nonzero constant through degree \(R\). Its
second logarithmic derivative is zero through degree \(R-2\):
differentiate twice, using the quotient rule and the inverse of
the unit series. By the product formula just given and the
definition of \(Q_S\), this is the series of
\(M^{-2}Q_S(P+z/M)\) through that degree. Hence \(Q_S\) has
order at least \(R-1=2d+1\) at \(P\).
A nonzero rational function on a smooth projective curve has zero
and pole divisors of equal degree. The lemma gives the
contradiction. Thus no nonzero coefficient vector is in the
kernel. By Theorem~\ref{thm:splitjets480}, each indicated row
replacement multiplies that row by the unit \(H_p^{m_r}\).
A block of \(p-1\) positive integers contains one positive
representative of each residue class modulo \(p-1\).
\end{proof}

\begin{corollary}[The span of the derivatives at one product shift]\leavevmode
\label{cor:fixedproductjets482}
In each coordinate \(l\) of an elliptic product fix a torsion level
\(M_l\ge2\), a mark \(t_l\) and a shift \(b_l\), with torsion
orders prime to the common good split prime \(p\). Suppose
\(p>2\max_l M_l^2+2\). Let \(v_{\boldsymbol a}\) be any fixed
array of vectors in a finite dimensional \(k\)-space, indexed by
\(\boldsymbol a\in\prod_l C_l[M_l]\). The vectors
\[
 \sum_{\boldsymbol a}v_{\boldsymbol a}
       \prod_l J_{r_l}(U_{b_l}U_{a_l}\Theta_{t_l}),
           \qquad 1\le r_l\le 2M_l^2+2 ,
\]
span exactly the span of the \(v_{\boldsymbol a}\). Each
coordinate admits the same replacements of positive orders as
in the theorem.
\end{corollary}
\begin{proof}
In each coordinate the rows in the theorem span the full dual
of the space indexed by \(C_l[M_l]\). Their tensor products span
the full dual on the product index set. Applying to this dual
the linear map sending a coordinate functional to the associated
\(v_{\boldsymbol a}\) gives the two asserted spans.
\end{proof}

\begin{corollary}[Ranks of the positive-derivative Hankel matrices]\leavevmode
\label{cor:fixedhankel482}
For a nonzero scalar coefficient vector in the theorem write
\(m_n=\sum_a u_a J_n(U_bU_a\Theta_t)\) for \(n\ge0\), and let
\(H=H_p\). There are unique coefficients \(w_z\in k\), indexed
by \(z^p=Hz\), with \(m_n=\sum_z w_z z^n\) for all \(n\ge0\)
(using \(z^0=1\)). At least one \(w_z\) with \(z\ne0\) is
nonzero. For every integer \(A\ge1\),
\[
 \operatorname{rank}(m_{A+i+j})_{0\le i,j\le p-2}
       =|\{z\ne0:w_z\ne0\}|.
\]
If \(H\) and all \(m_n\) belong to \(\mathbb F_p\), this rank is
at least the multiplicative order of \(H\) in \(\mathbb F_p^\times\).
\end{corollary}
\begin{proof}
The recurrence in Theorem~\ref{thm:splitjets480} defines the
moment functional \(Z^n\mapsto m_n\) on \(k[Z]/(Z^p-HZ)\).
Since \(H\ne0\), the polynomial has \(p\) distinct roots.
Evaluation at them identifies the quotient with \(k^p\), giving
the weights and their uniqueness. If their support were contained
in zero, every positive derivative would vanish, contrary to the
theorem. On the \(p-1\) nonzero roots the square matrix
\((z^i)_{0\le i\le p-2,z\ne0}\) is invertible. Multiplying it by
the diagonal matrix with entries \(w_z z^A\) and then by its
transpose gives the displayed Hankel matrix, proving the rank.
Over \(\mathbb F_p\), uniqueness gives \(w_{z^p}=w_z^p\).
Each nonzero root has Frobenius orbit of length exactly the
multiplicative order of \(H\), since \(z^p=Hz\).
\end{proof}

For the signed columns in Theorem~\ref{thm:signedcovariance478},
any torsion sums on a given product with the same derivative
symbol can be grouped into an array as in the corollary, retaining
all repeated indices in its coefficients. Its mixed derivatives
detect this entire array at the fixed shift. The designated Weyl
projections prescribe linear combinations of the derivatives. A
lower bound on their images still needs information about which
linear combinations the projected period supplies.


\subsection{Common split primes for signed theta derivatives}
\label{sec:commonsplit483}

At a split prime, a differential symbol multiplies the weights
in Corollary~\ref{cor:fixedhankel482} by its values at the
roots of \(Z^p-H_pZ\). To use this for a prescribed symbol, we
need many roots carrying weight on the same reduction. We will
choose a single prime for any given finite list of the marked
curves and torsion data. The coefficients within a signed sum
on their product can depend on all its torsion indices jointly.

\begin{lemma}[The recurrence scalar and the Frobenius trace]\leavevmode
\label{lem:frobtrace483}
Let the marked model in Lemma~\ref{lem:splitcmder480} reduce
over \(\mathbb F_p\), with \(y^2=f(x)=4x^3-g_2x-g_3\) and
\(\omega=dx/y\). Write \(a_p=p+1-\#C(\mathbb F_p)\). Then
\[
           H_p=[x^{p-1}]\,f(x)^{(p-1)/2}=a_p\pmod p .
\]
\end{lemma}
\begin{proof}
Write \(K=\mathbb F_p(C)\) and \(L=K^p\), so
\(1,x,\ldots,x^{p-1}\) is an \(L\)-basis of \(K\).
The \(L\)-linear operator \(\partial=y\,d/dx\) has
characteristic polynomial \(T^p-H_pT\).
Indeed it is annihilated by this separable polynomial by
Lemma~\ref{lem:splitcmder480}, and has rank \(p-1\).
Adjoin its eigenvalues to \(L\). This is a separable extension
and its compositum with the purely inseparable field \(K\) is
their tensor product, hence is a field. There is a nonzero
eigenvalue \(\lambda\), with eigenvector \(u\). The nonzero
elements \(1,u,\ldots,u^{p-1}\) have distinct eigenvalues
\(0,\lambda,\ldots,(p-1)\lambda\) by the derivation rule.
They exhaust the dimension and prove the polynomial formula.

For a field element \(g\) let \(\ell(g)\) be its coefficient
on \(x^{p-1}\) in the displayed basis. The adjugate of \(d/dx\)
maps \(g\) to the constant \(-\ell(g)\). Its only nonzero
cofactor comes from the entries \(1,\ldots,p-1\) on the
superdiagonal, whose product is \(-1\), with positive cofactor
sign. Multiplication by \(y\) has determinant \(y^p\), the norm
in the purely inseparable extension of degree \(p\). Its adjugate
is therefore \(y^p\) times multiplication by \(1/y\).
Taking the adjugates in reverse order gives
\(\operatorname{tr}\operatorname{adj}(\partial)=-y^p\ell(1/y)\);
here the trace of \(g\mapsto L_0(g)\,1\) is \(L_0(1)\) for
any \(L\)-linear functional \(L_0\).
This trace is the coefficient of \(T\) in the characteristic
polynomial, as \(p\) is odd. Since
\(1/y=f(x)^{(p-1)/2}/y^p\) and
\(3(p-1)/2<2p-1\), the only monomial of that numerator contributing
to \(x^{p-1}\) in the field basis has degree \(p-1\).
This proves the first equality.

There are \(1+f(u)^{(p-1)/2}\) points above \(x=u\) modulo
\(p\), including at a root of \(f\). Add the point at infinity.
The sum of \(u^j\) over \(u\in\mathbb F_p\) is zero modulo \(p\)
except for positive multiples of \(p-1\), where it is \(-1\).
The degree of \(f^{(p-1)/2}\) is less than \(2(p-1)\), giving
the asserted congruence for \(a_p\).
\end{proof}

\begin{theorem}[Simultaneously long orbits at a split prime]\leavevmode
\label{thm:commonsplit483}
Fix finitely many marked CM elliptic curves as above, with
any finite sets of algebraic coefficients and torsion translates
used in their theta functions. Take a finite Galois number field
\(J\) containing their CM fields, the data, the fields of the
algebraic translation ratios and scalar normalizations, all their
five-torsion points, and the fifth roots of unity. Fix
\(0<\delta<1/4\). For all sufficiently large \(X\), among the
rational primes \(p\in[X,2X]\) splitting completely in \(J\),
all but \(O(X^{1-2\delta})\) satisfy
\[
          \operatorname{ord}_{\mathbb F_p^\times}(H_p)
                     >X^{1/4-\delta}
\]
for every indicated curve at every place over \(p\).
The implied constant is absolute. In particular there are such
primes in each sufficiently distant interval \([X,2X]\) for the
fixed data.
\end{theorem}
\begin{proof}
Take \(X\) above the finitely many bad primes and the torsion
orders. At each completely split prime the residue field is
\(\mathbb F_p\), the group of rational points contains the full
five-torsion, and \(p\equiv1\bmod 5\). Consequently
\(a_p\equiv2\bmod 5\), excluding \(0,1,-1\).
Also \(|a_p|\le2\sqrt p\) by Hasse's inequality. Recall its
degree argument: for the \(p\)-power Frobenius \(T\), the quadratic
form on elliptic endomorphisms given by isogeny degree yields
\(\deg([m]-[n]T)=m^2-a_pmn+pn^2\). Its cross term follows from
\(\deg(1-T)=\#C(\mathbb F_p)\), since \(1-T\) is separable.
Nonnegativity for integral \(m,n\), and density of the ratios
\(m/n\) in \(\mathbb R\), give the bound.

Put \(B=\lfloor X^{1/4-\delta}\rfloor\). For each integer
\(a\) with \(2\le |a|\le2\sqrt{2X}\) form the nonzero integer
\(N_a=\prod_{j=1}^B(a^j-1)\).
Its absolute logarithm \(\log |N_a|\) is \(O(B^2\log X)\),
by summing \(\log(2|a|^j)\). Thus at most \(O(B^2)\) distinct
primes at least \(X\) divide it. Summing over the possible
integers \(a\) gives \(O(\sqrt X B^2)\) primes in total.
Outside this single exceptional set, each possible Frobenius
trace \(a_p\) has multiplicative order greater than \(B\)
modulo \(p\); for \(X\) large it is nonzero modulo \(p\).
The lemma identifies its residue with \(H_p\).
The Chebotarev density theorem for the identity conjugacy class
in \(\operatorname{Gal}(J/\mathbb Q)\) gives
\((1+o_J(1))X/([J:\mathbb Q]\log X)\) completely split primes
in \([X,2X]\) \cite[\S1A]{XIII-X-I-tz50}. This exceeds the exceptional
count for all sufficiently large \(X\).
\end{proof}

\begin{lemma}[The number of nonzero weights after scalar extension]\leavevmode
\label{lem:supportorbit483}
Fix a curve at one of these primes, with the data of
Theorem~\ref{thm:positivejets482} defined over the common
degree-one residue field, and \(p>2M^2+2\). Write \(d=M^2\)
and \(h=\operatorname{ord}_{\mathbb F_p^\times}(H_p)\).
For \emph{every} nonzero coefficient vector \(u\in k^d\),
the weight vector \(W(u)=(w_z)_{z\ne0,\ z^p=H_pz}\) of its
positive derivatives has at least \(\lceil h/d\rceil\)
nonzero entries.
\end{lemma}
\begin{proof}
The map \(W\) is linear and injective by
Corollary~\ref{cor:fixedhankel482} and
Theorem~\ref{thm:positivejets482}. The derivative sequences
of the individual basis translates are defined over
\(\mathbb F_p\). In fact the reference derivatives follow by
differentiating \(Dl=q_t,\ D\Theta_t=l\Theta_t\) from the
marked values \(l(0)=0,\Theta_t(0)=1\). The translates are
obtained by their algebraic ratios and the scalar multipliers,
all defined in \(J\). Thus coefficient Frobenius \(F:u\mapsto
u^{(p)}\) sends the support of \(W(u)\) to its permutation
\(z\mapsto z^p\), by uniqueness of the weights and conjugation
of all the moments.

The span of \(u,Fu,F^2u,\ldots\) is generated by its first
\(j\le d\) vectors. Once a power lies in the span of its
predecessors, that span is stable under the semilinear
automorphism \(F\), since its image has the same dimension.
The union of the supports of \(W\) on this span therefore
consists of entire Frobenius orbits, all of length \(h\).
It is nonempty by injectivity. It is already the union for the
first \(j\) vectors, each with the same support size as \(W(u)\).
Hence \(h\le j|\operatorname{supp}W(u)|\), as claimed.
\end{proof}

\begin{corollary}[Bounded-degree symbols on a fixed signed tensor]\leavevmode
\label{cor:projectsymbol483}
On a product of the indicated curves let
\(F=\sum_{\boldsymbol a}u_{\boldsymbol a}
\bigotimes_l U_{b_l}U_{a_l}\Theta_{t_l}\),
where \(a_l\in C_l[M_l]\), each \(b_l\) is fixed, and the
coefficient array has nonzero reduction. Take a common prime
as in the theorem with \(p>2\max_l M_l^2+2\). Put
\(\Delta_l=\lceil\operatorname{ord}(H_{p,l})/M_l^2\rceil\).
Let \(P\in k[Z_1,\ldots,Z_g]\) be nonzero of coordinate degrees
\(e_l<\Delta_l\).
For each multi-index \(\boldsymbol A\) with positive entries, the
square matrix on \(\boldsymbol i,\boldsymbol j\in\{0,\ldots,p-2\}^g\)
whose entry is the reduced mixed derivative
\[
              J_{\boldsymbol A+\boldsymbol i+\boldsymbol j}
                       \bigl(P(\partial)F\bigr)
\]
has rank at least \(\prod_l(\Delta_l-e_l)\). For fixed
characteristic-zero data, a single prime can be chosen to give
this nonvanishing for all nonzero reductions of symbols of any
prescribed bounded coordinate degrees.
\end{corollary}
\begin{proof}
In coordinate \(l\) let \(\mathcal W_l\) be the image of the
map \(W\) in the lemma, viewed as functions on the \(p-1\)
nonzero roots. The weight function \(G\) on the product of these
sets belongs to \(\bigotimes_l\mathcal W_l\), and is nonzero since
the tensor product of the injective maps \(W\) is injective.
A polynomial in the derivatives multiplies it pointwise by \(P\).
Roots with some zero coordinate make no contribution to the
displayed matrix since each \(A_l\) is positive.

We count the indices at which \(PG\ne0\). In one variable the
lemma supplies at least \(\Delta_1\) points where \(G\ne0\);
the polynomial can lose at most \(e_1\). For the induction choose
one value of the last variable for which the slice of \(G\) is
nonzero. It belongs to the tensor product of the preceding
spaces. Select any nonzero coefficient of \(P\) regarded as
a polynomial in the last variable. Apply the induction
hypothesis to this coefficient and the chosen slice. At least
\(\prod_{l<g}(\Delta_l-e_l)\) choices of the preceding coordinates
give nonzero values for both. Over each such choice, the last
fiber of \(G\) is a nonzero member of \(\mathcal W_g\), and
the last-variable polynomial is nonzero. The one-variable
bound gives at least \(\Delta_g-e_g\) points in that fiber.
This proves the product bound.

The matrix in the assertion is a tensor product of invertible
Vandermonde evaluation matrices on its left, their transposes on
its right, and between them the diagonal with entries
\(P(\boldsymbol z)G(\boldsymbol z)\prod_l z_l^{A_l}\).
This follows by expanding the moments in each coordinate.
Its rank is precisely the number just bounded. For fixed data
the theorem allows every \(\Delta_l\) to exceed the prescribed
degrees at one prime. A minor nonzero in the residue field is
also nonzero in characteristic zero for any integral lifts of
the data.
\end{proof}

The further differentiations in the matrix run through a box
of positive orders. In the period of
Theorem~\ref{thm:signedcovariance478} the final projection fixes
its own linear functional on the derivatives. The rank bound
proves survival under further differentiations of a fixed
symbol. Application to that final functional requires showing
which of these evaluations it retains.


\subsection{Integral elementary divisors of the theta derivative table}
\label{sec:jetsmith485}

We determine the integral image of all the input translates at
a single auxiliary level, with the output shifts and orders fixed
in advance. This gives the factorial cost in the marked
differential even when the highest order exceeds the residue
characteristic. Work over the integers \(O\) of a finite
coefficient extension at a prime \(p\ge5\). Use a good
Weierstrass model of the marked CM curve from
Section~\ref{sec:splitjets480}, with \(\omega=dx/y\) an integral
generator and \(p\nmid 2d_0\). We allow either splitting behavior
of \(p\) in the CM field. Fix \(M\ge2\) prime to \(p\) and a
finite nonempty set \(B\) of distinct shifts of order prime to
\(p\). For each \(b\in B\) fix a lift, the point \(P_b=b/M\) on
the cover, and an integer \(n_b\ge1\). Take \(O\) large enough
to contain these data and their translation factors. Put
\[
 d=M^2,\qquad n=\sum_{b\in B}n_b,\qquad
   D=\sum_{b\in B}n_b[P_b],\qquad
   P_*=Mt-\sum_{b\in B}n_bP_b\quad\text{in } C.
\]
The \(n\)-by-\(d\) matrix \(E\) has column \(a\in C[M]\) and row
\((b,j)\), \(0\le j<n_b\), with entry
\(J_j(U_bU_a\Theta_t)\). The original shifts act on the outside
of the input.

\begin{lemma}[Factorials in the marked local frame]\leavevmode
\label{lem:jetframe485}
At each \(P_b\) use a local parameter \(w_b\) of the smooth
integral curve, vanishing on the section \(P_b\), and trivialize
\(\mathcal L_M=[M]^*\mathcal O([t])\) by \(s_c\) with \(b+c\ne t\).
Let \(B_0\) map coefficients in the \(s_a\) to their ordinary
Taylor coefficients in \(w_b\) through degree \(n_b-1\) at every
\(P_b\). There is a block diagonal matrix
\(L\in\mathrm{GL}_n(O)\), lower triangular in each block, such that
\[
             E=L\,\operatorname{diag}_{(b,j)}(j!)\,B_0 .
\]
\end{lemma}
\begin{proof}
The proof of Lemma~\ref{lem:torsionbasis481} gives an integral
basis \(s_a(v)=(U_a\Theta_t)(Mv)\) after reduction also at the
present primes. Indeed its ingredients are the unit leading
term of a prime-to-\(p\) torsion ratio, the distinct deck
characters on \(C[M]\), and Riemann--Roch. All three apply at
good reduction when \(p\nmid M\). At \(P_b\) the section \(s_c\)
is a unit because its zero divisor is \([M]^*[t-c]\) and distinct
prime-to-\(p\) torsion sections stay distinct.

Write \(G_b(z)=(U_bU_c\Theta_t)(z)\). Its value at zero is a
unit and its undivided derivatives are integral. For the latter
assertion, \(q_t=-u_\Lambda-x(z+t)\) is regular at the origin
over \(O\) because \(t\) is a nonzero section of order two and the CM
formula in Lemma~\ref{lem:splitcmder480} for \(u_\Lambda\) has
unit denominator at \(p\nmid2d_0\). Integrating the identities
\(Dl=q_t,\ D\Theta_t=l\Theta_t\) as undivided derivative
sequences from \(l(0)=0,\Theta_t(0)=1\) preserves integrality
by the Leibniz rule. The ratio for the prime-to-\(p\) translate
is integral in an algebraic parameter at the origin, as proved
in Theorem~\ref{thm:trans57}, and its constant term here is a
unit. Composition of the two translations adds a root of unity.

For a section with local representative \(h=S/s_c\), its marked
derivatives are those of \(G_b(z)h(P_b+z/M)\), by the quotient
formula in Theorem~\ref{thm:shiftedjets481}. On the cover let
\(D_b=M^{-1}\partial_v\), so \(D_b^j w_b(P_b)\) is integral
in every positive order and \(D_b w_b(P_b)\) is a unit.
For the term \(w_b^j\) in \(h\), each undivided derivative of
order \(r\) of \(G_bw_b^j\) at the point is divisible by \(j!\).
In the repeated product rule assign the \(r\) labeled derivatives
to \(G_b\) and to the \(j\) ordered copies of \(w_b\). A term
survives only if every copy of \(w_b\) receives a derivative.
Permuting these copies acts freely on the assignments that
survive and preserves their value. Each orbit thus contributes
\(j!\) times an integral element. For \(r<j\) there is no such
assignment, and for \(r=j\) the quotient by \(j!\) is
\(G_b(0)(D_bw_b(P_b))^j\), a unit. The case \(j=0\) uses the
single empty family of copies and the same diagonal formula.
Truncation of \(h\) is legitimate since an order \(r\) derivative
depends only on its first \(r+1\) coefficients. This gives the
claimed matrices \(L\) directly.
\end{proof}

\begin{theorem}[The integral image at and above the interpolation degree]\leavevmode
\label{thm:jetsmith485}
Suppose \(d\ge n\). If \(d>n\), or if \(d=n\) and
\(P_*\ne0\), then
\[
             \operatorname{coker}_O E
                \simeq\bigoplus_{b\in B,\,0\le j<n_b} O/(j!).
\]
If \(d=n\) and \(P_*=0\), choose a row \(m=(b_*,n_{b_*}-1)\)
with \(n_{b_*}=\max_b n_b\). In this case
\[
             \operatorname{coker}_O E
                \simeq O\ \oplus
                   \bigoplus_{(b,j)\ne m,\ 0\le j<n_b}O/(j!).
\]
\end{theorem}
\begin{proof}
The matrix \(B_0\) reduces to restriction of all sections of
\(\mathcal L_M\) to the divisor \(D\) on the special fiber,
in the chosen Taylor bases. Its kernel over that field is
\(H^0(\mathcal L_M(-D))\). When \(d>n\) this has dimension
\(d-n\) by Riemann--Roch, so restriction is onto. If \(d=n\)
the class of \(\mathcal L_M(-D)\) in \(\mathrm{Pic}^0(C)=C\)
is \(P_*\). In fact the zeros of \(s_0\) are \(v_0+C[M]\) for
\(Mv_0=t\), and their group sum is \(M^2v_0=Mt\).
The sum of \(C[M]\) is zero by pairing inverses and, for even
\(M\), summing the four two-torsion points. Subtracting \(D\)
gives the formula, on either fiber. All the points involved have
order prime to \(p\), so \(P_*\ne0\) remains nonzero on reduction.
A nontrivial degree-zero bundle has no section. Thus \(B_0\)
is surjective over the residue field in this case as well.
In both cases some \(n\)-row square minor is a unit over \(O\).
The lemma now gives the first cokernel.

Suppose \(d=n\) and \(P_*=0\). Over the fraction field the
bundle \(\mathcal L_M(-D)\) is trivial. Choose a section
\(S_*\) of \(\mathcal L_M\) with divisor \(D\), multiplying by
a scalar to make its coordinates in the \(s_a\) integral with
some coordinate a unit. It is nonzero on reduction by the basis
lemma. In the local frames above,
\(S_*/s_c=w_b^{n_b}u_b(w_b)\) with \(u_b(0)\) a unit.
Indeed the lower coefficients vanish in characteristic zero,
and the reduction already has zeros of total multiplicity
at least \(n=d\) at the distinct reductions. Each therefore
has exactly the prescribed multiplicity.

On a tuple \((h_b\bmod w_b^{n_b})_b\) define the integral
linear form
\[
       R((h_b)_b)=
           \sum_b\operatorname{res}_{P_b}
                     \frac{h_b\,\omega}{w_b^{n_b}u_b(w_b)} .
\]
The indicated coefficient of \(w_b^{-1}dw_b\) depends only
on the given jet. The coefficient of the last Taylor entry at
each \(b\) is a unit since both \(u_b(0)\) and the coefficient
of \(dw_b\) in \(\omega\) are units. For \(h_b=S/s_c\) the
differentials in this sum come from \((S/S_*)\omega\). It has
no poles elsewhere, so its residues sum to zero. The image of
\(B_0\) lies in \(\ker R\). Over the residue field its rank is
\(d-1\) by the one-dimensional kernel computed from the trivial
degree-zero bundle. Since \(R\) is primitive and splits off one
coordinate, \(\ker R\) reduces onto the residue kernel of its
reduction. Nakayama's lemma gives
\(\operatorname{im}B_0=\ker R\) over \(O\).

Write \(R(h)=\sum_i R_i h_i\) and \(w_i=j!\) for row \(i=(b,j)\).
Eliminate coordinate \(m\) using the unit \(R_m\). After applying
\(\operatorname{diag}(w_i)\), the other rows on the free
coordinates of this kernel are diagonal with entries \(w_i\),
and the last row has entries \(-w_m R_i/R_m\). Each \(w_i\)
divides \(w_m\) by the choice of \(m\). Integral row additions
therefore clear the entire last row. Multiplying by \(L\) does
not change the cokernel. This proves the second formula.
\end{proof}

\begin{corollary}[Determinantal ideals after a fixed signed projection]\leavevmode
\label{cor:projectedminors485}
Use finitely many coordinates \(l\) in the first case of the
theorem. Let \(A\) be any specified integral linear map from the
array of their derivative values to a free target, incorporating
the desired signs and projections. Write \(E_\otimes=\bigotimes_lE_l\),
\(L_\otimes=\bigotimes_l L_l\), and
\(w_i=\prod_l j_l!\) for the row \(i=((b_l,j_l))_l\).
If \(I_s\) denotes the ideal generated by all minors of size \(s\),
then
\[
 I_s(AE_\otimes)
      =\sum_{|T|=s}
        \left(\prod_{i\in T}w_i\right)
                  I_s\bigl((A L_\otimes)_{*,T}\bigr).
\]
The sum uses subsets of the row index set of the derivative
array. One common input level can ensure the strict inequality
\(M^2>\sum_b n_{l,b}\) in every coordinate. Given a prime
\(P_0\) distinct from the coefficient primes under study, take
the least power \(M\ge2\) of \(P_0\) greater than
\(\sqrt{\max_l\sum_b n_{l,b}}\).
\end{corollary}
\begin{proof}
The tensor product of the surjective Taylor maps \(B_{0,l}\)
has an integral right inverse, by tensoring choices of right
inverses for the free targets. Appending this surjection to a
matrix leaves every determinantal ideal equal, by Cauchy--Binet
first for the product and then for its product with the right
inverse. In the resulting matrix \(AL_\otimes\operatorname{diag}(w_i)\)
each minor using column set \(T\) acquires exactly the product
displayed. This proves the formula, including vanishing ideals.
The stated power gives \(M^2>\sum_b n_{l,b}\) directly.
\end{proof}

For a product slot in the signed period, the formula applies when
its fixed class actions are expressed by these outer shifts and
the fixed linear map on the resulting jets. It computes the
ideal over all columns of the common input level. Determining
which minors survive on the prescribed input columns, including
the original vacuum \(a=0\), remains part of the integral period
comparison of Section~\ref{sec:missing}.


\subsection{The archimedean determinant of the same elliptic jets}
\label{sec:thetaalternant486}

Fix a complex embedding of the marked model in
Theorem~\ref{thm:jetsmith485}. Write \(C=\mathbb C/\Lambda\),
\(c=\pi/\operatorname{area}(\Lambda)\), and use the metric
\(\|f(v)\|=\exp(-c|v|^2/2)|f(v)|\) in degree one.
Let \(h\) be an odd theta section with simple zero at zero,
in the same Hermitian frame. Normalize its positive scalar
so that \(g(v)=\|h(v)\|\) satisfies
\(\int_C\log g(v)\,dv=0\), using Haar probability. The phase
of \(h\) can be chosen freely. Its sole logarithmic singularity
is integrable. Put
\(\gamma=\lim_{v\to0}g(v)/|v|=|h'(0)|>0\).
Thus \(g\) is the multiplicative Arakelov Green function with
one argument zero. Write \(\Theta_t(0)=1\) for the even marked
theta section of divisor \([t]\), \(M\ge2,\ d=M^2\), and order
the sections \(s_a(v)=(U_a\Theta_t)(Mv)\), \(a\in C[M]\).
Here and below a point in a holomorphic formula comes with a
chosen lift. Norms on the torus do not require lifts.

\begin{lemma}[Distribution over torsion]\leavevmode
\label{lem:green486}
For every integer \(m\ge2\),
\[
  \prod_{a\in C[m]}g(v+a)=g(mv),\qquad
  \prod_{a\in C[m]\setminus\{0\}}g(a)=m .
\]
In particular the sum of \(\log g(a)\) over the points of exact
order \(m\) is \(\Lambda(m)\), the von Mangoldt function.
Also \(\|\Theta_t(v)\|=g(v-t)/g(t)\).
\end{lemma}
\begin{proof}
Away from its zero the logarithm of the theta norm is the log
absolute value of a holomorphic nonvanishing function minus
\(c|v|^2/2\). Hence the two sides of the first identity have
logarithms with the same Laplacian and matching simple zeros.
Their log difference extends harmonically across those zeros.
It is constant on the compact torus and has mean zero because
translation and multiplication by \(m\) preserve Haar probability.
Letting \(v\to0\) gives the second assertion by the definition of
\(\gamma\). These are the elliptic projection and torsion product
formulas \cite[Thm.~3.2, Cor.~4.5]{XIII-X-dejong486}. Partition the
nonzero \(m\)-torsion by exact order and apply M\"obius inversion
to \(\log m=\sum_{e\mid m}\Lambda(e)\), with order one omitted.
For the last assertion the two norms have the same curvature
and divisor. Their ratio is constant by the same harmonic
argument. Evaluation at zero and \(g(-t)=g(t)\) fix it.
\end{proof}

\begin{theorem}[The alternating determinant in the marked frame]\leavevmode
\label{thm:thetaalternant486}
There is a constant \(\xi\) of absolute value one, independent
of \(v_1,\ldots,v_d\), such that
\[
 \det\bigl(s_a(v_i)\bigr)_{i,a}
  =\xi\,\frac{M^{d/2}}{g(t)^d}\,
     (U_{-Mt}h)\!\left(\sum_{i=1}^d v_i\right)
                 \prod_{i<j}h(v_j-v_i).
\]
The constant is determined, including its phase, by evaluation
at any ordered fiber \(v_0+C[M]\) for which \(d v_0\ne Mt\)
on \(C\).
\end{theorem}
\begin{proof}
First take distinct points \(v_i\) on the torus. Their evaluation
matrix is singular exactly when a nonzero section of
\(\mathcal L_M=[M]^*\mathcal O([t])\) vanishes at all of them.
The \(s_a\) form a basis in characteristic zero by the deck
character calculation of Lemma~\ref{lem:torsionbasis481}.
The degree of \(\mathcal L_M\) is \(d\) and its divisor class
has group sum \(Mt\), as in Theorem~\ref{thm:jetsmith485}.
Thus singularity for distinct points is equivalent to
\(\sum_i v_i=Mt\) on \(C\).

Consider the norm of the determinant, using the pullback
degree-\(d\) metric in every row. Fix all but \(v_i\) in general
position. As a section in \(v_i\), it vanishes at the other
\(d-1\) points by equality of rows in the same fiber, and at
\(Mt-\sum_{j\ne i}v_j\) by the divisor criterion. For general
fixed points these are \(d\) distinct zeros. The section is
nonzero for other \(v_i\), so each zero is simple. The product
\[
           g\!\left(\sum_i v_i-Mt\right)
                        \prod_{i<j}g(v_j-v_i)
\]
has these same zeros and curvature in \(v_i\). The logarithm
of the ratio of norms is harmonic across the zeros and constant
in this variable. This argument applies to every variable.
Equivalently the ratio on the locus of nonzero evaluations is
locally constant on the generic open set. It has a single value
there, since the complement of finitely many proper analytic
divisors in a connected complex torus product is connected
(a path can be displaced locally away from real codimension two).
Continuity gives the same norm identity at every tuple.

To evaluate its constant use a fiber \(v_0+C[M]\) as in the
statement. All the row fibers are canonically identified by
pullback from \(Mv_0\). In these identifications column \(a\)
is its value at \(v_0\) times a character of \(C[M]\), and
the \(d\) characters are distinct. The Fourier matrix has
determinant of modulus \(d^{d/2}\), since its product with its
adjoint is \(dI\). The product of the norms at \(v_0\) equals
\[
      g(t)^{-d}\prod_{a\in C[M]}g(Mv_0+a-t)
            =g(t)^{-d}g(dv_0-Mt)
\]
by the lemma and the unitary translation formula.
The group sum over the fiber is \(dv_0\). The product over its
unordered pairs of \(g\) of their differences is \(M^{d/2}\).
Indeed its square counts each nonzero difference exactly \(d\)
times and \(g(-a)=g(a)\). The lemma supplies the product over
the nonzero differences. Dividing gives the asserted constant
in norm, since \(d^{d/2}/M^{d/2}=M^{d/2}\).

The marked translation satisfies
\(\|U_wh(v)\|=g(v+w)\). Moreover
\[
 \left|\sum_i v_i\right|^2+\sum_{i<j}|v_j-v_i|^2
                         =d\sum_i |v_i|^2 .
\]
Thus the Gaussian metric factors on the two holomorphic sides
agree exactly on the universal cover. Their holomorphic quotient
has constant absolute value wherever defined. The open mapping
theorem makes it constant on a nonempty open set, and the
identity theorem gives the asserted formula everywhere, also
at repeated points. A nonsingular fiber fixes \(\xi\).
\end{proof}

\begin{corollary}[The absolute determinant for prescribed multiplicities]\leavevmode
\label{cor:confluent486}
Take pairwise distinct outer shifts \(b\) and integers
\(n_b\ge1\) with \(\sum_b n_b=d\). Use \(P_b=b/M\) and the
matrix \(E\) of Theorem~\ref{thm:jetsmith485}, and put
\(u=\sum_b\binom{n_b}{2}\). Its exact determinant norm is
\[
 \begin{aligned}
 |\det E|
   ={}&\frac{M^{d/2}}{g(t)^d}
         \left(\frac{\gamma}{M}\right)^u
         \left(\prod_b\prod_{j=0}^{n_b-1}j!\right)
         g\!\left(\sum_b n_bP_b-Mt\right)\\
      &{}\times\prod_{b<b'}g(P_{b'}-P_b)^{\,n_bn_{b'}} .
 \end{aligned}
\]
Any ordering of the shifts gives the same product. In particular
the very same group sum governs vanishing at the complex
embedding and the exceptional relation in the integral theorem.
\end{corollary}
\begin{proof}
In each cluster indexed by \(b\), replace \(P_b\) by the \(n_b\)
points \(P_b+\epsilon t_j\) for distinct fixed complex numbers
\(t_j\). Use the Hermitian translation at \(P_b\) as a holomorphic
local frame with unit norm at that point. In this frame the
representative of \(s_a\) at \(w\) is
\((U_bU_a\Theta_t)(Mw)\), by the explicit translation formula
with metric parameter \(dc\) on the cover. Expand each row in
\(\epsilon t_j\). Alternation forces distinct powers within
each cluster. Dividing the determinant by the product within
each cluster of \(\epsilon(t_j-t_i)\), then passing to the
limit, gives the determinant with consecutive Taylor-coefficient
rows. The coefficient of order \(j\) is
\(M^j J_j(U_bU_a\Theta_t)/j!\). This limiting identity also
holds if its determinant is zero. The additional metric factors
in the local frames tend to one. On the product of norms in the
theorem each internal pair contributes \(\gamma\) after division,
all other factors tending to the displayed values. The number
of internal pairs, and the sum of all the row orders, both equal
\(u\). This gives the formula with all factorials.
\end{proof}

\begin{corollary}[Fixed signed rows and selected torsion columns]\leavevmode
\label{cor:fixedminor486}
Let \(A\subset C[M]\) have size \(s\ge1\). Choose \(s\) fixed
linear functionals \(L_i\) on the \(s_a\), each given by a
finite signed linear combination of marked derivative
evaluations at prescribed shifts, with arbitrary orders. Every
minor \(\det(L_i(s_a))_{1\le i\le s,\,a\in A}\)
is obtained by applying an explicit product
of row functionals to the single holomorphic product in the
theorem, including its one common phase. This uses the prescribed
evaluations, with no further derivative orders in those rows.
\end{corollary}
\begin{proof}
For a row of order \(j\) at \(b=MP_b\), apply in its own
variable the functional
\[
 f\longmapsto M^{-j}\left.\frac{d^j}{dw^j}
   \left(e^{-dcw\overline{P_b}-dc|P_b|^2/2}f(P_b+w)\right)
                                                   \right|_{w=0}.
\]
The local translation identity just proved gives the desired
\(J_j\). Extend linearly with the specified signs. For \(s=d\)
apply these functionals in independent variables to the determinant
identity and use its multilinearity.

For \(s<d\) select a nonsingular fiber \(v_0+C[M]\). Identify
its row fibers by the pullback action and choose one common
frame, writing \(\lambda_z\) for evaluation there at \(v_0+z\).
Set \(s_a^0=\lambda_0(s_a)\ne0\), and let \(\chi_a\) be its
actual deck character, so \(\lambda_z(s_a)=\chi_a(z)s_a^0\).
For each \(a'\notin A\) append the evaluation row
\[
        \frac{1}{d\,s_{a'}^0}\sum_{z\in C[M]}
                                  \chi_{a'}(z)^{-1}\lambda_z .
\]
Character orthogonality makes this row one at \(a'\) and zero
on all other columns. Order these appended rows and permute
the full columns with the associated determinant sign. Expansion
along the appended rows leaves the requested minor. Each
\(\lambda_z\) in the holomorphic frame is a scalar times point
evaluation, determined by the common pullback identification
(equivalently normalize it so that \(\lambda_z(s_0)=\lambda_0(s_0)\)).
This specifies all the functionals occurring on the product.
\end{proof}

The formulas give the finite and archimedean sizes of a whole
elliptic interpolation determinant in the same normalization.
They also retain the cancellations in prescribed rows of one
elliptic slot. For the full period, different slots are coupled
by the class sum and the final projections in
Theorem~\ref{thm:signedcovariance478}. An upper estimate and
the required rank for that coupled determinant on the interval
of Section~\ref{sec:missing} remain to be proved.


\subsection{The rational norm of the marked elliptic determinant}
\label{sec:jetnorm487}

We now sum the archimedean identity of
Corollary~\ref{cor:confluent486} over the embeddings of the
actual coefficient field. Write \(K\) for the imaginary quadratic
field, assume complex multiplication by its full ring of integers
\(\mathcal O_K\), and use the same marked differential
\(\omega=dx/y\). For a proper nonzero ideal \(\mathfrak a\) put
\[
 \begin{aligned}
  \Phi_K(\mathfrak a)&=|(\mathcal O_K/\mathfrak a)^\times|,\\
  \Lambda_K(\mathfrak a)&=
     \begin{cases}\log N\mathfrak q,&
               \mathfrak a=\mathfrak q^e,\quad \mathfrak q
                    \text{ prime},\ e\ge1,\\
               0,&\text{otherwise},\end{cases}
  & b_K(\mathfrak a)&=\frac{\Lambda_K(\mathfrak a)}
                                    {2\Phi_K(\mathfrak a)} .
 \end{aligned}
\]
For a nonzero torsion point \(P\), abbreviate
\(b_K(\operatorname{Ann}_{\mathcal O_K} P)\) to \(b_K(P)\).

\begin{lemma}[The Green function in the original differential]\leavevmode
\label{lem:cmslope487}
At every complex embedding the coefficient \(\gamma\) of
Lemma~\ref{lem:green486} satisfies
\(\gamma^{12}=|g_2^3-27g_3^2|\) in the Weierstrass equation
\(y^2=4x^3-g_2x-g_3\) with \(dz=dx/y\). In particular
\(\gamma=1\) on the discriminant-one model of
Proposition~\ref{prop:ellnorm56}.
\end{lemma}
\begin{proof}
Let \(t_1,t_2,t_3\) be the nonzero two-torsion points and put
\(e_i=x(t_i)\). The elliptic function \(x-e_i\) has divisor
\(2[t_i]-2[0]\) and leading coefficient one at its double
pole in the coordinate \(z\). Cancellation of zeros, poles and
curvatures makes the logarithm of the ratio between
\(|x(v)-e_i|\) and \(g(v-t_i)^2/g(v)^2\) constant off their
divisors and harmonic across them. Comparison at zero gives
\[
   |x(v)-e_i|=\frac{\gamma^2}{g(t_i)^2}
                               \frac{g(v-t_i)^2}{g(v)^2}.
\]
Evaluate at \(t_j\) for \(i\ne j\), using
\(t_j-t_i=t_k\) at the third index. Multiplication over the
three unordered pairs gives
\(\prod_{i<j}|e_i-e_j|=\gamma^6/(\prod_i g(t_i))^2=\gamma^6/4\)
by Lemma~\ref{lem:green486} at \(m=2\). The cubic discriminant
identity is \(g_2^3-27g_3^2=16\prod_{i<j}(e_i-e_j)^2\).
Taking absolute values proves the claim.
\end{proof}

\begin{lemma}[The arithmetic mean of a torsion Green value]\leavevmode
\label{lem:cmorbit487}
Let \(P\ne0\) be a torsion point on \(C\). Take a number
field \(J\) containing \(K\) over which the curve, its CM action
and \(P\) are defined. If \(g_\sigma\) denotes the normalized
function of Lemma~\ref{lem:green486} on \(C^\sigma\), then
\[
      \frac{1}{[J:\mathbb Q]}\sum_{\sigma:J\hookrightarrow\mathbb C}
                \log g_\sigma(P^\sigma)=b_K(P).
\]
\end{lemma}
\begin{proof}
For each ideal class take its complex torus \(C_I=\mathbb C/I\),
and for each integral ideal \(\mathfrak a\) use the quotient
isogeny \(f:C_I\longrightarrow C_{\mathfrak a^{-1}I}\).
Its kernel is \(C_I[\mathfrak a]\) of order \(N\mathfrak a\).
Use any nonzero complex coordinates on the source and target.
The product of \(g_{C_I}(v+Q)\) over the kernel equals
\(g_{C_{\mathfrak a^{-1}I}}(f(v))\). In fact the isogeny
pulls back probability area on the target to \(N\mathfrak a\)
times probability area on the source, so the logarithms have
equal Laplacian off their matching simple zeros. Their log
difference is harmonic on the torus and has mean zero, just
as in Lemma~\ref{lem:green486}.

For a torus \(C_L=\mathbb C/L\) put
\(h(C_L)=\log(\gamma_{C_L}\sqrt{\operatorname{area}(L)})\);
this depends only on its complex isomorphism class since
scaling a lattice scales the derivative and the square root
of the area inversely. Let \(v\) tend to zero in the product
identity. The tangent factor satisfies
\(|f'(0)|^2\operatorname{area}(I)
       =N\mathfrak a\,\operatorname{area}(\mathfrak a^{-1}I)\)
in compatible coordinates. It follows that
\[
 \sum_{0\ne Q\in C_I[\mathfrak a]}\log g_{C_I}(Q)
       =\tfrac12\log N\mathfrak a
             + h(C_{\mathfrak a^{-1}I})-h(C_I).
\]
Average over all ideal classes. The last two terms cancel since
multiplication by \(\mathfrak a^{-1}\) permutes the classes.
Partition the sum on the left by exact annihilator. Unique
factorization of ideals gives
\(\log N\mathfrak a=\sum_{\mathfrak b\mid\mathfrak a}
\Lambda_K(\mathfrak b)\), with the unit ideal contributing zero.
M\"obius inversion therefore makes the class average of the sum
over exact annihilator \(\mathfrak a\) equal to
\(\tfrac12\Lambda_K(\mathfrak a)\).

For the passage to embeddings we use the main theorem of complex
multiplication for elliptic curves in its reciprocity form
\cite[Part II]{XIII-X-IV-cm37}. For the CM type given by the chosen
embedding of \(K\), the finite ideles act on the completed ideal
lattice and its torsion by multiplication, up to simultaneous
isomorphism of the marked pair. Reciprocity identifies these
isomorphism classes with the Galois orbit over \(K\).
The ideal part reaches every lattice class. At a fixed class,
the integral idele units map onto
\((\mathcal O_K/\mathfrak a)^\times\) by local reduction.
The module \(\mathfrak a^{-1}I/I\) is free of rank one over
\(\mathcal O_K/\mathfrak a\), as is seen separately at each
prime power and then by the Chinese remainder theorem.
Thus the unit group acts transitively on its
\(\Phi_K(\mathfrak a)\) generators, exactly the points of
the stipulated annihilator. Passing to isomorphism classes
at each lattice identifies the action of \(\mathcal O_K^\times\).
Every generator has the same stabilizer, consisting of units
congruent to one modulo \(\mathfrak a\), and \(g\) is invariant
under these isomorphisms by its normalization. Consequently the
uniform class and generator average just computed is the uniform
average on the Galois orbit. Embeddings of \(J\) extending the
chosen embedding of \(K\) give equal multiplicities on that
orbit, by passing to a normal closure. The other embedding uses
\(\overline{\mathfrak a}\), for which \(b_K\) has the same value.
Division by the number of generators proves the asserted mean.
\end{proof}

\begin{theorem}[The rational norm of the complete elliptic jet determinant]\leavevmode
\label{thm:jetnorm487}
Use the square matrix \(E\) of Corollary~\ref{cor:confluent486}
with the original discriminant-one differential. Assume
\(P_*=Mt-\sum_b n_bP_b\ne0\). Take a coefficient field \(J\)
containing all the points, the CM action, the algebraic
translation data and the entries of \(E\). With
\(u=\sum_b\binom{n_b}{2}\), one has
\[
 \begin{aligned}
 \frac{\log|N_{J/\mathbb Q}\det E|}{[J:\mathbb Q]}
   ={}& (d/2-u)\log M+\sum_b\sum_{j=0}^{n_b-1}\log(j!)
              -d\,b_K(t)+b_K(P_*)\\
      &{}+\sum_{b<b'} n_bn_{b'}\,b_K(P_b-P_{b'}).
 \end{aligned}
\]
For a general good algebraic scaling of the differential add
\(\frac{u}{12[J:\mathbb Q]}\log|N_{J/\mathbb Q}(g_2^3-27g_3^2)|\)
to the right. If \(P_*=0\) the determinant itself is zero.
\end{theorem}
\begin{proof}
We first justify using the same marked table after each
embedding. Phases have to be compatible in a determinant.
On the \(M\)-cover, write \(s_a/s_0=R_a(Mv)\). Theorem~\ref{thm:trans57}
shows as germs at zero that \(\sigma(R_a(Mv))\) is a single
root of unity \(\epsilon_a\) times the transported torsion ratio
for \(a^\sigma\). Here the origin parameter is taken with
respect to the transported differential. Both ratios on the
cover are rational functions, so this equality holds globally.
For each \(b\) choose \(c\) with \(b+c\ne t\). The ratio of
the two row germs for \(a,c\) is \((s_a/s_c)(P_b+z/M)\).
After transport it is the corresponding ratio multiplied by
\(\epsilon_a/\epsilon_c\), choosing the lift of \(b^\sigma\)
to be \(M\) times the chosen lift of \(P_b^\sigma\).
The denominator germ \(U_bU_c\Theta_t\) is itself a root of
unity times \(U_{b+c}\Theta_t\), whose algebraic transport has
a single phase in every derivative order by the same theorem.
Thus the matrix after embedding differs from the analytically
transported matrix by unit phases on its columns and on each
row block, as well as a possible permutation of columns.
Their determinant absolute values agree.

Take logarithms of the formula of
Corollary~\ref{cor:confluent486}. Distinct \(b\)'s give distinct
\(P_b\)'s, \(t\) is nonzero, and the assumed condition makes
every factor nonzero. Average each term over the embeddings.
Lemma~\ref{lem:cmorbit487} applies individually to \(t,P_*\)
and all the differences, on the same field \(J\). Additivity
of the logarithm needs no independence between these points.
Lemma~\ref{lem:cmslope487} evaluates the remaining tangent
factor, also for the alternate scaling. The singular case
follows at once from the complex formula.
\end{proof}

\begin{corollary}[Prime-by-prime mean content and the total pair contribution]\leavevmode
\label{cor:localjetnorm487}
For a rational prime \(p\) define \(b_{K,p}(P)=0\) unless
\(\operatorname{Ann}(P)=\mathfrak p^e\) for a prime
\(\mathfrak p\mid p\) and \(e\ge1\). In this case put
\(b_{K,p}(P)=f(\mathfrak p/p)/(2\Phi_K(\mathfrak p^e))\).
On the discriminant-one model with \(\det E\ne0\), the exact
mean of its valuations \(v_\lambda\), normalized by \(v_\lambda(p)=1\),
is
\[
\begin{aligned}
 \frac{1}{[J:\mathbb Q]}\sum_{\lambda\mid p}
      [J_\lambda:\mathbb Q_p]\,v_\lambda(\det E)
   ={}&(d/2-u)v_p(M)+\sum_b\sum_{j=0}^{n_b-1}v_p(j!)\\
     &{}-d\,b_{K,p}(t)+b_{K,p}(P_*)\\
     &{}+\sum_{b<b'}n_bn_{b'}\,b_{K,p}(P_b-P_{b'}).
\end{aligned}
\]
Furthermore, if the differences \(P_b-P_{b'}\) belong to
\(C[\mathfrak m]\), with integral nonzero ideal \(\mathfrak m\),
and \(n_{\max}=\max_b n_b\), then
\[
 0\ \le\ \sum_{b<b'}n_bn_{b'}b_K(P_b-P_{b'})
              \ \le\ \tfrac14 d\, n_{\max}\log N\mathfrak m .
\]
Equality on the right holds for a full coset of
\(C[\mathfrak m]\) with equal multiplicities.
\end{corollary}
\begin{proof}
Every term in the rational norm formula is a finite rational
linear combination of logarithms of rational primes. The
coefficient of \(\log p\) in \(b_K(P)\) is \(b_{K,p}(P)\).
The logarithm on the left has coefficient
\(v_p(N_{J/\mathbb Q}\det E)/[J:\mathbb Q]\), equal to the
displayed local mean by the norm formula. These coefficients
are unique, since clearing denominators in a rational relation
among logarithms of distinct primes gives an integer
factorization of one. This includes primes dividing \(M\) or the
shift orders.

For fixed \(b\) the differences with \(b'\ne b\) are distinct.
For any nonunit ideal \(\mathfrak a\mid\mathfrak m\) at most
\(\Phi_K(\mathfrak a)\) of them have that exact annihilator.
Multiplying by \(b_K(\mathfrak a)\) and summing over the
divisors of \(\mathfrak m\) gives at most
\(\tfrac12\log N\mathfrak m\). Bound each \(n_{b'}\) by
\(n_{\max}\), multiply by \(n_b\), sum using \(\sum_b n_b=d\),
and divide by two to pass to unordered pairs. All terms are
nonnegative. A full coset with equal multiplicities attains
every counting bound.
\end{proof}

These are arithmetic norm and valuation identities for the same
elliptic block over the period field and any larger coefficient
field. They use only its listed torsion points, their
multiplicities and the original differential. The signed sums
over several edges and the prescribed projection in the GRH
argument have a different determinant. Their interaction and the
uniform cohomology estimate of Section~\ref{sec:missing} remain
to be established.


\subsection{Green energy and signed elliptic determinants}
\label{sec:greenenergy488}

The absolute norm of a single determinant is given by
Theorem~\ref{thm:jetnorm487}. In a sum of determinants the largest
summand can depend on the embedding. We first bound the product over
point differences uniformly over configurations at each embedding.

\begin{theorem}[Weighted elliptic Green energy]\leavevmode
\label{thm:greenenergy488}
Let \(g\) be the multiplicative Green function of
Lemma~\ref{lem:green486} on a complex torus of area \(A\) in its
chosen coordinate. Put
\[
 \widehat\gamma=\sqrt A\,\lim_{v\to0}g(v)/|v|,\qquad
 L_C=\log\pi+1+\gamma_{\rm E}-2\log\widehat\gamma ,
\]
where \(\gamma_{\rm E}\) is Euler's constant. At distinct points
\(z_i\) take positive real weights \(w_i\), with
\(N=\sum_iw_i\) and \(B=\sum_iw_i^2\). Then
\[
 \sum_{i<j} w_iw_j\log g(z_i-z_j)
       \le {B\over4}\bigl(\log(N^2/B)+L_C\bigr),
 \qquad
 \sup_{z\in C}\log g(z)\le\tfrac12(\log 2+L_C).
\]
\end{theorem}
\begin{proof}
Rescale the coordinate by \(A^{-1/2}\), so that the lattice
\(\Lambda\) now has covolume one. The function \(G=\log g\) has
mean zero and \(G(z)=\log|z|+\log\widehat\gamma+o(1)\) at zero.
In this coordinate its distributional Laplacian is
\(2\pi(\delta_0-1)\). Indeed the logarithm of a holomorphic
simple zero contributes \(2\pi\delta_0\), and the log of the metric
contributes the Laplacian of \(-\pi|z|^2/2\).
For \(t>0\) let \(G_t=H_t*G\), with the periodic heat kernel
\[
 H_t(z)=\frac{1}{4\pi t}\sum_{\lambda\in\Lambda}
                              e^{-|z+\lambda|^2/(4t)}.
\]
The kernel has integral one by unfolding the lattice sum.
The Fourier coefficient of \(G_t\) at a nonzero dual vector
\(\xi\) is \(-e^{-4\pi^2t|\xi|^2}/(2\pi|\xi|^2)\), by applying
the Laplacian to the Fourier character
\(e^{2\pi i\langle \xi,z\rangle}\). The constant coefficient is
zero. The series is absolutely convergent. Each nonconstant
term of \(\sum_{i,j}w_i w_j G_t(z_i-z_j)\) is this negative
coefficient times the squared modulus of a sum of characters.
The double sum is therefore nonpositive.

We have \(\partial_tG_t=2\pi(H_t-1)\). At \(z\ne0\), integration
from time zero and \(H_t\ge0\) give \(G(z)\le G_t(z)+2\pi t\).
At zero the contribution of the central Gaussian to the heat
average of \(\log|z|\) is
\(\tfrac12\log(4t)-\tfrac12\gamma_{\rm E}\). This follows by
the radial substitution \(r^2/(4t)=y\) and
\(\int_0^\infty e^{-y}\log y\,dy=-\gamma_{\rm E}\).
For the small-time limit one may restrict to a disk about zero
smaller than half the shortest period. The heat mass outside
the disk decays exponentially, also against the integrable
periodic logarithmic singularity. Inside the disk the continuous
remainder tends in heat average to \(\log\widehat\gamma\).
Integrating the displayed time derivative from this limit gives
\[
 G_t(0)=\tfrac12\log(4t)-\tfrac12\gamma_{\rm E}
       +\log\widehat\gamma-2\pi t
       +2\pi\int_0^t {1\over4\pi s}
                         \sum_{0\ne\lambda\in\Lambda}
                          e^{-|\lambda|^2/(4s)}\,ds .
\]
The integral is nonnegative. If \(E\) is the pair sum in
the theorem, the nonpositive double sum and the off-diagonal
inequality give \(2E\le-BG_t(0)+2\pi t(N^2-B)\).
Substitute the lower bound for \(G_t(0)\) and take
\(t=B/(4\pi N^2)\). This proves the first assertion. Two points
with unit weights give the second at every nonzero difference.
At zero \(g\) vanishes.
\end{proof}

For the following results use the discriminant-one model with
CM by \(\mathcal O_K\). Write \(d_0=|\operatorname{disc}K|\),
\(D_{\xi_{d_0}}=L'(1,\xi_{d_0})/L(1,\xi_{d_0})\), and set
\[
 L_K=1+\log\!\left({\sqrt{d_0}\over2\pi}\right)+D_{\xi_{d_0}} .
\]
Use \(\mathbb E_\sigma\) for the uniform mean over all complex
embeddings of the common coefficient field \(J\) of
Theorem~\ref{thm:jetnorm487}. Lemma~\ref{lem:cmslope487} gives
\(\widehat\gamma_\sigma=\sqrt{A_\sigma}\), and
Theorem~\ref{thm:mean56} gives \(\mathbb E_\sigma L_{C^\sigma}=L_K\).

\begin{theorem}[A signed sum with prescribed row multiplicities]\leavevmode
\label{thm:signedcomplete488}
Fix \(M\ge2,\ d=M^2\), the nonzero mark \(t\), and positive
multiplicities \(n_1,\ldots,n_m\) of sum \(d\). Let
\(u=\sum_i\binom{n_i}{2}\), \(B=\sum_i n_i^2\). Form any finite
list of the complete marked matrices of
Corollary~\ref{cor:confluent486}. Use these multiplicities in each
matrix and any distinct torsion shifts within that matrix. For
\(Z=\sum_\lambda c_\lambda\det E_\lambda\ne0\), \(c_\lambda\in J\),
\[
 \begin{aligned}
 \mathbb E_\sigma\log|\sigma Z|
 \le{}&\mathbb E_\sigma\log\sum_\lambda|\sigma c_\lambda|
       +(d/2-u)\log M +\sum_i\sum_{j=0}^{n_i-1}\log(j!)
       -d\,b_K(t)\\
     &+\tfrac B4\bigl(\log(d^2/B)+L_K\bigr)
                         +\tfrac12(\log2+L_K).
 \end{aligned}
\]
\end{theorem}
\begin{proof}
At each embedding apply the confluent formula to all the
transported matrices. The transport preserves their absolute
determinants by the row and column phase calculation in
Theorem~\ref{thm:jetnorm487}. The tangent factor is one before
the factor \(M^{-u}\). The weighted energy bound with
\(z_i=P_i=b_i/M,\ w_i=n_i\) bounds the log of the pair product.
The supremum bound controls the remaining factor at
\(\sum n_iP_i-Mt\). Both bounds hold simultaneously for every
tuple of the given shape. A singular summand contributes zero.
The triangle inequality now bounds \(|\sigma Z|\) by
\(\sum|\sigma c_\lambda|\) times the common upper bound.
Take logarithms and average. Lemma~\ref{lem:cmorbit487} gives
\(\mathbb E_\sigma\log g_\sigma(t^\sigma)=b_K(t)\).
\end{proof}

\begin{corollary}[Arbitrary derivative orders on the complete input basis]\leavevmode
\label{cor:arbitraryorders488}
Keep the same CM curve, mark, and input columns \(a\in C[M]\).
Choose nonnegative orders \(j_1,\ldots,j_d\). In summand
\(\lambda\) let row \(i\) be the order-\(j_i\) derivative at zero
of \(U_{b_{\lambda i}}U_a\Theta_t\) in column \(a\), with any
torsion shifts \(b_{\lambda i}\), including repetitions. Denote
the resulting matrix by \(F_\lambda\). Define \(C(0)=1\) and
\(C(j)=j!(e/j)^{j/2}\) for \(j\ge1\). If
\(Z=\sum_\lambda c_\lambda\det F_\lambda\ne0\), then
\[
\begin{aligned}
 \mathbb E_\sigma\log|\sigma Z|\le{}&
 \mathbb E_\sigma\log\sum_\lambda|\sigma c_\lambda|
 +\tfrac d2\log M-d\,b_K(t)
 +\tfrac d4(\log d+L_K)+\tfrac12(\log2+L_K)\\
 &+\sum_i\log C(j_i)
 +\tfrac12\Big(\sum_i j_i\Big)
   \left(D_{\xi_{d_0}}-\gamma_{\rm E}
                         -\log(2\pi/\sqrt{d_0})\right).
\end{aligned}
\]
\end{corollary}
\begin{proof}
Use a different complex displacement \(z_i\) in each physical row
before differentiating, and put \(c=\pi/A_\sigma\).
The determinant of these holomorphic rows has its weight
\(\exp(-c\sum_i|z_i|^2/2)\) in the translated frame.
Its metric norm equals the norm of the undifferentiated
alternant at the \(d\) points \((b_{\lambda i}+z_i)/M\) on the
cover. Indeed its metric is the pullback metric, and translating
each row is an isometry with a scalar common to all its columns.
Apply Theorem~\ref{thm:thetaalternant486} and the energy bound
with \(d\) unit weights. The determinant vanishes at collisions,
so the bound extends to the entire polydisk. Its weighted norm
is at most
\[
 M^{d/2}g_\sigma(t^\sigma)^{-d}
  \exp\!\left\{ \tfrac d4(\log d+L_{C^\sigma})
                     +\tfrac12(\log2+L_{C^\sigma})\right\}.
\]
For \(j_i>0\) use the Cauchy circle \(|z_i|=\sqrt{j_i/c}\);
for \(j_i=0\) evaluate at zero. Removing the Gaussian and applying
the derivative formula on these circles multiplies the bound by
\(\prod_i C(j_i)c^{j_i/2}\). Differentiating once in each
independent row variable to its specified order gives the
claimed determinant by multilinearity. The transport proof of
Theorem~\ref{thm:jetnorm487} applies to the whole row germ and
hence to every chosen order. Apply the triangle inequality and
average as before, using Theorem~\ref{thm:mean56} for
\(\mathbb E_\sigma\log c\).
\end{proof}


\subsection{Theta derivatives at ordinary primary torsion}
\label{sec:ordinaryprimary492}

Let \(C=\mathbb C/\Lambda\) carry the marked model of
Theorem~\ref{thm:germ57}. Write \(O\) for its origin, \(t\) for its
nonzero point of order two, \(x=\wp,\ y=\wp'\), and
\(\partial=d/dz\) for the vector field dual to \(dx/y\).
Fix a good split prime \(p\ge5\) prime to the CM discriminant,
a place \(\lambda\) above it, and a valuation \(v(p)=1\).
Work over finite extensions at that place as required and denote
reduction in the algebraic closure \(k\) of the residue field by a bar.
The special fiber \(C_k\) is ordinary. For example, the nonzero
scalar in \(\partial^p=H_p\partial\) of
Lemma~\ref{lem:splitcmder480} is the Hasse invariant in this
differential. Put \(J_n(F)=F^{(n)}(0)\), without factorials.

\begin{lemma}[An additive character of the ordinary torsion]\leavevmode
\label{lem:addext492}
The differential determines an injective homomorphism
\(\lambda_C:C_k[p](k)\longrightarrow(k,+)\) with the following
description. In characteristic zero put
\[
 Z(P)=\zeta(z)-u_\Lambda z-(\pi/A)\bar z,\qquad P=[z]\ne O.
\]
This is independent of the representative \(z\).
Let \(P\) be torsion and \(m'p^\nu P=O\), with \(p\nmid m'\)
and \(\nu\ge1\). If \(\bar P\ne O\), then \(p^\nu Z(P)\) is
integral and
\[
 \overline{p^\nu Z(P)}
     =-\frac{1}{\overline{m'}}\,\lambda_C([m'p^{\nu-1}]\bar P).
\]
\end{lemma}
\begin{proof}
We give the integral group in which multiplication computes this
residue. Enlarge the local ring to contain \(C[2]\). For
\(d\in C[2]\), take an affine-line coordinate \(s_d\) over the open
set on which \(P+d\ne O\). Set \(h_0=0\) and
\[
 h_d(P)=\frac{y(P)}{2(x(P)-x(d))}\quad(d\ne O).
\]
Glue the coordinates by \(s_d-s_e=h_d-h_e\). The only possible
poles of \(h_d\) lie at \(d\) and \(O\).
At \(O\) every \(h_d\) for \(d\ne O\) has principal part
\(-1/T\) in \(T=-2x/y\). Distinct points of order two remain
distinct in the residue field. Thus all the transition functions
are regular on their overlaps over the integer ring. This gives a
smooth additive torsor \(\pi:G\to C\).

For inputs above \(P,Q\), take coordinates with
\(P+d,Q+e,P+Q+d+e\ne O\) and add by
\[
  s_{d+e}(\text{sum})=s_d(\text{first})+s_e(\text{second})
                         +B(P+d,Q+e),\qquad
  B(U,V)=\frac{y(U)-y(V)}{2(x(U)-x(V))}.
\]
Such charts cover each geometric fiber. One shift is excluded
when choosing \(d\) and at most two when choosing \(e\).
On this open set the chord is regular. Near \(U=V\) its other
expression has numerator \(4(x(U)^2+x(U)x(V)+x(V)^2)-g_2\)
and denominator \(2(y(U)+y(V))\). The only common exceptional
case on the diagonal has \(U+V=O\), already excluded.
The zeta addition formula reads \(Z(U+V)=Z(U)+Z(V)+B(U,V)\).
Oddness gives \(Z(d)=0\) at a nonzero point of order two.
Thus the complex points have coordinates \(s_d=Z(P+d)+a\);
the law adds the two constants \(a\). This proves compatibility
and the group identities on the generic fiber, and hence on the
integral charts by density. Inversion is given by
\((P,s_d)\mapsto(-P,-s_d)\). The identity over \(O\) has
\(s_d=0\) for \(d\ne O\), and the kernel is the additive group.
The section of abstract complex groups with \(a=0\) supplies the
unique torsion lift of each torsion point: multiplication by a
nonzero integer is bijective on the characteristic-zero kernel.
In particular these lifts at torsion points are algebraic,
by the same uniqueness applied over \(\overline{\mathbb Q}\).

The reduced torsor admits no regular section. A section in the
chart \(d=0\) would give a rational function \(f\) regular off
\(O\). For \(d\ne O\) its coordinate \(f+h_d\) must be regular
at \(O\). Hence \(f\) would have precisely one simple pole.
Riemann--Roch in genus one excludes such a function.

On the special fiber \([p]_G\) kills the additive kernel and so
descends to a homomorphism \(f:C_k\to G_k\). One can construct
it by applying \([p]_G\) to any local section of the torsor.
The differential of \([p]_G\) is zero and \(\pi\) is smooth,
so \(df=0\). It follows that \(f=s\circ F\), where \(F:C_k\to
C_k^{(p)}\) is relative Frobenius. Indeed the pullbacks of
rational coordinates through \(f\) lie in \(k(C_k)^p\), the
kernel of the invariant derivation. On regular local rings of
the smooth curve a rational \(p\)-th power regular at a point
has a regular root there. This also proves regularity of \(s\);
the homomorphism identity follows after pullback by \(F\).
The equality \(\pi f=[p]_{C_k}\) gives \(\pi s=V\), with
\(V:C_k^{(p)}\to C_k\) the Verschiebung. On an ordinary
elliptic curve \(V\) is a separable isogeny of degree \(p\).
If \(s\) killed its kernel, it would descend through \(V\):
the translates by the kernel are precisely the fibers of this
finite \'etale quotient. The descended map would be a regular
section of \(\pi\), which the pole calculation rules out.
Consequently \(s\) is injective on the cyclic kernel.
Now \(F\) sends \(C_k[p](k)\) bijectively onto \(\ker V(k)\).
Taking \(sF\) on those points gives the asserted \(\lambda_C\)
in the additive coordinate over \(O\).

Finally choose an integral point \(Q_0\) of \(G\) above \(P\)
with \(s_0=0\), possible since \(\bar P\ne O\).
Its torsion lift has coordinate \(Z(P)\). Subtracting the two
points gives this scalar in the kernel. Multiplication by
\(m'p^\nu\) gives
\(p^\nu Z(P)=-[m'p^\nu]_G Q_0/m'\) in that kernel.
The numerator on the right is integral and its reduction is
\(f([m'p^{\nu-1}]\bar P)\), as claimed.
\end{proof}

\begin{theorem}[Leading residues at every prescribed derivative order]\leavevmode
\label{thm:residuejet492}
Let \(M=p^\nu m_0,\ \nu\ge1,\ p\nmid m_0\), and \(M[w]=O\).
Put \(P=t+[w]\). Take \(m'=2m_0\) and set
\[
  L_\nu(w)=-\frac{\lambda_C([m'p^{\nu-1}]\bar P)}
                          {\overline{m'}} .
\]
This value depends additively on \(\overline{[w]}\) and is
independent of the prime-to-\(p\) multiple used to kill the
prime-to-\(p\) parts of \(t\) and \([w]\).
For \(F_w=U_w\Theta_t\) and all \(n\ge0\),
\[
 p^{\nu n}J_n(F_w)\ \text{is integral},\qquad
 \overline{p^{\nu n}J_n(F_w)}
       =\overline{J_0(F_w)}\,L_\nu(w)^n,\qquad 0^0=1.
\]
If the \(p\)-primary part of \(\overline{[w]}\) has order
\(p^\nu\), then \(v(J_n(F_w))=-\nu n\) for every \(n\ge0\).
\end{theorem}
\begin{proof}
Multiplication by \(m'\) removes the prime-to-\(p\) part.
On the remaining \(p\)-torsion after multiplication by
\(p^{\nu-1}\), the scalar \(\overline{m'}\) is invertible.
This proves independence and additivity. In particular \(L_\nu\)
is nonzero precisely when the indicated primary part has full order.

The value \(J_0\) is integral and is a unit if \(\bar P\ne O\).
To see this even when \(p\mid M\), write
\((F_w/\Theta_t)^M=c h\) as in Theorem~\ref{thm:trans57},
where \(h\) has horizontal divisor \(M([t-[w]]-[t])\) and
vertical order zero on the good model. The product of its \(M\)
translates by \([w]\) is a constant of valuation zero, giving
\(v(c)=0\). The pole section \(t\) misses \(O\); the potential
zero section \(t-[w]=-P\) misses \(O\) in the reduction when
\(\bar P\ne O\). Evaluate at \(O\) and use \(\Theta_t(0)=1\).

When \(P\ne O\), the logarithmic derivatives \(b_j\) of \(F_w\)
at zero satisfy
\[
 b_1=Z(P),\qquad b_2=-u_\Lambda-x(P),\qquad
 b_j=-\partial^{j-2}x(P)\quad(j\ge3).
\]
Indeed \(F_w\) is proportional to the Hermitian translate of
\(e^{-u_\Lambda z^2/2}\sigma(z)\) by a lift of \(P\), so the
translation exponential contributes exactly the antiholomorphic
term in \(b_1\). The formula for \(u_\Lambda\) in
Theorem~\ref{thm:germ57} gives integrality here.
If \(\bar P\ne O\), then \(b_j\) is integral for \(j\ge2\),
and Lemma~\ref{lem:addext492} gives the residue of \(p^\nu b_1\).

Suppose next that \(\bar P=O\) but \(P\ne O\). The point is in
the formal group, with \(0<v(T(P))\le1/(p-1)\).
For the upper bound, the strict formal logarithm obtained by
integrating the unit differential has form
\(T+\sum_{j\ge2}a_jT^j/j\) with integral \(a_j\).
For \(v(T)>1/(p-1)\) the first term has strictly smallest
valuation, since \(v_p(j)\le(j-1)/(p-1)\); the series converges
and is additive there by its formal identity. This neighborhood
is stable under the integral group law. Hence it contains no
nonzero torsion. In the chart \(d\ne O\) the torsion lift has
coordinate \(Z(P)+h_d(P)\). The integral lift with \(s_d=0\)
in the same argument as in the lemma gives
\(\overline{p^\nu(Z(P)+h_d(P))}=\lambda_C(O)=0\).
The Laurent series of \(h_d\) has integral coefficients and
pole order one at \(O\). Thus \(p^\nu h_d(P)\) is in the
maximal ideal. Also \(T^j\partial^{j-2}x\) is integral in the
formal coordinate for \(j\ge2\), as \(\partial\) is an integral
unit times \(d/dT\). It follows that \(p^{\nu j}b_j\) is in the
maximal ideal for every \(j\ge1\).

For either case \(P\ne O\), apply the integer polynomial identity
for \(J_n(F_w)/J_0(F_w)\) in the \(b_j\). Every monomial has
weight \(n\) with \(b_j\) of weight \(j\), and the sole monomial
in \(b_1\) alone is \(b_1^n\), with coefficient one. This follows
by differentiating \(F_w'=F_w(\log F_w)'\) repeatedly. After
multiplication by \(p^{\nu n}\), all terms with \(j\ge2\)
vanish in the residue field. If \(P=O\) in characteristic zero,
\([w]=t\), \(J_0=0\), and all the derivatives are integral by
Theorem~\ref{thm:trans57} at this prime-to-\(p\) translate.
This covers the remaining case. At full primary order \(\bar P\ne O\)
and both factors in the residue formula are nonzero, giving the
claimed valuation.
\end{proof}

\begin{corollary}[Leading reduction of a signed homogeneous test]\leavevmode
\label{cor:residuecolumn492}
In a fixed basis write one of the signed columns of
Theorem~\ref{thm:projectedmatrices489} as
\[
  Y_i=\sum_{a,\alpha}v_{i,a,\alpha}
                    \prod_h J_{\alpha_h}(F_{a_h}),\qquad
  |\alpha|=k_i .
\]
Use the actual shifts (with their chosen lifts) and a common
annihilator \(M\) as above. Choose a scalar \(c_i\ne0\) making
all \(c_iv_{i,a,\alpha}\) integral at the given place. Put
\(P_{i,a}(X)=\sum_\alpha c_i v_{i,a,\alpha}X^\alpha\).
The integral column \(c_i p^{\nu k_i}Y_i\) has reduction
\[
  \sum_a\Big(\prod_h\overline{J_0(F_{a_h})}\Big)
           \overline{P_{i,a}}\big((L_\nu(a_h))_h\big).
\]
On conjugate factors use the corresponding transported differential
and reduction. If a square minor of these reduced columns is
nonzero, the corresponding minor of the \(Y_i\) has valuation
\(-\sum_i(v(c_i)+\nu k_i)\), with the sum over its columns.
\end{corollary}
\begin{proof}
Distribute \(p^{\nu k_i}\) over the individual derivatives using
homogeneity. Apply Theorem~\ref{thm:residuejet492} factor by
factor and sum in the residue field. Taking a determinant commutes
with reduction. A nonzero residue gives valuation zero before
undoing the column scalings.
\end{proof}
The formula keeps the finite character coefficients and the original
projection. To use it in Section~\ref{nav:I-28} one needs a
nonzero minor, or control of the higher terms if the displayed
residue vanishes.


\subsection{Division functions for normalized theta translates}
\label{sec:divisionphase493}

Keep the marked elliptic curve, the lattice, and the normalization
\(\Theta_t(0)=1\) of Theorem~\ref{thm:germ57}. The exact scalar in
a torsion translation matters when several translations are added.
We recover it by using the evenness of \(\Theta_t\) on two opposite
points of a division cover. This is the symmetric translation
construction in the algebraic theory of theta functions; see
\href{https://arxiv.org/abs/math/0610163}{K. Bannai and
S. Kobayashi, \emph{Algebraic theta functions and \(p\)-adic
interpolation of Eisenstein--Kronecker numbers}}.

\begin{theorem}[A rational expression with its translation scalar]\leavevmode
\label{thm:divisionphase493}
Fix \(w\in M^{-1}\Lambda,\ M\ge1\). Put \(N=2M\),
\(b=[w/2]\in C[N]\), and \(d=[w/(2N)]\in C[N^2]\).
Choose a rational function \(r_b\) on \(C\) with divisor
\([N]^*([t-b]-[t])\). It can be chosen over any characteristic-zero
field containing the marked model and the point \(b\).
For the fixed complex lift \(w\) there is an identity of
meromorphic functions
\[
 \frac{U_w\Theta_t(z)}{\Theta_t(z)}
       =\frac{r_b(d+[z/N])}{r_b(-d-[z/N])}.
\]
In particular the marked model, \(u_\Lambda\), and \(C[N^2]\)
generate a field containing every derivative of every such
translate, with the relative scalars of different translates.
\end{theorem}
\begin{proof}
The divisor has degree zero. Under \(\operatorname{Pic}^0(C)=C\)
its class is \(-Nb=O\) since multiplication by \(N\) is its own
dual. Its principality over the field of definition follows by
base change for global sections of a line bundle: the space
of functions with the required divisor inequality has dimension
one after extension to an algebraic closure, hence already over
the base. Equality of divisors follows from degree zero.

Write \(b_0=w/2\) for the complex lift and
\(R(z)=U_{b_0}\Theta_t(z)/\Theta_t(z)\).
The transformation rule under a period \(\lambda\) is
\[
 R(z+\lambda)=\exp(2\pi i E(\lambda,b_0))R(z),
       \qquad E(v_1,v_2)=\operatorname{Im}(\bar v_1v_2)/A .
\]
Indeed division of the two theta multipliers cancels their
signs and leaves the exponent
\((\pi/A)(\bar\lambda b_0-\lambda\bar b_0)\).
Thus \(R(Nv)\) is periodic with the divisor defining \(r_b\).
It equals \(c\,r_b([v])\) for some nonzero constant \(c\).
By evenness of \(\Theta_t\) and the two translation exponentials,
\(R(-z-b_0)R(z)=1\). Also \(R(z+b_0)R(z)=
U_w\Theta_t(z)/\Theta_t(z)\) by multiplication of the collinear
translation factors. Substitution gives the displayed quotient,
where \(c\) cancels. This computation is valid as an identity
of meromorphic functions, including a common zero or pole in its
right side.

For derivatives, expand at \(d\) and \(-d\) by translation on
the algebraic curve and use the parameter \(T=-2x/y\) at \(O\).
The power series \(T(z/N)\) is defined over the model field:
integrate \(dx/y\), which has leading term \(dT\), then invert
the series and substitute \(z/N\). Division in the Laurent
series field handles any common vanishing orders. The resulting
quotient is holomorphic at zero because \(\Theta_t(0)=1\).
The derivatives of \(\Theta_t\) are in the marked model field
with \(u_\Lambda\) by the recurrence of Theorem~\ref{thm:germ57}.
This proves all field assertions with the specified lift.
\end{proof}

We choose the convention
\[
 e_N([s_0],[q_0])=\exp(2\pi iN E(s_0,q_0)),
       \qquad Ns_0,Nq_0\in\Lambda ,
\]
for the Weil pairing. It is perfect and alternating. Its
values are algebraic \(N\)-th roots of unity. The multiplier
calculation in the proof also identifies it without a choice of
a primitive root.

\begin{corollary}[Field automorphisms of a fixed translate]\leavevmode
\label{cor:deckphase493}
Let \(B\) contain the marked model, \(u_\Lambda\), and \(C[N]\).
For \(\sigma\in\operatorname{Gal}(B(C[N^2])/B)\) and \(j\ge0\),
\[
  \sigma\big(J_j(U_w\Theta_t)\big)=
       e_N(\sigma d-d,b)^2 J_j(U_w\Theta_t).
\]
\end{corollary}
\begin{proof}
Take \(r_b\) over \(B\). For \(s=[s_0]\in C[N]\), the equality
\(R(Nv)=c\,r_b([v])\) and the multiplier for the period \(Ns_0\)
give \(r_b(v+s)=e_N(s,b)r_b(v)\).
Now \(\sigma d-d\in C[N]\). In the rational expression of
Theorem~\ref{thm:divisionphase493}, the numerator and denominator
move by opposite points. The ratio therefore acquires the
square of this pairing. All base coefficients of the germ of
\(\Theta_t\) are fixed. Differentiating proves the assertion
also at a zero. The scalar is a character of \(\sigma\), since
the difference \(\sigma d-d\) is additive in \(\sigma\) when its
values are fixed by the group.
\end{proof}

\begin{theorem}[Galois decomposition at the prescribed tests]\leavevmode
\label{thm:fixedsectors493}
Consider any chosen exterior product \(\Xi=Y_1\wedge\cdots\wedge Y_s\)
of the prescribed columns in Theorem~\ref{thm:projectedmatrices489}.
Choose a coordinate basis for their common space and a field \(B\)
containing the entries of the finite coefficient vectors in that
basis, the model data just used, and \(C[N]\) for every elliptic
factor. Take \(M\) to kill all shifts and \(N=2M\).
Let \(H\) be any subgroup of the Galois group over \(B\) of the
compositum of the \(N^2\)-division fields of those curves.

In the full multilinear expansion of \(\Xi\), assign a character
to a term by multiplying the characters of Corollary~\ref{cor:deckphase493}
over all its translated factors. Write \(\Xi_\chi\) for the sum
of the terms assigned to \(\chi\), over the original signed
coefficients. Then
\[
  \Xi_\chi=\frac1{|H|}\sum_{\sigma\in H}
                \chi(\sigma)^{-1}\sigma\Xi,\qquad
  \Xi\ne0\ \Longleftrightarrow\
          \Xi_\chi\ne0\ \text{for some }\chi .
\]
For any embedding into \(\mathbb C\) and any fixed Hermitian
norm on the exterior space,
\[
       \frac1{|H|}\sum_{\sigma\in H}\|\sigma\Xi\|^2
                       =\sum_\chi\|\Xi_\chi\|^2 .
\]
At a finite place with valuation \(v\), for any one minor \(z\)
and the corresponding coordinates \(z_\chi\), put
\(m_H=\min_{\sigma\in H}v(\sigma z)\).
If \(z\ne0\), then
\[
             m_H-v(|H|)\ \le\ \min_\chi v(z_\chi)
                                  \ \le\ m_H .
\]
\end{theorem}
\begin{proof}
Relative to the lower field the action on \(C[N^2]\) is contained
in the matrices \(1+ND\) modulo \(N^2\). They commute, including
when \(N\) is composite. This holds on each factor, so \(H\) is
finite abelian. The lower field contains the roots of unity
occurring in the characters, by the Weil pairing on its rational
\(N\)-torsion. Since each scalar derivative transforms by its
stated character and each finite coefficient is fixed in the
chosen basis, every term of the expansion is an eigenvector.
Complex or algebraic character orthogonality gives the projector
formula, the equivalence, and, on expanding a squared norm, the
sum of squares. For the finite-place bounds use the nonarchimedean
triangle inequality first on each projector, then on
\(\sigma z=\sum_\chi\chi(\sigma)z_\chi\). All character values
are units there.
\end{proof}

The group in this theorem acts on algebraic numbers at the
given tests. The decomposition preserves each coupled class
index in the period. Terms carrying the same character still
enter a single signed sum. The theorem gives a determinant
criterion even if only a proper subgroup of the division group
acts on the coefficients.


\subsection{Independence of prescribed derivatives at split primes}
\label{sec:indepder494}

Let \(p^\nu\Vert M,\ \nu\ge1\), for a good split prime \(p\ge5\).
Write \(N=2M\). Fix an embedding into \(\mathbb C\), let \(\iota\)
be complex conjugation, and fix a place \(\mathfrak P\) of
\(\overline{\mathbb Q}\) above \(p\). Consider a finite list of
the marked CM curves with the same imaginary quadratic field
\(K\). Let \(B/\mathbb Q\) be a finite Galois extension containing
\(K\), their marked models and differentials, the constants
\(u_\Lambda\), and the fields of definition of their CM actions.
Assume each curve has good reduction at the places above \(p\)
and \(p\nmid e(B_{\mathfrak P}/\mathbb Q_p)\).
Denote by \(B_0\) the compositum of \(B\), the maximal rational
abelian extension \(\mathbb Q^{\rm ab}\), and the fields of the
\(N\)-torsion of all the curves. Allowing the abelian extension
here keeps finite characters of arbitrarily large conductor fixed.

These conditions on \(B\) allow the discriminant-one models of
the manuscript at every split \(p\ge5\) prime to the CM discriminant.
Indeed for full CM by \(\mathcal O_K\) the classical theorem of
complex multiplication places \(j\) in the Hilbert class field of
\(K\). That field is Galois over \(\mathbb Q\) and unramified at
such \(p\). Adjoin \(\mu_6\) and all the conjugates of the roots
in
\[
          g_2^3=j/1728,\qquad 27g_3^2=g_2^3-1 .
\]
The degrees above the Hilbert class field have only 2 and 3
as prime divisors. Good reduction follows from the unit
discriminant. The full two-torsion is unramified at these good
primes. Each CM endomorphism descends over the resulting field:
Galois conjugation preserving its tangent multiplier preserves
the endomorphism, since the tangent representation of CM in
characteristic zero is faithful. The expression
in \eqref{eq:quasi57} for \(u_\Lambda\) uses a Galois-stable
kernel and hence already belongs to the same field. Taking all
marked conjugates gives the required Galois field. This choice
does not require tame ramification for the subsequent rational
abelian coefficients adjoined in \(B_0\).

\begin{lemma}[Opposite inertia actions]\leavevmode
\label{lem:splitinertia494}
For each \(t_0\in\mathbb Z/p^\nu\mathbb Z\) there is an
automorphism \(h\) of \(\overline{\mathbb Q}\) fixing \(B_0\).
It is trivial on the prime-to-\(p\) torsion of every curve on the
list. On each \(C[p^{2\nu}]\) its actions on the connected and
the \'etale CM lines at \(\mathfrak P\), respectively, are
\(u=1+p^\nu t_0\) and \(u^{-1}\) modulo \(p^{2\nu}\).
\end{lemma}
\begin{proof}
The CM idempotents in
\(\mathcal O_K\otimes\mathbb Z_p=\mathbb Z_p\times\mathbb Z_p\)
split the \(p\)-divisible group of a good ordinary curve into
two height-one groups. They extend over the good model since
the CM endomorphisms do. The summand acting on the tangent line
is connected in the special fiber, and the summand with zero
tangent space is \'etale. At the conjugate place the two
idempotents are interchanged, as their tangent multipliers come
from the two embeddings at \(p\). Inertia over \(B\) fixes the
\'etale torsion. Its scalar action on the other rank-one Tate
module is the \(p\)-adic cyclotomic character, by the determinant
identity from the Weil pairing. Inertia also fixes prime-to-\(p\)
torsion by smooth good reduction.

Absolute inertia at \(\mathfrak P\) maps onto
\((\mathbb Z/p^{2\nu}\mathbb Z)^\times\) under the cyclotomic
character because the corresponding local cyclotomic extension
is totally ramified. Intersecting inertia with the absolute
Galois group of \(B\) cuts its index by
\(e(B_{\mathfrak P}/\mathbb Q_p)\), prime to \(p\).
Consequently the image of this intersection contains the whole
\(p\)-Sylow subgroup. Choose \(g\) in it with cyclotomic value
\(u\) at this level. The element \(g'=\iota g\iota^{-1}\) lies
in inertia at the conjugate place, fixes \(B\) by normality,
and has the same cyclotomic value. Thus \(h=g(g')^{-1}\) acts
in the stated way on both lines, fixes the prime-to-\(p\)
torsion, and fixes \(C[N]\) since \(u\equiv1\bmod p^\nu\).
It is a commutator, so it fixes \(\mathbb Q^{\rm ab}\). The same
\(g\) works on every curve.
\end{proof}

On each curve use a basis \((E_2,F_2)\) of the two lines at
level \(p^{2\nu}\), connected first. Take it so that at level
\(p^\nu\) the basis \((E_1,F_1)=([p^\nu]E_2,[p^\nu]F_2)\)
has Weil pairing a fixed primitive root \(\zeta_{p^\nu}\), with
the sign convention of Corollary~\ref{cor:deckphase493}.
For \(a=[w]\) write \(x_a,y_a\) for the coordinates of its
\(p\)-primary part in this lower basis.

\begin{theorem}[The character of a prescribed derivative]\leavevmode
\label{thm:labelder494}
For \(h\) of Lemma~\ref{lem:splitinertia494} the exact action on
every \(M\)-torsion translate, including its choice of lift, is
\[
      h\big(J_j(U_w\Theta_t)\big)
          =\zeta_{p^\nu}^{\,t_0 x_a y_a}J_j(U_w\Theta_t)
                 \qquad(j\ge0).
\]
In a product of derivatives the label in the exponent is
\(\sum_a x_a y_a\), counted with multiplicity. At several
distinct good split primes satisfying the same hypotheses
the labels can be separated simultaneously over \(B_0\).
\end{theorem}
\begin{proof}
Use \(d=[w/(2N)]\) and \(b=[w/2]=[N]d\) of
Theorem~\ref{thm:divisionphase493}. Write the \(p\)-part of \(d\)
as \(X E_2+Y F_2\). With \(N'=N/p^\nu\) we have
\[
 (h-1)d=t_0X E_1-t_0Y F_1,\qquad
 b_p=N'(X E_1+Y F_1),\qquad
 (x_a,y_a)=2N'(X,Y)\pmod {p^\nu}.
\]
Here \(h\) fixes every other primary part, and
\((1+p^\nu t_0)^{-1}=1-p^\nu t_0\bmod p^{2\nu}\).
Distinct primary parts pair trivially. On \(C[p^\nu]\) the
pairing \(e_N\) is \(e_{p^\nu}^{N'}\), as follows directly from
the alternating form defining it. Hence
\[
 e_N((h-1)d,b)^2
   =\zeta_{p^\nu}^{\,4t_0(N')^2XY}
   =\zeta_{p^\nu}^{\,t_0x_a y_a}.
\]
Corollary~\ref{cor:deckphase493} proves the derivative formula.
Characters multiply on products. Each of the automorphisms
constructed for another prime acts trivially on all the
prime-to-that-prime torsion. Thus their restrictions commute on
the common division field and the independent choices of \(t_0\)
give the direct product of label characters there.
\end{proof}

\begin{corollary}[Independent derivatives of arbitrarily prescribed orders]\leavevmode
\label{cor:indepder494}
For one of the curves choose \(p^\nu\) shifts \(w_u\) with
\(p\)-parts \((u,1)\), one for every \(u\in\mathbb Z/p^\nu\mathbb Z\).
Their prime-to-\(p\) components killed by \(M\), their lifts, and
their derivative orders \(j_u\ge0\) can be chosen independently.
Then the \(p^\nu\) numbers \(J_{j_u}(U_{w_u}\Theta_t)\) are
linearly independent over \(B_0\). More generally, in a signed
period with coefficient vectors over \(B_0\), a component with
an isolated nonzero term having a specified tuple of labels
cannot be cancelled by terms with a different tuple.

For a pure terminal frequency of
Lemma~\ref{lem:terminalfourier454}, supported on one CM half in
a locally integral CM product basis, its label at the terminal
prime is zero.
\end{corollary}
\begin{proof}
At \(\mathfrak P\) the \(y\)-coordinate specializes injectively
on the \'etale line. Thus \((u,1)\) specializes with full primary
order. Every stated derivative is nonzero by
Theorem~\ref{thm:residuejet492}. Applying the powers of an
automorphism for \(t_0=1\) to a relation over \(B_0\) gives the
Fourier matrix \((\zeta_{p^\nu}^{\,vu})_{v,u}\) on the summands.
Its inverse extracts each summand, proving independence. The
same character projector proves the assertion on isolated terms.
At the terminal place a pure frequency is in the dual position
component. In each scalar factor one of its split coordinates
is therefore zero. This gives the last claim and explains why
these terminal frequencies alone furnish no distinct labels.
\end{proof}


\subsection{Cyclic terminal Fourier projections on an elliptic curve}
\label{sec:cyclicneighbor495}

Let \(\Lambda\) be a scaled fractional ideal of \(\mathcal O_K\) in
the marked elliptic model, and let \(p\mathcal O_K=\mathfrak p
\bar{\mathfrak p}\) be a good split prime \(p\ge5\). Take the
frequency subgroup \(\mathfrak p^{-1}\Lambda/\Lambda\).
Choose a positive basis \(\omega_1,\omega_2\) of \(\Lambda\)
such that \(w=2\omega_1/p\) generates this subgroup and
\(\omega_1\in\bar{\mathfrak p}^{\,2}\Lambda\).
The proof below shows that this choice is possible for any
prescribed nonzero generator modulo \(\Lambda\).
The factor two gives \(pw\in2\Lambda\). For
\(k\in\mathbb F_p\), \(\zeta_p=\exp(2\pi i/p)\), put
\[
 S_k(z)=\frac1p\sum_{j=0}^{p-1}\zeta_p^{-kj}
                           U_{jw}\Theta_t(z).
\]

\begin{lemma}[The elliptic curve and the scalar of the Fourier sum]\leavevmode
\label{lem:cyclicneighbor495}
Put \(\Gamma=\mathbb Z\omega_1+\mathbb Z p\omega_2\) and
\(\Lambda'=\mathbb Z(\omega_1/p)+\mathbb Z p\omega_2\).
They are \(\bar{\mathfrak p}\Lambda\) and
\(\mathfrak p^{-1}\bar{\mathfrak p}\Lambda\), respectively.
On \(\mathbb C/\Lambda'\) there is a nonzero two-torsion mark
\(t'\) such that \(S_0=c\Theta'_{t'}\), where
\(\Theta'_{t'}(0)=1\) and \(c=S_0(0)\ne0\).
For each \(k\) choose \(\lambda_k\in\Lambda\) satisfying
\(\exp(2\pi iE(\lambda_k,w))=\zeta_p^{-k}\).
Then \(S_k=\epsilon_t(\lambda_k)^{-1}U_{\lambda_k}S_0\), with
\(\epsilon_t\) the original period sign. For \(k\ne0\) the
class of \(\lambda_k\) has exact order \(p\) in the
\(\bar{\mathfrak p}\)-torsion of \(\mathbb C/\Lambda'\).
In the common variable \(z\) there are differences \(e\) and
\(e'\) between distinct Weierstrass two-torsion \(x\)-coordinates
on the original and the new curve such that
\[
                          c^4=p^{-2}\,e'/e .
\]
\end{lemma}
\begin{proof}
The subgroup of \(\Lambda\) with \(E(\gamma,w)\) integral is
\(\Gamma\), by the displayed basis. It is also
\(\bar{\mathfrak p}\Lambda\): the adjoint of a CM scalar for
the alternating form is its complex conjugate, and this ideal
lattice has index \(p\) and annihilates the frequency under
the pairing. Adjoining \(w\) to \(\Gamma\) gives \(\Lambda'\).
The extra congruence in the choice of the basis is possible
with the prescribed generator. Indeed one can
lift its half at the first prime modulo \(p^2\), take zero at
the other, and lift the resulting primitive pair modulo \(p^2\)
to a unimodular integral basis. Elementary matrices generate
\(\operatorname{SL}_2(\mathbb Z/p^2\mathbb Z)\) by elimination
at a unit entry. With this choice \(w\) belongs to the claimed
fractional ideal, and the two superlattices of \(\Gamma\) have
the same index \(p\). In particular
\(\Lambda\cap\Lambda'=\Gamma\), and \(\Lambda/\Gamma\) is the
stated opposite CM line.

All \(U_{jw}\Theta_t\) are sections on the common cover by
\(\Gamma\). The operator \(U_\lambda\) for \(\lambda\in\Lambda\)
acts there on the \(j\)-th section as
\(\epsilon_t(\lambda)\exp(2\pi ijE(\lambda,w))\), by the
multiplier in Theorem~\ref{thm:divisionphase493}. This gives
the asserted formula for \(S_k\). Also \(U_w S_0=S_0\) because
\(U_{pw}\Theta_t=\Theta_t\). Thus \(S_0\) has the principal
polarization on \(\Lambda'\), which has the same area as
\(\Lambda\). It is even. It is nonzero since the \(p\) terms
in its sum have distinct eigencharacters under \(\Lambda/\Gamma\).
Its unique zero in this degree-one line is simple, fixed by
negation and away from zero by evenness. This gives \(t'\).
The quotient by the marked Hermitian theta with the same
divisor has periodic absolute value, since its remaining
period multipliers are unitary. It is constant by Liouville's
theorem.

To compute that constant put \(\tau=\omega_2/\omega_1\),
\(s=z/\omega_1\), and \(A=\operatorname{Im}(\bar\omega_1\omega_2)\).
For the even half-characteristic \((a,b)\) of the original mark
write
\[
 \vartheta_{a,b}(s,\tau)=
    \sum_{n\in\mathbb Z}\exp\{\pi i(n+a)^2\tau+
                                      2\pi i(n+a)(s+b)\}.
\]
The normalized Hermitian theta is this function divided by its
value at zero and multiplied by
\(\exp(\pi\bar\omega_1 z^2/(2A\omega_1))\).
This follows from the period multipliers of the series on
\(\omega_1,\omega_2\); it has the given even characteristic
and value one. Translation along \(2j\omega_1/p\) translates
the argument of the series by \(2j/p\) with the same exterior
Gaussian at \(z\). The average over \(j\) retains precisely
\(n+a=p(m+a)\), since \(p\) is odd. Consequently
\[
 c=\frac{\vartheta_{a,pb}(0,p^2\tau)}
          {\vartheta_{a,b}(0,\tau)} .
\]
The Jacobi identities identify the fourth power of each even
theta constant with the square of its first full period divided
by \(\pi^2\), times a difference of the Weierstrass roots.
These are the half-period formulas in
\href{https://dlmf.nist.gov/23.6.i}{NIST, Weierstrass functions
and theta functions, \S23.6(i)}, together with
\(\vartheta_3^4=\vartheta_2^4+\vartheta_4^4\).
Reducing a half-characteristic modulo integers introduces only
a sign here. The two first full periods are \(\omega_1\) and
\(\omega_1/p\), proving the identity in the common differential.
\end{proof}

\begin{theorem}[Integral derivatives of a complete cyclic projection]\leavevmode
\label{thm:cyclicunit495}
Use the choice of generator to precision \(p^2\) made in
Lemma~\ref{lem:cyclicneighbor495}. Let \(b_0\) lift any
prime-to-\(p\) torsion point of \(\mathbb C/\Gamma\).
At a place over \(p\) where the original differential is a good
invariant differential, use \(v(p)=1\). If the original frequency
line \(C[\mathfrak p]\) is connected in the reduction, then
\[
 v(c)=0,\qquad
 v\big(J_n(U_{b_0}S_k)\big)=0\quad (k\ne0,\ n\ge0),
 \qquad v\big(J_n(U_{b_0}S_0)\big)\ge n .
\]
At the other orientation \(v(c)=-1\), and every \(k,n\) satisfies
\(v(J_n(U_{b_0}S_k))\ge-1-n\).
\end{theorem}
\begin{proof}
The two isogenies
\[
    C_\Gamma=\mathbb C/\Gamma\ \longrightarrow\
    C=\mathbb C/\Lambda,\qquad
    C_\Gamma\ \longrightarrow\ C'=\mathbb C/\Lambda'
\]
have degree \(p\) and respective kernels the
\(\bar{\mathfrak p}\)-line and the \(\mathfrak p\)-line. One
can start algebraically with the quotient by a CM torsion
subgroup and its dual; the maps and the differential pulled
back from \(C\) are defined over a number field. The curves
have good models after finite extension since they have CM
and integral \(j\). For invariant differentials of good models
the pullback multiplier of an isogeny with \'etale reduction is
a unit. For these degree-\(p\) isogenies with connected kernel
the dual has \'etale kernel, so the two multipliers have
valuations one and zero by their composition \([p]\).

Write \(\alpha\,dz\) for a good invariant differential of
\(C'\), where \(dz\) denotes the common analytic variable.
At the connected frequency orientation the isogeny to \(C\)
is \'etale and the one to \(C'\) has connected kernel. Since
\(dz\) on \(C\) is good, \(v(\alpha)=1\). At the other
orientation the roles are reversed and \(v(\alpha)=-1\).
Choose a discriminant-one model of \(C'\); its differential
has this same valuation relative to \(dz\).
In the normalized coordinate \(z'=\alpha z\), every difference
between distinct two-torsion roots is a unit, because the
discriminant is a unit and \(p\ge5\). Returning to \(z\)
multiplies the root differences by \(\alpha^2\).
The lemma therefore gives \(4v(c)=-2+2v(\alpha)\).

The section \(U_{b_0}S_k\) is \(c\) times the translate of
the normalized theta on \(C'\) by \(b_0+\lambda_k\), up to a
root of unity from the period sign and the translation law.
All the translations are rational in the common real period
lattice, so the composition phases are indeed roots of unity.
Passing to \(z'\) multiplies its \(n\)-th derivative in the
original variable by the factor \(\alpha^n\) relative to
the derivative in \(z'\). The equality of the Hermitian
translations under this variable change follows from
\(\operatorname{area}(\alpha\Lambda')=|\alpha|^2A\).
At the connected frequency orientation and \(k\ne0\), the
point \([\lambda_k]\) on \(C'\) has full \'etale order \(p\).
The added prime-to-\(p\) shift preserves this order. By
Theorem~\ref{thm:residuejet492}, the derivative in \(z'\)
has valuation \(-n\). At \(k=0\) it is integral by
Theorem~\ref{thm:trans57}. In the opposite orientation the
only primary shift is connected on \(C'\), so
Theorem~\ref{thm:torsion312} gives integrality in \(z'\)
in every order. Combining these facts with the two values
of \(v(\alpha)\) and \(v(c)\) proves the assertions.
\end{proof}

\clearpage
\section{References}\label{nav:bibliography-XIII}
\begingroup\def\section*#1{}\begin{thebibliography}{99}\small
\bibitem{XIII-X-IV-cm37} J. S. Milne, \emph{Complex Multiplication}, version 0.10, 2020, Part II. \href{https://www.jmilne.org/math/CourseNotes/CM.pdf}{Arithmetic theory of CM abelian varieties}.
\bibitem{XIII-X-IV-silver37} J. H. Silverman, \emph{The Arithmetic of Elliptic Curves}, second edition, Springer, 2009, Chapters IV and VII. \href{https://link.springer.com/book/10.1007/978-0-387-09494-6}{Formal groups and reduction}.
\bibitem{XIII-X-IV-av43} D. Mumford, \emph{Abelian Varieties}, Oxford University Press, 1970, \S17, p.~163. See also J. S. Milne, \href{https://jmilne.org/math/xnotes/AVs.pdf}{\emph{Abelian Varieties}, Remark 7.3}.
\bibitem{XIII-X-IV-nist43} NIST, \emph{Digital Library of Mathematical Functions}, Chapter 20. \href{https://dlmf.nist.gov/20.2}{Quasiperiodicity, definitions and zeros} and \href{https://dlmf.nist.gov/20.5}{Product formulas}.
\bibitem{XIII-X-IV-tz49} J. Thorner and A. Zaman, \emph{An explicit bound for the least prime ideal in the Chebotarev density theorem}, Algebra Number Theory \textbf{11} (2017), 1135--1197. \href{https://arxiv.org/abs/1604.01750}{Theorem 1.2}.
\bibitem{XIII-X-IV-dual49} C.-L. Chai and F. Oort, \emph{Moduli of abelian varieties and \(p\)-divisible groups}, Clay Math. Proc. \textbf{8} (2009), 439--536. \href{https://www2.math.upenn.edu/~chai/papers_pdf/CMIPchaioort.pdf}{Duality theorem for abelian schemes, Section 10}.
\bibitem{XIII-X-IV-gan} W. T. Gan, \emph{Automorphic forms and the theta correspondence}, Arizona Winter School lectures. \href{https://arxiv.org/abs/2303.14918}{Lecture notes}.
\bibitem{XIII-X-IV-kudla29} S. S. Kudla, \emph{Notes on the local theta correspondence}, European School of Group Theory, 1996. \href{https://www.math.toronto.edu/skudla/castle.pdf}{Chapter III, Section 6, pp.~53--54}.
\bibitem{XIII-X-I-tz50} J. Thorner and A. Zaman, \emph{A unified and improved Chebotarev density theorem}, Algebra Number Theory \textbf{13} (2019). \href{https://msp.org/ant/2019/13-5/ant-v13-n5-p02-p.pdf}{Theorem 1.4 and Theorem 3.1}.
\bibitem{XIII-X-dejong486} R. de Jong, \emph{On the Arakelov theory of elliptic curves}, Enseign. Math. (2) \textbf{51} (2005), 179--201. \href{https://arxiv.org/pdf/math/0312359}{Theorem 3.2 and Corollary 4.5}.
\end{thebibliography}\endgroup
\end{document}
